% English translation of Hermann Weyl: Die Idee der Riemannschen Fläche (1913).
% Translated for the ReadTheMasters project under CC-BY 4.0.
% Made from our own transcription (original.tex), not from any existing translation.
% Every math expression, symbol, \origpage{N} marker, and structural anchor is
% reproduced unchanged; only the prose is translated.
%
\documentclass{article}
\usepackage{readmasters}

\begin{document}
\origpage{v}

\section*{Preface and Introduction.}

The present work reproduces the principal content of a course of lectures delivered by me at the University of Göttingen in the winter semester of 1911/12, whose primary purpose was to develop the fundamental ideas of Riemann's function theory in a form that fully satisfies all modern requirements of rigor. Such a rigorous presentation---which, especially in establishing the fundamental concepts and theorems of analysis situs that enter into function theory, does not appeal to intuitive plausibility, but rather provides set-theoretically exact proofs---has not existed hitherto. The scientific work that remained to be done here may perhaps not be valued especially highly as an achievement. Nevertheless, I believe I can claim to have sought, with earnestness and conscientiousness, the \emph{simplest and most appropriate methods} leading to the intended goal; and in many places I have had to strike out on paths different from those that have become traditional in the literature since the appearance of C.~Neumann's classic book on ``Riemann's Theory of Abelian Integrals'' (1865). To a far greater extent than appears from the citations, I have been aided by the fundamental topological investigations of \emph{Brouwer} published in recent years, whose conceptual sharpness and concentration one must admire; and I secretly hope that something of the spirit that animates the works of this researcher has also come alive in this book of mine.

It was formerly customary, and has remained customary so far as I see in all presentations of the theory of Riemann surfaces up to now, to take over the notion of a \emph{curve} as it seems to be given to our sensory intuition without conceptual fixation, and to make naive use of those properties which force themselves upon us in this notion with a kind of intuitive evidence (e.g., the theorem that a curve has two banks). But ``intuitive evidence'' by no means relieves us---today there can no longer be any doubt about this---of the necessity of providing proofs for these very truths that are ultimately grounded in the axioms of arithmetic; at the very least, such proofs become necessary as soon as those fluid intuitions have expanded into general abstract concepts (as the method of mathematics as an exact science entails) and have, as it were, congealed within them. For it is certain that the general mathematical concept of a ``continuous curve'' covers much to which we find nothing corresponding in our intuition.

\origpage{vi}
And a rigorous set-theoretic foundation for the topological concepts and theorems in question for Riemann's function theory is all the more necessary since the ``points'' of which the basic configurations (curves and surfaces) consist here are not spatial points in the ordinary sense, but can be arbitrary mathematical objects of another kind (e.g., function elements). --- But to evade the difficulties that the general concept of curve brings with it through specialization, by restricting oneself to continuously differentiable or analytic curves, or even to polygons, is within analysis situs certainly an inadmissible procedure; for this discipline demands a concept of curve that behaves invariantly under arbitrary one-to-one continuous point transformations.

It cannot be denied: the discovery of the generality, extending far beyond all our intuitions, of such concepts as ``function,'' ``curve,'' etc. on the one hand, and the need for logical rigor on the other---profitable, indeed necessary, though they were for our science---have nonetheless also provoked unhealthy symptoms in the development of mathematics today. A portion of that mathematical production which strives to pursue these concepts into their subtlest nuances and---distortions, or seeks to grasp them in their widest outlines, having evaporated into the void or seeped away into side channels, has lost touch with the living stream of science. The idea of the Riemann surface also demands, if we wish to do justice to modern rigorous demands with respect to exactness, a wealth of abstract and subtle concepts and considerations for its presentation. But one need only sharpen one's gaze a little to recognize that this entire wide-meshed logical web (in which the beginner will perhaps become entangled) is not what essentially matters: it is only the \emph{net} with which we haul up the actual \emph{idea}---which in its essence is simple, grand, and divine---from the τόπος ἄτοπος, as Plato says,---like a pearl from the sea---to the surface of our intellectual world. But to grasp the kernel that this network of subtle and painstaking concepts envelops,---that which constitutes the life, the true substance, the inner value of the theory---for this a book (and even a teacher) can offer only scanty hints; here each individual must wrestle anew for understanding on their own.

One still encounters here and there the view that the Riemann surface is nothing more than an ``image,'' than a (very valuable, \emph{very suggestive}, one admits) means for visualizing and illustrating the multi-valuedness of functions. This view is fundamentally mistaken. The Riemann surface is an indispensable objective constituent of the theory; it is its very foundation. Nor is it something that is distilled a posteriori more or less artificially out of the analytic functions, but must throughout

\origpage{vii}
be regarded as the prius, as the native soil on which functions can first grow and thrive. It must admittedly be conceded that Riemann himself somewhat obscured this true relation of functions to the Riemann surface by the form of his presentation---perhaps only because he did not wish to impose all too alien notions on his contemporaries; obscured this relation also by speaking only of those multi-sheeted covering surfaces spreading out over the plane with individual winding points, of which one still thinks primarily today when Riemann surfaces are mentioned, and did not make use of the more general notion (first developed to transparent clarity later by \emph{Klein}), whose characteristic feature may be stated as this: that in it the connection to the plane of an independent complex variable, as well as the connection to three-dimensional point space in general, is fundamentally severed. And yet there can be no doubt that only in Klein's conception do Riemann's fundamental ideas come into their own in their natural simplicity, their living and penetrating power. On this conviction the present work is based.

In detail, its contents are structured in the following manner. In the \emph{first chapter}, three things are treated:

1.\ a precise discussion of the relationship of Weierstrass's concepts of ``\emph{analytic function}'' and ``\emph{analytic configuration}'' to the idea of the Riemann surface (§§~1---7);

2.\ a rigorous fixation of the concept of \emph{surface} in general and of the \emph{Riemann surface} in particular (§§~4---7);

3.\ an exact foundation of those theorems of analysis situs that are unconditionally necessary for constructing Riemann's function theory (§§~8---11). In these topological considerations, the concept of \emph{covering surface}, so important for uniformization theory, plays a far greater role (and, as I believe, rightly so) than is otherwise assigned to it.

The subject of \emph{Chapter~II} is above all the \emph{fundamental problem} of Riemann's function theory: for a given Riemann surface, to find the associated functions, especially for the case of a closed Riemann surface. The existence proofs are provided here along the path contemplated by Riemann himself with the aid of the so-called \emph{Dirichlet principle} (§§~12---15). The viability of this path was shown, as is well known, by \emph{Hilbert}, who supported with a reliable proof the Dirichlet minimal principle that was formerly held to be evident, but was then challenged by Weierstrass. By taking into account the works of other authors building upon Hilbert, and through a hitherto unpublished approach due to the author, substantially simplified in comparison to Riemann and Hilbert, we have succeeded in giving this chain of reasoning so transparent a form that in terms of simplicity it is also equal, if not

\origpage{viii}
superior, to the alternating method devised by \emph{Schwarz} and \emph{C.~Neumann}, which has hitherto alone been presented in textbooks. It demonstrates its great scope by the fact that it can be transferred without modification to open Riemann surfaces; a circumstance that benefits above all uniformization theory. §§~16---18 then bring an outline of the systematic \emph{theory of functions on a closed Riemann surface} following upon the existence theorems, but give only so much of this theory as is necessary to make clear the function-theoretic fruitfulness of Riemann's fundamental idea and to demonstrate the commanding role that the genus, originating in analysis situs, plays in the realm of functions. In doing so, contrary to custom, I have adhered to the normalization of Abelian integrals originally resulting from the proof of the existence theorems, in which real and imaginary parts must be separated from each other, since this normalization has the advantage of being independent of any dissection of the Riemann surface. Finally, the concluding sections (§§~19---21) are devoted to the \emph{theory of uniformization}, boldly sketched by \emph{Klein} and \emph{Poincaré}, and recently placed on a broad foundation by \emph{Koebe}. We thereby enter the temple in which the deity (if I may use this image) is restored to itself from the earthly bondage of its individual realizations: in the symbol of the \emph{two-dimensional non-Euclidean crystal}, the prototype of Riemann surfaces itself becomes visible, (so far as this is possible) pure and freed from all obscurities and contingencies. It was therefore clear that the decisive results of uniformization theory belonged in this book.

In the preparation of the manuscript, a write-up of my lectures prepared by Mr.~\emph{Frankfurther} for the local mathematical reading room was of great use to me; for his careful and devoted work I would like to express my sincerest thanks to Mr.~Frankfurther here as well. Besides what is reproduced here, that course of lectures also contained an outline of the theory of elliptic functions, expounded with greater completeness on algebraic functions, and presented the Riemann-Weierstrass solution of Jacobi's inversion problem with the aid of the $\vartheta$-series. Having hardly any fundamental importance for the idea of the Riemann surface, I have set these matters aside for now. Nor was it with an entirely light heart that I refrained from bridging from Riemann's function-theoretic standpoint to the curve-theoretic conception of \emph{Clebsch}, \emph{Brill}, and \emph{Noether}; but care has been taken, particularly through the citations, to open access for the reader in various directions to that literature in which information on this and on many other further developments of Riemann's theory may be found. The citations at the same time contain, albeit incompletely and in faint outline, a history of the guiding ideas that govern this most important part of function theory.

It will perhaps be found that I have been

\origpage{ix}
all too lavish with the \emph{coining of new words}. Yet I believe I have nowhere offended against the principle that an already established terminology must not be disturbed. Beyond that, however, I believe every author has not only the right, but the duty---if indeed their book possesses enough inner substance to signify anything as a whole---to give the material to be presented its own distinct form, its fitting and adequate expression. Even if some inconvenience may perhaps result from this for someone wishing to orient themselves quickly regarding the content of individual sections or to compare it with the train of thought of other works. However, through the fact that technical terms are emphasized by bold type at the place of their introduction, and through a detailed subject index located on the final pages of the book, sufficient care will, I hope, be taken for clarity as well.

Of \emph{factual prerequisites} I need not presuppose much; certainly no more than, e.g., the first volume of Osgood's well-known textbook of function theory (2nd ed., Leipzig: B.~G. Teubner, 1912) contains; I assume that the reader is familiar with the simplest examples of multi-sheeted Riemann surfaces and the symbolism of group theory. The reading does, however, I believe, require a not inconsiderable degree of abstract mathematical training; one must know how never to lose sight, amidst the seemingly complicated conceptual constructions and trains of thought, of the inwardly simple fundamental ideas around which the whole is centered.

The most important of these fundamental ideas, in so far as they do not stem directly from Riemann's creations, are due to the man to whom I have been permitted to dedicate this book in sincere and deep veneration. Privy Councillor \emph{Klein}, despite the burden of other work and despite his impaired health, did not fail to discuss the entire material with me in frequent oral conversations; for his remarks, which on several occasions prompted me to replace my original presentation with a more correct and appropriate one, I am indebted to him with the greatest gratitude. For suggestions for improvement of diverse kinds and assistance with proofreading, I also have to thank most warmly my friends, Messrs.~\emph{Koebe}, \emph{Groß}, \emph{Bieberbach}, and \emph{Weitzenböck}. How much of what I have to offer in the way of novelty in what follows goes back to earlier conversations with Koebe, I am no longer able to determine today. It was Dr.~\emph{Groß} (who assisted me in proofreading with boundless care) from whom the impetus originated in the first place for me to prepare for print my lectures on Riemann's function theory from the winter semester of 1911/12, which he attended as one of my auditors. My thanks go finally to the publishers B.~G. Teubner for the well-known care they have bestowed upon the book thus produced in typography and production.

\emph{Göttingen}, April 1913.

Hermann Weyl.

\origpage{x}

\section*{Table of Contents.}

First Chapter.

Concept and Topology of Riemann Surfaces.

Page

§~1. Weierstrass's Concept of an Analytic Function --- 1

§~2. Concept of an Analytic Configuration --- 5

§~3. Relationship of the Concepts ``Analytic Function'' and ``Analytic Configuration'' to Each Other --- 12

§~4. Concept of a Surface --- 16

§~5. Examples of Surfaces --- 25

§~6. Analytic Configurations, Considered as Surfaces --- 30

§~7. Concept of a Riemann Surface --- 34

§~8. Schlichtartig Surfaces --- 43

§~9. Covering Surfaces. Simply Connected Surfaces. Monodromy Theorem and Cauchy's Integral Theorem --- 47

§~10. One-Sidedness and Two-Sidedness of Surfaces. The Residue Theorem --- 56

§~11. Integral Functions. Genus. Canonical Dissection --- 68

Second Chapter.

Functions on Riemann Surfaces.

§~12. The Dirichlet Integral --- 78

§~13. On the Poisson Integral --- 82

§~14. Approach to the Proof of the Existence Theorems. Construction of the Elementary Differentials --- 91

§~15. Proof of the Dirichlet Principle --- 100

§~16. Relations Between the Differentials on a Closed Riemann Surface --- 108

§~17. Single-Valued Functions as a Subclass of Additive and Multiplicative Functions. Riemann-Roch Theorem and Abel's Theorem --- 117

§~18. The Algebraic Function Field --- 134

§~19. Uniformization --- 141

§~20. Riemann Surfaces and Non-Euclidean Motion Groups. Fundamental Domains. Poincaré's $\Theta$-Series --- 148

§~21. Conformal Mapping of a Riemann Surface onto Itself --- 159

Index of Technical Terms --- 166

Corrections and Additions --- 170

\origpage{1}
\section*{First Chapter. Concept and Topology of Riemann Surfaces.}

\subsection*{§ 1. Weierstrass's Concept of an Analytic Function.}

If $z$ is a complex variable and $a$ a fixed complex number, then following Weierstrass one designates as a \textbf{function element with center} $a$ any series proceeding in positive integral powers of $z - a$
\[ \mathfrak{P}(z - a) = A_0 + A_1(z - a) + A_2(z - a)^2 + \cdots, \tag{1} \]
which does not converge only for $z = a$. Otherwise, the coefficients $A_0$, $A_1$, $A_2$, \dots\ may be arbitrary complex numbers. The domain of convergence of such a power series consists either of the entire complex $z$-plane or of the interior of a definite circle $|z - a| < r$ [$r > 0$], the ``\textbf{circle of convergence}'', and a part${}^{1)}$ of the points situated on the periphery $|z - a| = r$ of this circle.

In the interior of its circle of convergence (by which, if applicable, the entire plane is also to be understood as a circle of radius $r = \infty$), such a function element represents a regular analytic function in Cauchy's sense. Conversely, it is known from the elements of function theory that a single-valued regular-analytic function in the neighborhood $|z - a| < r$ of a place $a$ can be expanded into a power series (1) convergent in this neighborhood, provided that the neighborhood belongs entirely to the domain of regularity of the analytic function. To be sure, only a circular piece of the whole function field is then brought to representation by the power series.

If one proceeds from the power series, then the endeavor must be directed toward extending the definition of the analytic function, which initially is given by the power series (1) only within the circle of convergence, over further domains of the $z$-plane in such a way that in this extension it does not forfeit its analytic character. The means for this is Weierstrass's principle of analytic

\textbf{1)} The word ``part'' (of a set) is here to be understood such that both the whole set and the ``empty set'' (containing no element) fall under this concept.

\origpage{2}
continuation.${}^{1)}$ It turns out that the plan of conquering as wide a domain of the $z$-plane as possible for the analytic function to be defined can be carried out in only one single way. But the single-valuedness of the function is in general lost in the process of analytic continuation. In this one must by no means see a defect; rather, it is a great advantage that in this way multiple-valued analytic functions also become amenable to exact treatment.

If $b$ is a value of $z$ belonging to the interior of the circle of convergence $|z - a| < r$, then by rearranging the series (1) in powers of $z - b$ there arises, as is well known, a new power series
\[ \mathfrak{Q}(z - b) = B_0 + B_1(z - b) + B_2(z - b)^2 + \cdots, \tag{2} \]
which converges at least in the circle $|z - b| < r - |b - a|$; its circle of convergence may, however, very well have a larger radius than $r - |b - a|$. Since in any case (1) and (2) agree in their values in the domain common to both their circles of convergence, (2) then provides us with an extension of the definition of our analytic function beyond the original region. We shall say that (2) is an \textbf{immediate analytic continuation} of (1). The general process of (mediate) \textbf{analytic continuation} consists in applying that of immediate analytic continuation not merely once, but an arbitrary finite number of times in succession---in an analogous way, say, to how in projective geometry a general projective mapping can be defined as the successive execution of an arbitrary number of immediate projective, that is, perspective mappings.

The \textbf{analytic continuation} can be carried out \textbf{along a given curve} $\mathfrak{c}$. This shall mean the following. Let there be given a curve emanating from the point $z = a$, that is, let to each real value $\lambda$ in the interval $0 \leqq \lambda \leqq 1$ there be assigned in a continuous manner a point $z_\lambda$ of the complex $z$-plane, in particular to $\lambda = 0$ the point $z_0 = a$, to $\lambda = 1$ a certain point $z_1 = c$. Furthermore, to each value of the parameter $\lambda$ let there also be assigned a function element $\mathfrak{P}_\lambda$ with center $z_\lambda$; let $\mathfrak{P}_0$ be the given function element (1), and let this condition moreover be fulfilled: If $\lambda_0$ is any $\lambda$-value, $\mathfrak{c}'$ an arbitrary sub-arc of $\mathfrak{c}$ lying entirely in the interior of the circle of convergence of $\mathfrak{P}_{\lambda_0}$ (defined by an inequality of the form $|\lambda - \lambda_0| \leqq \varepsilon$; $\varepsilon$ a positive constant),

\textbf{1)} Cf.\ the memoir on the ``Definition analytischer Funktionen einer Veränderlichen vermittelst algebraischer Differentialgleichungen'', written in the year 1842 and first published in the Mathematische Werke of Weierstrass, Vol.\ 1 (1894), especially pp.\ 83---84; further, the first pages in Riemann's ``Theorie der Abelschen Funktionen'' (1857) [Werke, 2nd ed., pp.\ 88---89]. Encyklopädie II B 1 (article by Osgood), No.\ 13.

\origpage{3}
then the function elements $\mathfrak{P}_\lambda$ belonging to the $\lambda$-values satisfying this inequality all arise from $\mathfrak{P}_{\lambda_0}$ by immediate analytic continuation. Then we say that it has been possible to continue the prescribed element (1) analytically along $\mathfrak{c}$, and call $\mathfrak{P}_1$ that function element with center $c$ which arises from (1) by analytic continuation along $\mathfrak{c}$. Conversely, the element $\mathfrak{P}_0$ then proceeds from $\mathfrak{P}_1$ by analytic continuation along the curve $\mathfrak{c}$ traversed in reverse.

\emph{Analytic continuation along a given curve $\mathfrak{c}$ is, if at all, possible in only one way.} For if one had two different analytic continuations
\[ \mathfrak{P}_\lambda, \quad \mathfrak{P}_\lambda^*, \]
then let $\lambda_0$ be the lower bound of all those $\lambda$-values for which the corresponding power series $\mathfrak{P}_\lambda$, $\mathfrak{P}_\lambda^*$ do not agree. Let $\mathfrak{P}_{\lambda_0}$, $\mathfrak{P}_{\lambda_0}^*$ both converge in the circle $|z - z_{\lambda_0}| < r_0$. One can then demarcate by an inequality of the form $|\lambda - \lambda_0| \leqq \varepsilon$ an arc $\mathfrak{c}_0$ on $\mathfrak{c}$ which lies entirely in the circle described about $z_{\lambda_0}$ with radius $\frac{1}{2}r_0$. The $\mathfrak{P}_\lambda$, $\mathfrak{P}_\lambda^*$ belonging to the points of this arc arise by immediate analytic continuation from $\mathfrak{P}_{\lambda_0}$, $\mathfrak{P}_{\lambda_0}^*$, and since among the $\lambda$-values satisfying that inequality there are found those for which $\mathfrak{P} \neq \mathfrak{P}^*$, we must also have $\mathfrak{P}_{\lambda_0} \neq \mathfrak{P}_{\lambda_0}^*$. Therefore $\lambda_0$ cannot be $= 0$. If then $\lambda_1$ is a value $> 0$ and $> \lambda_0 - \varepsilon$, but $< \lambda_0$, then $\mathfrak{P}_{\lambda_0}$, $\mathfrak{P}_{\lambda_0}^*$ arise by immediate analytic continuation from $\mathfrak{P}_{\lambda_1} = \mathfrak{P}_{\lambda_1}^*$. This is a contradiction. This shows in particular that the terminal element $\mathfrak{P}_1$ is uniquely determined by the initial element and the curve of continuation.

We have further to consider that by applying immediate analytic continuation a finite number of times---where the occurring element centers are successive points of the curve $\mathfrak{c}$---one can arrive from the initial element at the terminal element. Indeed: if one of the function elements $\mathfrak{P}_\lambda$ has the radius of convergence $r = \infty$, then this holds also for all others, in particular for the initial element, and the terminal element already arises from it by a single application of immediate analytic continuation. Otherwise, to each $\lambda$ there belongs a finite radius of convergence $r_\lambda$ of $\mathfrak{P}_\lambda$. If $\lambda_1$ is any value of the parameter $\lambda$, then I claim that we shall have
\[ |r_\lambda - r_{\lambda_1}| \leqq |z_\lambda - z_{\lambda_1}|, \]
so long as $\lambda$ differs sufficiently little from $\lambda_1$, and from this it follows that $r_\lambda$ depends continuously on $\lambda$ ($0 \leqq \lambda \leqq 1$). For if $\lambda_2$ ($> \lambda_1$) is such a $\lambda$-value that the entire arc of the curve $z_\lambda$ ($\lambda_1 \leqq \lambda \leqq \lambda_2$) lies in the interior of the circle described about $z_{\lambda_1}$ with radius $\frac{1}{2}r_{\lambda_1}$, then $\mathfrak{P}_{\lambda_2}$ arises from $\mathfrak{P}_{\lambda_1}$ by immediate analytic continuation, and we therefore have
\[ r_{\lambda_2} \geqq r_{\lambda_1} - |z_{\lambda_2} - z_{\lambda_1}|. \]

\origpage{4}
Since the entire arc $z_\lambda$ ($\lambda_1 \leqq \lambda \leqq \lambda_2$) is, however, also contained in the circle of convergence of $\mathfrak{P}_{\lambda_2}$, the converse holds as well:
\[ r_{\lambda_1} \geqq r_{\lambda_2} - |z_{\lambda_1} - z_{\lambda_2}|. \]
From the continuity thus established, it follows that $r_\lambda$ possesses a \emph{positive} minimum $r^0$. If we now choose the values $0 = \lambda_0 < \lambda_1 < \lambda_2 < \dots < \lambda_n = 1$ such that $|z_\lambda - z_{\lambda_h}| < r^0$ remains for $\lambda_h \leqq \lambda \leqq \lambda_{h+1}$ [$h = 0, 1, \dots n-1$], then in the sequence $\mathfrak{P}_{\lambda_0}$, $\mathfrak{P}_{\lambda_1}$, \dots\ $\mathfrak{P}_{\lambda_n}$ each function element arises from the preceding one by immediate analytic continuation.

If the continuation of the initial element along $\mathfrak{c}$ is impossible, then there is on the curve a definite point, the \textbf{``critical point''}, at which the procedure necessarily reaches its limit. Formulated more precisely: there is a threshold $\lambda = \Lambda_0$ of the following kind. If $\lambda_0 < \Lambda_0$, then the process of analytic continuation along the sub-curve $z = z_\lambda$ ($0 \leqq \lambda \leqq \lambda_0$) can be carried out, but this ceases from $\lambda_0 = \Lambda_0$ onward. Here, of course, $\Lambda_0 = 1$ is also possible, that is, only the endpoint of $\mathfrak{c}$ may be the critical point; on the other hand, $\Lambda_0 > 0$ always.

One further theorem on analytic continuation is of importance. If one has two curves
\[ z = z_1(\lambda), \quad z = z_2(\lambda), \]
which lead from the same point $a$ $\{ = z_1(0) = z_2(0) \}$ to the same endpoint $c$ and thereby always remain in sufficient proximity to each other, then the analytic continuation, if it can be performed along the first curve, can also be carried out along the second, and yields the same terminal element. The condition that the curves shall remain in sufficient proximity to each other means: there exists a positive number $\delta$ such that, whenever the inequality $|z_1(\lambda) - z_2(\lambda)| < \delta$ is satisfied for all $\lambda$, the assertion of our theorem holds. The proof follows directly from the fact that the terminal element can be obtained from the initial element by applying the process of immediate analytic continuation a finite number of times.

We are now in a position to state the general Weierstrassian definition of the analytic function as follows: \emph{An \textbf{analytic function} is the totality} $\mathbf{G}$ \emph{of all those function elements that can arise from a given function element by analytic continuation.}

Every function element of $\mathbf{G}$ can be obtained from every such element by analytic continuation.

Two analytic functions $\mathbf{G}_1$, $\mathbf{G}_2$, for which it can be shown that they have a single function element in common, are identical throughout, that is, every element of $\mathbf{G}_1$ is also contained in $\mathbf{G}_2$ and conversely.

\origpage{5}
If
\[ \mathfrak{P}(z - a) = A_0 + A_1(z - a) + A_2(z - a)^2 + \cdots \]
is a function element belonging to $\mathbf{G}$, then the number $A_0$ is called a \textbf{value} of the analytic function $\mathbf{G}$ at the point $z = a$.

To be sure, this Weierstrassian conception of the multiple-valued analytic function as a \emph{totality of function elements} has at first glance something artificial about it. When one speaks of $\sqrt{z}$ or $\lg z$, one hardly imagines the totality of those power series that are capable of representing pieces of these multiple-valued functions. Nevertheless, the Weierstrassian definition, to which simplicity and precision cannot be denied, proves itself a firm starting point for analytic function theory. Through gradual elaboration of the Weierstrassian view, we shall subsequently arrive at the Riemannian conception, in which the independent variable $z$, just like the dependent variable $u$ hitherto represented by a totality $\mathbf{G}$ of function elements, appear as single-valued analytic functions of a parameter---a parameter, to be sure, that in general varies not in a complex plane, but on a certain two-dimensional manifold, the so-called Riemann surface.

First, however, following Weierstrass, we must extend the concept of the analytic function to that of the analytic configuration.

\subsection*{§ 2. Concept of an Analytic Configuration.}

From an analytic function the analytic configuration arises by considering the function not merely, as hitherto, at those places where it behaves regularly, but by adjoining the places where it possesses a branch point of finite order or a pole (or both at once). We obtain the rigorous formulation of the concept of an analytic configuration when we extend the previous concept of a function element in an appropriate manner${}^{1)}$.

The function element (1) : $u = \mathfrak{P}(z - a)$ can be represented with the aid of a complex parameter $t$ as follows:
\[ z = a + t, \quad u = \mathfrak{P}(t) = A_0 + A_1 t + A_2 t^2 + \cdots \]
By relinquishing herein the privileged role played by $z$, and moreover permitting also finitely many negative powers of $t$, we arrive at the more general definition:

Let
\[ z = P(t), \quad u = Q(t) \]
be any two series proceeding in integral powers of $t$

\textbf{1)} See Weierstrass, Vorlesungen über die Theorie der Abelschen Transzendenten (edited by G.~Hettner and I.~Knoblauch), Werke, Vol.~4, pp.~16---19.

\origpage{6}
and containing only finitely many negative powers of $t$, of such a kind that in a certain neighborhood $|t| < r$ ($r$ a positive constant) of the origin of the $t$-plane 1) both series converge and 2) the same pair of values $(z, u)$ never arises for two distinct $t$-values of this neighborhood: then we say that this pair of power series \textbf{defines a function element}.

It is not our meaning that by ``function element'' we understand directly the pair of power series $P(t)$, $Q(t)$; rather, we regard these two series merely as one \emph{representation} of the intended function element, which besides this one admits infinitely many equally valid representations. Concerning the transition from one such representation to another, we make the following natural agreements.

If in both $P(t)$ and $Q(t)$ one replaces the parameter $t$ by a power series
\[ t(\tau) = c_1 \tau + c_2 \tau^2 + \cdots, \]
then let $P(t)$ transform into the power series $\Pi(\tau)$, and $Q(t)$ into $\mathrm{K}(\tau)$. We presuppose that $t(\tau)$ is convergent in a certain neighborhood of $\tau = 0$ and that the first coefficient $c_1 \neq 0$; then with the aid of a certain positive constant $\varrho$ we can demarcate about $\tau = 0$ a neighborhood $|\tau| < \varrho$ such that in it $t(\tau)$ 1) converges and remains in absolute value $< r$ and 2) always assumes two distinct values at two distinct places $\tau$. In this neighborhood $|\tau| < \varrho$, $\Pi$ and $\mathrm{K}$ then also converge, and at two distinct places $\tau_1$, $\tau_2$ of this neighborhood one never simultaneously has $\Pi(\tau_1) = \Pi(\tau_2)$, $\mathrm{K}(\tau_1) = \mathrm{K}(\tau_2)$. We call the pair $\Pi, \mathrm{K}$ \textbf{equivalent} to the original $P, Q$, no matter what the coefficients $c_1, c_2, \dots$ of the substituted series $t(\tau)$ may be, provided only that convergence holds and $c_1 \neq 0$. This latter presupposition entails that conversely $P(t), Q(t)$ can be obtained from $\Pi(\tau), \mathrm{K}(\tau)$ by substituting for $\tau$ a certain power series in $t$:
\[ \tau = \gamma_1 t + \gamma_2 t^2 + \cdots \quad \Bigl(\gamma_1 = \frac{1}{c_1} \neq 0\Bigr). \]

The relation of equivalence is thus a mutual one. Moreover, every pair of power series is evidently equivalent to itself, and if two pairs of power series are equivalent to a third, they are also equivalent to each other. These facts entitle us to regard equivalent pairs of power series as \emph{representations of the same}, and non-equivalent ones as \emph{representations of distinct} function elements. Or expressed differently: \emph{Two pairs of power series, each defining a function element, define the same function element if and only if they are equivalent.}

\origpage{7}
We rely here on a mode of explanation that must often be employed elsewhere in mathematics as well, and which has its psychological roots in the capacity of our mind for \emph{abstraction}. This type of definition rests on the following general principle: If between things of any domain of operations there is defined a relation $\sim$ \emph{having the character of an equivalence} [i.e., a relation satisfying the rules:
\[ \begin{gathered} \text{1.\ } a \sim a; \quad \text{2.\ from } a \sim b \text{ follows } b \sim a; \\ \text{3.\ from } a \sim c, \ b \sim c \text{ follows } a \sim b], \end{gathered} \]
then it is possible to regard each of the objects $a$ of that original domain of operations as a \emph{representative} of a thing $\alpha$ in such a way that two objects $a, b$ appear as representatives of the same thing $\alpha$ if and only if they are equivalent to each other in the sense of the relation $\sim$. This principle will in particular always be applicable when in the objects $a, b, \dots$ we are interested only in those properties which are invariant under the relation $\sim$. Its application has the effect of replacing an unwieldy mode of expression by a shorter one, which takes account of the prevailing point of view of interest in the investigation by automatically shedding everything that is \emph{inessential} from the standpoint of this interest in the investigated objects. I mention here two examples of such a ``\emph{definition by abstraction}''.

1.\ Of two parallel straight lines one says they have the same \emph{direction}, of two non-parallel ones, they have different directions. The original objects ($a$) are the straight lines, the relation having the character of an equivalence is parallelism; what is required is to assign to each straight line a something, its ``direction'', in such a way that the identity of the assigned ``directions'' corresponds to the parallelism of the straight lines.

2.\ A ``motion'' (of a point) is given when the position of the moving point $p$ is given at each moment $\lambda$ of a certain time interval $\lambda_0 \leqq \lambda \leqq \lambda_1$: $p = p(\lambda)$. If one has two such motions $p = p(\lambda)$, $q = q(\mu)$, then one says that these motions traverse the same ``\emph{path}'' if and only if $\lambda$, the time parameter of the first motion, can be expressed as a continuous monotonically increasing function of the time parameter $\mu$ of the second motion: $\lambda = \lambda(\mu)$, in such a way that the first motion thereby transforms into the second: $p(\lambda(\mu)) \equiv q(\mu)$. Here it is the concept of ``path'' that is to be defined in this way${}^{1)}$.

From this digression we return to our extended concept of a function element. Among all equivalent representations

\textbf{1)} This concept means something different from the \emph{point set} consisting of all points passed through during the motion. What is involved here is the same difference as that between the path traversed by a pedestrian (which, as long as the pedestrian marches, is \emph{in statu nascendi}) and the (long existing) path \emph{on which} he marches.

\origpage{8}
of one and the same function element, we shall attempt to single out a definite, simplest possible one as the ``\emph{normal representation}''. In doing so, we distinguish several cases. If $P(t)$ contains no negative powers of $t$:
\[ z = P(t) = a + a_1 t + a_2 t^2 + \cdots \]
and if $a_1 \neq 0$, then we can introduce $z - a = \tau$ as a new representing parameter and obtain
\[ z = a + \tau, \quad u = \mathrm{K}(\tau) \tag{3} \]
as a new representation of the same function element. If $Q(t)$ also contains no negative powers, then the same holds for $\mathrm{K}$, and we have before us such a function element as we considered in §~1, which we now wish to designate for distinction as \textbf{regular function elements}. (3) is the sought-for normal representation. A regular function element admits only a single normal representation, and two regular function elements are therefore certainly \emph{distinct if their normal representations are distinct.} Only on the basis of these circumstances are we genuinely justified in designating our present concept of a function element as an extension of the one taken as basis in §~1.

If no powers of $t$ with negative exponent occur in the expansion of $z$, but if we assume more generally that
\[ z = a + a_\mu t^\mu + a_{\mu+1} t^{\mu+1} + \cdots \quad (a_\mu \neq 0), \]
so that $a_\mu$ ($\mu \geqq 1$) is the first coefficient distinct from $0$ besides the constant term, then for $t$ one can substitute such a power series in $\tau$
\[ t = c_1 \tau + c_2 \tau^2 + \cdots \quad (c_1 \neq 0) \]
that
\[ z = a + \tau^\mu \]
holds, and we obtain a representation of the form
\[ z = a + \tau^\mu, \quad u = \mathrm{K}(\tau). \tag{3*} \]
Indeed: if $\sqrt[\mu]{a_\mu}$ is a definite one of the $\mu$ roots, then as is well known there exists a unique power series $\gamma_1 + \gamma_2 t + \gamma_3 t^2 + \cdots$ with leading term $\gamma_1 = \sqrt[\mu]{a_\mu}$, whose $\mu^{\text{th}}$ power is $= a_\mu + a_{\mu+1} t + \cdots$. By solving
\[ \tau = \gamma_1 t + \gamma_2 t^2 + \cdots \]
for $t$, one obtains the desired result. But according to which of the $\mu$ roots $\sqrt[\mu]{a_\mu}$ one takes, one obtains $\mu$ distinct representations (3*): they all arise from one of them by replacing $\tau$ by $\tau \cdot \zeta$, where $\zeta$ denotes an arbitrary $\mu^{\text{th}}$ root of unity. A function element given by
\[ z = a' + \tau^{\mu'}, \quad u = \mathrm{K}'(\tau) \quad [\mu' \text{ integral and positive}] \tag{4} \]

\origpage{9}
can be identical with the one represented by (3*) only if $a' = a$; $\mu' = \mu$, and if (4) agrees coefficient for coefficient with one of the $\mu$ normal representations arising from (3*) by the substitutions $\tau \mid \tau \cdot \zeta$. In particular, the integer $\mu$ is accordingly characteristic for the function element in question and independent of its particular representation; we say that the element is \textbf{ramified of the $(\mu - 1)^{\text{th}}$ order}; in the case $\mu = 1$, which we considered above, it is called \textbf{unramified}.

If negative powers of $t$ occur in the expansion of $z$, let $t^{-\nu}$ be the lowest:
\[ z = a_{-\nu} \cdot t^{-\nu} + a_{-\nu+1} t^{-\nu+1} + \cdots \quad (a_{-\nu} \neq 0). \]
One can replace $t$ by such a convergent power series in $\tau$: $t = c_1 \tau + \cdots$ ($c_1 \neq 0$) that
\[ z = \tau^{-\nu} \quad \text{and then} \quad u = \mathrm{K}(\tau) \tag{3**} \]
results. This is now the normal representation; when $\nu > 1$, it is not uniquely determined by the element, but there are precisely $\nu$ distinct ones, all of which proceed from one by replacing $\tau$ by $\tau \cdot \zeta$, where $\zeta$ is an arbitrary $\nu^{\text{th}}$ root of unity. Here too one speaks of a ramification of the $(\nu - 1)^{\text{th}}$ order.

In deriving the normal representations (3), (3*), (3**), we have singled out the variable $z$ over $u$, as in §~1. Alongside the regular function elements have stepped the irregular ones, alongside the unramified the ramified. Now by transferring the concepts of immediate and mediate analytic continuation to arbitrary (including irregular) function elements, we arrive without further ado at the definition of the \textbf{analytic configuration}.

Let $\mathfrak{e}$ be a function element and
\[ z = P(t), \quad u = Q(t) \tag{5} \]
any representation of it, $|t| < r$ ($r > 0$) any circle in which this representation is \textbf{valid} (that is, in which $P(t)$, $Q(t)$ converge and $P(t_1) = P(t_2)$, $Q(t_1) = Q(t_2)$ never simultaneously occur for $t_1 \neq t_2$, $|t_1| < r$, $|t_2| < r$). For each value $t_0$ for which $|t_0| < r$, we can then rearrange the series $P(t)$, $Q(t)$ in powers of $t' = t - t_0$ and thus obtain a new pair of power series $P'(t')$, $Q'(t')$ and with it a new function element $\mathfrak{e}_{t_0}$. Of all the function elements $\mathfrak{e}_{t_0}$ thus belonging to the various $t_0$ situated in the circle $|t_0| < r$, we say that they form an \textbf{analytic neighborhood} of the original element $\mathfrak{e} (= \mathfrak{e}_0)$, and this designation shall be applied whichever representation of the element $\mathfrak{e}$ by a parameter $t$ and whichever circle $|t| < r$ in which this representation is valid was chosen. The analytic neighborhood described above, determined by the representation (5) together with the inequality $|t| < r$,

\origpage{10}
we call a \textbf{$t$-neighborhood} with respect to the representing parameter, where a brief designation is desired.

If one has two distinct representations of the same function element $\mathfrak{e}$:
\[ z = P(t), \quad u = Q(t); \quad \text{resp.} \quad z = \Pi(\tau), \quad u = \mathrm{K}(\tau) \]
(which transform into one another by the substitution
\[ t = t(\tau) = c_1 \tau + c_2 \tau^2 + \cdots \quad [c_1 \neq 0] \]
into one another), then the following important theorem holds:

\emph{Every $t$-neighborhood of $\mathfrak{e}$ contains a $\tau$-neighborhood of $\mathfrak{e}$, and therefore conversely every $\tau$-neighborhood contains a $t$-neighborhood.}

Proof: Let the $t$-neighborhood be determined by $|t| < r$. We then choose a positive number $\varrho$ such that in $|\tau| < \varrho$ the power series $t(\tau)$ converges, remains in absolute value $< r$, and assumes no value twice. This inequality $|\tau| < \varrho$ determines a $\tau$-neighborhood, which I claim is contained in the original $t$-neighborhood. By rearranging $\Pi(\tau), \mathrm{K}(\tau)$ in powers of $\tau' = \tau - \tau_0$ ($|\tau_0| < \varrho$) let there arise the representation $\Pi'(\tau'), \mathrm{K}'(\tau')$ of the element $\mathfrak{e}_{\tau_0}$, and by rearranging $P(t), Q(t)$ in powers of $t' = t - t_0$ [$t_0 = t(\tau_0)$, $|t_0| < r$] the element $(P'(t'), Q'(t')) = \mathfrak{e}_{t_0}$; then $\mathfrak{e}_{\tau_0} = \mathfrak{e}_{t_0}$. For by rearranging $t(\tau)$ in powers of $\tau' = \tau - \tau_0$, one obtains
\[ t'(\tau') = t(\tau) - t_0 = c_1' \tau' + c_2' \tau'^2 + \cdots, \tag{6} \]
and one evidently has
\[ P'(t'(\tau')) = \Pi'(\tau'), \quad Q'(t'(\tau')) = \mathrm{K}'(\tau'); \]
for in a certain neighborhood of the point $\tau_0$ of the complex $\tau$-plane, $\Pi'(\tau - \tau_0)$, $P'(t(\tau) - t_0)$, for instance, are regular functions that agree in value with $\Pi(\tau)$. But in the substitution (6), which transforms the pair $P'(t'), Q'(t')$ into $\Pi'(\tau'), \mathrm{K}'(\tau')$, the first coefficient
\[ c_1' = \left[ \frac{dt(\tau)}{d\tau} \right]_{\tau = \tau_0} \neq 0, \]
since otherwise the function $t(\tau)$ in the vicinity of $\tau = \tau_0$ would always assume one and the same value $t$ at at least two places. The identity of $\mathfrak{e}_{\tau_0}, \mathfrak{e}_{t_0}$ is thereby proved.

The concept of \textbf{analytic neighborhood} introduced here corresponds, in a form more expedient for further formulations, to the concept of immediate analytic continuation utilized in its place in §~1. In this concept of analytic neighborhood it also becomes apparent that the irregular elements are to be regarded merely as exceptional places compared to the regular ones; for the analytic neighborhood of any element, if only taken sufficiently small, always consists (apart perhaps from this element itself) exclusively of regular function elements. In order to see this, one need only choose the circle $|t| < r$ determining the sought-for neighborhood of $\mathfrak{e} = (P(t), Q(t))$ so

\origpage{11}
small that throughout it (except perhaps for $t = 0$) one has $\frac{dP}{dt} \neq 0$.

An \textbf{analytically connected sequence of function elements} is given when to each real value $\lambda$ of the interval $0 \leqq \lambda \leqq 1$ there is assigned a function element $\mathfrak{e}(\lambda)$ such that the following condition is fulfilled: If $\lambda_0$ is any value of the parameter $\lambda$, $\mathfrak{e}_0 = \mathfrak{e}(\lambda_0)$ the associated function element, $\mathfrak{U}_0$ an arbitrary analytic neighborhood of $\mathfrak{e}_0$, then there always exists a positive number $\varepsilon$ such that $\mathfrak{e}(\lambda)$ belongs to $\mathfrak{U}_0$ as long as $|\lambda - \lambda_0| \leqq \varepsilon$. The analytically connected sequence \textbf{connects} the initial element $\mathfrak{e}(0)$ with the terminal element $\mathfrak{e}(1)$.

\emph{An \textbf{analytic configuration} is a totality} $\mathbf{G}$ \emph{of function elements with the following properties:}

\emph{1.\ Any two function elements belonging to} $\mathbf{G}$ \emph{can be connected with each other by an analytically connected sequence consisting entirely of function elements belonging to} $\mathbf{G}$.

\emph{2.\ } $\mathbf{G}$ \emph{cannot be extended by the addition of function elements in any way such that the extended totality still possesses property 1.}

From 2.\ it follows in particular that along with a function element $\mathfrak{e}$, every analytic neighborhood of $\mathfrak{e}$ also belongs to $\mathbf{G}$.

We can bring these definitions closer to ourselves through an analogy. For in function theory the Weierstrassian function element evidently assumes the same role that the \emph{point} as spatial element plays in geometry, say of three-dimensional space. We therefore place the function element in analogy with the point, and, as of three-dimensional point space, we then speak of the ``space (that is: nothing other than the totality) of function elements''; to be sure, since the determination of a function element depends on infinitely many continuously variable determining constituents, we must ascribe infinitely many dimensions to this space. All concepts concerning the \emph{continuity} of three-dimensional space can be traced back to the single one: ``\emph{neighborhood} of a point'' (as the neighborhood of a point we consider, say, the interior of every sphere described about this point as center). In the space of function elements, what we called above ``analytic neighborhood'' shall furnish the analogue of the ordinary concept of ``neighborhood'' in point space. Then the \emph{analytically connected sequence of function elements} corresponds quite precisely to the \emph{continuous curve} in point space. The space of function elements possesses, according to these stipulations, a substantially different structure from ordinary three-dimensional space: for whereas point space constitutes a single connected whole (any two of its points can be connected by a continuous curve),

\origpage{12}
the infinite-dimensional space of function elements decomposes into infinitely many (two-dimensional) ``layers''; each such layer by itself is a continuously (that is, here: analytically) connected whole, but the individual layers have no connection with each other at any place; these ``layers'' are none other than the analytic configurations. In Euclidean point space, one designates as a ``\emph{domain}'' a point set of such a nature that 1.\ to each point of the set there exists a neighborhood belonging entirely to the set, and 2.\ any two points of the set can be connected by a continuous curve consisting entirely of points of the set. Accordingly, in the space of function elements the analytic configuration would be designated as a \emph{domain incapable of further extension}. The will to conceive analytic configurations on the basis of the discussed analogy as \emph{two-dimensional manifolds} immediately leads us into the midst of the Riemann-Kleinian modes of visualization. The ordinary Riemann surfaces presented in elementary textbooks do indeed have precisely this significance: to represent each function element of the analytic configuration by a single point of the surface in such a way that analytically connected sequences of function elements of the configuration appear as continuous curves on the Riemann surface.

Before we can pursue further the trains of thought stimulated by this, however, we must still orient ourselves concerning the relationship of the two concepts of ``analytic function'' and ``analytic configuration''.

\subsection*{§ 3. Relationship of the Concepts ``Analytic Function'' and ``Analytic Configuration'' to Each Other.}

A first remark is this: If one has succeeded, in the manner depicted in §~1, in continuing a regular function element regularly along a given curve $z = z(\lambda) [0 \leqq \lambda \leqq 1]$ [whereby to each value $\lambda$ there is assigned a regular function element $\mathfrak{e}(\lambda)$, in which the expansion of $z$ in powers of the representing parameter $t$ begins with the constant term $z(\lambda)$], then one has thereby obtained an analytically connected sequence of function elements in the sense of §~2. In truth: if $\lambda_0$ is any $\lambda$-value, $\mathfrak{U}_0$ an arbitrary analytic neighborhood of $\mathfrak{e}(\lambda_0) = \mathfrak{e}_0$, then on account of the regularity of $\mathfrak{e}_0$ one can choose the variable $z' = z - z(\lambda_0)$ as the representing parameter of $\mathfrak{e}_0$, and there then exists a $z'$-neighborhood of $\mathfrak{e}_0$, defined say by $|z'| < r_0$, which is contained entirely in $\mathfrak{U}_0$. If about $z(\lambda_0)$ one demarcates by the inequality $|\lambda - \lambda_0| \leqq \varepsilon (\varepsilon > 0)$ an arc for which throughout $|z(\lambda) - z(\lambda_0)| < r_0$, then the $\mathfrak{e}(\lambda)$ corresponding to the points of this arc arise from $\mathfrak{e}_0$ by rearranging in powers of $z - z(\lambda)$, and accordingly belong to the $z'$-neighborhood of

\origpage{13}
$\mathfrak{e}_0$ determined by $|z'| < r_0$ and thus also to the neighborhood $\mathfrak{U}_0$. --- The converse, that an analytically connected sequence of regular function elements is obtained from the initial element by the process of analytic continuation described in §~1, requires no proof.

\emph{In an analytically connected sequence $\mathfrak{e}(\lambda) [0 \leqq \lambda \leqq 1]$ of function elements, there always occur only finitely many irregular (in particular also only finitely many ramified) function elements.} For otherwise those values $\lambda$ to which irregular elements correspond would possess an accumulation point $\lambda_0$, and there would lie in every analytic neighborhood of $\mathfrak{e}(\lambda_0)$ infinitely many irregular elements, which is impossible.

Let the expansion of $z$ in the representation of the function element $\mathfrak{e}(\lambda)$ begin with the constant term $z(\lambda)$ [for the sake of convenience we exclude the possibility that it begins with negative powers of the representing parameter $t$]. We wish to consider the case where ramified function elements occur among the $\mathfrak{e}(\lambda)$; let $\mathfrak{e}(0)$ be regular, but $\lambda_0 (< 1)$ the smallest $\lambda$-value for which $\mathfrak{e}(\lambda)$ is irregular; let $\mu - 1$ be the order of ramification of $\mathfrak{e}(\lambda_0)$. The elements $\mathfrak{e}(\lambda)$ from $\lambda = 0$ up to $\lambda_0$ (excl.) are then obtained uniquely by analytic continuation along the given curve $z = z(\lambda)$. But this changes from $\lambda_0$ onward; beyond $\mathfrak{e}(\lambda_0)$, analytic continuation along the given curve can be continued in precisely $\mu$ distinct ways: the trunk divides into $\mu$ branches, along one of which the given analytically connected sequence runs. Each of these branches can in its further course ramify once again, and the next ramification of one branch can perfectly well occur at a different place from that for the other branch; and so forth. A branch can also naturally, if one reaches a critical place that means not merely a pole or branch point of finite order, cease entirely before one has reached $\lambda = 1$ on it. From this description one may recognize how felicitously the intuitive conception underlying Riemann's designation ``branch point'' captures the essence of the matter. For the proof of our assertions the following simple consideration suffices.

Let
\[ z = z(\lambda_0) + t^\mu, \quad u = Q(t) \]
be one of the $\mu$ normal representations of $\mathfrak{e}(\lambda_0)$, which may be valid for $|t| < r$. Let a $\lambda_1 > \lambda_0$ be chosen such that $|z(\lambda) - z(\lambda_0)| < r^\mu$ remains for $\lambda_0 \leqq \lambda \leqq \lambda_1$. Then there exist in the $t$-plane $\mu$ curve branches emanating from the origin $t = 0$, all of which remain in the circle $|t| < r$ and transform by the mapping $z = z(\lambda_0) + t^\mu$ into one and the same curve arc $z = z(\lambda) [\lambda_0 \leqq \lambda \leqq \lambda_1]$; these $\mu$ curve branches proceed from one another by rotation about the origin by $\frac{2\pi}{\mu}, \frac{4\pi}{\mu}, \dots$. By continuing $Q(t)$ along each of these curve branches (immediate analytic

\origpage{14}
continuation!), one obtains the $\mu$ possible continuations of $\mathfrak{e}(\lambda_0)$ along the prescribed curve beyond $\lambda = \lambda_0$ up to $\lambda = \lambda_1$. That each of these continuations can be carried through to the end ($\lambda = 1$) as an analytically connected sequence is not asserted and is also in general not correct.

The circumstance that in every analytically connected sequence of function elements only finitely many irregular ones occur makes it possible to \emph{bypass} these irregular elements. Two function elements that can be connected by an analytically connected sequence at all can thus also be connected by one which (apart perhaps from the initial and terminal elements) consists exclusively of regular elements. For the proof we assume for the sake of simplicity that the analytically connected sequence $\mathfrak{e}(\lambda)$ contains \emph{only one} irregular element $\mathfrak{e}(\lambda_0) [0 < \lambda_0 < 1]$. Let
\[ z = P(t), \quad u = Q(t) \]
be a representation of it valid for $|t| < r$; here let $r$ be chosen so small that the elements of the $t$-neighborhood $\mathfrak{U}_0$ of $\mathfrak{e}(\lambda_0)$ determined by $|t| < r$ are, apart from $\mathfrak{e}(\lambda_0)$ itself, all regular. One can take two values $\lambda_1 < \lambda_0$ and $\lambda_2 > \lambda_0$ such that all $\mathfrak{e}(\lambda) [\lambda_1 \leqq \lambda \leqq \lambda_2]$ belong to $\mathfrak{U}_0$. $\mathfrak{e}(\lambda_1)$ arises from the underlying representation of $\mathfrak{e}(\lambda_0)$ by rearranging it in powers of $t - t_1$ (where $t_1$ is a certain point belonging to the circle $|t| < r$), $\mathfrak{e}(\lambda_2)$ by rearranging in powers of $t - t_2$. The two points $t_1$, $t_2$ in the $t$-plane can be connected within the circle $|t| < r$ by a curve that does not pass through the origin. By assigning to each point $t_0$ of this curve that function element which one obtains by rearranging $P(t)$, $Q(t)$ in powers of $t - t_0$, one succeeds in connecting $\mathfrak{e}(\lambda_1)$ with $\mathfrak{e}(\lambda_2)$ by an analytically connected sequence consisting of purely regular function elements.

From the proved facts it follows:

1.\ \emph{All the regular function elements of an analytic configuration constitute a single analytic function.}

2.\ \emph{Every analytic function consists of all the regular function elements of an analytic configuration uniquely determined by the function.}

To this is added the further theorem:

3.\ \emph{The irregular function elements of an analytic configuration are present only in countable quantity.}

We conduct the proof with the aid of the theorem stated by Poincaré and Volterra${}^{1)}$,

\textbf{1)} Poincaré, Rendiconti del Circolo matematico di Palermo, Vol.~2 (1888), pp.~197---200. Volterra, Atti della Reale Academia dei Lincei, Ser.~4, $\mathrm{IV}_2$, p.~355.

\origpage{15}
\emph{that in an analytic configuration there are at most countably infinitely many regular function elements $u = \mathfrak{P}(z - a)$ with prescribed center $z = a$.} For all these elements can be generated from one of them, $\mathfrak{P}_1$, by continuing $\mathfrak{P}_1$ in a regular manner along curves in the $z$-plane that emanate from $a$ and return thither. For every such curve, however, one can construct a polygonal path consisting of finitely many rectilinear segments, which runs in such proximity to the curve that regular analytic continuation along the polygonal path is likewise possible and leads to the same terminal element as continuation along the curve. In doing so, one can furthermore ensure that the vertices of this polygonal path possess rational coordinates relative to $a$; let these relative coordinates $z - a$ be:
\[ \frac{n_1' + i n_1''}{n_1}, \quad \frac{n_2' + i n_2''}{n_2}, \dots, \frac{n_h' + i n_h''}{n_h} \qquad (i = \sqrt{-1}), \]
where in each case $n_f'$, $n_f''$, $n_f$ are integers without common divisor and $n_f > 0$ [$f = 1, 2, \dots, h$]. To this polygonal path we assign the number
\[ \sum_{f=1}^h |n_f'| + \sum_{f=1}^h |n_f''| + \sum_{f=1}^h n_f = N \]
to it.
Among the polygonal paths emanating from $a$ and returning to $a$ whose vertices differ from $a$ by rational complex numbers, there are certainly only finitely many to which the same number $N$ belongs. If I choose $N$ successively $= 3, 4, 5, \dots$, I thereby arrange all these polygonal paths into an enumerated sequence. Each of these polygonal paths determines either one or no function element with center $a$, according to whether the regular continuation of $\mathfrak{P}_1$ along the polygonal path can be accomplished or not. But in this way one also securely obtains all regular elements belonging to the analytic configuration with center $a$, whose countability is thereby proved.

Instead of the $z$-plane I can make use of the $z$-sphere proceeding from it by stereographic projection, on which $z = \infty$ is also represented by a single point. If
\[ z = a + t^\mu, \ u = Q(t); \quad \text{resp.} \quad z = t^{-\mu}, \ u = Q(t) \]
is an irregular element of a given analytic configuration in its normal representation, which may be valid for $|t| < r$, then for each value $z_0$ ($\neq a$ resp.\ $\infty$) satisfying the condition
\[ |z - a| < r^\mu \quad \text{resp.} \quad |z| > r^{-\mu} \tag{7} \]
there exist precisely $\mu$ regular function elements $u = \mathfrak{P}(z - z_0)$ with center $z_0$ belonging to the $t$-neighborhood of the irregular element determined by $|t| < r$. We say briefly: the irregular element

\origpage{16}
``induces'' at the point $z_0$ those $\mu$ regular elements. The inequality (7) signifies on the sphere a spherical cap, which I now prefer, however, to replace by that largest spherical cap $K$ not extending beyond it whose center lies at $a$ (resp.\ $\infty$); let the radius of $K$ be called $\varkappa$. If I connect $z_0$, lying in $K$, with the center within the spherical cap by a great-circle arc, then the irregular element arises from each of the $\mu$ elements induced at $z_0$ by analytic continuation along this circular arc from $z_0$ to $a$ (resp.\ $\infty$). If one therefore has two distinct irregular elements with the same center $a$ resp.\ $\infty$, it is excluded that they induce two identical regular elements at a point $z_0$. If one has two irregular elements $\mathfrak{e}_1$, $\mathfrak{e}_2$ with distinct centers and we denote the associated spherical caps $K$ by $K_1$, $K_2$, their radii by $\varkappa_1$, $\varkappa_2$, then let both spherical caps shrink about their center in such a way that they possess only half the radius $\frac{1}{2}\varkappa_1$, $\frac{1}{2}\varkappa_2$. Let, say, $\varkappa_1 \geqq \varkappa_2$. If these two smaller spherical caps $k_1$, $k_2$ overlap, then let $z_0$ denote a point common to them. I claim then that none of the elements induced by $\mathfrak{e}_1$ at $z_0$ is identical with an element induced there by $\mathfrak{e}_2$. For the center of $K_2$ lies within the spherical cap $K_1$, as does the great-circle arc connecting $z_0$ with this center within $K_2$. By continuing one of the regular elements induced by $\mathfrak{e}_1$ at $z_0$ along this circular arc leading to the center of $K_2$, one obtains one of the regular elements induced there by $\mathfrak{e}_1$, but never the irregular element $\mathfrak{e}_2$. Our assertion is thereby established.

In the manner now that we have assigned the (small) spherical caps $k_1$, $k_2$ to the irregular elements $\mathfrak{e}_1$, $\mathfrak{e}_2$ here, we assign a spherical cap $k$ to each irregular element of the given analytic configuration. If the irregular elements were not present merely in countable quantity, there would have to exist a rational point $z_0$ contained in more than merely countably many spherical caps $k$. Each of the irregular elements to which these spherical caps belong induces at the point $z_0$ at least one regular element belonging to the analytic configuration that has $z_0$ as center, and indeed, as we have just demonstrated, distinct irregular elements always induce distinct ones. But that is impossible, since there can be no more than countably many regular function elements with center $z_0$ in the given analytic configuration.

\emph{The analytic configuration accordingly differs from the analytic function only in that countably many irregular elements have been adjoined.}

\subsection*{§ 4. Concept of a Surface.}

It was already mentioned at the close of §~2 that an analytic configuration gains greatly in intuitive clarity when one succeeds

\origpage{17}
in representing each element of the configuration by a point of a spatial surface $\mathfrak{F}$ in such a way that these representing points together simply cover the surface $\mathfrak{F}$, and each analytically connected sequence of elements of the analytic configuration appears as a continuous curve on $\mathfrak{F}$. The problem of discovering such a surface $\mathfrak{F}$ that visualizes the analytic configuration can, to be sure, from a purely objective standpoint be rejected as an inappropriate formulation of the question, since three-dimensional space intrinsically has nothing whatever to do with analytic configurations, and one evidently refers to it not for logical-mathematical reasons at all, but because it stands particularly close to our sensory spatial intuition; one might designate it as an anthropomorphism contrary to scientific principle to force in this way alien notions upon things, instead of taking them as they are, merely to satisfy our need for ``images and parables''. These reproaches of the pure logician do not touch us, however, if we pursue the other conception, also already touched upon, according to which \emph{the analytic configuration itself appears as a two-dimensional manifold}, to which all those concepts of continuity that encounter us in ordinary geometry can be transferred. On the contrary: it would mean overlooking one of the most essential sides of the subject if one did not make use of this notion.

The concept of ``two-dimensional manifold'' or of ``surface'' shall therefore for us not be bound to the idea of the spatial point, but shall receive a much more general abstract meaning. Whenever any totality of things (which will assume the role of ``points'') is given at all, and by definition there exists between them a continuous connection similar to that between the points of a plane, we speak of a two-dimensional manifold. But since all continuity concepts can be traced back to the single one of neighborhood, there belong to the definition of a two-dimensional manifold two things:

1.\ Specification of those things which are to count as ``points'' of the manifold;

2.\ A definition of the concept of ``neighborhood''.

In precise formulation: when do we wish to say that a \textbf{two-dimensional manifold $\mathfrak{F}$ is given}? When the following is the case:

\emph{Given a totality of things called} \textbf{``points of the manifold $\mathfrak{F}$''}. \emph{For each point $\mathfrak{p}$ of the manifold $\mathfrak{F}$ certain sets of points of the manifold are defined as} \textbf{``neighborhoods of $\mathfrak{p}$ on $\mathfrak{F}$''}. \emph{Every neighborhood $\mathfrak{U}_0$ of a point $\mathfrak{p}_0$ of the manifold on $\mathfrak{F}$ must contain $\mathfrak{p}_0$ itself and admit a one-to-one mapping onto the interior points of an ordinary Euclidean circle}

\origpage{18}
\emph{$K_0$ (where $\mathfrak{p}_0$ transforms into the center of the circle) of the following nature:}

1.\ \emph{if $\mathfrak{p}$ is any point of $\mathfrak{U}_0$ and $\mathfrak{U}$ a neighborhood of $\mathfrak{p}$ on $\mathfrak{F}$ consisting only of points of $\mathfrak{U}_0$, then the image of $\mathfrak{U}$ (projected into $K_0$ by that mapping) contains the image point of $\mathfrak{p}$ in its interior; that is, about the image point $p$ of $\mathfrak{p}$ a circular disk $k$ can be described such that every point of $k$ is the image of a point of $\mathfrak{U}$;}

2.\ \emph{if $K$ is the interior of any circle situated entirely in $K_0$ with center $p$, then there always exists a neighborhood $\mathfrak{U}$ of $\mathfrak{p}$ on $\mathfrak{F}$ whose image lies entirely in $K$.}

We now explain briefly how on the basis of the concept of neighborhood all continuity concepts can be transferred from the ordinary plane to arbitrary two-dimensional manifolds.

If a set $\mathfrak{E}$ of points of the manifold $\mathfrak{F}$ is given, then the point $\mathfrak{p}$ of $\mathfrak{F}$ is called an \textbf{accumulation point} (\textbf{limit point}) of the set $\mathfrak{E}$ if in every neighborhood of $\mathfrak{p}$ on $\mathfrak{F}$ there lie infinitely many points of $\mathfrak{E}$. --- A point belonging to the set $\mathfrak{E}$ that is not an accumulation point of $\mathfrak{E}$ is called an \textbf{isolated point} in $\mathfrak{E}$. On the other hand, $\mathfrak{p}$ is an \textbf{interior point} of $\mathfrak{E}$ if there exists a neighborhood of $\mathfrak{p}$ on $\mathfrak{F}$ all of whose points belong to $\mathfrak{E}$. --- $\mathfrak{E}$ is called \textbf{closed} if every accumulation point of $\mathfrak{E}$ belongs to $\mathfrak{E}$.

A continuous \textbf{curve} $\gamma$ on $\mathfrak{F}$ is given when to each real number $\lambda$ of the interval $0 \leqq \lambda \leqq 1$ there is assigned a point $\mathfrak{p}(\lambda)$ on $\mathfrak{F}$ in such a way that the condition of continuity is fulfilled. This states: If $\lambda_0$ is any value of the parameter $\lambda$, $\mathfrak{U}_0$ any neighborhood of the point $\mathfrak{p}(\lambda_0)$ on $\mathfrak{F}$, then there always exists a positive number $\varepsilon$ of such a kind that for all $\lambda$ satisfying the condition $|\lambda - \lambda_0| \leqq \varepsilon$, the point $\mathfrak{p}(\lambda)$ belongs to $\mathfrak{U}_0$. --- The curve $\gamma$ \textbf{connects} the \textbf{initial point} $\mathfrak{p}(0)$ with the \textbf{terminal point} $\mathfrak{p}(1)$.

A point set $\mathfrak{E}$ on $\mathfrak{F}$ is called a \textbf{domain} if 1) every point of $\mathfrak{E}$ is an interior point of this point set and 2) any two points of $\mathfrak{E}$ can be connected by a continuous curve on $\mathfrak{F}$ consisting exclusively of points of $\mathfrak{E}$. --- \emph{We restrict the concept of two-dimensional manifold $\mathfrak{F}$ utilized hitherto by the further requirement that $\mathfrak{F}$ itself shall be a domain (that is, consist of one piece).}

If to each point $\mathfrak{p}$ of a certain point set $\mathfrak{E}$ on $\mathfrak{F}$ there is assigned a real or complex number $f(\mathfrak{p})$, then a \textbf{function} $f$ is defined in $\mathfrak{E}$; the number $f(\mathfrak{p})$ is the \textbf{value} of this function at the point $\mathfrak{p}$. If $\mathfrak{p}_0$ denotes an arbitrary point of $\mathfrak{E}$, then $f$ is called \textbf{continuous} at the point $\mathfrak{p}_0$ if for every positive number $\varepsilon$ there exists a neighborhood $\mathfrak{U}_\varepsilon$ of $\mathfrak{p}_0$ on $\mathfrak{F}$ such that $|f(\mathfrak{p}) - f(\mathfrak{p}_0)| < \varepsilon$ remains for all $\mathfrak{p}$ belonging both to $\mathfrak{E}$ and to $\mathfrak{U}_\varepsilon$.

\origpage{19}
If $\mathfrak{F}_1$, $\mathfrak{F}_2$ denote two two-dimensional manifolds and $\mathfrak{E}_1$ a point set on $\mathfrak{F}_1$, and if furthermore to each point $\mathfrak{p}^{(1)}$ of $\mathfrak{E}_1$ there is assigned a point $\mathfrak{p}^{(2)}$ of $\mathfrak{F}_2$, then an \textbf{application} [mapping] of $\mathfrak{E}_1$ \textbf{into} $\mathfrak{F}_2$ is thereby given. The set of image points on $\mathfrak{F}_2$ is designated as the \textbf{image set} of $\mathfrak{E}_1$. This mapping is called \textbf{continuous} if the following holds: If $\mathfrak{p}_0^{(1)}$ is any point of $\mathfrak{E}_1$, $\mathfrak{p}_0^{(2)}$ its image point on $\mathfrak{F}_2$, $\mathfrak{U}^{(2)}$ any neighborhood of $\mathfrak{p}_0^{(2)}$ on $\mathfrak{F}_2$, then there always exists a neighborhood $\mathfrak{U}^{(1)}$ of $\mathfrak{p}_0^{(1)}$ on $\mathfrak{F}_1$ such that the image point of every point belonging simultaneously to $\mathfrak{E}_1$ and $\mathfrak{U}^{(1)}$ is a point of $\mathfrak{U}^{(2)}$. --- The continuous image of a closed set is again closed. --- If $\mathfrak{p}^{(1)}(\lambda)\,[0 \leqq \lambda \leqq 1]$ is a continuous curve on $\mathfrak{F}_1$ all of whose points belong to the set $\mathfrak{E}_1$ continuously mapped into a second manifold $\mathfrak{F}_2$, then one can define a continuous curve on $\mathfrak{F}_2$ by assigning to each value of the parameter $\lambda$ that point $\mathfrak{p}^{(2)}$ on $\mathfrak{F}_2$ which proceeds from the point $\mathfrak{p}^{(1)}(\lambda)$ under the given mapping of $\mathfrak{E}_1$. This curve is called the \textbf{image} of the given one. --- If a closed set $\mathfrak{E}_1$ (on $\mathfrak{F}_1$) is mapped in a one-to-one and continuous manner onto a subset $\mathfrak{E}_2$ of $\mathfrak{F}_2$, then the \emph{inverse mapping} (which comes about by assigning to each point $\mathfrak{p}^{(2)}$ of $\mathfrak{E}_2$ that single point $\mathfrak{p}^{(1)}$ of $\mathfrak{E}_1$ whose image is $\mathfrak{p}^{(2)}$) is also a continuous mapping.

Alongside the concept of continuous mapping in the ordinary sense, which comes into consideration mainly for closed sets, there appears a concept of continuity tailored in an analogous manner to domains consisting only of interior points, and for which a name has hitherto been lacking; I propose the designation ``domain-continuous'' [gebiets-stetig]. The mapping of a set $\mathfrak{E}_1$ on $\mathfrak{F}_1$ consisting only of interior points is called \textbf{domain-continuous} if the image $\mathfrak{B}_0^{(2)}$ of every neighborhood $\mathfrak{U}_0^{(1)}$ situated entirely in $\mathfrak{E}_1$ of an arbitrary point $\mathfrak{p}_0^{(1)}$ belonging to $\mathfrak{E}_1$ always contains the image point $\mathfrak{p}_0^{(2)}$ of $\mathfrak{p}_0^{(1)}$ in its interior. A one-to-one and invertibly domain-continuous mapping is also continuous in the ordinary sense. A fundamental theorem proved by L.~E.~J.~\emph{Brouwer}${}^{1)}$ states that the converse of this also holds: every one-to-one continuous mapping of a set consisting only of interior points is, together with its inverse, domain-continuous. By this theorem the concept of domain-continuity is rendered superfluous once again. But since in what follows we nowhere have occasion to make use of this \emph{difficult-to-prove theorem}, I should nonetheless like to retain here the word ``domain-continuous''.

\textbf{1)} Mathematische Annalen Vol.~70 (1911), pp.~161---165; Vol.~71 (1912), pp.~305---313; Vol.~72 (1912), pp.~55---56. \emph{Brouwer's} investigations relate to $n$-dimensional domains. The two-dimensional case, which alone comes into consideration for us, was already settled earlier with the aid of the so-called ``Jordan curve theorem'' by \emph{Schoenflies} (Göttinger Nachrichten 1899, pp.~282---290), \emph{Osgood} (ibid.\ 1900, pp.~94---97), and \emph{F.~Bernstein} (ibid.\ 1900, pp.~98---102).

\origpage{20}
--- The requirements 1. and 2. that we imposed on p.~18 in the definition of a two-dimensional manifold on the concept of neighborhood can now be expressed briefly as follows: \emph{Every neighborhood of a point $\mathfrak{p}_0$ on $\mathfrak{F}$ can be mapped in a one-to-one and invertibly domain-continuous manner onto the interior of a circle (where $\mathfrak{p}_0$ transforms into the center of the circle).}

The definition of what is to be regarded as a ``neighborhood'' of a point on a manifold $\mathfrak{F}$ can be modified up to a certain degree without our wishing thereby to regard the manifold $\mathfrak{F}$ as altered. For if we do not touch the definition of those things which are to count as points of $\mathfrak{F}$, but replace the original (``first'') definition of the concept of neighborhood by a second one that likewise satisfies the requirement just newly formulated, and if then the theorem holds: Every point set to be designated as a neighborhood of $\mathfrak{p}_0$ according to the second definition contains a point set that is to be addressed as a neighborhood of $\mathfrak{p}_0$ according to the first definition, and conversely, whatever the point $\mathfrak{p}_0$ on $\mathfrak{F}$ may be---then we wish to agree to regard the manifold determined by the second definition as not distinct from that determined by the first. In this agreement there is present once again a ``definition by abstraction'', the proof of justification of which we can leave to the reader. Moreover, none of the continuity concepts enumerated above is in any way affected by the fact that the first definition of neighborhood is replaced by the second: a point set that is to be designated as closed or as a domain according to the first definition remains so according to the second as well; correspondingly it stands with continuous functions, continuous curves, continuous mappings, etc.

That branch of mathematics which deals with the continuity properties of two- (and multi-) dimensional manifolds is designated as \textbf{analysis situs} or \textbf{topology}. This discipline has played an important role in function theory since \emph{Riemann}, and in the following sections we shall have to occupy ourselves still more thoroughly with the topology of two-dimensional manifolds. Two manifolds must be regarded as \textbf{equivalent in the sense of analysis situs} if they can be mapped onto each other point for point in a one-to-one and invertibly domain-continuous manner. All mutually equivalent two-dimensional manifolds may be conceived as realizations of one and the same ``ideal manifold'', in which the individual features by which these various realizations differ from one another are obliterated (definition by abstraction). The continuity concepts collected above are pure concepts of analysis situs, since they behave invariantly under arbitrary one-to-one and domain-continuous mappings. By

\origpage{21}
regarding the analytic configuration as a two-dimensional manifold, we are led to bring the analysis situs properties of this manifold into the foreground as the most incisive and primitive. It seems to me a requirement of objectivity to establish these properties then also only with the aid of such concepts and methods as belong to analysis situs. This requirement was satisfied by \emph{Riemann} in his construction of function theory.

A sub-domain $\mathfrak{G}$ on a manifold $\mathfrak{F}$ is itself a two-dimensional manifold $\mathfrak{G}$ if as a neighborhood of a point of $\mathfrak{G}$ on $\mathfrak{G}$ we consider every neighborhood of this point on $\mathfrak{F}$ contained entirely in $\mathfrak{G}$. Every curve on $\mathfrak{F}$ belonging entirely to $\mathfrak{G}$ is then also a continuous curve on $\mathfrak{G}$, etc.

In order to be able to state essential theorems on two-dimensional manifolds, it seems to be necessary to subject this concept to a further restriction, whose purpose is to ensure the applicability of the ``\emph{method of exhaustion}'' to the manifolds to be investigated. For the exhaustion of a planar domain one most simply uses \emph{triangles} (\emph{method of triangulation}). We shall accordingly henceforth consider only such manifolds as can be triangulated in an analogous sense to a planar domain, and this restriction will be unexceptionable for us because we shall subsequently be able to provide the proof of the possibility of triangulation for every analytic configuration.

A planar triangle can most easily be described with the aid of the homogeneous barycentric coordinates belonging to its vertices. We thereby arrive at the following definition. If to each triple $(\xi_1, \xi_2, \xi_3) \ne (0, 0, 0)$ of three real, \emph{non-negative} numbers there is assigned a point $\mathfrak{p}$ of the manifold $\mathfrak{F}$ such that 1) the same point $\mathfrak{p}$ is assigned to two such triples if and only if both number triples have the same ratio $(\xi_1 : \xi_2 : \xi_3)$, and 2) the condition of continuity is fulfilled, then we say that a \textbf{triangle} $\Delta$ is defined on $\mathfrak{F}$. $\xi_1 : \xi_2 : \xi_3$ is called the \textbf{coordinate ratio} of $\mathfrak{p}$ in $\Delta$. --- The condition of continuity will state: If $(\xi_1^0, \xi_2^0, \xi_3^0)$ is any one of the triples in consideration, $\mathfrak{p}^0$ the assigned point, $\mathfrak{U}_0$ an arbitrary neighborhood of $\mathfrak{p}^0$ on $\mathfrak{F}$, then there exists a positive number $\varepsilon$ of such a kind that for all triples in consideration satisfying the inequalities
\[
|\xi_1 - \xi_1^0| < \varepsilon, \quad |\xi_2 - \xi_2^0| < \varepsilon, \quad |\xi_3 - \xi_3^0| < \varepsilon
\]
the image point $\mathfrak{p}$ belongs to that neighborhood $\mathfrak{U}_0$.

To the definition of a triangle there belongs accordingly not merely the specification of those points of $\mathfrak{F}$ that are to belong to it, but it counts as completely determined only when to each of its points there is assigned a coordinate ratio of three non-negative numbers. The

\origpage{22}
points $1:0:0$, $0:1:0$, $0:0:1$ will be designated as the three \textbf{vertices} of the triangle; the three \textbf{edges} of the same are furnished by $\xi_1 = 0$, resp.\ $\xi_2 = 0$, resp.\ $\xi_3 = 0$.

The requirement of the possibility of a triangulation we shall now have to formulate [in closest connection with \emph{Brouwer's} fundamental works${}^{1)}$ and in some agreement with the \emph{Dehn-Heegaard} theory developed in the Encyklopädie${}^{2)}$] as follows. Let there be defined on a manifold $\mathfrak{F}$ finitely or infinitely many triangles $\Delta$ (the ``elementary triangles'' of the triangulation), such that every point of the manifold belongs to at least one of the triangles $\Delta$ and moreover the following conditions are fulfilled:

\rmfigure{figures/fig-1.png}{Fig.~1.}{Interior point, edge point, vertex of a triangulation}

1) If $\mathfrak{p}$ is a point belonging to the $\Delta$-triangle $\Delta_0$ not situated on the edges of $\Delta_0$, then neither $\mathfrak{p}$ nor any point of a certain neighborhood of $\mathfrak{p}$ belongs to any of the triangles $\Delta$ distinct from $\Delta_0$.

2) If $\mathfrak{p}$ is situated on an edge $k$ of $\Delta_0$, but is not a vertex of $\Delta_0$, then there exists another $\Delta$-triangle $\Delta_1$ to which $\mathfrak{p}$ additionally belongs; $\Delta_0$, $\Delta_1$ have all points on $k$, but no others, in common; neither $\mathfrak{p}$ nor any point of a certain neighborhood of $\mathfrak{p}$ belongs to any $\Delta$-triangle distinct from $\Delta_0$ and $\Delta_1$.

3) If $\mathfrak{p}$ is a vertex of a triangle $\Delta$, then there exist finitely many $\Delta$-triangles: $\Delta_0, \Delta_1, \dots, \Delta_h$ to which the point $\mathfrak{p}$ belongs; in all these triangles $\mathfrak{p}$ is a vertex, and they hang together in a single cycle in such a way that $\Delta_0$ has with $\Delta_1$ precisely one edge in common, $\Delta_1$ with $\Delta_2$ an edge, $\dots$, finally $\Delta_h$ again with $\Delta_0$ an edge${}^{3)}$; neither $\mathfrak{p}$ nor any point of a certain neighborhood of

\textbf{1)} See especially L.~E.~J.~\emph{Brouwer}, Über Abbildung von Mannigfaltigkeiten, Math.\ Ann.\ Vol.~71 (1912), pp.~97 ff. Further J.~\emph{Hadamard}, ``Sur quelques applications de l'indice de Kronecker'' in the 2nd volume of J.~\emph{Tannery}, Introduction à la théorie des fonctions d'une variable, 2nd ed.\ (Paris 1910). The conception of an arbitrary closed spatial surface as a polyhedron is found clearly developed in \emph{Möbius}, especially in the posthumous papers ``Zur Theorie der Polyeder und der Elementarverwandtschaft'' (1861), Werke Vol.~II, pp.~517 ff.

\textbf{2)} Encyklopädie III A B 3, pp.~153 ff.

\textbf{3)} We designate this configuration as a \textbf{triangle star}.

\origpage{23}
$\mathfrak{p}$ belongs to a triangle $\Delta$ not appearing in the cycle $\Delta_i$ ($i = 0, 1, \dots, h$):

\emph{then the manifold is} \textbf{partitioned into elementary triangles} $\Delta$.

In the future we take into consideration only such two-dimensional manifolds as can be partitioned into elementary triangles; for such manifolds I shall use the shorter designation \textbf{surface}. It will emerge presently what great significance the possibility of triangulation possesses for analysis situs on a surface.

If one has partitioned a surface $\mathfrak{F}$ into elementary triangles $\Delta$, then an ordered sequence of finitely many (mutually distinct) of these triangles $\Delta$:
\[
\Delta_1, \Delta_2, \dots, \Delta_n
\]
will be designated as a (\textbf{simple}) \textbf{chain} of triangles if $\Delta_j$ always has an edge in common with $\Delta_{j+1}$; the chain \emph{``connects''} $\Delta_1$ with $\Delta_n$. (If $\Delta_n$ again also has an edge $k$ in common with $\Delta_1$, then the chain is \emph{closed} across the edge $k$.) The circumstance that any two points of $\mathfrak{F}$ can be connected by a continuous curve on $\mathfrak{F}$ has the consequence that any two elementary triangles $\Delta$ can be connected with each other by a simple chain of triangles $\Delta$. First of all, namely, there can be only finitely many triangles $\Delta$ having points in common with a continuous curve $\mathfrak{p} = \mathfrak{p}(\lambda)\,[0 \leqq \lambda \leqq 1]$. If there were infinitely many, one would obtain an infinite sequence of curve points $\mathfrak{p}(\lambda_j)\,[j = 1, 2, \dots]$, no two of which belong to the same triangle; if $\lambda_\infty$ is then an accumulation point of the $\lambda_j$, points of infinitely many triangles would have to be encountered in every neighborhood of $\mathfrak{p}(\lambda_\infty)$: according to the requirements 1)---3) to be imposed on every triangulation, this is impossible. In order now to connect two given triangles $\Delta$, $\Delta_0$ and $\Delta^0$, by a chain, I proceed as follows: I choose any point in the interior of $\Delta_0$ and a point in the interior of $\Delta^0$ and connect both on $\mathfrak{F}$ by a continuous curve $\mathfrak{p} = \mathfrak{p}(\lambda)\,[0 \leqq \lambda \leqq 1]$. I observe when this curve \emph{for the last time} leaves the triangle $\Delta_0$---I express myself as though $\lambda$ signifies time---, that is, I seek the largest value $\lambda = \lambda_0$ for which $\mathfrak{p}(\lambda)$ belongs to $\Delta_0$. The point $\mathfrak{p}(\lambda_0)$ will lie on an edge $k$ of $\Delta_0$. If it is not a vertex (\emph{first case}), then the curve now crosses over into that triangle which borders on $\Delta_0$ along $k$: this $\Delta_1$ will be the next link in the chain to be constructed. If, on the other hand, $\mathfrak{p}(\lambda_0)$ is a vertex, I form the entire triangle star about the vertex $\mathfrak{p}(\lambda_0)$, which I can regard as a single closed chain $\Delta_0\,\Delta_1 \cdots \Delta_h$. If $\Delta^0$ occurs among its triangles, I am already finished; otherwise I observe when the curve leaves the star for the last time. The point $\bar{\mathfrak{p}}_0$ through which this occurs is certainly not a point

\origpage{24}
of $\Delta_0$ (for $\Delta_0$ was already left for the last time earlier); rather, in the sequence $\Delta_1, \dots, \Delta_h$ let the triangle $\Delta_g$ be the first (and, if $\bar{\mathfrak{p}}_0$ is not a vertex, also the only one) to which the point $\bar{\mathfrak{p}}_0$ belongs. Then I take $\Delta_1, \dots, \Delta_g$ as the next links of the chain to be constructed. (In the first case I write $g = 1$.) Now I proceed with $\Delta_g$ just as initially with $\Delta_0$ and obtain the next link or links of the sought-for chain; to the last of the links thus newly obtained I apply the same continuation procedure once more, and so on. All the links I obtain are triangles having points in common with the curve. But since the decisive moment for the continuation is always that when the curve leaves the triangle or star \emph{for the last time}, I also never return to an already obtained triangle. Consequently the procedure, since only finitely many triangles are available altogether, must terminate at some point; but that will occur only when I have penetrated to the endpoint of the curve, that is, to the triangle $\Delta^0$. By the described construction, $\Delta_0$ is accordingly indeed connected with $\Delta^0$ by a simple chain.

From this we conclude further that the elementary triangles $\Delta$ can form only a finite or countably infinite set. We proceed from a triangle $\Delta_0$, then form those finitely many triangles $\Delta$ that abut the edges and vertices of $\Delta_0$, then those finitely many that abut the free edges and vertices of the triangles admitted up to now, and so forth. We thereby arrange the $\Delta$ into a countable sequence; that every triangle $\Delta$ is ultimately included follows precisely from the theorem on the possibility of chain connection.

A point set on the surface $\mathfrak{F}$ possessing no accumulation point is finite or countable. For each triangle $\Delta$ can evidently receive only finitely many points of such a set. On manifolds for which we do not presuppose partitionability into elementary triangles, the last theorem is not in general correct. Were the set of irregular elements in an analytic configuration not countable, the analytic configuration could certainly not be triangulated. Conversely, this fact that an analytic configuration in truth contains only countably many irregular elements will become for us in §~6 the essential starting point for the actual triangulation of an arbitrary analytic configuration.

If on $\mathfrak{F}$ there exists no infinite point set without accumulation point at all, then $\mathfrak{F}$ is called \textbf{closed}. A set without accumulation point is obtained by choosing arbitrarily a point in the interior of each elementary triangle $\Delta$ of a definite partition of $\mathfrak{F}$. Hence: \emph{A closed surface can be partitioned into finitely many, but never into infinitely many elementary triangles. But also conversely: An} \textbf{open} (\emph{= un-}

\origpage{25}
\emph{closed) surface can be partitioned into infinitely many, but never into finitely many elementary triangles.} For if the number of elementary triangles $\Delta$ is finite, every infinite point set on $\mathfrak{F}$ has an accumulation point, since infinitely many points of such a set must fall upon at least one of the triangles $\Delta$.

\subsection*{§ 5. Examples of Surfaces.}

1.\ The \emph{Euclidean plane} (open).

2.\ The \emph{interior of a square} (open).

\rmfigure{figures/fig-2.png}{Fig.~2.}{Triangulation of the plane and of the surface T (p.~28)}
\rmfigure{figures/fig-3.png}{Fig.~3.}{Triangulation of the interior of a square}

By the accompanying figures a definite triangulation of these surfaces is to be indicated. Here each triangle is to be made into a ``triangle on $\mathfrak{F}$'' in the sense of §~4 by representing its points through the homogeneous barycentric coordinates belonging to the vertices of the triangle.

3.\ The \emph{projective plane}. It can be defined arithmetically as the totality of ratios $\xi_1 : \xi_2 : \xi_3$ formed from any three real numbers $(\xi_1, \xi_2, \xi_3) \ne (0, 0, 0)$. The concept of neighborhood is to be understood such that a variable ``point'' $\xi_1 : \xi_2 : \xi_3$ converges to the fixed point $\xi_1^0 : \xi_2^0 : \xi_3^0$ if and only if
\[
\frac{(\xi_2\,\xi_3^0 - \xi_3\,\xi_2^0)^2 + (\xi_3\,\xi_1^0 - \xi_1\,\xi_3^0)^2 + (\xi_1\,\xi_2^0 - \xi_2\,\xi_1^0)^2}{[(\xi_1^0)^2 + (\xi_2^0)^2 + (\xi_3^0)^2] [\xi_1^2 + \xi_2^2 + \xi_3^2]}
\]
tends to 0. A triangulation (into ten elementary triangles) is indicated by Figure~4; it is immediately evident how the individual triangles can be conceived as ``triangles'' in our general sense with the aid of the projective triangular coordinates belonging to each of them. The projective plane is \emph{closed}.

4.\ The \emph{surface of the regular octahedron} decomposes in a natural way into eight triangles which, referred to their barycentric coordinates,

\origpage{26}
can be conceived as elementary triangles of a triangulation. Closed.

5.\ By projection of the regular octahedron from its center onto the surface of the circumscribed sphere, one obtains the sphere as a manifold partitioned into eight elementary triangles (the octants). Sphere and octahedron are equivalent in the sense of analysis situs.

\rmfigure{figures/fig-4.png}{Fig.~4.}{Triangulation of the projective plane}
\rmfigure{figures/fig-5.png}{Fig.~5.}{Octahedron}

6.\ The \emph{Möbius band}.${}^{1)}$ It arises from a long, narrow rectangle when one twists it by $180^\circ$ in three-dimensional space and bends it together so that the narrow ends of the rectangle come into coincidence in reversed orientation. This spatial surface can be represented analytically, referred to rectangular coordinates $xyz$, most simply with the aid of two real parameters $\varrho$, $\varphi$ as follows${}^{2)}$:

\rmfigure{figures/fig-6.png}{Fig.~6.}{Möbius band}
\[
\left\{
\begin{aligned}
x &= (a - \varrho \sin \tfrac{1}{2}\varphi) \cos \varphi \\
y &= (a - \varrho \sin \tfrac{1}{2}\varphi) \sin \varphi \\
z &= \varrho \cos \tfrac{1}{2}\varphi.
\end{aligned}
\right.
\]
$\varphi$ is unrestricted in variation, $\varrho$ bound by the condition $|\varrho| < h$; $a$ and $h$ denote positive constants, of which $h$ must be the smaller. In the sense of analysis situs the Möbius band is evidently identical with the following manifold $\mathfrak{B}$:

In a Euclidean plane with rectangular coordinates $\varrho$, $\varphi$ we consider the strip $|\varrho| < 1$ and the discrete group $\varGamma$ of paddle motions${}^{3)}$ $(\varrho, \varphi) \to (\varrho'\,\varphi')$ defined by
\[
\varrho' = (-1)^n \varrho, \quad \varphi' = \varphi + 2n\pi \quad [n = 0, \pm 1, \pm 2, \cdots]
\]
A

\textbf{1)} Möbius, Werke, Vol.~2, pp.~484--485 and pp.~519--521.

\textbf{2)} The surface represented by these equations is admittedly not a developable.

\textbf{3)} The word ``paddle motion'' [Paddelbewegung] (borrowed from Low German usage) is intended to indicate that the group is obtained by repeatedly flipping the plane about an axis, alternately to the left and to the right, but at the same time shifting it forward each time (by the distance $2\pi$) in the direction of that axis.

\origpage{27}
point system $S$ in this plane is called a \textbf{system of points equivalent with respect to $\varGamma$} (or more briefly: a \textbf{point system belonging to $\varGamma$}) if by every motion of the group $\varGamma$ the points of $S$ transform into each other, but also every point of $S$ can be carried into every other by a motion of this

\rmfigure{figures/fig-7a.png}{Fig.~7\,a.}{Point on B with neighborhood and angle}

group. A point system belonging to $\varGamma$ that belongs to the strip $|\varrho| < 1$ we call a ``point'' of the manifold $\mathfrak{B}$ to be defined. If $S_0$ is such a ``point'' of $\mathfrak{B}$ and if about the individual (Euclidean) points of the point system $S_0$ one describes purely congruent circles (Fig.~7\,a) lying entirely within the strip, then those point systems $S$ belonging to $\varGamma$ which are represented by a point in each of these circles form a ``neighborhood'' of $S_0$ on $\mathfrak{B}$. If of an ordinary point belonging to the point system $S$ belonging to $\varGamma$ one says it ``\emph{lies over}'' the ``point'' $S$ of $\mathfrak{B}$, then the parallel strip appears as a ``\emph{covering surface}'' over $\mathfrak{B}$, which covers $\mathfrak{B}$ infinitely-many-sheeted without, however, being \emph{ramified} at any place relative to $\mathfrak{B}$. In the parallel strip we apply the Euclidean measurement of angles; by conceiving it as a covering surface of $\mathfrak{B}$, this measurement of angles transfers directly to $\mathfrak{B}$. In this, however, the measure of angle becomes indeterminate not merely up to integral multiples of $2\pi$, but also with respect to its sign (Fig.~7\,a). This circumstance is connected with a property of the Möbius band designated as its ``\emph{one-sidedness}'' (without leaving the surface and without climbing over its boundary, one can pass from one side to the other), and which we shall examine more closely later on account of its function-theoretic significance. --- A triangulation of $\mathfrak{B}$ and therewith of the Möbius band is indicated by Fig.~7\,b.

\rmfigure{figures/fig-7b.png}{Fig.~7\,b.}{Triangulation of B}

\rmfigure{figures/fig-8.png}{Fig.~8.}{Construction of the torus}

7.\ The \emph{torus} arises (Fig.~8) when a circle of radius $r$ is rotated about an axis situated in its plane that does not meet the circle; let $R (> r)$ be the perpendicular distance of the center of the circle from the axis. The rectangular coordinates $xyz$ of the points on the surface can be represented with the aid of two real parameters $\sigma$, $\varphi$, whose geometric significance is evident, in the form

\origpage{28}
\[
\left\{
\begin{aligned}
x &= (R + r \cos \sigma) \cos \varphi, \\
y &= (R + r \cos \sigma) \sin \varphi, \\
z &= r \sin \sigma
\end{aligned}
\right.
\]
To the torus as a surface in three-dimensional Euclidean space there belongs a natural (Euclidean) measurement of angles. If a curve on the torus
\[
\varphi = \varphi(\lambda), \quad \sigma = \sigma(\lambda)
\]
is given, where $\frac{d\varphi}{d\lambda}, \frac{d\sigma}{d\lambda}$ are continuous and throughout $\left(\frac{d\varphi}{d\lambda}\right)^2 + \left(\frac{d\sigma}{d\lambda}\right)^2 \ne 0$, then the differential quotient of the arc length $s = s(\lambda)$ of this curve is calculated from:
\[
\left(\frac{ds}{d\lambda}\right)^2 = (R + r \cos \sigma)^2 \left(\frac{d\varphi}{d\lambda}\right)^2 + r^2 \left(\frac{d\sigma}{d\lambda}\right)^2.
\]
If one introduces a new parameter $\psi$ instead of $\sigma$:
\[
\psi = \int_0^\sigma \frac{d\sigma}{\frac{R}{r} + \cos \sigma}, \tag{8}
\]
one can write instead
\[
ds^2 = e(d\varphi^2 + d\psi^2), \text{ where } e = (R + r \cos \sigma)^2. \tag{9}
\]
By (8), $\psi$ is defined as a single-valued, continuously differentiable function of $\sigma$, and conversely $\sigma$ as an identical function of $\psi$. The values $(\varphi, \psi)$ corresponding to a definite point on the torus are determined only up to integral multiples of
\[
2\pi, \text{ resp. } \int_0^{2\pi} \frac{d\sigma}{\frac{R}{r} + \cos \sigma} = \frac{2\pi r}{\sqrt{R^2 - r^2}}
\]
We are thus led to consider in a Euclidean plane with rectangular coordinates $\varphi, \psi$ the translation group
\[
\left.
\begin{gathered}
\varphi' = \varphi + 2m\pi, \quad \psi' = \psi + \frac{2\pi r}{\sqrt{R^2 - r^2}}\,n \\
\begin{pmatrix}
m = 0, \pm 1, \pm 2, \cdots \\
n = 0, \pm 1, \pm 2, \cdots
\end{pmatrix}
\end{gathered}
\right\} \varGamma
\]
and to address a system of points equivalent with respect to $\varGamma$, a so-called \textbf{``point lattice''}, as a ``point'' of a new manifold $\mathfrak{T}$, over which the plane then spreads out as an infinitely-many-sheeted, but nowhere ramified covering surface. The Euclidean measurement of angles of the plane transfers directly to the manifold $\mathfrak{T}$, and there does not result from it (as in the case of the Möbius

\origpage{29}
band) an uncertainty regarding the sign of the angle. Onto this manifold $\mathfrak{T}$ the torus is mapped in a one-to-one and invertibly domain-continuous manner by means of its representation through the parameters $\varphi$, $\psi$; but not only that: the mapping is also \emph{conformal}. Equation (9) states, namely, that the ratio of an element of arc $ds$ on the torus to the image element $(\sqrt{d\varphi^2 + d\psi^2})$ in the $\varphi\psi$-plane is $= \sqrt{e} = R + r \cos \sigma$, thus depending only on the place where the element of arc is located, but not on its direction. This conformal mapping provides the foundation for the theory of analytic functions on the torus. --- The torus is \emph{closed}; a triangulation is given by Fig.~2, when the rectangle drawn in heavy lines there, whose sides are parallel to the coordinate axes, possesses side lengths $2\pi$ and $\frac{2\pi r}{\sqrt{R^2 - r^2}}$. A ``point'' of $\mathfrak{T}$ is also marked in this figure by zero circles.

8.\ The possible \emph{simultaneous positions of two hands} moving on a dial are the points of a manifold on which one describes a continuous curve when one moves the two hands simultaneously in any continuous manner. This manifold is evidently closed and equivalent to the torus.

9.\ In order for a surface in the sense of analysis situs to be given, it suffices to characterize the triangles of which it consists (after it has been triangulated in a definite way) by any symbols (e.g., numbers) and to specify the incidence relations of these triangles. Most simply this is done by numbering the vertices and then recording each triangle in the form $(efg)$ [or $(feg)$ or $\cdots$], where $e, f, g$ are the numbers of its vertices.${}^{1)}$ The schema thus arising will satisfy the condition that each vertex $e_0$ occurs in only a single finite cycle of triangle symbols:
\[
(e_0 e_1 e_2), \quad (e_0 e_2 e_3), \quad \cdots, \quad (e_0 e_{h-1} e_h), \quad (e_0 e_h e_1)
\]
Furthermore, one will in no way be able to divide the vertex numbers into two classes such that each triangle symbol exhibits either three vertices from the one or three vertices from the other class. Every schema of symbols $(efg)$ satisfying these conditions also defines a definite surface in the sense of analysis situs (together with a definite triangulation of the same). This abstract mode of description is especially suited for \emph{closed} surfaces, since for these the schema consisting of only finitely many symbols can be completely written down. The construction of all admissible finite schemata is a purely combinatorial task. The question naturally

\textbf{1)} Möbius, Werke Vol.~II, pp.~478\,f.

\origpage{30}
arises immediately as to when two such finite schemata represent one and the same surface (in different triangulations); the further developments on analysis situs contained in this chapter will furnish important contributions toward answering this question.

The schema of the \emph{tetrahedron} and therewith of the \emph{sphere} reads, for example:
\[
(123) \quad (234) \quad (341) \quad (412).
\]
In a different, octahedral triangulation [cf.\ Fig.~5] the sphere can be described by the table
\[
\begin{gathered}
( \text{I}\,1\,2 ) \quad ( \text{I}\,2\,3 ) \quad ( \text{I}\,3\,4 ) \quad ( \text{I}\,4\,1 ) \\
( \text{II}\,1\,2 ) \quad ( \text{II}\,2\,3 ) \quad ( \text{II}\,3\,4 ) \quad ( \text{II}\,4\,1 )
\end{gathered}
\]
The \emph{projective plane} in the triangulation drawn above [Fig.~4] possesses the following schema${}^{1)}$:
\[
\begin{gathered}
( 1\ 2\ \text{I} ) \quad ( 2\ 3\ \text{II} ) \quad ( 3\ 1\ \text{III} ) \\
( 1\ 2\ \text{II} ) \quad ( 2\ 3\ \text{III} ) \quad ( 3\ 1\ \text{I} ) \\
( 1\ \text{II}\ \text{III} ) \quad ( 2\ \text{III}\ \text{I} ) \quad ( 3\ \text{I}\ \text{II} ) \\
( \text{I}\ \text{II}\ \text{III} ).
\end{gathered}
\]

10.\ Every \emph{analytic configuration} provides an example of a ``surface'' in our sense, if we regard the function elements belonging to the analytic configuration as ``points'' and let the concept of ``neighborhood'' required for the definition of a surface coincide with the concept of ``analytic neighborhood'' treated in §~2. It is clear first of all that under this interpretation every analytic configuration becomes a two-dimensional manifold; only it remains to be proved that this manifold is always also capable of a definite triangulation.

\subsection*{§ 6. Analytic Configurations, Considered as Surfaces.}

A triangle $\Delta$ on an arbitrary two-dimensional manifold $\mathfrak{F}$ can be represented point by point by a Euclidean triangle $D$ with vertices 1, 2, 3 in such a way that to each point of $\Delta$, which according to p.\ 21 possesses a definite coordinate ratio $\xi_1 : \xi_2 : \xi_3$ (in $\Delta$), one associates that point of $D$ whose homogeneous barycentric coordinates with respect to the vertices 1, 2, 3 possess the same value $\xi_1 : \xi_2 : \xi_3$. And conversely, a point set $\Delta$ on $\mathfrak{F}$ that is mapped in a one-to-one and continuous manner onto a Euclidean triangle $D$ may, by virtue of this mapping, be addressed as a triangle on $\mathfrak{F}$. Every straight line segment located in $D$ is then the image of a curve traversing entirely in $\Delta$

\textbf{1)} The triangles standing in the first two rows extend through infinity.

\origpage{31}
on $\mathfrak{F}$, which we designate as an \textbf{``elementary segment''} in $\Delta$, and every Euclidean triangle $D^*$ lying entirely in $D$ is the image of a triangle $\Delta^*$ on the manifold; we call $\Delta^*$ a \textbf{``sub-triangle''} of $\Delta$.

Let $\mathfrak{F}$ be partitioned in a definite way into elementary triangles. Another partition $\zeta'$ of the same manifold into elementary triangles is called a \textbf{subdivision} of the first, $\zeta$, if every elementary triangle of the partition $\zeta'$ is a sub-triangle of an elementary triangle of the partition $\zeta$. If $\zeta$ is a given partition into elementary triangles, one can very easily obtain such a sequence of triangulations $\zeta_1, \zeta_2, \zeta_3, \cdots$ that 1) $\zeta_1 = \zeta$, 2) each $\zeta_{n+1}$ is a subdivision of the preceding $\zeta_n$, and 3) the partition $\zeta_n$ becomes \emph{arbitrarily fine} as $n$ increases indefinitely. The latter shall mean: if $\mathfrak{p}$ is an arbitrary point of $\mathfrak{F}$, and $\mathfrak{U}$ is an arbitrary neighborhood of $\mathfrak{p}$ on $\mathfrak{F}$, then there is always an index $n$ such that the elementary triangle (or exceptionally: the elementary triangles) of $\zeta_n$ in which $\mathfrak{p}$ lies is (are) itself (themselves) contained entirely in $\mathfrak{U}$. If countably infinitely many points $\mathfrak{p}_2, \mathfrak{p}_3, \cdots$ were given on $\mathfrak{F}$, we can furthermore arrange this sequence of increasingly fine subdivisions such that $\mathfrak{p}_2$ is contained among the vertices occurring in the partition $\zeta_2$, $\mathfrak{p}_3$ among the vertices of the partition $\zeta_3$, and so on. We shall make use of this remark immediately.

First it shall be shown that \emph{every domain $\mathfrak{G}$ on a surface $\mathfrak{F}$ (which according to p.\ 21 can itself be conceived as a two-dimensional manifold) is capable of a triangulation if this holds for $\mathfrak{F}$.} One starts from a triangulation $\zeta = \zeta_1$ of $\mathfrak{F}$ and constructs for it a sequence of subdivisions $\zeta_2, \zeta_3, \dots$ whose mesh size falls below any bound with increasing index. Let $\varGamma_1$ be those triangles of $\zeta_1$ that lie entirely within $\mathfrak{G}$; $\varGamma_2$ those triangles of $\zeta_2$ that lie entirely in $\mathfrak{G}$ but are not parts of $\varGamma_1$; $\varGamma_3$ those triangles of $\zeta_3$ that lie entirely in $\mathfrak{G}$ but are neither parts of triangles $\varGamma_1$ nor of $\varGamma_2$; and so on. All the triangles $\varGamma_1, \varGamma_2, \varGamma_3, \dots$, which are everywhere distinct in their interior points, but together fill out the entire domain $\mathfrak{G}$, may be called by a common name ``the triangles $\varGamma$''. \emph{On the edges of a triangle $\varGamma$ there lie only finitely many vertices of triangles $\varGamma$.} For if there lay infinitely many, they would have a condensation point $\mathfrak{p}$. To $\mathfrak{p}$ belongs a neighborhood $\mathfrak{U}$ lying entirely in $\mathfrak{G}$, and there is found an index $n$ such that all triangles of the partition $\zeta_n$ that contain $\mathfrak{p}$ lie entirely in $\mathfrak{U}$ and thus entirely in $\mathfrak{G}$. But in the interior of the pair of triangles (or star of triangles) of $\zeta_n$ containing $\mathfrak{p}$, there can lie no vertex of triangles $\varGamma$, apart perhaps from $\mathfrak{p}$ itself. This is a contradiction. --- Our construction is now continued as follows: an individual triangle $\varGamma$, let it be called $\varGamma^0$, can be partitioned into

\origpage{32}
finitely many sub-triangles in such a way that the vertices of these sub-triangles are precisely all the vertices of the triangles $\varGamma$ lying on the edges of $\varGamma^0$ (Fig.~9). If we carry out this partitioning process not only with $\varGamma^0$, but with every triangle $\varGamma$, we thereby obtain the desired partition of $\mathfrak{G}$ into elementary triangles.

\rmfigure{figures/fig-9.png}{Fig.~9.}{Partition into sub-triangles}

By considerations of an entirely analogous kind, the proof announced at the conclusion of the preceding section also succeeds, namely that every analytic configuration $\mathbf{G}$, conceived as a two-dimensional manifold, can be partitioned into elementary triangles. We operate with a $z$-sphere, to whose points we imagine assigned in the known manner (by stereographic projection) the values of the complex variable $z$ including $\infty$. A function element belonging to $\mathbf{G}$:
\[
u = \text{power series in } \sqrt[\mu]{z - a}
\]
we call a point of $\mathbf{G}$ ``lying over the point $z = a$ of the $z$-sphere''; if $u$ proceeds in powers of $\frac{1}{\sqrt[\nu]{z}}$, the function element in question lies over the point $\infty$ of the $z$-sphere. In accordance with this mode of expression, $\mathbf{G}$ appears as a certain covering surface over the $z$-sphere, as a configuration of the kind that Riemann introduced in his famous dissertation${}^{1)}$ for the investigation of multiple-valued analytic functions, and which since then, under the name of \emph{``Riemann surface''}, has dominated the theory of these functions.

We mark on the sphere the countably many points $a_2, a_3, \dots$ over which lie ramified function elements (\textbf{``branch points''}) of $\mathbf{G}$. Let $\zeta_1$ denote the partition of the sphere into its eight octants, which according to § 5, Example 5) shall count as triangles on the sphere; let $\zeta_1, \zeta_2, \zeta_3, \dots$ be a sequence of triangulations of the sphere beginning with $\zeta_1$ of such a nature that $\zeta_{n+1}$ is always a subdivision of $\zeta_n$, $a_n$ appears as a vertex of the partition $\zeta_n$, and $\zeta_n$ becomes arbitrarily fine as $n$ increases indefinitely. We consider a triangle $D$ on the sphere that occurs in any of these partitions $\zeta_n$. A point set $\Delta$ on $\mathbf{G}$ will be designated as a triangular piece of $\mathbf{G}$ lying over $D$ if the relation
\[
\text{''}\textit{to }\Delta\textit{ belonging point }\mathfrak{p} \to \textit{spherical point over which }\mathfrak{p}\textit{ lies}\text{''}
\]
is a one-to-one continuous mapping between the points of $\Delta$ and $D$. $\Delta$ then appears, in consequence of this mapping, as a triangle defined on $\mathbf{G}$.

\textbf{1)} \emph{Werke}, 2nd ed., pp.\ 7--9.

\origpage{33}
We now first seek out those triangles $\Delta_1^*$ on $\mathbf{G}$ that lie over a spherical triangle of the partition $\zeta_1$ and contain no branch point; of such triangles $\Delta_1^*$ there will be finitely or infinitely many, but perhaps none at all exists. Next we seek out those triangles $\Delta_2^*$ on $\mathbf{G}$ that lie over a triangle of the partition $\zeta_2$, that possess \emph{at most one} branch point and indeed one lying over $a_2$, and that are not parts of triangles $\Delta_1^*$; next those triangles $\Delta_3^*$ on $\mathbf{G}$ that possess the following properties: each $\Delta_3^*$ must lie over a spherical triangle of the partition $\zeta_3$, may contain at most a single branch point and this must then lie over $a_2$ or over $a_3$, and $\Delta_3^*$ must not be part of a triangle $\Delta_1^*$ or $\Delta_2^*$; and so on. The triangles $\Delta^*\,(\Delta_1^*, \Delta_2^*, \Delta_3^*, \dots)$, everywhere distinct in their interior points, together fill out the entire surface $\mathbf{G}$. In fact: if $\mathfrak{e}$ is any element of $\mathbf{G}$ lying over the point $a$ of the $z$-sphere, and if $k$ is a spherical cap around $a$ in which the normal representation of $\mathfrak{e}$ is valid, but $n$ is an index such that all triangles $D_n$ of the partition $\zeta_n$ that contain $a$ lie entirely in $k$, and $a$, in case it is a branch point, coincides with one of the points $a_2, a_3, \dots, a_n$: then over each of these triangles $D_n$ there lies at least one triangle on $\mathbf{G}$ that contains $\mathfrak{e}$. These triangles on $\mathbf{G}$ are triangles $\Delta_n^*$ if they were not already contained as parts in triangles $\Delta_1^*$ or $\Delta_2^*$ $\dots$ or $\Delta_{n-1}^*$.

Without the introduction of new vertices, one can now carry out with each triangle $\Delta^*$ such a division into finitely many triangles that the resulting triangles $\Delta$ also acquire this property: every point that is a vertex for \emph{one} $\Delta$ is likewise a vertex in all (finitely many) triangles $\Delta$ to which it belongs. To see that we have thereby obtained a partition of $\mathbf{G}$ into elementary triangles $\Delta$, one notes the following easily verifiable facts:

Two triangles $\Delta$ lying over one and the same triangle of the $z$-sphere have either no point in common or a single vertex, which is then a branch point.

Two triangles $\Delta$ lying over two different triangles of the $z$-sphere with a common edge have either no point in common, or a vertex that is then a branch point, or an entire edge lying over the common edge of the two spherical triangles.

Two triangles $\Delta$ lying over two spherical triangles without a common point or with only one common vertex have either no point or a vertex in common.

If $\mathfrak{e}$ is an unramified vertex of a triangle $\Delta$, then the triangles $\Delta$ containing $\mathfrak{e}$ form a single cycle that lies, triangle for triangle, over a star of the $z$-sphere. If $\mathfrak{e}$ is a branch point of $(\mu - 1)^{\text{th}}$ order $(\mu \geqq 2)$, then it is a vertex for all triangles $\Delta$ containing $\mathfrak{e}$; these group themselves around $\mathfrak{e}$ in a single cycle, and

\origpage{34}
indeed, every $\mu$ triangles of the cycle lie over one and the same triangle of the $z$-sphere.

\subsection*{§ 7. Concept of the Riemann Surface.}

If $\mathfrak{e}$ is an element belonging to $\mathbf{G}$ and
\[ z = P(t), \quad u = Q(t) \]
a representation of it valid for $|t| < r\;[r > 0]$, then to each $t,\,|t| < r$, there belongs an element $\mathfrak{e}_t$ that arises by rearranging the expansion of $\mathfrak{e}$:
\[ z = P(t + t'), \quad u = Q(t + t') \]
in powers of $t'$. This transition $t \to \mathfrak{e}_t$ is a one-to-one domain-continuous mapping of the disc $|t| < r$ onto a certain neighborhood of the point $\mathfrak{e}$ on $\mathbf{G}$: the parameter $t$ thus appears as a continuous function on $\mathbf{G}$ defined in this neighborhood; we designate it as a \textbf{``local uniformizing variable''} belonging to the point $\mathfrak{e}$. Every other local uniformizing variable $\tau$ belonging to $\mathfrak{e}$ is representable for sufficiently small $t$ in the form
\[ \tau = \gamma_1 t + \gamma_2 t^2 + \cdots \qquad (\gamma_1 \neq 0) \]
A complex-valued function $f$ defined in a domain of $\mathbf{G}$ is said to be \textbf{regular-analytic} at a point $\mathfrak{e}$ of this domain if in a certain neighborhood of $\mathfrak{e}$ it can be represented as a regular power series in the local uniformizing variable $t$ belonging to $\mathfrak{e}$:
\[ f = a_0 + a_1 t + a_2 t^2 + \cdots. \]
Which local uniformizing variable $t$ denotes here is immaterial; if such a representation by \emph{one} local uniformizing variable is possible, then $f$ is also representable in the same manner by every other local uniformizing variable (of course always only in a certain neighborhood of $\mathfrak{e}$ that still depends on the choice of the local uniformizing variable).

If the expansion of $z\;[z = P(t)]$ in the representation of the function element $\mathfrak{e}$ contains no negative powers of $t$, then the constant initial term $z_0$ of this expansion depends solely on $\mathfrak{e}$, but not on the representation of the element $\mathfrak{e}$; $z_0$ will be designated as the \textbf{value} of the complex variable $z$ at the ``point'' $\mathfrak{e}$. If that expansion begins with negative powers of $t$, then the same holds for every representation of the function element $\mathfrak{e}$, and the \textbf{value} of $z$ at the point $\mathfrak{e}$ is $=\infty$. \emph{In this way $z$ appears as a single-valued regular-analytic function on the ``surface'' $\mathbf{G}$ defined everywhere except for isolated points of infinity.} Quite analogously, we can conceive $u$ as a regular-analytic function uniquely defined on the whole of $\mathbf{G}$ except for poles. As the independent argument in both functions there figures not a complex variable in the ordinary sense (i.\,e.\ not a point that \emph{varies in a planar domain}), but a variable point on the ``Riemann surface'' $\mathbf{G}$.

\origpage{35}
For the assertion that $z$ and $u$ are analytic functions on the surface $\mathbf{G}$, it is essential that $\mathbf{G}$ \emph{is not merely given as a surface in the sense of analysis situs.} For on a surface of which only analysis situs properties are taken into consideration, one can indeed speak of continuous functions, but not of ``continuously differentiable'', ``analytic'' (or even ``entire rational'') functions or the like. In order to be able to pursue analytic function theory on a surface $\mathfrak{F}$ in a manner analogous to that in the complex plane, a definition must rather be provided (in addition to the definition of the surface) by which the meaning of the expression ``analytic function on the surface'' is established in such a way that all theorems on analytic functions in the plane that are valid ``in the small'' carry over to this more general concept. Theorems valid ``in the small'' are understood here as those whose correctness is always asserted only for a certain neighborhood of a point, about whose size the theorem itself gives no information. By such an explanation of the expression ``analytic function on $\mathfrak{F}$'', the surface $\mathfrak{F}$ becomes a \textbf{Riemann surface}. This conception of the concept of the Riemann surface, first developed in geometric form by \emph{F.~Klein} in his treatise ``On Riemann's Theory of Algebraic Functions and Their Integrals'',${}^{1)}$ is more general than that which \emph{Riemann} himself employed in his foundational works on the theory of analytic functions. But there can be no doubt that only with this generalized formulation do Riemann's ideas emerge in their full simplicity and power. The foundation for this, moreover, was laid by \emph{Riemann} himself through the conceptual formulations concerning $n$-dimensional manifolds developed in his Habilitation lecture,${}^{2)}$ and it may well be taken as certain that for Riemann the thoughts spun out in that lecture stood in close connection with his function-theoretic investigations, although these connections were not explicitly emphasized by him.

\textbf{1)} Leipzig 1882. See further \emph{Klein}, Neue Beiträge zur Riemannschen Funktionentheorie, Math.\ Ann.\ Vol.\ 21 (1883), §§ 1---3 [pp.\ 146---151]. Surfaces closed by boundary identification as carriers of analytic functions occur already earlier (see \emph{Riemann}, Art.\ 12 of the ``Theorie der Abelschen Funktionen'', Werke p.\ 121; \emph{H.~A.~Schwarz} in his fundamental paper on the integration of the partial differential equation $\frac{\partial^2 u}{\partial x^2} + \frac{\partial^2 u}{\partial y^2} = 0$ from the year 1870, Gesammelte mathematische Abhandlungen Vol.\ II, p.\ 161; \emph{Dedekind}, Crelle's Journal Vol.\ 83, 1877, pp.\ 274\,ff. Freely situated surfaces in space were first brought in, admittedly only for analysis situs investigations, by \emph{Tonelli} (1875; Atti dei Lincei, ser.\ II, t.\ 2) and \emph{Clifford} (1876; Mathematical Papers, pp.\ 249\,ff.). \emph{Klein} himself received, as he communicates in the preface to his treatise ``Über Riemanns Theorie'' (p.\ IV), the first impulse toward his conception through an occasional oral remark of \emph{Prym} (1874). In Klein's work it is always only a matter of \emph{closed} configurations. The most general concept is found explicitly probably first in \emph{Koebe}, cf.\ e.\,g.\ Göttinger Nachrichten 1908, pp.\ 338---339 (footnote).

\textbf{2)} ``Über die Hypothesen, welche der Geometrie zu Grunde liegen'', Werke, 2nd ed., pp.\ 272---287.

\origpage{36}
\textbf{General Definition of the Concept of the Riemann Surface.}

\emph{If a surface $\mathfrak{F}$ is given, and if in addition it is defined for each point $\mathfrak{p}_0$ of $\mathfrak{F}$ and each function $f(\mathfrak{p})$ on $\mathfrak{F}$ present in any neighborhood of $\mathfrak{p}_0$ when $f(\mathfrak{p})$ shall be called regular-analytic at the point $\mathfrak{p}_0$, then there is thereby given a \textbf{Riemann surface} ${}^{\mathfrak{R}}\mathfrak{F}$, whose points are regarded as the points of $\mathfrak{F}$. But that definition must satisfy the following conditions:}

\emph{1.~If $\mathfrak{p}_0$ is any point of $\mathfrak{F}$, then there exists a function $t(\mathfrak{p})$ that is regular-analytic not only at the point $\mathfrak{p}_0$ (where it possesses the value $0$) but also at all points $\mathfrak{p}$ of a certain neighborhood of $\mathfrak{p}_0$ on $\mathfrak{F}$, and which maps this neighborhood one-to-one and domain-continuously into the complex $t$-plane; such a function is called a \textbf{local uniformizing variable} for $\mathfrak{p}_0$.}

\emph{2.~If $f(\mathfrak{p})$ is any function regular-analytic at the point $\mathfrak{p}_0$ and $t(\mathfrak{p})$ is a local uniformizing variable belonging to $\mathfrak{p}_0$, then there always exists a neighborhood $\mathfrak{U}_0$ of $\mathfrak{p}_0$ in which $f(\mathfrak{p})$ can be represented as a regular power series in $t(\mathfrak{p})$}
\[ f(\mathfrak{p}) = a_0 + a_1 t(\mathfrak{p}) + a_2 (t(\mathfrak{p}))^2 + \cdots \tag{10} \]
\emph{.}

From these requirements it follows: If $\tau$, alongside $t$, is another local uniformizing variable belonging to $\mathfrak{p}_0$, then in a certain neighborhood of $\mathfrak{p}_0$ a representation
\[ \tau = \gamma_1 t + \gamma_2 t^2 + \cdots \]
must be valid. But since conversely $t$ must also be expressible in an analogous manner by $\tau$, necessarily $\gamma_1 \neq 0$. In order to demonstrate the analyticity of a function $f(\mathfrak{p})$ at the point $\mathfrak{p}_0$, it therefore always suffices to prove the existence of a representation (10) by \emph{a single} local uniformizing variable $t(\mathfrak{p})$ belonging to $\mathfrak{p}_0$.

If any two Riemann surfaces $\mathfrak{F}_1$, $\mathfrak{F}_2$ are given, then a mapping that maps $\mathfrak{F}_1$ point by point, one-to-one and domain-continuously onto $\mathfrak{F}_2$ in such a way that every function regular-analytic at any point of $\mathfrak{F}_1$ passes by this mapping into a function on $\mathfrak{F}_2$ that is regular-analytic at the image point, is called a \textbf{conformal mapping}. The reason for this designation will soon become evident. Two Riemann surfaces that can be mapped conformally onto one another are to be regarded as \textbf{(conformally) equivalent} and merely as different realizations of one and the same ideal Riemann surface. Only those properties that are invariant under conformal mapping can ever count as \textbf{intrinsic properties of a Riemann surface}, which therefore, if they belong to \emph{one} Riemann surface $\mathfrak{F}$, also attach to every Riemann surface equivalent to it. All analysis situs qualities naturally belong to these intrinsic properties of a Riemann surface.

\origpage{37}
Every sub-domain of a Riemann surface is itself a Riemann surface. Every local uniformizing variable for a point $\mathfrak{p}$ maps a certain neighborhood of $\mathfrak{p}$ conformally onto a planar domain. Here the plane is likewise to be conceived as a Riemann surface, and indeed in the way that corresponds to the elementary concept of an analytic function in the complex Gaussian number plane.

We have discussed above in what sense an analytic configuration can be viewed as a Riemann surface. But the concepts ``analytic configuration'' and ``Riemann surface'' do not coincide. Through an analytic configuration $(z, u)$ there is given to us not merely a Riemann surface, but simultaneously two functions $z$ and $u$ on it, regular except for poles. Here $z$ and $u$ satisfy the following condition: there exist no two distinct points $\mathfrak{p}_1^0, \mathfrak{p}_2^0$ on the Riemann surface, associated local uniformizing variables $t_1$, resp.\ $t_2$, and two series $P(t), Q(t)$ proceeding in integral powers of $t$ of such a kind that
\[ \begin{gathered}
z = P(t_1), \quad u = Q(t_1) \text{ in the neighborhood of } \mathfrak{p}_1^0, \\
z = P(t_2), \quad u = Q(t_2) \text{ in the neighborhood of } \mathfrak{p}_2^0
\end{gathered} \]
is satisfied. For an arbitrary Riemann surface one always obtains an analytic configuration by distinguishing on it any two functions $z, u$ regular except for poles that satisfy the condition just formulated. If, however, there exists \emph{one} such pair of functions $z, u$ at all, it can also always be chosen in infinitely many different ways; e.\,g.\ instead of $z, u$ I can use any two linear combinations of $z, u$:
\[ \begin{gathered}
z' = a z + b u, \quad u' = A z + B u \\
[a, b;\, A, B \text{ constant};\; a B - b A \neq 0].
\end{gathered} \]
\emph{That to every given Riemann surface there really belongs a pair of functions $(z, u)$, i.\,e.\ an analytic configuration,} is a fundamental fact of Riemann's function theory, the proof of which for closed Riemann surfaces will be provided in Chap.\ II of this work with the help of the \emph{Thomson}--\emph{Dirichlet} principle employed by \emph{Riemann} for the same purpose.

\origpage{38}
\emph{We now enumerate several invariant concepts that concern the behavior of functions and curves on Riemann surfaces.}

If the expansion of a function $f(\mathfrak{p})$ regular at the point $\mathfrak{p}_0$ in powers of the local uniformizing variable $t$ belonging to $\mathfrak{p}_0$ begins with the term $a_m t^m\,(a_m \neq 0)$, then $\mathfrak{p}_0$ is a \textbf{zero of the $m^{\text{th}}$ order} of $f$. We also say: $f$ has at the place $\mathfrak{p}_0$ the \textbf{order} $m$. If $\tau(\mathfrak{p})$ is any other local uniformizing variable for $\mathfrak{p}_0$,
\[ t = c_1 \tau + c_2 \tau^2 + \cdots \quad (c_1 \neq 0), \]
then the expansion of $f$ in powers of $\tau$ begins with the term $a_m c_1^m \cdot \tau^m$; from this follows the ``invariance'' of the order $m$ of a zero. If a function $f(\mathfrak{p})$ permits in a certain neighborhood of the point $\mathfrak{p}_0$, apart from the point $\mathfrak{p}_0$ itself, an expansion
\[ f = \frac{a_{-n}}{t^n} + \cdots + \frac{a_{-1}}{t} + a_0 + a_1 t + \cdots \quad [a_{-n} \neq 0,\, n \text{ integer and } > 0], \]
then $f$ has at the place $\mathfrak{p}_0$ a \textbf{pole of $n^{\text{th}}$ order}; we also say: $f$ has at the place $\mathfrak{p}_0$ the \textbf{order} $-n$. Proof of invariance as above. If the function $f(\mathfrak{p})$ has at the place $\mathfrak{p}_0$ the order $k\,(\gtreqqless 0)$ and $g(\mathfrak{p})$ has the order $l$, then $f \cdot g$ has at $\mathfrak{p}_0$ the order $k + l$, and $\frac{f}{g}$ the order $k - l$.

If a single-valued function $f$ is regular-analytic in the neighborhood of the point $\mathfrak{p}_0$, apart from this point itself, then one can either assign to it at the point $\mathfrak{p}_0$ such a value that it is regular at $\mathfrak{p}_0$ as well, or it has a pole at $\mathfrak{p}_0$, or it has there an \textbf{essential singularity}. In the last case $f(\mathfrak{p})$ comes arbitrarily close to every value in every neighborhood of $\mathfrak{p}_0$.

If $z, u$ are two functions regular except for poles in a domain $\mathfrak{G}$ of the Riemann surface, then every rational expression $R(z, u)$ of $z, u$ is also regular except for poles in $\mathfrak{G}$. Namely, by substituting for $z, u$ their expansions in powers of the local uniformizing variable $t$ belonging to a point of the domain, one obtains for $R(z, u)$ a series proceeding in integral powers of $t$ that contains at most finitely many negative powers. Thus, even where an indeterminate expression of the form $\frac{0}{0}$ initially arises for $R$ by direct substitution of the \emph{values} of $z$ and $u$, there are in truth no places of indeterminacy present.

A real function $U$ is called a \textbf{harmonic} or \textbf{potential function} at a place $\mathfrak{p}_0$ if there exists a function regular-analytic at this point with whose real part $U$ agrees in a certain neighborhood of $\mathfrak{p}_0$. If $U$ is harmonic at all points of a domain, but not $=\mathrm{const.}$, then $U$ can have a maximum (greatest value) or minimum (least value) at no point of this domain. \emph{Therefore on a closed Riemann surface, apart from constants, there exists no everywhere harmonic and a fortiori no everywhere regular-analytic (single-valued)}

\origpage{39}
\emph{function.} --- For a real function $U$ to be \textbf{continuously differentiable} at the place $\mathfrak{p}_0$, $U$, which in a certain neighborhood of $\mathfrak{p}_0$ can be conceived as a function of the local uniformizing variable $t = x + iy$ ($x, y$ real) belonging to $\mathfrak{p}_0$, must possess continuous first partial derivatives $\frac{\partial U}{\partial x}$, $\frac{\partial U}{\partial y}$ for sufficiently small values of $|t|$. Which local uniformizing variable $t$ is drawn upon in applying this criterion is of course immaterial. Similarly one can speak of $2\,\text{times}$, $3\,\text{times}$, \ldots continuously differentiable functions.

A curve $\mathfrak{p} = \mathfrak{p}(\lambda)\;[0 \leqq \lambda \leqq 1]$ can be represented in the form $t = t(\lambda)$ in such a neighborhood of one of its points $\mathfrak{p}(\lambda_0) = \mathfrak{p}_0$ that is mapped one-to-one and conformally onto a domain of the $t$-plane by the local uniformizing variable $t$ belonging to $\mathfrak{p}_0$. If $t(\lambda)$, for real $\lambda$ lying sufficiently close to $\lambda_0$ and belonging to the interval $(01)$, agrees with a convergent power series $b_1 (\lambda - \lambda_0) + b_2 (\lambda - \lambda_0)^2 + \cdots$ in which $b_1 = \left(\frac{dt}{d\lambda}\right)_{\lambda = \lambda_0} \neq 0$, then the curve is called \textbf{analytic for $\lambda = \lambda_0$}. The proof of invariance is trivial. By using as local uniformizing variable the parameter $\tau$ defined by $t = b_1 \tau + b_2 \tau^2 + \cdots$, a certain piece of the curve containing $\mathfrak{p}(\lambda_0)$ appears in the $\tau$-plane as a segment of the real axis. By an \textbf{analytic curve} simply, one is to understand one that is analytic for all values of the parameter $\lambda$ running from $0$ to $1$.

If it is only known that the derivative $\frac{dt}{d\lambda}$ exists and is continuous in $\lambda$ for values $\lambda$ lying sufficiently close to $\lambda_0$ and in the interval $(01)$, but possesses a value $\neq 0$ for $\lambda = \lambda_0$, then we call the curve \textbf{continuously differentiable for $\lambda_0$}. If
\[ \mathfrak{p} = \overset{1}{\mathfrak{p}}(\lambda), \quad \mathfrak{p} = \overset{2}{\mathfrak{p}}(\lambda) \tag{11} \]
are two curves emanating from the same point $\mathfrak{p}_0 = \overset{1}{\mathfrak{p}}(0) = \overset{2}{\mathfrak{p}}(0)$, continuously differentiable for $\lambda = 0$, and
\[ t = \overset{1}{t}(\lambda), \quad t = \overset{2}{t}(\lambda) \quad [0 \leqq \lambda \leqq \lambda_1] \]
are the images of the initial arcs of those two curves in the plane of the local uniformizing variable $t$ belonging to $\mathfrak{p}_0$, then we shall also designate the angle $\vartheta$ which these arcs form with one another at the origin of the $t$-plane, and which is determined up to integral multiples of $2\pi$ by the equations
\[ \begin{gathered}
\left(\frac{d\overset{1}{t}}{d\lambda}\right)_{\lambda = 0} = r_1 e^{i\,\vartheta_1}, \quad \left(\frac{d\overset{2}{t}}{d\lambda}\right)_{\lambda = 0} = r_2 e^{i\,\vartheta_2} \\
[r_1, r_2 > 0; \quad \vartheta_1, \vartheta_2 \text{ real}] \\
\vartheta = \vartheta_2 - \vartheta_1.
\end{gathered} \]
as the \textbf{angle} of the two curves (11) emanating from $\mathfrak{p}_0$ on the

\origpage{40}
Riemann surface. This measure of angle is an invariant because the transition from one local uniformizing variable $t$ to another $\tau$ is mediated by an \emph{angle-preserving} mapping of a domain of the $t$-plane containing the origin onto an identical such domain of the $\tau$-plane. \emph{On a Riemann surface there thus exists an invariant measure of angle.}${}^{1)}$ The concept of ``conformal'' mapping explained above coincides, after the introduction of this measure of angle, with the concept of one-to-one \emph{angle-preserving} mapping.

A surface free of singularities in Euclidean space, such as a sphere or torus, on which (after agreement on the direction of rotation) a definite natural Euclidean angle measurement exists, can be conceived in a unique way as a Riemann surface such that the angle measurement belonging to it as a Riemann surface in the manner just described agrees with the natural one. Here the possibility of mapping a neighborhood of each point of that surface conformally (in the Euclidean sense) onto a planar domain must be established.${}^{2)}$

For the \emph{sphere}, stereographic projection provides a function $z$ regular on the entire sphere except for a single pole of 1st order, which takes every value including $\infty$ once and only once. Every other function regular except for poles on the sphere is a rational function of $z$, so that the theory of functions regular except for poles on the sphere essentially coincides with that of rational functions of one variable $z$. The choice of the independent variable $z$ is of course arbitrary to a certain degree; for in addition to $z$, any function $z'$ that arises from $z$ through a linear transformation
\[ z' = \frac{az + b}{cz + d} \quad [a, b, c, d \text{ constant}; \quad ad - bc \neq 0] \]
is also suitable as such. The position and multiplicity of the zeros and poles of an analytic function free of essential singularities can be prescribed arbitrarily on the sphere, provided only that the total order of the zeros agrees with the total order of the poles.

Matters turn out quite differently on the \emph{torus} (cf.\ Chap.\ II). As the intrinsic reason for the great differences that exist between the behavior of functions on the torus on the one hand, and functions on the sphere on the other, one can point almost everywhere to the single circumstance (which lies not in the realm of function theory, but of analysis situs!) that on the torus a system of two closed curves meeting each other at one point can be drawn ($\varphi = 0$ and $\psi = 0$ in the notation of § 5, Example 7) which does not dissect the torus --- whereas such a system does not exist on the sphere. If we map the torus conformally onto the point-lattice manifold $\mathfrak{T}$ (§ 5), the functions free of essential singularities on the torus appear as those single-valued functions of the complex variable $w = \varphi + i\psi$, regular except for poles, which \emph{are doubly periodic with periods $2\pi$, $\frac{2\pi ir}{\sqrt{R^2 - r^2}}$, that is, as elliptic functions} (with purely imaginary period ratio $= \frac{ir}{\sqrt{R^2 - r^2}}$).

\textbf{1)} That in a certain sense such a measurement of length also exists is a deep-lying fact from the theory of ``uniformization''. Cf.\ §§ 19, 20 of this book.

\textbf{2)} Cf.\ \emph{L.~Lichtenstein}, Beweis des Satzes, daß jedes hinreichend kleine, im wesentlichen stetig gekrümmte, singularitätenfreie Flächenstück auf einen Teil einer Ebene zusammenhängend und in den kleinsten Teilen ähnlich abgebildet werden kann. Abhandlungen der K.~Preuß.\ Akademie der Wissenschaften vom Jahre 1911, Anhang. --- Difficulties arise if the spatial surface possesses corners. On this problem cf.\ \emph{H.~A.~Schwarz} in the paper already cited, Gesammelte Abhandlungen Vol.\ 2, p.\ 161; it found its solution through \emph{Koebe}, Göttinger Nachrichten 1908, pp.\ 359---360, and \emph{R.~König}, Mathematische Annalen Vol.\ 71, 1912, pp.\ 184---205.

\origpage{41}
Two tori are always equivalent as surfaces in the sense of analysis situs. However, they cannot in general be mapped conformally onto one another; rather, this is possible if and only if the ``modulus'' $\frac{r\sqrt{R^2 - r^2}}{R^2}$ has the same value for both tori. The word \textbf{``modulus''} has here this meaning: If in a family of Riemann surfaces, all of which are mutually equivalent in the sense of analysis situs, a number is somehow assigned to each of them in such a way that two surfaces of the family are in any case assigned the same number if the two surfaces are \emph{conformally} equivalent, then that number, regarded in its dependence on the Riemann surfaces of the family, is called a modulus of the family. The fact that equivalence in the sense of analysis situs, generally speaking, does not entail conformal equivalence of Riemann surfaces is of fundamental importance. --- We represent the points of the first torus to ourselves as the point lattices of a $w_1$-plane with periods $2\pi$ and $2\pi i a_1 = 2\pi i \frac{r_1}{\sqrt{R_1^2 - r_1^2}}$, and the points of the second torus as the point lattices of a $w_2$-plane corresponding to the periods $2\pi$ and $2\pi i a_2\;[a_2 > 0]$. If a conformal mapping of one torus onto the other is given, then to each lattice $w_1$ there is assigned a lattice $w_2$ in such a way that if we impart to the lattice $w_1$ an infinitesimal displacement $dw_1$, the image lattice $w_2$ experiences an infinitesimal displacement $dw_2$ whose ratio $\frac{dw_2}{dw_1}$ to $dw_1$ depends only on the lattice $w_1$, but not on the \emph{direction} of the displacement $dw_1$; regarded in its dependence on the lattice $w_1$, this ratio is a regular-analytic function on the first torus, hence must be $=\mathrm{const.} = A$. From this it follows that the conformal mapping must be representable by a formula
\[ w_2 = A w_1 + B \quad [A, B \text{ constants}] \tag{12} \]

\origpage{42}
in the sense that from (12), if one substitutes for $w_1$ all the points of a $w_1$-lattice, one always obtains all the points $w_2$ of the corresponding $w_2$-lattice. A simple lattice consideration of number-theoretic nature teaches that this is possible only if $A = \pm 1$ or $= \pm\,\frac{i}{a_1}$; in the first case $a_2 = a_1$, in the second $a_2 = \frac{1}{a_1}$. In any case, therefore, for the possibility of conformal mapping of the two tori the equation
\[ a_1 + \frac{1}{a_1} = a_2 + \frac{1}{a_2} \]
is the necessary (and obviously also sufficient) condition. This agrees with our assertion, since
\[ \frac{r}{\sqrt{R^2 - r^2}} + \frac{\sqrt{R^2 - r^2}}{r} = \frac{R^2}{r\sqrt{R^2 - r^2}} \]

We conclude this section with some general remarks on \emph{the idea of the Riemann surface}. The fundamental idea underlying its introduction is by no means restricted to complex function theory. A function of two real variables $x, y$ is a \emph{function in the plane}; but it is certainly just as justified to investigate functions on the sphere, on the torus, or on a surface in general, as precisely in the plane. As long as one is concerned admittedly only with the behavior of functions ``in the small'' --- and to this most considerations of analysis refer ---, the concept of a function of two real variables is general enough, since the neighborhood of every point of a two-dimensional manifold can be represented by $x, y$ (or $x + iy$). As soon as one proceeds, however, to the investigation of the behavior of functions \emph{``in the large''}, the functions in the plane form an important but \emph{special case among infinitely many other equally legitimate ones}; \emph{Riemann} and \emph{Klein} taught us not to remain standing at this special case. Applied to complex function theory, this means: \emph{before proceeding to the study of any species of functions, that surface which provides the domain of variability of the independent argument must always first be defined; then it must be explained what ``analytic function'' shall mean on this surface, whereby the surface becomes a Riemann surface; and only then can one set to work on the functions themselves.} Accordingly, one has to observe properties of three different levels in functions: 1.~and these are the most incisive: the analysis situs qualities of the Riemann surface on which the functions exist; 2.~the intrinsic properties of this Riemann surface not belonging to the domain of analysis situs (e.\,g.\ a definite value of a ``modulus''); 3.~those properties (such as position

\origpage{43}
and order of zeros and poles) by which functions differ with respect to their behavior on the same Riemann surface. From this standpoint the Weierstrassian concept of the analytic configuration plays only a secondary role: it comes about only through combining two functions existing on one and the same Riemann surface. It is a natural step, instead of only two functions, then also to consider simultaneously perhaps three or four or even more functions of the same variable point on a Riemann surface. Geometrically speaking, taking this step means nothing other than: passing from the study of planar analytic ``curves'' to that of curves in three-dimensional, four-dimensional, etc., space.

The single-valued functions on a Riemann surface $\mathfrak{F}$, regular everywhere except for poles, we shall mostly designate briefly as the \textbf{``functions on the surface''}. For \emph{closed Riemann surfaces} we shall obtain in Chap.~II an overview of all these functions. But one can also envisage the following more general class of functions belonging to $\mathfrak{F}$: Let a function element on $\mathfrak{F}$ be given (i.\,e.\ a power series proceeding in integral powers of the local uniformizing variable $t$ belonging to a point $\mathfrak{p}_0$ of $\mathfrak{F}$); one can then attempt to continue this function element analytically along all possible paths on $\mathfrak{F}$ in a manner analogous to that described for the case of the plane in §~1. If such a continuation is uniquely possible along all paths, i.\,e.\ such that one encounters at most poles, but never points beyond which continuation is impossible altogether (``natural boundaries'') or becomes multiple-valued (branching relative to $\mathfrak{F}$), then the resulting function, as the example $w = \varphi + i\psi$ on the torus shows, need by no means be single-valued; rather, in general, such a function obtained by analytic continuation will become single-valued only on a certain \emph{covering surface} spreading without boundaries and without branchings over $\mathfrak{F}$. Now in many questions, notably in \emph{uniformization theory}, it is of great importance to extend the realm of single-valued functions on $\mathfrak{F}$ to the more comprehensive realm of all (finitely or infinitely multiple-valued) functions that are unramified on $\mathfrak{F}$ and without essential singularities and natural boundaries. This imposes upon us the obligation to consider, in this chapter, the remainder of which is devoted to the questions of analysis situs fundamental to function theory, alongside each surface $\mathfrak{F}$ the \emph{covering surfaces} belonging to it.

\subsection*{§ 8. Schlichtartig Surfaces.}

A surface $\mathfrak{F}$ in a definite triangulation $\zeta$ I denote by $\mathfrak{F}_\zeta$. As an \textbf{elementary segment on} $\mathfrak{F}_\zeta$ counts every elementary segment in $\Delta$ contained entirely in a triangle $\Delta$ of $\zeta$ (see §~6, p.~31). By virtue of

\origpage{44}
one of the two ends of the elementary segment being designated as the initial point and the other as the terminal point, this segment is \textbf{directed}. Finitely many elementary segments $\sigma_1\,\sigma_2 \cdots \sigma_n$ form a \textbf{polygonal path} if the terminal point of $\sigma_h$ always coincides with the initial point of $\sigma_{h+1}\;[h = 1, 2, \dots, n-1]$. If any two segments of a polygonal path have a point in common (and only one) only when they succeed one another, then the polygonal path \textbf{does not self-intersect}: it is a \textbf{``simple''} polygonal path. If the terminal point of $\sigma_n$ coincides with the initial point of $\sigma_1$, then the polygonal path is \textbf{closed}; $\sigma_1$ is then the segment succeeding $\sigma_n$ (cyclic ordering). A closed polygonal path in which two segments have a point in common only when they succeed one another is designated as a \textbf{polygon}.

\emph{Two points of a domain $\mathfrak{G}$ on $\mathfrak{F}_\zeta$ can always be connected by a simple polygonal path running entirely in $\mathfrak{G}$.} Namely, one can first connect the two points by a curve $\gamma$ running entirely in $\mathfrak{G}$. From the partition $\zeta$ one can produce such a fine subdivision $\zeta'$ that all elementary triangles of $\zeta'$ that have points in common with $\gamma$ lie in $\mathfrak{G}$. By the construction indicated on pp.~23\,f., one obtains a simple chain of triangles of the partition $\zeta'$ that connects that triangle of $\zeta'$ in which the initial point of $\gamma$ lies with the triangle containing the terminal point of $\gamma$. The chain of triangles can then be replaced immediately by a simple polygonal path that traverses the triangles of the chain successively in one segment each. --- Furthermore: If $\sigma$ is any polygonal path on $\mathfrak{F}_\zeta$, one can specify such a subdivision $\zeta'$ of $\zeta$ that the segments of which $\sigma$ consists are edges of the partition $\zeta'$ (or composed of edges of $\zeta'$), so that $\sigma$ is an \textbf{edge path} on $\mathfrak{F}_{\zeta'}$.

\rmfigure{figures/fig-10.png}{Fig.~10.}{Simple polygonal path with adjacent triangles}

Of a closed set $\mathfrak{E}$ on $\mathfrak{F}$ we say that it \textbf{does not dissect} $\mathfrak{F}$ if the points of $\mathfrak{F}$ that do not belong to $\mathfrak{E}$ constitute a single domain. \emph{A simple polygonal path $\sigma$ does not dissect $\mathfrak{F}_\zeta$.} For the proof we make such a subdivision $\zeta'$ of $\zeta$ that $\sigma$ appears as an edge path. The triangles of $\zeta'$ abutting $\sigma$ can then be numbered in such a way: 1, 2, 3, \dots, that any two consecutive triangles in this numbering have an edge in common not belonging to $\sigma$. By the figure the beginning of such a numbering is indicated${}^{1)}$. I assert that any two points $\mathfrak{p}$ and $\mathfrak{q}$ of $\mathfrak{F}$ can be connected with one another by a continuous curve not meeting $\sigma$. I connect $\mathfrak{p}$ and

\textbf{1)} One will admittedly not always be able to avoid here that the same triangle occasionally occurs twice or several times with different numbers.

\origpage{45}
$\mathfrak{q}$ by a curve $\gamma$ on $\mathfrak{F}$. If $\gamma$ meets the polygonal path $\sigma$, let $\mathfrak{p}'$ be the first, $\mathfrak{q}'$ the last intersection point, $\mathfrak{p}''$ a point on $\gamma$ so shortly before $\mathfrak{p}'$ that $\mathfrak{p}''$ lies in one of the triangles abutting $\sigma$; $\mathfrak{q}''$ a point on $\gamma$ so shortly after $\mathfrak{q}'$ that $\mathfrak{q}''$ is also still situated in such a triangle. With the help of the above-mentioned triangle numbering I can then connect $\mathfrak{p}''$ and $\mathfrak{q}''$ by a polygonal path that does not meet $\sigma$ and runs through the triangles abutting $\sigma$ in the order of their numbering. This polygonal path forms, together with the arc $\mathfrak{p}\mathfrak{p}''$ and the arc $\mathfrak{q}''\mathfrak{q}$ of $\gamma$, a curve as we desire.

If $\mathfrak{E}$ is a closed set that does not dissect $\mathfrak{F}_\zeta$, and $\sigma$ is an elementary segment on $\mathfrak{F}_\zeta$ that has the two endpoints, but otherwise no point, in common with $\mathfrak{E}$, then the union $\mathfrak{E} + \sigma$ dissects the surface $\mathfrak{F}$ either not at all or into two domains. For the proof I can imagine $\sigma$ as the common edge of two triangles $\Delta_1$, $\Delta_2$ of $\zeta$. Let $\mathfrak{p}_0$ be a point of $\sigma$ that is neither of the endpoints, $\mathfrak{p}_0\mathfrak{q}_1$ a small elementary segment leading into the interior of $\Delta_1$ having no point in common with $\mathfrak{E}$, and $\mathfrak{p}_0\mathfrak{q}_2$ a similar segment leading into the interior of $\Delta_2$. I assert: every point $\mathfrak{p}$ of $\mathfrak{F}$ not belonging to $\mathfrak{E} + \sigma$ can be connected without crossing $\mathfrak{E} + \sigma$ either with $\mathfrak{q}_1$ or with $\mathfrak{q}_2$. I first connect $\mathfrak{p}$ with $\mathfrak{q}_1$ without crossing $\mathfrak{E}$ by the curve $\gamma$. If I thereby meet $\sigma$ for the first time at the point $\mathfrak{p}_1$ (which can be neither of the endpoints of $\sigma$!), I shall be able to specify a point $\mathfrak{p}_1'$ shortly before $\mathfrak{p}_1$ such that the partial arc $\mathfrak{p}_1'\mathfrak{p}_1$ of $\gamma$ lies either entirely in $\Delta_1$ or entirely in $\Delta_2$. According as the one or the other is the case, I can reach from this arc, by means of a single elementary segment in $\Delta_1$ resp.\ $\Delta_2$ that meets neither $\mathfrak{E}$ nor $\sigma$, the segment $\mathfrak{p}_0\mathfrak{q}_1$, resp.\ $\mathfrak{p}_0\mathfrak{q}_2$. --- In particular, $\mathfrak{F}$ is not dissected by $\mathfrak{E} + \sigma$ at all if $\mathfrak{q}_1$ can be connected with $\mathfrak{q}_2$ (i.\,e.\ a point on the \emph{one} with a point on the \emph{other} ``bank'' of $\sigma$) without crossing $\mathfrak{E} + \sigma$.

The combination of the facts proven in the two preceding paragraphs yields the result that \emph{every polygon dissects the surface $\mathfrak{F}_\zeta$ into at most two domains}; for a polygon can be conceived as having arisen by addition of a single segment ($\sigma$) to a simple polygonal path ($\mathfrak{E}$). If $\mathfrak{F}_\zeta$ is indeed dissected by every polygon, then the surface $\mathfrak{F}$ is called (after Koebe) \textbf{schlichtartig}. That there exist non-schlichtartig surfaces is shown by the example of the torus.

\emph{Theorem: If a surface $\mathfrak{F}$ is schlichtartig in the triangulation $\zeta\,\mathrm{I}$, then it is also schlichtartig if it is partitioned into elementary triangles in any other way $\zeta\,\mathrm{II}$.}

We distinguish the concepts relating to the one and the other partition, such as ``segment'', ``polygon'', and the like, by the addition of a I, resp.\ II. If $\mathfrak{F}$ were not schlichtartig in the triangulation $\zeta\,\mathrm{II}$, there would exist a polygon II, $\pi_{\mathrm{II}}$, that does not dissect $\mathfrak{F}$. Let $\sigma_{\mathrm{II}}$ be a segment of $\pi_{\mathrm{II}}$,

\origpage{46}
which we can assume to run, apart perhaps from its endpoints, in the interior of a triangle $\Delta_{\mathrm{II}}$${}^{1)}$. We cross (Fig.~11) $\sigma_{\mathrm{II}}$ by an elementary segment II lying entirely in the interior of this triangle $\Delta_{\mathrm{II}}$, $\tau_{\mathrm{II}} = \mathfrak{q}'\mathfrak{q}''$, which may meet $\sigma_{\mathrm{II}}$ at $\mathfrak{p}$. $\mathfrak{q}'$ can be connected with $\mathfrak{q}''$ without crossing $\pi_{\mathrm{II}}$ by a continuous curve, and hence also by a simple polygonal path I, which we call $\varSigma_{\mathrm{I}}$. I may assume that $\varSigma_{\mathrm{I}}$, apart from its two endpoints, has no point in common with the segment $\tau_{\mathrm{II}}$.${}^{2)}$

\rmfigure{figures/fig-11.png}{Fig.~11.}{To the invariance proof of the schlichtartig property}

Since $\varSigma_{\mathrm{I}}$ does not dissect the surface and the polygon $\pi_{\mathrm{II}}$ (if I omit from it a small segment containing the point $\mathfrak{p}$) connects a point on one bank of $\tau_{\mathrm{II}}$ with a point on the other bank without crossing $\varSigma_{\mathrm{I}} + \tau_{\mathrm{II}}$, $\mathfrak{F}$ is not dissected\ednote{Misprint in the original: durcb instead of durch.} by $\varSigma_{\mathrm{I}} + \tau_{\mathrm{II}}$. Now I consider a segment $\sigma_{\mathrm{I}}$ of $\varSigma_{\mathrm{I}}$, of which I can again assume that it is not an edge of a triangle I, cross it in the triangle I in which it lies by a small segment I, $\tau_{\mathrm{I}}$, and connect the endpoints of $\tau_{\mathrm{I}}$ by a simple polygonal path $T_{\mathrm{I}}$ that does not cross $\varSigma_{\mathrm{I}} + \tau_{\mathrm{II}}$. I may assume that $T_{\mathrm{I}}$ has only the endpoints in common with $\tau_{\mathrm{I}}$; $\pi_{\mathrm{I}} = T_{\mathrm{I}} + \tau_{\mathrm{I}}$ is then a polygon I that does not dissect $\mathfrak{F}$, since the curve $\varSigma_{\mathrm{I}} + \tau_{\mathrm{II}}$ connects two points lying opposite each other on the two banks of $\sigma_{\mathrm{I}}$ without crossing $\pi_{\mathrm{I}}$. Thus $\mathfrak{F}$ cannot have been schlichtartig in the triangulation $\zeta\,\mathrm{I}$ either.${}^{3)}$

\textbf{1)} Namely, if all segments of $\pi_{\mathrm{II}}$ are edges, we can replace an arbitrary one of them, $\sigma_{\mathrm{II}}$, by a once-broken polygonal path running in one of the two triangles abutting $\sigma_{\mathrm{II}}$, without $\pi_{\mathrm{II}}$ losing the property of not dissecting $\mathfrak{F}$.

\textbf{2)} For otherwise let $\overset{*}{\mathfrak{q}}'$ be that point where $\varSigma_{\mathrm{I}}$, traversed from $\mathfrak{q}'$ to $\mathfrak{q}''$, meets the segment $\mathfrak{p}\mathfrak{q}'$ for the last time, and $\overset{*}{\mathfrak{q}}''$ that point where after this moment $\varSigma_{\mathrm{I}}$ meets the segment $\mathfrak{p}\mathfrak{q}''$ for the first time, while $\varSigma_{\mathrm{I}}^*$ is the part of $\varSigma_{\mathrm{I}}$ traversed between these two events: then $\varSigma_{\mathrm{I}}^*$ has only the endpoints in common with the segment II, $\overset{*}{\tau}_{\mathrm{II}} = \overset{*}{\mathfrak{q}}'\overset{*}{\mathfrak{q}}''$, and we could replace $\varSigma_{\mathrm{I}}$ by $\varSigma_{\mathrm{I}}^*$, and $\tau_{\mathrm{II}}$ by $\overset{*}{\tau}_{\mathrm{II}}$.

\textbf{3)} We have thereby proven a piece of the ``Jordan curve theorem'', which asserts that every simple closed curve on a schlichtartig surface (in particular in the plane) dissects this surface. In addition to the first (incomplete) proof by C.~Jordan himself in his \emph{Cours d'analyse}, 2nd ed., Vol.~I, pp.~91---99, see especially Brouwer, Math.~Ann., Vol.~69 (1910), pp.~169---175.

\origpage{47}

\subsection*{§ 9. Covering Surfaces. Simply Connected Surfaces. Monodromy Theorem and Cauchy's Integral Theorem.}

Even sharper than the concept of the schlichtartig property is that of \emph{simple connectivity}, which stands in close relation to the construction of covering surfaces. If $\mathfrak{F}$ is a given surface, the \textbf{``base surface''}\ednote{Misprint in the original: fi ligature instead of fl ligature (Grundfiäche instead of Grundfläche).}, we will call the surface $\bar{\mathfrak{F}}$ a \textbf{covering surface over $\mathfrak{F}$} if to each point $\bar{\mathfrak{p}}$ on $\bar{\mathfrak{F}}$ there is assigned a single point $\mathfrak{p}$ on $\mathfrak{F}$ as the \textbf{``trace point''} of $\bar{\mathfrak{p}}$; we then also say that $\bar{\mathfrak{p}}$ \textbf{lies over} $\mathfrak{p}$. This assignment shall, in any case when we designate $\bar{\mathfrak{F}}$ as \textbf{unramified}${}^{1)}$ (relative to $\mathfrak{F}$), satisfy the following condition: if $\bar{\mathfrak{p}}_0$ is an arbitrary point of $\bar{\mathfrak{F}}$, there is always a neighborhood of $\bar{\mathfrak{p}}_0$ on $\bar{\mathfrak{F}}$ that is \emph{related by that assignment one-to-one and domain-continuously onto a domain of $\mathfrak{F}$}. Let $\bar{\mathfrak{F}}$ denote an unramified covering surface over $\mathfrak{F}$ in this sense. Let $\mathfrak{p} = \mathfrak{p}(\lambda)\;[0 \leqq \lambda \leqq 1]$ be a curve $\gamma$ on $\mathfrak{F}$ emanating from the point $\mathfrak{p}_0 = \mathfrak{p}(0)$, and $\bar{\mathfrak{p}}_0$ a point on $\bar{\mathfrak{F}}$ over $\mathfrak{p}_0$. Then (cf.\ §~1) two cases can occur:

either: there is a unique continuous curve $\bar{\mathfrak{p}} = \bar{\mathfrak{p}}(\lambda)\;[0 \leqq \lambda \leqq 1]$ on $\bar{\mathfrak{F}}$ emanating from $\bar{\mathfrak{p}}_0$, such that for every value of the parameter $\lambda$ the point $\bar{\mathfrak{p}}(\lambda)$ lies over $\mathfrak{p}(\lambda)$;

or: there exists a threshold $\varLambda_0\;(0 < \varLambda_0 \leqq 1)$, such that over every partial arc $0 \leqq \lambda \leqq \lambda_0$ of $\gamma$ for which $\lambda_0 < \varLambda_0$ there does lie one (and only one) curve on $\bar{\mathfrak{F}}$ emanating from $\bar{\mathfrak{p}}_0$, but this is no longer the case as soon as $\lambda_0 \geqq \varLambda_0$. This latter case can be conceived as follows: if one traces from $\bar{\mathfrak{p}}_0$ the continuous change of a point $\bar{\mathfrak{p}}$ on $\bar{\mathfrak{F}}$ whose trace point on $\mathfrak{F}$ describes the curve $\gamma$, one encounters, before the end is reached, over the point $\mathfrak{p}(\varLambda_0)$ of $\mathfrak{F}$ a boundary of the covering surface.

If only the first case always occurs, whatever the point $\bar{\mathfrak{p}}_0$ on $\bar{\mathfrak{F}}$ may be and whatever the curve $\gamma$ emanating from the trace point $\mathfrak{p}_0$ of $\bar{\mathfrak{p}}_0$ may be, we shall accordingly have to designate the covering surface as \textbf{unbounded}. If over each point of $\mathfrak{F}$ there lies a single point of the unramified unbounded covering surface $\bar{\mathfrak{F}}$, then $\bar{\mathfrak{F}}$ is \textbf{single-sheeted} and not essentially different from $\mathfrak{F}$. If to a surface $\mathfrak{F}$ there belong no other unramified unbounded covering surfaces than single-sheeted ones, then $\mathfrak{F}$ is called \textbf{simply connected}${}^{2)}$.

\textbf{1)} The word does not really express completely what it is intended to say according to the definition; it would be more correct (but also more cumbersome) to say ``unramified and unpleated''.

\textbf{2)} This definition emphasizes that property of simply connected surfaces which is the decisive one for function-theoretic applications.

\origpage{48}
E.\,g.\ the interior $\mathfrak{K}$ of a circle in the Euclidean number plane is a simply connected surface. In order to show this, we consider an arbitrary unramified unbounded covering surface $\bar{\mathfrak{K}}$ over $\mathfrak{K}$. Let $\bar{\mathfrak{o}}$ be a point on $\bar{\mathfrak{K}}$ lying over the center $\mathfrak{o}$ of $\mathfrak{K}$. By seeking out, for each straight line segment $\mathfrak{o}\mathfrak{p}$ in $\mathfrak{K}$ emanating from $\mathfrak{o}$, that curve on $\bar{\mathfrak{K}}$ beginning at $\bar{\mathfrak{o}}$ (and ending say at $\bar{\mathfrak{p}}$) of which that segment is the trace, we obtain for each point $\mathfrak{p}$ of $\mathfrak{K}$ a definite point $\bar{\mathfrak{p}}$ of $\bar{\mathfrak{K}}$ lying over it. If we can show that this assignment $\mathfrak{p} \to \bar{\mathfrak{p}}$ is domain-continuously invertible, then the simple connectivity of $\mathfrak{K}$ is thereby proven; for then all curves on $\bar{\mathfrak{K}}$ emanating from $\bar{\mathfrak{o}}$ whose trace lines run from $\mathfrak{o}$ to $\mathfrak{p}$ must terminate in the \emph{same} point $\bar{\mathfrak{p}}$, and $\bar{\mathfrak{K}}$ is single-sheeted.

To provide that proof, however, we proceed as follows: Let the point $\bar{\mathfrak{q}}$ describe on $\bar{\mathfrak{K}}$, starting from $\bar{\mathfrak{o}}$ and ending at $\bar{\mathfrak{p}}$, that curve of which the straight segment $\sigma = \mathfrak{o}\mathfrak{p}$ is the trace in $\mathfrak{K}$, and let $\mathfrak{q}$ be at each moment the trace point of $\bar{\mathfrak{q}}$. To each point $\mathfrak{q} = \mathfrak{q}_0$ we can assign a circle $k_{\mathfrak{q}_0}$ with center $\mathfrak{q}_0$ such that there exists a domain $\bar{\mathfrak{G}}$ on $\bar{\mathfrak{K}}$ containing $\bar{\mathfrak{q}}_0$ which, by virtue of the assignment ``\emph{point on} $\bar{\mathfrak{K}} \to$ \emph{trace point in} $\mathfrak{K}$'', is mapped one-to-one and domain-continuously onto $k_{\mathfrak{q}_0}$. As long as $\bar{\mathfrak{q}}$ moves on that arc of curve whose trace is the segment $\sigma\mathfrak{q}_0$ cut out from $\sigma$ by $k_{\mathfrak{q}_0}$, $\bar{\mathfrak{q}}$ lies in this domain $\bar{\mathfrak{G}}$. According to the so-called Heine-Borel theorem, which belongs to the foundations of infinitesimal analysis${}^{1)}$, we can choose finitely many points of the segment $\sigma$ (cf.\ Fig.~12)
\[ \mathfrak{q}_0 = \mathfrak{o}, \mathfrak{q}_1, \mathfrak{q}_2, \dots, \mathfrak{q}_{n-1}, \mathfrak{q}_n = \mathfrak{p} \quad (\text{in this order}) \]
such that the associated intervals $\sigma\mathfrak{o}, \sigma\mathfrak{q}_1, \sigma\mathfrak{q}_2, \dots, \sigma\mathfrak{p}$ cover the entire segment $\mathfrak{o}\mathfrak{p}$ in such a way that $\mathfrak{q}_{h+1}$ is always situated in the interior of the circle $k_{\mathfrak{q}_h}\;[h = 0, 1, \dots, n-1]$. Now let $\sigma' = \mathfrak{o}\mathfrak{p}'$ be a straight line segment from $\mathfrak{o}$ whose endpoint $\mathfrak{p}'$ lies in $k_{\mathfrak{p}}$ and which otherwise runs so close to $\mathfrak{o}\mathfrak{p}$ that it contains a point $\mathfrak{q}_1'$ lying simultaneously in the interior of $k\mathfrak{o}$ and $k_{\mathfrak{q}_1}$, a point $\mathfrak{q}_2'$ that

\rmfigure{figures/fig-12.png}{Fig.~12.}{Simple connectivity of the circular disc}

\textbf{1)} Cf.\ e.\,g.\ Lebesgue, \emph{Leçons sur l'intégration}, Paris 1904, pp.\ 104---105.

\origpage{49}
lies simultaneously in the interior of $k_{\mathfrak{q}_1}$ and $k_{\mathfrak{q}_2}$, and so on, finally a point $\mathfrak{q}_{n-1}'$ belonging simultaneously to the interior of $k_{\mathfrak{q}_{n-2}}, k_{\mathfrak{q}_{n-1}}$. We draw the straight segments $\mathfrak{q}_1\mathfrak{q}_1', \mathfrak{q}_2\mathfrak{q}_2', \dots, \mathfrak{q}_{n-1}\mathfrak{q}_{n-1}', \mathfrak{p}\mathfrak{p}'$. That curve on $\bar{\mathfrak{K}}$ emanating from $\bar{\mathfrak{o}}$ whose trace is the triangle $\mathfrak{o}\mathfrak{q}_1\mathfrak{q}_1'\mathfrak{o}$ is closed, since this triangle lies entirely in $k\mathfrak{o}$. That curve emanating from $\bar{\mathfrak{q}}_1$ whose trace is the quadrilateral $\mathfrak{q}_1\mathfrak{q}_2\mathfrak{q}_2'\mathfrak{q}_1'$ closes, because the quadrilateral lies entirely in $k_{\mathfrak{q}_1}$. From these two facts it follows that the curve emanating from $\bar{\mathfrak{o}}$ whose trace describes the triangle $\mathfrak{o}\mathfrak{q}_2\mathfrak{q}_2'\mathfrak{o}$ closes. Continuing in the same manner, we arrive at the result that the curve beginning at $\bar{\mathfrak{o}}$ whose trace is the triangle $\mathfrak{o}\mathfrak{p}\mathfrak{p}'\mathfrak{o}$ is closed; that therefore the curve beginning at $\bar{\mathfrak{p}}$ whose trace is the segment $\mathfrak{p}\mathfrak{p}'$ terminates in the \emph{same} point $\bar{\mathfrak{p}}'$ as the line emanating from $\bar{\mathfrak{o}}$ whose trace point describes the segment $\sigma' = \mathfrak{o}\mathfrak{p}'$. The desired proof is thereby accomplished.

It follows from this: \emph{If $\bar{\mathfrak{F}}$ is an unramified unbounded covering surface over an arbitrary base surface $\mathfrak{F}$, $\mathfrak{p}_0$ is a point on $\mathfrak{F}$ with neighborhood $\mathfrak{U}$, and $\bar{\mathfrak{p}}_0$ is a point of $\bar{\mathfrak{F}}$ lying over $\mathfrak{p}_0$, then there exists a unique domain $\bar{\mathfrak{U}}$ on $\bar{\mathfrak{F}}$ containing $\bar{\mathfrak{p}}_0$ that appears mapped one-to-one and domain-continuously onto $\mathfrak{U}$ by virtue of the assignment ``point of $\bar{\mathfrak{F}} \to$ trace point on $\mathfrak{F}$''.} In quite an analogous way one recognizes that if $\Delta$ is an elementary triangle containing the point $\mathfrak{p}_0$ of a definite triangulation $\zeta$ of $\mathfrak{F}$, there lies over $\Delta$ a point set $\bar{\Delta}$ containing $\bar{\mathfrak{p}}_0$ which is mapped one-to-one and continuously onto $\Delta$ by that assignment, and which transforms into a triangle on $\bar{\mathfrak{F}}$ by assigning to each point of $\bar{\Delta}$ the same coordinate ratio $\xi_1 : \xi_2 : \xi_3$ that corresponds to the trace point in $\Delta$. From this it emerges: \emph{If one constructs for each elementary triangle $\Delta$ of $\zeta$ all the triangles $\bar{\Delta}$ lying over it on $\bar{\mathfrak{F}}$, one obtains a triangulation $\bar{\zeta}$ of $\bar{\mathfrak{F}}$.}

Furthermore we can conclude: \emph{If for two curves $\gamma', \gamma''$ on $\mathfrak{F}$ connecting the same points $\mathfrak{p}_0$, $\mathfrak{p}_1$ with each other, one determines those two curves on $\bar{\mathfrak{F}}$ emanating from a point $\bar{\mathfrak{p}}_0$ over $\mathfrak{p}_0$ of which $\gamma', \gamma''$ are the trace lines, both lead to the same endpoint $\bar{\mathfrak{p}}_1$ over $\mathfrak{p}_1$ if $\gamma', \gamma''$ run sufficiently close to each other.} The latter condition is formulated more precisely as follows: $\gamma'$ and $\gamma''$ shall permit partition into finitely many consecutive partial arcs
\[ \gamma_1', \gamma_2', \dots, \gamma_n'; \quad \text{resp.} \quad \gamma_1'', \gamma_2'', \dots, \gamma_n'' \]
such that $\gamma_h'$ and $\gamma_h''$ always belong entirely to the neighborhood of a suitable point $\mathfrak{q}_h$ on $\mathfrak{F}$. In particular, if $\gamma'$ is given, $\gamma''$ can be constructed as a polygonal path on $\mathfrak{F}$ (after $\mathfrak{F}$ is triangulated in a definite way) in such a way that the curves $\gamma'$ and $\gamma''$ lead to the same endpoint on \emph{every} unramified,

\origpage{50}
unbounded covering surface $\bar{\mathfrak{F}}$, provided only that they both emanate from the same initial point on $\bar{\mathfrak{F}}$.

An unramified, unbounded covering surface $\bar{\mathfrak{F}}$ over $\mathfrak{F}$ is customarily designated as \textbf{regular} if it never occurs that of two curves on $\bar{\mathfrak{F}}$ having the same trace line in $\mathfrak{F}$, one is closed and the other unclosed. If $\bar{\mathfrak{F}}$ is regular and $\bar{\mathfrak{p}}_0$, $\bar{\mathfrak{p}}_0'$ are any two points on $\bar{\mathfrak{F}}$ that \textbf{``cover each other''} (i.\,e.\ have the same trace point in $\mathfrak{F}$), then there exists a unique one-to-one and domain-continuously invertible mapping of $\bar{\mathfrak{F}}$ onto itself in which each point $\bar{\mathfrak{p}}$ passes into a point $\bar{\mathfrak{p}}'$ covering it, and in particular $\bar{\mathfrak{p}}_0$ into $\bar{\mathfrak{p}}_0'$. These \textbf{``deck transformations of $\bar{\mathfrak{F}}$ into itself''} form a \emph{group} $\varGamma$, and in $\varGamma$, conceived as an abstract group, the relation of $\bar{\mathfrak{F}}$ to $\mathfrak{F}$, in so far as it possesses analysis situs character, finds pure and complete expression. To demonstrate the existence of the deck transformation
\[ T\colon \bar{\mathfrak{p}} \to \bar{\mathfrak{p}}' \]
connect $\bar{\mathfrak{p}}_0$ with $\bar{\mathfrak{p}}$ on $\bar{\mathfrak{F}}$ by a curve $\bar{\gamma}$ and draw that curve $\bar{\gamma}'$ on $\bar{\mathfrak{F}}$ emanating from $\bar{\mathfrak{p}}_0'$ whose trace line on $\mathfrak{F}$ agrees with the trace line of $\bar{\gamma}$. Assign the endpoint $\bar{\mathfrak{p}}'$ of $\bar{\gamma}'$ to $\bar{\mathfrak{p}}$. $\bar{\mathfrak{p}}'$ is independent of the choice of the curve $\bar{\gamma}$ connecting $\bar{\mathfrak{p}}_0$ with $\bar{\mathfrak{p}}$ on account of the presupposed regularity of $\bar{\mathfrak{F}}$. For if $\bar{\gamma}$, $\bar{\gamma}_1$ are two such curves, then the curve $\bar{\gamma} - \bar{\gamma}_1$ consisting of $\bar{\gamma}$ and the backwardly traversed line $\bar{\gamma}_1$ is closed. Therefore the curve $(\bar{\gamma} - \bar{\gamma}_1)'$ emanating from $\bar{\mathfrak{p}}_0'$ whose trace line coincides with that of $\bar{\gamma} - \bar{\gamma}_1$ is also \emph{closed}; i.\,e.\ $\bar{\gamma}'$ and $\bar{\gamma}_1'$ lead to the same endpoint $\bar{\mathfrak{p}}'$. --- It is clear that the representation of a covering surface $\bar{\mathfrak{F}}$ by an abstract group $\varGamma$ in the described manner is possible only for \emph{regular} covering surfaces.

Among all unramified unbounded covering surfaces belonging to a given surface $\mathfrak{F}$, there is one that is the ``strongest''${}^{1)}$; it is characterized by the statement: A curve $\tilde{\gamma}$ on it is closed only if all curves situated on any unramified unbounded covering surfaces of $\mathfrak{F}$ having the same trace curve on $\mathfrak{F}$ as $\tilde{\gamma}$ are closed. By virtue of this property the \textbf{``universal covering surface''}\ednote{Misprint in the original: fi ligature instead of fl ligature (Überlagerungsfiäche instead of Überlagerungsfläche).} $\tilde{\mathfrak{F}}$ is \emph{regular}, and the group of its deck transformations, regarded as an abstract group, is an analysis situs invariant of the base surface $\mathfrak{F}$. In addition, $\tilde{\mathfrak{F}}$ is simply connected. For if $\tilde{\mathfrak{F}}^*$ were a more than single-

\textbf{1)} Poincaré, \emph{Bulletin de la société mathématique de France}, Vol.\ 11, 1883, pp.\ 113---114.

\origpage{51}
sheeted unramified and unbounded covering surface over $\tilde{\mathfrak{F}}$, one could draw on $\tilde{\mathfrak{F}}^*$ an unclosed line whose trace curve on $\tilde{\mathfrak{F}}$ is closed. But since $\tilde{\mathfrak{F}}^*$ is also an unramified unbounded covering surface over $\mathfrak{F}$, this possibility contradicts outright the characteristic property of $\tilde{\mathfrak{F}}$.

The universal covering surface can be defined, for example, as follows: Each curve $\gamma$ emanating from a fixed point $\mathfrak{p}_0$ of $\mathfrak{F}$ defines a ``point of $\tilde{\mathfrak{F}}$'', of which we say that it lies over the endpoint of $\gamma$; two such curves $\gamma, \gamma'$ then define \emph{the same point on $\tilde{\mathfrak{F}}$} if and only if on every unramified unbounded covering surface over $\mathfrak{F}$ two curves emanating from the same point whose trace lines are $\gamma, \gamma'$ always terminate in the same point. --- Let $\gamma_0$ be a curve on $\mathfrak{F}$ from $\mathfrak{p}_0$ to $\mathfrak{p}$ defining the point $\tilde{\mathfrak{p}}$ on $\tilde{\mathfrak{F}}$, and $\mathfrak{U}$ a neighborhood of $\mathfrak{p}$ on $\mathfrak{F}$. If to $\gamma_0$ I attach all possible curves $\gamma$ emanating from $\mathfrak{p}$ and running in $\mathfrak{U}$, I say that the points on $\tilde{\mathfrak{F}}$ defined by all the resulting paths of curves $\gamma_0 + \gamma$ form a ``neighborhood'' $\tilde{\mathfrak{U}}$ of $\tilde{\mathfrak{p}}$. Over each point of $\mathfrak{U}$, because $\mathfrak{U}$ is simply connected, there lies a single point of $\tilde{\mathfrak{U}}$, and consequently our concept of ``neighborhood'' satisfies the requirements to be placed on it (§~4). Every triangulation of $\mathfrak{F}$ carries over immediately to $\tilde{\mathfrak{F}}$, so that we are entitled to designate $\tilde{\mathfrak{F}}$ as a ``surface''${}^{1)}$.

\emph{Every simply connected surface is schlichtartig.} For if on a triangulated surface $\mathfrak{F}_\zeta$ there lies a polygon $\pi$ consisting of edges that does not dissect it, one can always construct over $\mathfrak{F}$ an unramified unbounded, two-sheeted covering surface in the following way${}^{2)}$: To each elementary triangle $\Delta$ of $\zeta$ one assigns two ``triangles lying over $\Delta$'' $\Delta_1, \Delta_2$. If $\Delta', \Delta''$ are two triangles of $\zeta$ connected along an edge not belonging to the polygon $\pi$, then along the corresponding edge $\Delta_1'$ and $\Delta_1''$ shall be connected with each other, and likewise $\Delta_2'$ with $\Delta_2''$. If, however, the common edge of $\Delta', \Delta''$ belongs to the polygon $\pi$, we join together along the corresponding

\textbf{1)} If $\mathfrak{F}$ is closed and accordingly partitioned into finitely many elementary triangles, the construction of $\tilde{\mathfrak{F}}$ is only a matter of determining for each of the finitely many closed chains of triangles on $\mathfrak{F}$ whether it also closes on $\tilde{\mathfrak{F}}$, and this must naturally be decidable by a definite finite procedure. Cf.\ the genetic construction of $\tilde{\mathfrak{F}}$ in \emph{Koebe}, Über die Uniformisierung beliebiger analytischer Kurven II, Crelle's Journal Vol.\ 139, 1911, pp.\ 271---276.

\textbf{2)} The description of the covering surface is carried out according to the method generally described in §~5, Example 9.

\origpage{52}
edge $\Delta_1'$ with $\Delta_2''$ and $\Delta_2'$ with $\Delta_1''$. --- That every straight polygon in the Euclidean plane dissects it is a special case of our theorem, since the simple connectivity of the Euclidean plane can be seen in exactly the same way as the simple connectivity of the interior of the circle.

It is of importance to decide when a domain composed of finitely many triangles is simply connected. On a triangulated surface $\mathfrak{F}_\zeta$ let there be given a domain $\mathfrak{P}$ that consists exclusively of points of the elementary triangles $\Delta_1, \Delta_2, \dots, \Delta_n$ of $\zeta$, but also really contains all interior points of these triangles; an edge of these triangles shall either (apart perhaps from its endpoints) belong \emph{entirely to} $\mathfrak{P}$ [\textbf{interior edge}] or shall belong to $\mathfrak{P}$ with \emph{no point} [\textbf{boundary edge}]; a vertex from which only interior edges emanate shall always lie within $\mathfrak{P}$. Such a domain $\mathfrak{P}$ we call a \textbf{polyhedron}, and indeed an \textbf{open} one if it possesses at least one boundary edge, and a\ednote{Misprint in the original: eine instead of ein.} \textbf{closed} one if no boundary edges are present. $\mathfrak{P}$ can be a closed polyhedron only if $\mathfrak{F}$ is a closed surface and $\mathfrak{P}$ is identical with $\mathfrak{F}$. A simple polygonal path consisting of edges that belongs entirely to $\mathfrak{P}$ up to its two endpoints not lying within $\mathfrak{P}$ shall be called a \textbf{crosscut}. If $\mathfrak{P}$ is open and does not consist merely of a single triangle $\Delta_1$, there exist crosscuts in $\mathfrak{P}$ (e.\,g.\ one or two edges of a triangle $\Delta_i$ of which at least one edge is a boundary edge form such a crosscut). A crosscut dissects $\mathfrak{P}$, as the considerations in §~8 show, either not at all or into two domains, which are then likewise open polyhedra. \emph{If the open polyhedron $\mathfrak{P}$ is simply connected, every crosscut must dissect $\mathfrak{P}$;} this is recognized in exactly the same way as it was inferred above from simple connectivity that every closed polygon dissects. The converse also holds here: \emph{an open polyhedron that is dissected by every crosscut is simply connected. (A closed polyhedron that is dissected by every polygon consisting of edges is simply connected.)}

\rmfigure{figures/fig-13.png}{Fig.~13.}{Crosscuts}

\textbf{1.} Let $\mathfrak{P}$ be an open polyhedron satisfying the hypothesis of the theorem, and let it be dissected by the crosscut $\sigma$ into $\mathfrak{P}_1, \mathfrak{P}_2$. Then each of the two polyhedra $\mathfrak{P}_1, \mathfrak{P}_2$ also satisfies this hypothesis. For an arbitrary crosscut $\sigma_1$ of $\mathfrak{P}_1$ is namely either itself a crosscut $\sigma'$ in $\mathfrak{P}$ or can be completed to such a crosscut by adding one

\origpage{53}
or two polygonal paths that are parts of $\sigma$ (Fig.~13). Consider two triangles bordering each other along an edge belonging to $\sigma_1$, and in the interior of each of these triangles a point $\mathfrak{p}^*$, resp.\ $\mathfrak{p}^{**}$. Since $\mathfrak{P}$ is dissected by $\sigma'$, $\mathfrak{p}^*$ and $\mathfrak{p}^{**}$ cannot be connected within $\mathfrak{P}$ without crossing $\sigma'$ (p.~45). $\mathfrak{p}^*$, $\mathfrak{p}^{**}$ both lie in $\mathfrak{P}_1$, but according to what has been said they cannot be connected within $\mathfrak{P}_1$ without crossing $\sigma_1$; i.\,e.\ $\mathfrak{P}_1$ is dissected by $\sigma_1$.

\textbf{2.} We consider a covering surface $\bar{\mathfrak{P}}$ (unramified, unbounded) over $\mathfrak{P}$. Let $\Delta_1, \Delta_2$ be two triangles bordering each other along an edge $\sigma^0$ of $\sigma$, so that $\Delta_1$ belongs to $\mathfrak{P}_1$ and $\Delta_2$ to $\mathfrak{P}_2$; let $\bar{\Delta}_1, \bar{\Delta}_2$ be two triangles on $\bar{\mathfrak{P}}$ lying over $\Delta_1$, resp.\ $\Delta_2$, \emph{with a common edge}. If we now assume that the theorem to be proven has already been proven for polyhedra consisting of fewer triangles than $\mathfrak{P}$, we conclude that $\mathfrak{P}_1$ and $\mathfrak{P}_2$ are simply connected. Consequently we can assign to each triangle $\Delta$ of $\mathfrak{P}$ an overlying $\bar{\Delta}$ of $\bar{\mathfrak{P}}$ in such a way that $\bar{\Delta}_1$ appears assigned to the triangle $\Delta_1$, $\bar{\Delta}_2$ to the triangle $\Delta_2$, and to any two triangles $\Delta$ whose common edge is an interior edge of $\mathfrak{P}_1$ or $\mathfrak{P}_2$ there always correspond two triangles $\bar{\Delta}$ that likewise have an edge in common. If $\mathfrak{p}_0$ is an endpoint of $\sigma^0$ without being an endpoint of the crosscut $\sigma$, then over the star $\mathfrak{S}$ of elementary triangles of $\mathfrak{F}_\zeta$ grouping themselves around $\mathfrak{p}_0$ ($\mathfrak{S}$ belongs entirely to $\mathfrak{P}$) there lies a cycle of triangles $\bar{\mathfrak{S}}$ on $\bar{\mathfrak{P}}$ that contains $\bar{\Delta}_1, \bar{\Delta}_2$ and over each triangle $\Delta$ of $\mathfrak{S}$ (on account of the unramifiedness of $\bar{\mathfrak{P}}$) contains a unique triangle $\bar{\Delta}$, and here $\bar{\Delta}$ can be no other triangle than precisely that which was assigned to $\Delta$ by our assignment (see Figure~14).

\rmfigure{figures/fig-14.png}{Fig.~14.}{Star of elementary triangles around a vertex}

If $\sigma'$ is the edge of $\sigma$ following upon $\sigma^0$ at the vertex $\mathfrak{p}_0$, and $\Delta_1', \Delta_2'$ are the two triangles in $\mathfrak{P}_1$, resp.\ $\mathfrak{P}_2$ with the common edge $\sigma'$, then it follows that the triangles $\bar{\Delta}_1', \bar{\Delta}_2'$ assigned to them by us likewise have an edge in common. By thus concluding from edge to edge, traversing the crosscut $\sigma$ from $\sigma^0$ in both directions, we arrive at the insight that the triangles $\bar{\Delta}$ also cohere with one another across the crosscut $\sigma$. The single-sheetedness of $\bar{\mathfrak{P}}$, and thus the simple connectivity of $\mathfrak{P}$, is thereby proven. For for a polyhedron consisting of \emph{a single} triangle, the correctness of our theorem cannot be doubted.

The \emph{significance of covering surfaces for complex function theory} is evident already from the fact that every unramified, unbounded covering surface $\bar{\mathfrak{F}}$ belonging to a \emph{Riemann surface} $\mathfrak{F}$ is itself immediately

\origpage{54}
a Riemann surface. Namely, by assigning to all points of the covering surface lying over a point $\mathfrak{p}$ of $\mathfrak{F}$ the same function value as to the point $\mathfrak{p}$, one obtains from each local uniformizing variable $t(\mathfrak{p})$ for a point $\mathfrak{p}_0$ of $\mathfrak{F}$ a definite local uniformizing variable $t$ for each of the points of $\bar{\mathfrak{F}}$ lying over $\mathfrak{p}_0$, and obtains for each function on $\mathfrak{F}$ free of essential singularities a function of the same character on $\bar{\mathfrak{F}}$. But every function on an arbitrary unramified, unbounded covering surface $\bar{\mathfrak{F}}$ appears again as a single-valued function on the \emph{universal covering surface} $\tilde{\mathfrak{F}}$, so that considering the latter is suited to replace to a certain degree that of all other less comprehensive covering surfaces. Among the functions on $\tilde{\mathfrak{F}}$, the functions $f$ on the base surface $\mathfrak{F}$ are characterized by the property that they behave invariantly under the group of deck transformations of $\tilde{\mathfrak{F}}$; i.\,e.\ they satisfy the identity
\[ f(\tilde{\mathfrak{p}}) = f(\tilde{\mathfrak{p}}\,T), \]
if the assignment $T\colon \tilde{\mathfrak{p}} \to \tilde{\mathfrak{p}}\,T$ is any one of these deck transformations.

The function-theoretic exploitation of the simple connectivity of a surface rests on the following \emph{Monodromy Theorem}:

\emph{If $\mathfrak{F}$ is a simply connected surface,}
\[ z = a_{-m} t^{-m} + \cdot\cdot + a_0 + a_1 t + a_2 t^2 + \dots \quad [t = \text{local uniformizing variable for } \mathfrak{p}_0], \]
\emph{is a function element belonging to a point $\mathfrak{p}_0$ of $\mathfrak{F}$ as center, and if in the analytic continuation of $z$ along arbitrary paths on $\mathfrak{F}$ one never encounters other critical points than ordinary poles, then by this continuation one obtains a single-valued regular-analytic function on $\mathfrak{F}$, except for poles.}

Proof: In a manner analogous to how the universal covering surface was defined, one can explain an unramified, unbounded covering surface over $\mathfrak{F}$ that possesses the property that a curve on it whose trace line on $\mathfrak{F}$ runs from $\mathfrak{p}_0$ back to $\mathfrak{p}_0$ is closed if and only if continuation of the function element $z$ along the trace line leads back to the initial element. On account of the presupposed simple connectivity of $\mathfrak{F}$, this covering surface must be single-sheeted. The proof of the monodromy theorem is thereby already accomplished.

A special case of the monodromy theorem is \emph{Cauchy's Integral Theorem}. In order to formulate it generally, we must speak of ``differentials'' on a Riemann surface. While a ``function'' is characterized by having a definite value at each place of its domain of

\origpage{55}
definition, a \textbf{differential} $dz$ at a place $\mathfrak{p}$ is not endowed with a definite value $(dz)_t^\mathfrak{p}$ in itself, but only in relation to the differential $dt$ of each local uniformizing variable $t$ belonging to $\mathfrak{p}$; if $t, \tau$ are two local uniformizing variables belonging to the same place $\mathfrak{p}$, then one must always have
\[ (dz)_t^\mathfrak{p} = (dz)_\tau^\mathfrak{p} \cdot \left(\frac{d\tau}{dt}\right)_{t = 0} \]
A function $z$ regular-analytic in the domain $\mathfrak{G}$ possesses a differential $dz$ for which
\[ (dz)_t^\mathfrak{p} = \left(\frac{dz}{dt}\right)_{t = 0} \]
when $\mathfrak{p}$ is any point of $\mathfrak{G}$, and $t$ is any local uniformizing variable belonging to $\mathfrak{p}$. A harmonic function $u$ also gives occasion to a differential $dw$ --- according to the formula
\[ (dw)_t^\mathfrak{p} = \left(\frac{\partial u}{\partial x} - i \frac{\partial u}{\partial y}\right)_{t = 0} \quad [t = x + iy]. \]

By multiplication of a differential $dz$ with a function $f$ there arises a new differential $dZ$:
\[ (dZ)_t^\mathfrak{p} = f(\mathfrak{p}) \cdot (dz)_t^\mathfrak{p}. \]
If at a place $\mathfrak{p}\colon (dz)_t^\mathfrak{p} = 0$ for a local uniformizing variable $t$ belonging to $\mathfrak{p}$, this is the case for \emph{every} local uniformizing variable for $\mathfrak{p}$; $\mathfrak{p}$ is then a \textbf{zero} of $dz$. If $dZ, dz$ are two differentials defined in the same domain $\mathfrak{G}$ and $dz$ nowhere possesses a zero, then $\frac{dZ}{dz} = f$ is a function in $\mathfrak{G}$; for
\[ \frac{(dZ)_t^\mathfrak{p}}{(dz)_t^\mathfrak{p}} = f(\mathfrak{p}) \]
is independent of the choice of the local uniformizing variable $t$.

If there exists a neighborhood of the point $\mathfrak{p}_0$ (with local uniformizing variable $t$) in which $dz$ and $dt$ exist, $dt \neq 0$, and $\frac{dz}{dt}$ is regular-analytic, then $dz$ is called \textbf{regular-analytic} at the place $\mathfrak{p}_0$. If $\frac{dz}{dt}$ has a zero of $m^{\text{th}}$ order for $t = 0$, we also say of the differential $dz$ that it has a \textbf{zero of $m^{\text{th}}$ order} in $\mathfrak{p}_0$ (or is of order $m$); if $\frac{dz}{dt}$ has a pole of $n^{\text{th}}$ order ($dz$ is then defined in the neighborhood of $\mathfrak{p}_0$ with exclusion of this point itself), we say that $dz$ has a \textbf{pole of $n^{\text{th}}$ order} in $\mathfrak{p}_0$ (or is of order $-n$). These order numbers are independent of the choice of the local uniformizing variable $t$

\origpage{56}
. The same holds for the \textbf{residue}${}^{1)}$ of the differential $dz$ at the place $\mathfrak{p}_0$, i.\,e.\ the coefficient $A_{-1}$ of the expansion
\[ \frac{dz}{dt} = A_{-n} t^{-n} + \dots + A_{-1} t^{-1} + A_0 + A_1 t + A_2 t^2 + \dots \]
The differential of a function regular except for poles is itself, apart from poles, regular-analytic and nowhere has a residue different from $0$. The converse of this theorem, in so far as it is correct, forms the content of Cauchy's integral theorem in its general formulation.

\emph{If $dz$ is a differential regular except for poles in a simply connected domain $\mathfrak{G}$ that nowhere possesses a residue $\neq 0$, then there exists a single-valued function $z$, regular apart from poles, whose differential coincides with the given $dz$ in all of $\mathfrak{G}$.}

Proof: Let $\mathfrak{p}_0$ be a place in $\mathfrak{G}$ at which $dz$ behaves regularly, and with reference to a local uniformizing variable $t$ belonging to $\mathfrak{p}_0$ let the expansion
\[ \frac{dz}{dt} = A_0 + A_1 t + A_2 t^2 + \dots. \]
hold for sufficiently small $t$. Then one forms the function element
\[ z = A_0 t + A_1 \frac{t^2}{2} + A_2 \frac{t^3}{3} + \dots. \]
This element permits an analytic continuation on every curve $\gamma$ emanating from $\mathfrak{p}_0$ and running in $\mathfrak{G}$, in which one encounters no other critical points than poles (with the help of this analytic continuation we define the integral $\displaystyle\int\limits_\gamma dz$), and thus by the monodromy theorem one obtains a single-valued function $z$ in $\mathfrak{G}$ of the desired quality.

\subsection*{§ 10. One-Sidedness and Two-Sidedness of Surfaces. The Residue Theorem.}

Let there be given in the Euclidean plane a closed curve $\mathfrak{C}$ with a definite sense of traversal, furthermore a point $O$ not situated on $\mathfrak{C}$, and in the pencil of rays through $O$ (more briefly: at $O$) a definite sense of rotation $\mathfrak{s}$. If, while the variable point $P$ traverses the curve $\mathfrak{C}$ once in the prescribed sense, we trace the continuous variation of the angle $\varphi$ that the segment $OP$ forms with a fixed ray through $O$, then the difference
``value of $\varphi$ at the end --- value of $\varphi$ at the beginning of the traversal''
will be an integral multiple $2n\pi$ of $2\pi$; the integer $n$ is called the \textbf{order} of $O$ with respect to $\mathfrak{C}$, in symbols

\textbf{1)} One must insist by all means that the residue is something that belongs to a \emph{differential}, not to a \emph{function}.

\origpage{57}
\[ n = \underset{\mathfrak{C}}{\operatorname{ord}}(O). \]
This number $n$ depends, in addition to the point $O$ and the curve $\mathfrak{C}$ traversed in a definite manner, also on the sense of rotation $\mathfrak{s}$ at $O$; if we replace this by the opposite sense, $n$ changes its sign.

The possibility of establishing a sense of rotation on an arbitrary surface rests on the following fundamental theorem:

\emph{If a domain $\mathfrak{G}$ of the Euclidean plane is mapped one-to-one and domain-continuously onto an identical such domain $\mathfrak{G}'$, where the point $O'$ corresponds to the point $O$ of $\mathfrak{G}$; and if furthermore a definite sense of rotation $\mathfrak{s}$ is given at $O$, then one can establish a sense of rotation $\mathfrak{s}'$ at $O'$ in such a way that the order of the point $O$ with respect to every closed curve running in $\mathfrak{G}$ not passing through $O$ agrees with the order of $O'$ with respect to the image curve.}

For the proof we use a fixed circle $\mathfrak{k}$ in $\mathfrak{G}$ with center $O$ ($x, y$ denote Cartesian coordinates):
\[ x = a \cos 2\pi\lambda, \quad y = a \sin 2\pi\lambda \quad [0 \leqq \lambda \leqq 1]; \]
with respect to $\mathfrak{k}$, $O$ possesses the order $1$. We decide at $O'$ initially arbitrarily upon one of the two possible senses of rotation, and then call $n_0$ the order of $O'$ with respect to the image curve $\mathfrak{k}'$ of $\mathfrak{k}$. The proof divides into two parts:

\textbf{1.} If $\mathfrak{C}$ is an arbitrary closed curve in $\mathfrak{G}$ not passing through $O$, and $n = \underset{\mathfrak{C}}{\operatorname{ord}}(O)$, then
\[ \underset{\mathfrak{C}'}{\operatorname{ord}}(O') = n \cdot n_0. \]

\textbf{2.} $n_0$ is $= +1$ or $-1$.

\emph{Proof of 1.:} Let
\[ x = x(\lambda), \quad y = y(\lambda) \quad [0 \leqq \lambda \leqq 1] \]
be the curve $\mathfrak{C}$; we divide it into partial arcs
\[ 0 \leqq \lambda \leqq \lambda_1, \quad \lambda_1 \leqq \lambda \leqq \lambda_2, \dots, \quad \lambda_{r-1} \leqq \lambda \leqq 1, \]
in such a way that on each of these partial arcs the variations of the azimuth $\varphi$ remain less than $\frac{\pi}{2}$. Let $P_0, P_1, \dots, P_{r-1}$ be the curve points corresponding to the values $0, \lambda_1, \dots, \lambda_{r-1}$. If we thus trace the continuous change of the angle $\varphi = \varphi(\lambda)$ that $OP$ encloses with a fixed ray passing through $O$ while $P$ traverses the $h^{\text{th}}$ arc $\mathfrak{C}_h\colon \lambda_{h-1} \leqq \lambda \leqq \lambda_h$, there never occur two $\varphi$-values whose difference is in absolute value $\geqq \frac{\pi}{2}$. To $\mathfrak{C}_h$ we assign the circular arc $\mathfrak{k}_h\colon Q_{h-1} Q_h$ cut out from $\mathfrak{k}$ by the rays $OP_{h-1}$, $OP_h$ and represent it with the help

\origpage{58}
of a parameter $\lambda$ in such a way: $x = x_h(\lambda)$, $y = y_h(\lambda)$, that when $\lambda$ increases monotonically from $\lambda_{h-1}$ to $\lambda_h$, the assigned point $(xy)$ traverses that circular arc monotonically from $Q_{h-1}$ to $Q_h$. The connecting line segment of a point $P$ on $\mathfrak{C}_h$ with a point $Q$ on $\mathfrak{k}_h$ never contains the point $O$; for then the rays $OP, OQ$ would enclose the angle $\pi$ with each other, which is impossible.

\rmfigure{figures/fig-15.png}{Fig.~15.}{Arc Ch and circular arc kh}

\[ \mathfrak{K} = \mathfrak{k}_1 + \mathfrak{k}_2 + \dots + \mathfrak{k}_r \]
is a closed curve that traverses the circle $\mathfrak{k}$ piecewise monotonically and in total $n$ times. The image curve $\mathfrak{K}'$ of $\mathfrak{K}$ accordingly traverses $\mathfrak{k}'$ piecewise monotonically and in total $n$ times; consequently $O'$ has with respect to $\mathfrak{K}'$ the order $n \cdot n_0$.

If $\mu$ is a number $\geqq 0$ and $\leqq 1$, the equations
\[ \begin{aligned} x &= \mu x(\lambda) + (1 - \mu) x_h(\lambda) \\ y &= \mu y(\lambda) + (1 - \mu) y_h(\lambda) \end{aligned} \quad (\lambda_{h-1} \leqq \lambda \leqq \lambda_h) \quad h = 1, 2, \dots, r \]
represent a closed curve $\mathfrak{C}_\mu$ that does not pass through $O$, whose image curve $\mathfrak{C}_\mu'$ therefore does not pass through $O'$. If we denote by $n_\mu'$ the order of $O'$ with respect to $\mathfrak{C}_\mu'$, then in consequence, and because $\mathfrak{C}_\mu'$ depends continuously on $\mu$, $n_\mu'$ varies continuously with $\mu$, but must always be an integer and is accordingly constant for $0 \leqq \mu \leqq 1$. For $\mu = 1$, $\mathfrak{C}_\mu = \mathfrak{C}$, and for $\mu = 0$ it agrees with the curve $\mathfrak{K}$. Therefore we must have
\[ \underset{\mathfrak{C}'}{\operatorname{ord}}(O') = \underset{\mathfrak{K}'}{\operatorname{ord}}(O') = n \cdot n_0 \]

\emph{Proof of 2.:} If $\mathfrak{E}$ is a closed set in the Euclidean plane and $P$ is any point not belonging to $\mathfrak{E}$, we call the collection of points connectable with $P$ by continuous curves not meeting $\mathfrak{E}$ a \textbf{domain determined by $\mathfrak{E}$}. Two domains determined by $\mathfrak{E}$ are either completely identical or have not a single point in common. The \textbf{boundary} of a domain is formed by all points that, without belonging to the domain itself, are condensation points of points of the domain. The boundary of a domain $\mathfrak{G}$ determined by $\mathfrak{E}$ is a closed subset $\mathfrak{E}^*$ of $\mathfrak{E}$. $\mathfrak{G}$ is then also one of the domains determined by $\mathfrak{E}^*$.

From this it emerges for $\mathfrak{k}'$: The closed curve $\mathfrak{k}'$ determines two domains in the plane; the one, $\mathfrak{I}$, which we call the interior, consists of all image points of the points situated within $\mathfrak{k}$; the other, the ``exterior'', $\mathfrak{A}$, contains all distant points of the plane. As ``distant'' are to count those points of the plane that lie outside a circle around $O'$ containing the entire curve $\mathfrak{k}'$ in its interior.

\origpage{59}
Every point of $\mathfrak{k}'$ belongs both to the boundary of $\mathfrak{I}$ and of $\mathfrak{A}$. A partial arc of $\mathfrak{k}'$ determines only a single domain.

At each point of the plane we take as basis the same sense of rotation for determining its order with respect to $\mathfrak{k}'$. If one then lets a point move continuously in the plane without meeting $\mathfrak{k}'$, its order with respect to $\mathfrak{k}'$ must change continuously and thus remain constant. Both for all points of $\mathfrak{A}$ and for all points of $\mathfrak{I}$ there exists therefore a constant order; this is for $\mathfrak{A}$: $= 0$, since the distant points certainly have order $0$ with respect to $\mathfrak{k}'$. Of the points of $\mathfrak{I}$ we assert that their order is $= \pm 1$.

Through $O'$ we draw a straight line. If we trace this from $O'$ to both sides, we must in both directions meet a point of $\mathfrak{k}'$ for the first time. The two points $E, F$ of $\mathfrak{k}'$ thus obtained divide this curve into two partial arcs. One of them, $\mathfrak{k}'_1$, forms together with the straight line segment $\overrightarrow{FE}$ a closed curve $\mathfrak{C}_1$; the other, $\mathfrak{k}'_2$, with the segment $\overrightarrow{EF}$ a second closed curve $\mathfrak{C}_2$. We cross the segment $EF$ at the point $O'$ by a small segment $O_1 O_2$ perpendicular to $EF$ lying entirely in $\mathfrak{I}$. The order of a point $P$ with respect to $\mathfrak{C}_1$ obviously jumps by $\pm 1$ when $P$, traversing the segment $O_1 O_2$ monotonically, passes $O'${}${}^{1)}$:
\[ \underset{\mathfrak{C}_1}{\operatorname{ord}} (O_1) - \underset{\mathfrak{C}_1}{\operatorname{ord}} (O_2) = \pm 1. \]

\rmfigure{figures/fig-16.png}{Fig.~16.}{Domain division by the curve k'}

Since $\mathfrak{C}_1$, consisting of a partial arc of $\mathfrak{k}'$ not dissecting the plane and a straight line segment, can dissect the plane into at most two domains $\mathfrak{I}_1, \mathfrak{A}_1$, $O_1$ or $O_2$ must belong to that one of these two domains $\mathfrak{A}_1$ in which the distant points of the plane lie; one of the two points $O_1, O_2$, say $O_2$, must therefore have order $0$ with respect to $\mathfrak{C}_1$; then $O_1$ has with respect to $\mathfrak{C}_1$ the order $\pm 1$ and lies in the other domain $\mathfrak{I}_1$ determined by $\mathfrak{C}_1$.

Similarly one sees that $\mathfrak{C}_2$ dissects the plane into two domains $\mathfrak{I}_2, \mathfrak{A}_2$, of which let $\mathfrak{A}_2$ contain the distant points of the plane, and furthermore that one of the following two cases must occur:
\[ 1) \quad \underset{\mathfrak{C}_2}{\operatorname{ord}} (O_1) = 0, \quad \underset{\mathfrak{C}_2}{\operatorname{ord}} (O_2) = \pm 1; \quad O_1 \text{ lies in } \mathfrak{A}_2; \]

\textbf{1)} For the angle subtended by the segment $\overrightarrow{FE}$ as seen from $P$ converges to $+\pi$ or $-\pi$ according as $P$ approaches $O'$ from $O_1$ or from $O_2$; but the total angular change suffered by a ray passing through $P$ whose other endpoint describes the curve $\mathfrak{k}'_1$ depends continuously on $P$, even at the place $P = O'$.

\origpage{60}
\[ 2) \quad \underset{\mathfrak{C}_2}{\operatorname{ord}} (O_1) = \pm 1, \quad \underset{\mathfrak{C}_2}{\operatorname{ord}} (O_2) = 0; \quad O_2 \text{ lies in } \mathfrak{A}_2. \]
In Case 1) there would result, in agreement with our assertion,
\[ \underset{\mathfrak{C}}{\operatorname{ord}} (O_1) = \underset{\mathfrak{C}_1}{\operatorname{ord}} (O_1) + \underset{\mathfrak{C}_2}{\operatorname{ord}} (O_1) = \pm 1. \]
We now show that 2) is impossible.

The points of $\mathfrak{k}'_2$, if we disregard the endpoints $E, F$, lie in $\mathfrak{A}_1$. For in arbitrary proximity to a point $Q$ of $\mathfrak{k}'_2$ there are points of $\mathfrak{A}$, and these are connectable with the distant points of the plane by continuous curves not meeting $\mathfrak{k}'$, thus running entirely in $\mathfrak{A}$, and therefore also not meeting the segment $EF$ and $\mathfrak{k}'_1 + EF = \mathfrak{C}_1$. In arbitrary proximity to $Q$ there are thus found points belonging to $\mathfrak{A}_1$, and since the point $Q$ does not lie on $\mathfrak{C}_1$, $Q$ itself must therefore lie in $\mathfrak{A}_1$. $\mathfrak{I}_1$ consists of all points $J$ connectable with $O_1$ by curves $\mathfrak{c}$ that do not meet $\mathfrak{C}_1$. $\mathfrak{c}$ itself lies entirely in $\mathfrak{I}_1$, and since $\mathfrak{k}'_2$ apart from its endpoints lies in $\mathfrak{A}_1$, $\mathfrak{c}$ also does not meet $\mathfrak{k}'_2$, hence also not $\mathfrak{k}'_2 + EF = \mathfrak{C}_2$. If that second case occurred, $O_1$ would lie simultaneously in $\mathfrak{I}_2$, and, since $\mathfrak{c}$ does not meet the curve $\mathfrak{C}_2$, $J$ would also lie in $\mathfrak{I}_2$. Every point of $\mathfrak{I}_1$ would accordingly be a point of $\mathfrak{I}_2$. Similarly one recognizes the converse; accordingly $\mathfrak{I}_1$ would be identical with $\mathfrak{I}_2$. Apart from the endpoints, $\mathfrak{k}'_2$ does not belong to the boundary of $\mathfrak{I}_1$, nor does $\mathfrak{k}'_1$ belong to the boundary of $\mathfrak{I}_2$. If now $\mathfrak{I}_1 = \mathfrak{I}_2$, the boundary of $\mathfrak{I}_1$ could thereafter consist only of points of the segment $EF$. But that is absurd, since the segment $EF$ does not dissect the plane. The untenability of assumption 2) is proven.

The sense of rotation $\mathfrak{s}'$ established by the fundamental theorem at the beginning of this section shall be called the \textbf{image sense of rotation} of $\mathfrak{s}$ under the mapping $S'$ of the domain $\mathfrak{G}$ onto $\mathfrak{G}'$: $\mathfrak{s}$ \textbf{``passes'' under the mapping} $S'$ \textbf{``over into''} $\mathfrak{s}'$. If $\mathfrak{G}$ has been mapped by a one-to-one and domain-continuous mapping $S'$ onto $\mathfrak{G}'$, and by another such mapping $S''$ onto $\mathfrak{G}''$, then the image sense of rotation $\mathfrak{s}' = \mathfrak{s} S'$ of $\mathfrak{s}$ generated by $S'$ passes into the image sense of rotation $\mathfrak{s}'' = \mathfrak{s} S''$ generated by $S''$ by the mapping $S'^{-1} S''$ that transforms $\mathfrak{G}'$ into $\mathfrak{G}''$.

These facts suffice to recognize the possibility of establishing a sense of rotation $\mathfrak{s}$ at an arbitrary point $\mathfrak{p}_0$ of a given manifold $\mathfrak{F}$. In every one-to-one and domain-continuous mapping $S$ of a neighborhood of $\mathfrak{p}_0$ onto a domain of the Euclidean plane, in which $\mathfrak{p}_0$ passes into $O$, $\mathfrak{s}$ will manifest itself in an image sense of rotation $\mathfrak{s} S$ prevailing at $O$. If one has two such mappings $S$ and $T$, there always exists a neighborhood of $\mathfrak{p}_0$ of which both mappings provide a one-to-one and domain-continuous image in the Euclidean plane. The image senses of rotation $\mathfrak{s} S$, $\mathfrak{s} T$ must be of such a kind that $\mathfrak{s} S$ passes into $\mathfrak{s} T$ by the mapping $S^{-1} T$.

\origpage{61}
(The explicit specification of a sense of rotation $\mathfrak{s}$ can of course take place only ``figuratively'' by giving for some definite mapping $S$ the image sense of rotation $\mathfrak{s} S$ in the Euclidean plane).

If at any two points of the plane $O_1, O_2$ one takes as basis the same sense of rotation $\mathfrak{s}_1 = \mathfrak{s}_2$ for determining the order of these points, this makes itself noticeable in that $O_1, O_2$ certainly always possess the same order with respect to a closed planar curve $\mathfrak{C}$ if $O_1, O_2$ can be connected by a curve not meeting $\mathfrak{C}$. If $O_1, O_2$ lie in a domain $\mathfrak{G}$ that is transferred into $\mathfrak{G}'$ by a one-to-one and domain-continuous mapping $S'$, this remark shows, when we apply it to curves $\mathfrak{C}$ lying in $\mathfrak{G}$, that from $\mathfrak{s}_1 = \mathfrak{s}_2$ there follows
\[ \mathfrak{s}_1 S' = \mathfrak{s}_2 S' \]
Thereafter the following definition arises naturally:

\emph{If at each point of a manifold $\mathfrak{F}$ a sense of rotation $\mathfrak{s}$ is established, it is called \textbf{continuous} at the point $\mathfrak{p}_0$ if there exists a neighborhood of $\mathfrak{p}_0$ such that, if this is mapped in any way one-to-one and domain-continuously onto a planar domain, the image sense of $\mathfrak{s}$ is the same at all points of this domain.}

And only when $\mathfrak{s}$ is everywhere continuous will one say that a \textbf{uniform sense of rotation} is defined on the manifold $\mathfrak{F}$. If such a ``uniform'' specification of a sense of rotation is possible on $\mathfrak{F}$, $\mathfrak{F}$ is called \textbf{two-sided} (examples: Euclidean plane, sphere, torus), otherwise \textbf{one-sided}${}^{1)}$ (examples: projective plane, Möbius strip). A closed curve on $\mathfrak{F}$ can be of two kinds: if I start from a definite sense of rotation $\mathfrak{s}$ at a point of the curve and continue it continuously along the curve, I arrive at the initial point either with the same or with the opposite sense of rotation. On two-sided surfaces there exist only curves of the first kind; on one-sided ones, curves of both the first and the second kind. From this the names are explained. That the sense of rotation reverses along a curve of the second kind manifests itself on a one-sided spatial surface such as the Möbius strip in that, wandering along this curve, I pass from one ``side'' of the surface to the other, so that the surface in truth does not have two separated ``sides'' at all. Over every one-sided surface $\mathfrak{F}$ there exists an unramified, unbounded, two-sheeted, two-sided covering surface on which a curve is closed if and only if its trace curve on $\mathfrak{F}$ is closed and furthermore of the first kind. The two sheets of this covering surface can be regarded as the one and the other ``side'' of $\mathfrak{F}$.

\textbf{1)} This definition of one-sidedness was given by \emph{Klein}, Math. Ann. Vol.~9 (1876), p.~479.

\origpage{62}
We now imagine the surface $\mathfrak{F}$, of which we wish to assume that it is two-sided and thus that on it a uniform sense of rotation is established, present in a definite triangulation $\zeta$. A triangle $\Delta$ of this triangulation with vertices $1, 2, 3$ is mapped onto a planar equilateral triangle $D$ by identifying the coordinate ratios $\xi_1\colon \xi_2\colon \xi_3$ corresponding to the points of $\Delta$ with the homogeneous barycentric coordinates (formed with respect to the vertices of $D$) of the image points in $D$. This mapping $S$ is domain-continuously invertible for the interior of $\Delta$ and $D$. $\mathfrak{s}$ thus generates at all interior points of $D$ a definite image sense of rotation $\mathfrak{s} S$, and indeed the same at all points; but this induces, in the manner evident from the figure, a definite sense of traversal on the boundary of $D$ and therewith a definite cyclic order of vertices or \textbf{indicatrix} of $D$ resp.\ $\Delta$: either
\[ (123) = (231) = (312) \]
or
\[ (213) = (132) = (321). \]

\rmfigure{figures/fig-17.png}{Fig.~17.}{Direction of rotation and indicatrix}

After having thus obtained for each triangle $\Delta$ of the triangulation a definite indicatrix, I consider two such triangles $\Delta_3$ with vertices $123$ and $\Delta_4$ with vertices $124$, which border each other along the edge $12$. I assert: in the traversal of $\Delta_3$ taking place with the prescribed indicatrix, the edge $12$ is traversed in the \emph{opposite} direction to that in the traversal of $\Delta_4$ prescribed by its indicatrix. I express this briefly by saying that the indicatrix of $\Delta_3$ is \textbf{coherent} with that of $\Delta_4$. Draw in the plane two equilateral triangles $D_3 = (123)$, $D_4 = (124)$ meeting at the edge $12$. $\Delta_3$ is mapped continuously onto $D_3$ by the coordinate ratios of its points (let the mapping be called $S_3$) and $\Delta_4$ onto $D_4$ (mapping $S_4$). Here the vertices designated with identical figures correspond to each other. The common edge of $\Delta_3, \Delta_4$ is mapped twice, both by $S_3$ and by $S_4$, onto the common edge of $D_3, D_4$. For this edge $S_4^{-1} S_3$ is a one-to-one continuous mapping in which each of the two ends passes into itself. I choose a point $\mathfrak{p}$ on the edge $12$ (which is neither of the endpoints), a point $\mathfrak{p}_3$ in the interior of $\Delta_3$, and a point $\mathfrak{p}_4$ in the interior of $\Delta_4$:
\[ p_3 = \mathfrak{p}_3 S_3, \quad p_4 = \mathfrak{p}_4 S_4; \quad p^3 = \mathfrak{p} S_3, \quad p^4 = \mathfrak{p} S_4. \]
One can easily map $D_4$ one-to-one and continuously onto itself (mapping $s_4$) in such a way that each segment $p^4 p$ ($p$ on the boundary of $D_4$) emanating from $p^4$ passes into the segment $p^3 p$, but on the edge $12$ $s_4$ agrees with $S_4^{-1} S_3$. By Figure 18 such a mapping is indicated. Then by that mapping $T_3$ which for $\Delta_3$ is: $= S_3$, and for $\Delta_4$: $= S_4 s_4$,

\origpage{63}
the interior of $\Delta_3 + \Delta_4$ is mapped one-to-one and domain-continuously onto the interior of the rhombus $D_3 + D_4$. Analogously let a mapping $T_4$ be constructed that agrees with $S_4$ for $\Delta_4$. $T_3^{-1} T_4$ is a one-to-one continuous mapping of the rhombus $D_3 + D_4$ into itself in which the pencil of rays through the point $p^3$ passes into the pencil of rays through the point $p^4$ and every point on the boundary of the rhombus remains fixed. Draw the image $\mathfrak{s} T_3$ of the sense of rotation $\mathfrak{s}$ prevailing at the points $\mathfrak{p}_3$ and $\mathfrak{p}$. $\mathfrak{s} T_3$ at $p^3$ is then the same sense of rotation as at $p_3$. Let the endpoint of a movable ray through the point $p_3$ or $p^3$ rotating around it with the sense of rotation $\mathfrak{s} T_3$ traverse the boundary of the rhombus say in the direction 13241. By the mapping $T_3^{-1} T_4$, $\mathfrak{s} T_3$ at the point $p^3$ passes into $\mathfrak{s} T_4$ at the point $p^4$. Accordingly, the endpoint of a movable ray through $p^4$ rotating around this point in the sense of rotation $\mathfrak{s} T_4$ obviously describes the rhombus likewise in the direction 13241. The same holds consequently also for a ray through $p_4$ that revolves around this point with the sense of rotation $\mathfrak{s} T_4$. But $\mathfrak{s} T_3 = \mathfrak{s} S_3$ at $p_3$ is decisive for the indicatrix of $D_3$, and $\mathfrak{s} T_4 = \mathfrak{s} S_4$ at $p_4$ decisive for the indicatrix of $D_4$; in the assumed case $(321)$ would be the indicatrix of $D_3$, and $(412)$ that of $D_4$.

\rmfigure{figures/fig-18.png}{Fig.~18.}{Abbildung eines Dreiecks $D_4$ auf sich selbst}

A uniform sense of rotation $\mathfrak{s}$ on the triangulated surface $\mathfrak{F}$ thus induces in each elementary triangle $\Delta$ of $\zeta$ an indicatrix, which for any two adjacent triangles cohere with each other. Our consideration yields, however, immediately also the converse result: If to each triangle $\Delta$ of $\zeta$ an indicatrix is assigned in such a way that they cohere for any two adjacent triangles, then there is thereby established on $\mathfrak{F}$ a uniform sense of rotation $\mathfrak{s}$ that induces in each $\Delta$ the indicatrix\ednote{Misprint in the original: Indikratix instead of Indikatrix.} in question.${}^{1)}$ A two-sided surface $\mathfrak{F}_\zeta$ can therefore be recognized by the fact that it is possible to assign to each $\Delta$ such an indicatrix that any two adjacent triangles possess indicatrices standing in coherence with each other.

\textbf{1)} By this property that one can impart coherent indicatrices to all elementary triangles, \emph{Möbius} (1865; Werke Vol.~II, p.~477 and 482) and also \emph{Brouwer} (Math. Ann. Vol.~71, p.~101) define two-sided surfaces. The proof that this property belongs to a surface independently of the manner of its triangulation is given by \emph{Brouwer}, Math. Ann. Vol.~71, p.~324 (footnote) in a different way from that done here; his proof is also valid for $n$-dimensional manifolds.

\origpage{64}
\emph{Example:} From the scheme of the projective plane given in §~5 its one-sidedness emerges purely combinatorially. The scheme read:
\[
\begin{gathered}
(\mathrm{I}\ \mathrm{II}\ \mathrm{III}) \\
/\qquad\qquad\qquad\quad |\quad\qquad\qquad\qquad \backslash \\
\begin{matrix}
(1\ \mathrm{II}\ \mathrm{III}) \\
/\quad\backslash \\
(1\ \mathrm{II}\ 2) \quad (1\ 3\ \mathrm{III})
\end{matrix}
\qquad
\begin{matrix}
(\mathrm{I}\ 2\ \mathrm{III}) \\
/\quad\backslash \\
(\mathrm{I}\ 2\ 1) \quad (3\ 2\ \mathrm{III})
\end{matrix}
\qquad
\begin{matrix}
(\mathrm{I}\ \mathrm{II}\ 3) \\
/\quad\backslash \\
(\mathrm{I}\ 1\ 3) \quad (2\ \mathrm{II}\ 3)
\end{matrix}
\end{gathered}
\]
In this notation each triangle of the last row stands in mediate coherence to $(\mathrm{I}\ \mathrm{II}\ \mathrm{III})$; as an incoherently written intermediate link there appears in each case a triangle of the second row. If the projective plane were two-sided, then in the scheme any two among the six triangles of the last row in which two identical figures appear would have to stand in coherence; in truth, however, three incoherences occur:
\[
\begin{array}{ccccc}
(1\ \mathrm{II}\ 2) & \qquad\qquad & (1\ 3\ \mathrm{III}) & \qquad\qquad & (3\ 2\ \mathrm{III}) \\
| & & | & & | \\
(\mathrm{I}\ 2\ 1) & & (\mathrm{I}\ 1\ 3) & & (2\ \mathrm{II}\ 3)
\end{array}
\]

[\emph{Addition:} In the general considerations carried out above we utilized the fact that one can map each pair of triangles of a triangulation one-to-one and continuously onto a rhombus (where vertex passes into vertex). In the same way one can map a star of triangles one-to-one and continuously onto a regular polygon of the Euclidean plane in such a way that the interior vertex of the star passes into the center of the polygon, and the boundary vertices of the star into the polygon vertices. If one partitions this polygon somehow into triangles $D$ in such a way that on the periphery of the polygon no new vertices arise, and the center of the polygon falls into the \emph{interior} of one of the triangles $D$, one has thereby altered the original triangulation of the surface in the interior of that star of triangles in such a way that a point which was formerly a vertex now comes to lie in the interior of an elementary triangle.]

If we proceed from a fixed elementary triangle $\Delta_0$ of the triangulated surface $\mathfrak{F}_\zeta$ to which we impart a definite indicatrix, and if $\Delta$ is any other elementary triangle belonging to the partition $\zeta$, one can connect $\Delta_0$ with $\Delta$ by a simple chain of triangles:
\[ \Delta_0, \Delta_1, \Delta_2, \cdots, \Delta_m = \Delta. \tag{K} \]
Impart to $\Delta_1$ the indicatrix coherent with the indicatrix of $\Delta_0$, to $\Delta_2$ that which coheres with that of $\Delta_1$, and so on; thus each triangle of the chain and finally also $\Delta$ itself receives a definite indicatrix. If
\[ \Delta_0' = \Delta_0, \Delta_1', \Delta_2', \cdots, \Delta_n' = \Delta \tag{K'} \]
is another simple chain connecting $\Delta_0$ with $\Delta$, we can

\origpage{65}
also in this chain continue the indicatrix of $\Delta_0$ according to the principle of coherence. If in \emph{every} chain we arrive at the same indicatrix of $\Delta$, and that whatever the terminal triangle $\Delta$ may be, then the surface must be two-sided. If it is one-sided, there will accordingly exist for a certain triangle $\Delta$ two chains $(K)$ and $(K')$ leading to different indicatrices for the terminal triangle $\Delta$. I note those triangles of the chain $(K)$ that also occur in $(K')$, in that order in which they appear in $(K)$:
\[ \Delta_0, \Delta_i, \Delta_k, \cdots, \Delta_m. \tag{13} \]
Let say $\Delta_l$ (perhaps only $\Delta_m$) be the first triangle among these to which in the chain $(K')$ a different indicatrix attaches than in the chain $(K)$, and $\Delta_k$ the triangle preceding $\Delta_l$ in the sequence (13). The piece of the chain $(K)$ connecting $\Delta_k$ with $\Delta_l$ provides, together with the piece of the chain $(K')$ leading from $\Delta_l$ to $\Delta_k$, a (non-self-intersecting) closed chain of triangles. In this the indicatrix reverses when continued according to the principle of coherence. On a one-sided surface the impossibility of an indicatrix satisfying the coherence condition must accordingly be decided already in a simple closed chain of triangles $K_0$.

If one draws in each triangle of such a chain $K_0$ an elementary segment $\sigma$ leading over from a point of the edge where the preceding triangle of the chain $K_0$ abuts to a point of that edge where the next succeeding triangle abuts, and indeed in such a way that these elementary segments $\sigma$ together form a closed polygon $\pi$, then $\pi$ has \emph{no separated banks}: if one imagines $\pi$ as a ditch along which a promenade runs, this promenade does not lead us back to the starting point after a single circuit, but terminates on the opposite bank of $\pi$.

\rmfigure{figures/fig-19.png}{Fig.~19.}{Dreieckskette, in der sich die Indikatrix umkehrt, und Polygon pi ohne getrennte Ufer}

This will become clear from the figure, in which the triangles of the chain $K_0$ are represented by Euclidean triangles and the heavy segments form the polygon $\pi$. The property of $\pi$ of possessing no \textbf{separated banks} is

\origpage{66}
strictly formulated as follows: Not a single one of those domains on $\mathfrak{F}$ containing $\pi$ is dissected by $\pi$. From this it emerges that schlichtartig surfaces are always two-sided.

Conversely: \emph{on a two-sided surface it can never happen that a polygon $\pi$ possesses no separated banks.} Replacing $\zeta$, if necessary, by a subdivision, we can assume that $\pi$ consists entirely of edges. Let the vertices of $\pi$ be, in this order,
\[ 1, 2, 3, 4, \cdots, n, n + 1 = 1, \cdots \]
Since the surface is to be two-sided, a definite indicatrix belongs to each triangle, and these cohere with one another. Among the two triangles with the edge $12$ there is one in whose traversal prescribed by the indicatrix the edge $\overrightarrow{12}$ is traversed in the direction of the arrow.

\rmfigure{figures/fig-20.png}{Fig.~20.}{Linkes Ufer eines Polygons}

Let it be called $\Delta^0$. Starting from $\Delta^0$, I form in the manner illustrated by the figure the triangles $\Delta^1, \Delta^2, \Delta^3, \cdots$, all bordering on $\pi$, of which every two consecutive ones have an edge in common not belonging to $\pi$. Those triangles of this chain possessing an edge $h, h + 1$ lying on $\pi$ carry such an indicatrix that in the traversal determined by it the edge $\overrightarrow{h, h + 1}$ is described in the direction of the arrow. \{One recognizes this in the figure e.\,g.\ for $\Delta^3$ by continuing the indicatrix of $\Delta^0$ by the condition of coherence in the chain $\Delta^0 \Delta^1 \Delta^2 \Delta^3$.\} After finitely many steps I come for the first time again to a triangle having $12$ as an edge. Its indicatrix must read: $(1\ 2\ *)$; it can therefore only be the triangle $\Delta^0$ and not the other triangle $\Delta_1^0$ with the edge $12$, since this possesses the indicatrix $(2\ 1\ *)$. The chain of triangles thus obtained from $\Delta^0$ forms the one bank of $\pi$, and the chain to be derived in a similar manner from the other triangle $\Delta_1^0$ with the edge $12$ forms the other bank.

\emph{Every Riemann surface is two-sided.}

Let $\mathfrak{p}_0$ be a point of a Riemann surface, and $t', t''$ two arbitrary local uniformizing variables for $\mathfrak{p}_0$ mapping a certain neighborhood of $\mathfrak{p}_0$ onto planar domains $\mathfrak{G}'$, resp.\ $\mathfrak{G}''$, where the point $t' = 0$, resp.\ $t'' = 0$ corresponds to the point $\mathfrak{p}_0$. Let these mappings be called $T', T''$. The mapping $T'^{-1} T''$ of $\mathfrak{G}'$ onto $\mathfrak{G}''$ is mediated by a formula
\[ t'' = \text{regular function of } (t') \text{ in } \mathfrak{G}'. \]
Let $\mathfrak{s}'$ be that sense of rotation at the origin of the $t'$-plane which carries the positive real axis through $90^\circ$ into the positive imaginary axis. Correspondingly let the sense of rotation $\mathfrak{s}''$ at the origin of the $t''$-plane be defined. I assert: by the mapping $T'^{-1} T''$, $\mathfrak{s}'$ passes into $\mathfrak{s}''$. This circumstance permits regarding $\mathfrak{s}'$ and $\mathfrak{s}''$ as the images generated by the

\origpage{67}
mappings $T', T''$ of a definite sense of rotation $\mathfrak{s}$ at $\mathfrak{p}_0$, which is thereby established in a manner independent of the choice of the local uniformizing variables. Since the sense of rotation $\mathfrak{s}$ thus determined for each point $\mathfrak{p}_0$ of the Riemann surface is also everywhere continuous --- there is then nothing more to prove about that ---, the two-sidedness of Riemann surfaces is established.

Around $t' = 0$ in $\mathfrak{G}'$ I describe a small circle $\mathfrak{k}'$:
\[ t' = a e^{i\varphi} \quad (0 \leqq \varphi \leqq 2\pi) \]
[$a > 0$ is the constant radius]. Let the image of $\mathfrak{k}'$ provided by $T'^{-1} T''$ in $\mathfrak{G}''$ be called $\mathfrak{k}''$. According to the definition of the concept ``order'',
\[ 2\pi i \cdot \underset{\mathfrak{k}''}{\text{ord.}} (t'' = 0) = \int_{\mathfrak{k}''} \frac{dt''}{t''} \]
There holds, however,
\[ \frac{dt''}{t''} = \left(\frac{1}{t'} + a_0 + a_1 t' + a_2 t'^2 + \dots \right) dt', \]
where the power series converges in a circle of the $t'$-plane containing $\mathfrak{k}'$ in its interior. Thus it follows
\[ \int_{\mathfrak{k}''} \frac{dt''}{t''} = \int_{\mathfrak{k}'} \frac{dt'}{t'} = 2\pi i, \]
and therewith it is proven that $\underset{\mathfrak{k}''}{\text{ord.}} (t'' = 0) = + 1$ and not $= - 1$, or that $\mathfrak{s}'$ passes into $\mathfrak{s}''$ by the mapping $T'^{-1} T''$.

If the Riemann surface $\mathfrak{F}$ is triangulated in any way, the ``positive'' sense of rotation $\mathfrak{s}$ just established induces a ``positive'' indicatrix in each triangle. If $dz$ is a differential regular at all points of an elementary triangle of the triangulated surface including the boundary, then the integral of $dz$ taken around the boundary of the triangle with positive indicatrix is equal to zero. If on the other hand $dz$ is regular in the entire triangle including the boundary, apart from a pole situated in the interior, then this integral is $= 2\pi i$ times the residue of $dz$ at that pole.

If $dz$ is regular except for poles on a closed Riemann surface, one can partition $\mathfrak{F}$ into elementary triangles in such a way that the poles fall into the interior of the triangles and never two or more poles into the interior of the same triangle (p.~64). If one integrates $\displaystyle\int dz$ around all triangles with positive indicatrix and adds, one obtains $2\pi i$ times the sum of all residues of $dz$. But if one traverses all triangle perimeters with positive indicatrix, then on account of coherence each edge is traversed twice, but in the opposite sense. Consequently that sum of integrals must on the other hand be $= 0$, and we obtain the theorem:

\emph{The sum of the residues of a differential regular except for poles on a closed Riemann surface is $0$.}

\origpage{68}

\subsection*{§ 11. Integral Functions. Genus. Canonical Dissection.}

The process of integration of differentials on a Riemann surface can be replaced for arbitrary surfaces by the following generalization.

A \emph{curve function} $F$ is defined on a surface $\mathfrak{F}$ if to each curve $\gamma$ on $\mathfrak{F}$ a number $F(\gamma)$ is assigned; $F(\gamma)$ is then designated as the \emph{value} of $F$ for the curve $\gamma$. If for two curves $\gamma', \gamma''$ the terminal point of $\gamma'$ coincides with the initial point of $\gamma''$, one can compose from them a single curve $\gamma = \gamma' + \gamma''$; if always
\[ F(\gamma' + \gamma'') = F(\gamma') + F(\gamma''), \]
then the curve function is called \emph{linear}. A linear curve function has for a closed curve a value that does not change when one shifts the initial point of the closed curve along it. If $F$ has the value $0$ for every closed curve, we write $F \sim 0$ ($F$ \emph{homologous} to $0$). There then exists a ``point function'' $f(\mathfrak{p})$ on the surface such that for every curve
\[ F(\gamma) = f(\mathfrak{p}_2) - f(\mathfrak{p}_1) \]
holds, where $\mathfrak{p}_1, \mathfrak{p}_2$ denote the initial and terminal points of $\gamma$. We consider only such linear curve functions which are everywhere $\sim 0$ ``in the small''; these shall be called \emph{``integral functions''}. There shall thus exist for each point of the surface a neighborhood of such a kind that for every closed curve running in this neighborhood $\gamma_0\colon F(\gamma_0) = 0$. For an integral function $F$ there holds always
\[ F(-\gamma) = - F(\gamma), \]
where $-\gamma$ denotes the path $\gamma$ traversed in the opposite sense:
\[ \gamma\colon \mathfrak{p} = \mathfrak{p}(\lambda) \; [0 \leqq \lambda \leqq 1]; \quad -\gamma\colon \mathfrak{p} = \mathfrak{p}(1 - \lambda) \; [0 \leqq \lambda \leqq 1]. \]
On a simply connected surface, in particular therefore on the universal covering surface, every integral function is $\sim 0$.

Integral functions can be multiplied by constant factors and added. Integral functions $F_1, \dots, F_m$ between which a homology
\[ c_1 F_1 + \dots + c_m F_m \sim 0 \]
exists with constant coefficients $c$ not all vanishing are called \emph{linearly dependent}. If there exist finitely many, say $h$, linearly independent integral functions $F_1, \dots, F_h$ of such a kind that every integral function $F$ is homologous to a linear combination of these $h$ integral functions with constant coefficients, then $F_1, \dots, F_h$ form a \emph{``basis''} of the linear family of inhomologous integral functions, and the number $h$ of basis functions (which is obviously the same for every basis) is designated as the \emph{``degree''} of this family. We call it also the \emph{``degree of con-}

\origpage{69}
\emph{nectivity''}${}^{1)}$ of the surface $\mathfrak{F}$ whose integral functions we consider. If there exists no finite basis for the integral functions, $\mathfrak{F}$ possesses an infinitely high degree of connectivity.

We consider in particular a polyhedron $\mathfrak{P}$ and set out to express its degree of connectivity $h$ by the number $d$ of triangles of $\mathfrak{P}$, and the numbers $e, k$ of its interior vertices and edges. In each of the finitely many triangles $\Delta$ of which $\mathfrak{P}$ is composed, we enter the barycenter $(1 : 1 : 1)$. If (as always in the following consideration) $\Delta_1, \Delta_2$ are any two of these triangles meeting along an interior edge of $\mathfrak{P}$, and $\mathfrak{s}_1, \mathfrak{s}_2$ their barycenters, then for every curve $\gamma$ running within $\Delta_1 + \Delta_2$ leading from $\mathfrak{s}_1$ to $\mathfrak{s}_2$ the integral function $F(\gamma)$ has the same value, which we wish to denote by $x_{\Delta_1 \Delta_2} [F]$. This is due to the fact that not only an individual triangle, but also the pair of triangles $\Delta_1 + \Delta_2$ is simply connected. It is always $x_{\Delta_1 \Delta_2} = - x_{\Delta_2 \Delta_1}$. If $\mathfrak{e}$ is a vertex situated within $\mathfrak{P}$, and $\Delta_1, \Delta_2, \dots, \Delta_r$ are the triangles grouping themselves around $\mathfrak{e}$ in cyclic order, then
\[ x_{\Delta_1 \Delta_2} + x_{\Delta_2 \Delta_3} + \dots + x_{\Delta_r \Delta_1} = 0 \tag{$\mathfrak{e}$} \]
must hold, because each star of triangles is also simply connected. (The left-hand side of this equation changes its sign if we reverse the cyclic ordering of the triangles $\Delta_i$, i.\,e.\ the indicatrix in the star of triangles.) If we prescribe the system of numbers $x_{\Delta_1 \Delta_2}$ in any way according to the stated conditions, one recognizes without difficulty that there always exists an integral function to which these $x_{\Delta_1 \Delta_2}$ belong in the indicated manner${}^{2)}$. An integral function $F$ is then and only then

\textbf{1)} This number corresponds to the \emph{``connectivity number''} of \emph{Riemann}, but for closed surfaces is lower by $1$ than this. Cf.\ \emph{Schläfli}, Crelle's Journal Vol.\ 76, p.\ 152, footnote, and \emph{Klein}, Math.\ Ann., Vol.\ 7, p.\ 550, footnote.

\textbf{2)} E.\,g.\ as follows: In an individual star of triangles I can assign to the triangles $\Delta_1, \dots, \Delta_r (\Delta_{r+1} = \Delta_1)$ numbers $g_1, \dots, g_r (g_{r+1} = g_1)$ such that
\[ x_{\Delta_1 \Delta_2} = g_2 - g_1, \quad x_{\Delta_2 \Delta_3} = g_3 - g_2, \dots, \quad x_{\Delta_r \Delta_1} = g_1 - g_r \]
holds. I define in the interior of this star a point function $f$ that in the interior of the triangle $\Delta_i\colon = g_i$, on the edge separating $\Delta_i$ from $\Delta_{i+1}$, apart from the endpoints, possesses the constant value $\frac{1}{2} \{g_i + g_{i+1}\}$, and at the common vertex of the triangles $\Delta_i$ becomes equal to the arithmetic mean of the $r$ numbers $g_i$. For a curve $\gamma = (\mathfrak{p}_1 \mathfrak{p}_2)$ running within this star I define $F(\gamma) = f(\mathfrak{p}_2) - f(\mathfrak{p}_1)$. In the choice of the $g_i$ there lies an arbitrariness in so far as I can increase them all by the same number; but that has no influence on $F(\gamma)$. If $\gamma$ lies in the interior of two stars of triangles simultaneously (which then have in common a pair of triangles in which $\gamma$ is situated), the value of $F(\gamma)$ also does not depend on which of the two stars of triangles I use for computation in the indicated manner. If $\gamma$ is an arbitrary curve, I can partition it into finitely many consecutive arcs $\gamma_1, \gamma_2, \dots, \gamma_n$ such that each $\gamma_i$ lies entirely in the interior of a star of triangles. I then set:
\[ F(\gamma) = F(\gamma_1) + F(\gamma_2) + \dots + F(\gamma_n). \]
According to this explanation there always results the same value, whatever partition into partial arcs $\gamma_j$ I may undertake. To compare the two values resulting for two different partitions and ascertain their agreement, I need only apply both partitions simultaneously. $F(\gamma)$ is an integral function as we desire.

\origpage{70}
$\sim 0$ if to each triangle $\Delta$ of $\mathfrak{P}$ there corresponds a number $g_\Delta$ such that for all pairs of triangles $\Delta_1, \Delta_2$ having an interior edge of $\mathfrak{P}$ in common,
\[ x_{\Delta_1 \Delta_2} [F] = g_{\Delta_2} - g_{\Delta_1} \]
holds.

If of any two oppositely equal numbers $x_{\Delta_1 \Delta_2}, x_{\Delta_2 \Delta_1}$ we retain always only one and assign it to the common edge of the triangles $\Delta_1, \Delta_2$, we have $k$ ``unknowns'' $x$; between them exist the $e$ linear homogeneous equations $(\mathfrak{e})$ corresponding to the individual interior vertices $\mathfrak{e}$ of $\mathfrak{P}$. Are these equations linearly independent of one another? Suppose there existed between their left-hand sides an identity with coefficients $y_{\mathfrak{e}}$. Let us consider an edge $\mathfrak{e}\mathfrak{f}$ whose two vertices $\mathfrak{e}, \mathfrak{f}$ lie in the interior of $\mathfrak{P}$. Let the two equations belonging to the vertices $\mathfrak{e}$ and $\mathfrak{f}$ be written in such form that they correspond to two coherent indicatrices of the two stars of triangles belonging to $\mathfrak{e}$ and $\mathfrak{f}$:
\[ \begin{aligned} (\mathfrak{e})\colon & \quad \Delta_1, \Delta_2, \Delta_3, \dots, \Delta_r \\ (\mathfrak{f})\colon & \quad \Delta'_1 = \Delta_2, \Delta'_2 = \Delta_1, \Delta'_3, \dots, \Delta'_s. \end{aligned} \]
The unknown $x_{\Delta_1 \Delta_2}$ then occurs only in these two equations and indeed with opposite signs, and it must therefore be $y_{\mathfrak{e}} = y_{\mathfrak{f}}$. If on the other hand $\mathfrak{e}\bullet$ is an edge of which one vertex $\mathfrak{e}$ lies in the interior, while the other $\bullet$ lies on the boundary of $\mathfrak{P}$, one concludes in the same way $y_{\mathfrak{e}} = 0$. If $\mathfrak{P}$ is open, one can lay from every interior vertex an edge path to the boundary $\mathfrak{e}\mathfrak{f}\dots\mathfrak{l}\bullet$ and finds
\[ y_{\mathfrak{e}} = y_{\mathfrak{f}} = \dots = y_{\mathfrak{l}} = 0. \]
If $\mathfrak{P}$ is closed, but one-sided, one can draw from $\mathfrak{e}$ an edge path returning to $\mathfrak{e}$ on which the indicatrix reverses: concluding along this from vertex to vertex, one obtains $y_{\mathfrak{e}} = - y_{\mathfrak{e}}$, thus likewise $y_{\mathfrak{e}} = 0$. If $\mathfrak{P}$ is closed and two-sided, furnish all stars of triangles on $\mathfrak{P}$ with indicatrices mutually coherent with each other, and write each of the equations $(\mathfrak{e})$ corresponding to this indicatrix. Then, since one can connect every vertex with every other by an edge path, one finds that all $y_{\mathfrak{e}}$ are equal to one another, and there really exists between the left-hand sides of the equations $(\mathfrak{e})$ the identity with coefficients $y_{\mathfrak{e}} = 1$. If we set $\varepsilon = 0$ for closed two-sided polyhedra, otherwise $\varepsilon = 1$, then the equations $(\mathfrak{e})$ have accordingly $k - e + 1 - \varepsilon$ linearly independent solutions.

\origpage{71}
If one arbitrarily assigns to each triangle $\Delta$ of $\mathfrak{P}$ a number $g_\Delta$, one always obtains through the general formula $x_{\Delta_1 \Delta_2} = g_{\Delta_2} - g_{\Delta_1}$ a solution of the equations in question; of the solutions arising in this way we shall say that they are $\sim 0$. The systems of numbers $\{g_\Delta\}$ form a linear family of degree $d$. Such a system yields $0$ for all $g_{\Delta_2} - g_{\Delta_1}$ if and only if $g_\Delta$ has the same value for all triangles $\Delta$ (we use here that $\mathfrak{P}$ is connected). The solutions of $(\mathfrak{e})$ that are $\sim 0$ thus form a linear family of degree $d - 1$. The excess
\[ (k - e + 1 - \varepsilon) - (d - 1) = (k - e - d + 2) - \varepsilon \]
gives the sought greatest number $h$ of linearly independent integral functions on $\mathfrak{P}$. \emph{It follows from the significance of $h$ uncovered by the above considerations that for closed surfaces this quantity is completely independent of the manner of their triangulation}${}^{1)}$; furthermore one must always have $h \geqq 0$.

That between $l$ closed paths $\gamma_1, \gamma_2, \dots, \gamma_l$ the \emph{homology}${}^{2)}$
\[ c_1 \gamma_1 + c_2 \gamma_2 + \dots + c_l \gamma_l \sim 0 \]
with numbers $c_i$ as coefficients holds, shall mean that for every integral function $F$ the equation
\[ c_1 F(\gamma_1) + c_2 F(\gamma_2) + \dots + c_l F(\gamma_l) = 0 \]
takes place. If the surface we are considering has the finite degree of connectivity $h$, then between $l > h$ closed paths there always exists such a homology with coefficients not all vanishing. One needs only to solve the $h$ homogeneous linear equations
\[ c_1 F_i(\gamma_1) + c_2 F_i(\gamma_2) + \dots + c_l F_i(\gamma_l) = 0 \quad [i = 1, \dots, h] \]
in which $F_1, \dots, F_h$ form a basis for the integral functions, in order to

\textbf{1)} Usually the \emph{connectivity number} (after \emph{Riemann}) is explained with the help of dissection of the given surface into a simply connected one, where within analysis situs one cannot avoid admitting arbitrary continuous curves as dissection lines; but the proofs that the number so defined is determined by the surface alone are not rigorous and seem difficult to bring into rigorous form. By our taking here as foundation a new definition standing in closest relation to function-theoretic applications (theory of Abelian integrals), this shortcoming has been avoided. One may well say that our procedure brings to light the essential core of the method applied by \emph{Weierstrass}, \emph{Hensel-Landsberg}, and others in the theory of algebraic functions, namely first to investigate the behavior of the \emph{integrals} and to draw from this conclusions about the \emph{paths of integration}, in a form freed from all function-theoretic contingencies. --- The equation $e + d - k = 2$ valid for closed, two-sided, simply connected polyhedra is designated as the \emph{Euler polyhedron formula}. See \emph{Euler}, Petrop.\ Novi Comm.\ 4 (1752---53), p.~109.

\textbf{2)} Cf.\ \emph{Poincaré}, Analysis situs, Journal de l'École polytechnique, Ser.\ 2, Vol.\ 1 (1895), p.~19.

\origpage{72}
determine such coefficients. There must accordingly exist a certain number $h' \; (\leqq h)$ of closed paths $\gamma_1, \dots, \gamma_{h'}$ such that between these no homology (with coefficients not all $0$) exists, but certainly between any $h' + 1$ closed paths. $\gamma_1, \dots, \gamma_{h'}$ then form a basis for the family of inhomologous closed paths, and $h'$ is the degree of this family. In an analogous manner as $h' \leqq h$ was found above, one proves $h \leqq h'$; accordingly $h' = h$.

\emph{The special nature of the objects (closed paths) of which our family consists entails that a basis $\gamma_1, \dots, \gamma_h$ can be found such that for every closed path a homology}
\[ \gamma \sim n_1 \gamma_1 + \dots + n_h \gamma_h \]
\emph{with integral coefficients $n_1, \dots, n_h$ takes place.} Let $\gamma'_1, \gamma'_2, \dots, \gamma'_h$ initially be an arbitrary basis of closed curves. To each closed line
\[ \gamma \sim r_1 \gamma'_1 + r_2 \gamma'_2 + \dots + r_h \gamma'_h \]
assign in a Cartesian $h$-dimensional space the point with coordinates $(r_1, r_2, \dots, r_h)$:
\[ \gamma \sim (r_1, r_2, \dots, r_h). \]
The system $G$ of points thus corresponding to all closed curves $\gamma$ forms a \emph{``lattice''}; i.\,e.\ with every point $(r_1, r_2, \dots, r_h)$ there belongs also $(-r_1, -r_2, \dots, -r_h)$ to $G$, and if $(r'_1, r'_2, \dots, r'_h)$, $(r''_1, r''_2, \dots, r''_h)$ are any two points belonging to $G$, then
\[ (r'_1 + r''_1, r'_2 + r''_2, \dots, r'_h + r''_h) \]
is also always contained in $G$. For from two closed paths $\gamma', \gamma''$ one can always obtain a closed path
\[ \gamma \sim \gamma' + \gamma'' \]
by traversing $\gamma'$ once from one of its points $\mathfrak{p}'$, then going from $\mathfrak{p}'$ along an arbitrary curve $\sigma$ to a point $\mathfrak{p}''$ of $\gamma''$, traversing $\gamma''$ from $\mathfrak{p}''$, and finally returning along $\sigma$ in the opposite direction to $\mathfrak{p}'$. In particular, all points with integral coordinates belong to $G$.

Lattice points do not, however, lie in arbitrary proximity to the origin $(0, 0, \dots, 0)$; rather there exists an integer $N$ such that the coordinates of each point belonging to $G$ transform into integers by multiplication with $N$. Proof: Since the coefficients of the equations $(\mathfrak{e})$ are integers, one can in particular specify $h$ \emph{integral} solutions
\[ \{ x_{\Delta_1 \Delta_2}^{(i)} \} \quad [i = 1, 2, \dots, h] \]
from which all others can be linearly composed in the sense of homology. To these $h$ solutions correspond $h$ linearly independent integral functions $F_i$, \emph{whose values for every closed curve}

\origpage{73}
\emph{are integers}${}^{1)}$. If we set the determinant, different from $0$,
\[ |F_i(\gamma'_j)|_{\substack{i = 1, 2, \dots, h \\ j = 1, 2, \dots, h}} = N, \]
our assertion emerges in view of the equations
\[ r_1 F_i(\gamma'_1) + r_2 F_i(\gamma'_2) + \dots + r_h F_i(\gamma'_h) = F_i(\gamma), \quad [i = 1, 2, \dots, h] \]
on whose right-hand sides stand integers.

Among all (\emph{finitely many}) lattice points of the form
\[ (r_1, 0, \dots, 0) \quad [0 < r_1 \leqq 1] \]
let
\[ \gamma_1 \sim (r_1^{(1)}, 0, \dots, 0) \]
be that for which $r_1$ has its smallest value; among the finitely many lattice points of the form
\[ (r_1, r_2, 0, \dots, 0) \quad [0 \leqq r_1 < r_1^{(1)}; \quad 0 < r_2 \leqq 1] \]
--- of which there certainly exist some --- furthermore let
\[ \gamma_2 \sim (r_1^{(2)}, r_2^{(2)}, 0, \dots, 0) \]
be that for which $r_2^{(2)}$ is as small as possible; among all lattice points
\[ (r_1, r_2, r_3, 0, \dots, 0) \quad [0 \leqq r_1 < r_1^{(1)}, \quad 0 \leqq r_2 < r_2^{(2)}; \quad 0 < r_3 \leqq 1] \]
\[ \gamma_3 \sim (r_1^{(3)}, r_2^{(3)}, r_3^{(3)}, 0, \dots, 0) \]
that for which $r_3^{(3)}$ turns out smallest; and so on. The closed curves $\gamma_1, \gamma_2, \gamma_3, \dots, \gamma_h$ then form a basis as we seek it${}^{2)}$. For if $\gamma \sim (r_1, r_2, \dots, r_h)$ is an arbitrary point belonging to $G$, one can successively determine the integers $n_h, n_{h-1}, \dots, n_1$ such that
\[ \begin{gathered} \gamma - (n_h \gamma_h + n_{h-1} \gamma_{h-1} + \dots + n_1 \gamma_1) \sim (\bar{r}_1, \bar{r}_2, \dots, \bar{r}_h) \\ 0 \leqq \bar{r}_i < r_i^{(i)} \quad (i = h, h - 1, \dots, 1) \end{gathered} \]
results. But then one concludes from the meaning of $\gamma_h, \gamma_{h-1}, \dots, \gamma_1$ successively
\[ \bar{r}_h = 0, \quad \bar{r}_{h-1} = 0, \dots, \quad \bar{r}_1 = 0, \]
and the final result is that asserted:
\[ \gamma \sim n_h \gamma_h + \dots + n_2 \gamma_2 + n_1 \gamma_1. \]

Only a basis $\{\gamma_i\}$ by which all closed curves can be expressed linearly, homogeneously, and \emph{integrally} in

\textbf{1)} In order to calculate the value $F(\gamma)$ of an integral function $F$ for a closed path $\gamma$, one can replace $\gamma$ by a closed polygonal path $\pi$ running in sufficient proximity to $\gamma$ (p.~49); we assume it such that, without passing through vertices, it crosses the edges at all meeting points. If $\pi$ passes at any such crossing point from the triangle $\Delta_1$ over into the triangle $\Delta_2$, then $F(\gamma) = F(\pi)$ is equal to the sum $\displaystyle\sum x_{\Delta_1 \Delta_2} [F]$ extended over all crossing points.

\textbf{2)} This concluding procedure is designated by \emph{Minkowski} as ``adaptation of a number lattice with respect to an included lattice''. Cf.\ \emph{Diophantische Approximationen} (Leipzig 1907), pp.~90---95.

\origpage{74}
the sense of homology shall from now on count as a real \emph{basis}. The transition from one such real basis to another is mediated by linear transformation with integral coefficients whose determinant must be $= \pm 1$.

Over the closed surface $\mathfrak{F}$ of connectivity degree $h$ there exists an unramified unbounded (regular) covering surface $\hat{\mathfrak{F}}$ on which a curve whose trace curve on $\mathfrak{F}$ is closed closes if and only if that trace curve is $\sim 0$. All integral functions on $\mathfrak{F}$, when regarded as integral functions on $\hat{\mathfrak{F}}$, become homologous to $0$; hence I call $\hat{\mathfrak{F}}$ the \emph{covering surface of the integral functions}. If $\gamma_1, \dots, \gamma_h$ is a basis for the closed paths on $\mathfrak{F}$, and we denote generally by $\hat{\mathfrak{p}} S_i$ that point on $\hat{\mathfrak{F}}$ at which one arrives if, starting from the arbitrary point $\hat{\mathfrak{p}}$ on $\hat{\mathfrak{F}}$, one traverses on $\hat{\mathfrak{F}}$ a path whose trace line in $\mathfrak{F}$ is closed and $\sim \gamma_i$, then the $S_i$ denote deck transformations of $\hat{\mathfrak{F}}$ into itself and together generate the entire group of these deck transformations, which is a commutative (Abelian) group:
\[ S_1^{n_1} S_2^{n_2} \dots S_h^{n_h} \]
[$n_1, n_2, \dots, n_h$ run independently of one another through all integers].

In the case of closed two-sided surfaces, \emph{Riemann} constructed for the function-theoretic applications to be discussed in Chap.~II a special basis of closed paths, the so-called \emph{canonical dissection}.${}^{1)}$ We base this construction on the following considerations.

On a simply connected polyhedron every integral function is $\sim 0$, thus its degree of connectivity $h = 0$. \emph{The converse of this theorem also holds.} For if, e.\,g., the open polyhedron $\mathfrak{P}$ for which $h = 0$ is presupposed were not simply connected, one could introduce a crosscut (consisting of $n$ edges) that does not dissect $\mathfrak{P}$. Thereby $\mathfrak{P}$ is transformed into a polyhedron $\mathfrak{P}'$ whose numbers are calculated from the formulas
\[ d' = d, \quad k' = k - n, \quad e' = e - (n - 1) \]
and for $\mathfrak{P}'$ there would result $h' = - 1$, which is impossible.

Let there be given on a closed two-sided surface $\mathfrak{F}_\zeta$ present in a definite triangulation $\zeta$ a polygon $\pi$ consisting of edges of $\zeta$ that does not dissect $\mathfrak{F}$. Then, since $\pi$ does not dissect the surface, one can draw a polygon $\pi'$ on $\mathfrak{F}_\zeta$ that crosses $\pi$ at a single place, but otherwise has no point in common with $\pi$. If we

\textbf{1)} \emph{Riemann}, Theorie der Abelschen Funktionen, Werke, 2nd ed., pp.~129---130.

\origpage{75}
make such a subdivision $\zeta^*$ of $\zeta$ that $\pi'$ also consists entirely of edges, then $\pi + \pi'$ is a closed point set of such a kind that it leaves undissected not only $\mathfrak{F}$, but every domain to which it belongs. The triangles of $\zeta^*$ having edges lying on $\pi + \pi'$ can namely be arranged in a single closed chain (where admittedly triangles can occur multiple times in the chain) in such a way that every two consecutive triangles of the chain have an edge in common not belonging to $\pi + \pi'$. In producing this arrangement use is made of the fact that $\pi$ as well as $\pi'$ possesses separated banks. $\pi + \pi'$ we call a \textbf{pair of loop cuts}.

\rmfigure{figures/fig-21.png}{Fig.~21.}{Pair of loop cuts}

If on $\mathfrak{F}_\zeta$ one has $q$ mutually non-intersecting pairs of loop cuts
\[ \pi_1 + \pi'_1, \quad \pi_2 + \pi'_2, \quad \cdots, \quad \pi_q + \pi'_q, \]
that do not dissect $\mathfrak{F}_\zeta$, one can connect an arbitrary point $\mathfrak{p}_0$ on the surface (not lying on these $q$ pairs of loop cuts) by simple polygonal paths $\sigma_1, \dots, \sigma_q$ with the intersection points of the pairs of loop cuts. We can choose the $\sigma$ such that any two of them do not intersect except at $\mathfrak{p}_0$, and that $\sigma_i$ has the intersection point, but otherwise no point, in common with the $i^{\text{th}}$ pair of loop cuts, but does not meet the remaining pairs of loop cuts at all: we have then a ``team'' of $q$ pairs of loop cuts attached by the ``reins'' $\sigma$ to the point $\mathfrak{p}_0$. If we make such a subdivision $\zeta^*$ of $\zeta$ that the entire cut system
\[ \Sigma = \sum_{i=1}^q (\pi_i + \pi'_i + \sigma_i) \]
consists of edges, we recognize that it leaves undissected not only $\mathfrak{F}$, but every domain that contains $\Sigma$ entirely (Fig.~22). By severing the connection of the triangles across the edges of the cut system, we obtain from the closed polyhedron an open polyhedron $\mathfrak{P}'$. The number of triangles $d$ has remained the same in this operation; the number $k - e$ has decreased by $2q - 1$. There must therefore (if $h$ denotes the degree of connectivity of $\mathfrak{F}$) be
\[ h - (2q - 1) \geqq 1 \quad \text{or} \quad 2q \leqq h \]
If still $2q < h$, then $\mathfrak{P}'$ is not simply connected. We then draw in $\mathfrak{P}'$ a crosscut $v$ that does not dissect $\mathfrak{P}'$, and connect two points situated opposite each other on the two banks of an edge belonging to $v$ by a polygon $\pi'_{q+1}$ in $\mathfrak{P}'$ not meeting $v$. We omit from $v$ at the beginning and end a small piece of the first resp.\ last edge, and connect the beginning and end of the thus mutilated

\origpage{76}
crosscut $v^-$ by a polygonal path $\tau$ within $\mathfrak{P}'$ that remains in such proximity to the cut system $\Sigma$ that it cannot meet $\pi'_{q+1}$ either. $v^- + \tau = \pi_{q+1}$ is then a closed polygon that does not dissect $\mathfrak{P}'$, since $\pi'_{q+1}$ connects the two banks of a segment belonging to $\pi_{q+1}$ without crossing $\pi_{q+1}$ within $\mathfrak{P}'$. $\pi_{p+1} + \pi'_{q+1}$\ednote{Misprint in the original: $\pi_{p+1}$ instead of $\pi_{q+1}$.} is a further pair of loop cuts having not a single point in common with the preceding ones and together with them still not dissecting $\mathfrak{F}$.

\rmfigure{figures/fig-22.png}{Fig.~22.}{Canonical dissection}

If the number $h > 2q$, it must consequently be $\geqq 2q + 2$. This shows that the number $h$ for two-sided closed surfaces can only be even. If we set $h = 2p$, there exists a team of $p$ pairs of loop cuts $\pi + \pi'$ attached with the help of reins $\sigma$ to a point $\mathfrak{p}_0$ of the surface. $p$ is called the \textbf{genus} of $\mathfrak{F}$. We can assume (passing eventually from the original triangulation $\zeta$ to a subdivision $\zeta^*$) that all these lines $\pi$, $\pi'$, $\sigma$ consist of edges; they \emph{yield the} \textbf{canonical dissection} \emph{of $\mathfrak{F}_{\zeta^*}$ into a simply connected polyhedron} $\mathfrak{F}_*$.

Let each triangle of $\mathfrak{F}_{\zeta^*}$ be furnished with an indicatrix according to the principle of coherence. Of two elementary triangles of the triangulation meeting along a directed edge $\overrightarrow{12}$, we call that whose $\text{indicatrix} = (12*)$ the \textbf{left}, the other the \textbf{right}. To the reins $\sigma$ we impart the direction from $\mathfrak{p}_0$ to the intersection points; to the $\pi$, $\pi'$ such a sense of traversal that $\pi'$ crosses the line $\pi$ from right to left (hence $\pi$ crosses $\pi'$ from left to right).

We consider an integral function $F$ on $\mathfrak{F}$. There exists on $\mathfrak{F}_*$ a point function $f$ such that for every curve $\gamma$ running in $\mathfrak{F}_*$ whose initial and terminal points shall be called $\mathfrak{p}_1$, $\mathfrak{p}_2$,
\[ F(\gamma) = f(\mathfrak{p}_2) - f(\mathfrak{p}_1) \]
holds. Let $\Delta'$, $\Delta''$ be two triangles having in common an edge $\varkappa$ belonging to the canonical cut system ($\Delta'$ lying on the left, $\Delta''$ on the right bank of the cut in question); then there exists in $\Delta' + \Delta''$ a function $f_0$ such that for every curve $\gamma = (\mathfrak{p}_1 \mathfrak{p}_2)$ running in $\Delta' + \Delta''$
\[ F(\gamma) = f_0(\mathfrak{p}_2) - f_0(\mathfrak{p}_1). \]
One must then have
\[ f = f_0 + c' \text{ in } \Delta', \quad f = f_0 + c'' \text{ in } \Delta'' \]

\origpage{77}
where $c'$, $c''$ are constants. $c'' - c'$ I call the \textbf{jump} of $F$ on $\varkappa$. If one considers a star of triangles grouping itself around a vertex situated on the canonical cut system that does not coincide with any of the intersection points, it turns out that the jump of $F$ possesses the same value for all edges belonging to the same one of the cuts $\pi$, $\pi'$, $\sigma$. Let it be
\[ = a_i \text{ for } \pi_i, \quad = a'_i \text{ for } \pi'_i, \quad = c_i \text{ for } \sigma_i. \]
If one considers the star of triangles grouping itself around the intersection point of the $i^{\text{th}}$ pair of loop cuts, there results $c_i = 0$. The reins $\sigma_i$ can accordingly be omitted. The numbers $a_i$, $a'_i$ are called the \textbf{periods} of the integral function at the cuts $\pi_i$, $\pi'_i$.

If $\gamma$ is an arbitrary closed curve, it can be replaced for the computation of $F(\gamma)$ by a polygon running in sufficient proximity to $\gamma$; let this be assumed such that it crosses the pairs of loop cuts at finitely many places that are not vertices of the triangulation. If $+ n_i$ is the number of times that that polygon crosses $\pi_i$ from left to right, diminished by the number of times that the same happens from right to left, and $- n'_i$ has the analogous meaning for $\pi'_i$, then obviously
\[ F(\gamma) = \sum_{i=1}^p n_i a_i - \sum_{i=1}^p n'_i a'_i; \]
on the other hand
\[ a_i = F(\pi'_i), \quad - a'_i = F(\pi_i). \]
The $\pi_i$, $\pi'_i$ thus form a real basis for the closed paths.

\emph{The genus $p$ is, as Möbius and Jordan}${}^{1)}$ \emph{have proven, the only analysis situs invariant of closed two-sided surfaces.} The theorem namely that two closed two-sided surfaces equivalent in the sense of analysis situs must possess the same $p$ can be inverted. It causes no great trouble to remodel Jordan's proof into one equal in rigor to the analysis situs considerations undertaken in this work; we do not wish, however, to enter into this. What decisive importance attaches to the genus of a closed Riemann surface for the theory of functions on this surface will emerge sufficiently from the function-theoretic theorems of the next chapter.

It is easy to specify closed two-sided surfaces in three-dimensional Euclidean space of relatively simple nature whose genus $p$ has a given value; e.\,g.\ a sphere provided with $p$ handles possesses the genus $p$.

\textbf{1)} Möbius, ``Theorie der elementaren Verwandtschaft'' (1863), Werke Vol.~II, pp.~435---471; C.~Jordan, Journal de mathématiques, Ser.~2, Vol.~11 (1866), p.~105. Further: W.~Dyck, Math.\ Ann.\ Vol.~32 (1888), p.~457.

\origpage{78}

\section*{Second Chapter. Functions on Riemann Surfaces.}

\subsection*{§ 12. The Dirichlet Integral.}

Let $\mathfrak{E}$ be a closed point set lying entirely in the finite in the Euclidean plane with rectangular coordinates $xy$. To determine the area of $\mathfrak{E}$, draw in the plane the square mesh of side length $\varepsilon$ produced by the parallels
\[ y = m\varepsilon \ (m \text{ arbitrary integer}) \]
to the $x$-axis and the parallels
\[ x = n\varepsilon \ (n \text{ arbitrary integer}) \]
to the $y$-axis. If then $i_\varepsilon$ is the number of squares of the mesh consisting entirely of interior points of $\mathfrak{E}$, while $a_\varepsilon\ (> i_\varepsilon)$ is the number of those squares of the mesh that have any points in common with $\mathfrak{E}$ at all, there exist${}^{1)}$ the limits
\[ \lim_{\varepsilon=0} \varepsilon^2 i_\varepsilon = I, \quad \lim_{\varepsilon=0} \varepsilon^2 a_\varepsilon = A. \]

If $I = A$ turns out to be the case, one calls this number the \textbf{content} (or area) of $\mathfrak{E}$. Those points of $\mathfrak{E}$ in whose arbitrary proximity points not belonging to $\mathfrak{E}$ are encountered form the \textbf{boundary} of $\mathfrak{E}$. For $\mathfrak{E}$ to possess a definite content, it is obviously necessary and sufficient that the boundary of $\mathfrak{E}$ possess content $0$.

A continuously differentiable curve, and all the more an analytic or piecewise analytic curve in the plane, possesses content $0$. Namely, the continuously differentiable curve $\gamma$ of length $l$ can be divided into $n + 1$ partial arcs each of length $\frac{l}{n + 1}$. If we then use a square mesh of edge length $\varepsilon = \frac{l}{n}$, such a single

\textbf{1)} C.~Jordan, \emph{Cours d'analyse}, 2nd ed., Vol.~1, p.~28.

\origpage{79}
partial arc can obviously have points in common with at most four squares of the mesh; in total there will thus be at most $4(n + 1)$ squares of the mesh that contain points of $\gamma$:
\[ a_\varepsilon \cdot \varepsilon^2 \leqq 4(n + 1) \frac{l^2}{n^2}, \]
from which the correctness of the assertion $\lim\limits_{\varepsilon=0} a_\varepsilon \cdot \varepsilon^2 = 0$ emerges.${}^{1)}$

If the closed point set $\mathfrak{E}$, of which we now assume once and for all that it possesses a content $J$, lies entirely in a domain $\mathfrak{G}$ in which a continuous function $f$ is defined, then the integral $\displaystyle\iint\limits_{\mathfrak{E}} f\,dx\,dy$ has a clear meaning. An approximate value of it is computed by multiplying the content $\varepsilon^2$ of each square of the $\varepsilon$-mesh belonging to $\mathfrak{E}$ by the value of $f$ at the center of this square and adding all these contributions originating from the individual squares. The number so ascertained converges to the value of the integral when one lets $\varepsilon$ tend to $0$. It is immaterial here whether one sums only over the squares of the mesh situated entirely within $\mathfrak{E}$, or in addition also takes into account some or all squares abutting the boundary of $\mathfrak{E}$. If $\mathfrak{E}$ is partitioned into finitely many closed sets $\mathfrak{E}_1, \mathfrak{E}_2, \cdots, \mathfrak{E}_n$ with definite content (such that the $\mathfrak{E}_i$ are everywhere distinct in their interior points, each $\mathfrak{E}_i$ is a part of $\mathfrak{E}$, but also every point of $\mathfrak{E}$ is contained in at least one of the $\mathfrak{E}_i$), then from the definition there immediately results the addition law
\[ \iint_{\mathfrak{E}} f\,dx\,dy = \sum_{i=1}^n \iint_{\mathfrak{E}_i} f\,dx\,dy. \]

If $\mathfrak{G}$ is mapped one-to-one and domain-continuously onto a domain $\mathfrak{G}'$ of the $x'y'$-plane:
\[ x' = x'(xy), \qquad y' = y'(xy) \tag{14} \]
where the derivatives
\[ \begin{Vmatrix} \dfrac{\partial x'}{\partial x} & \dfrac{\partial x'}{\partial y} \\[1ex] \dfrac{\partial y'}{\partial x} & \dfrac{\partial y'}{\partial y} \end{Vmatrix} \]
shall be continuous functions of $xy$ and their determinant $d$ everywhere $\neq 0$, then the transformation formula holds:
\[ \iint_{\mathfrak{E}} f\,dx\,dy = \iint_{\mathfrak{E}'} \frac{f}{d}\,dx'\,dy'. \]

\textbf{1)} The proof is due to C.~Jordan (\emph{Cours d'analyse}, 2nd ed., Vol.~1, p.~107).

\origpage{80}
Here $\mathfrak{E}'$ is the image of $\mathfrak{E}$, and in the integral on the right $\frac{f}{d}$ is to be expressed as a function of $x', y'$ by virtue of the transformation formulas.

If $u$ is a continuously differentiable function in $\mathfrak{G}$, one is accustomed (incidentally with little historical justification) to designate
\[ \mathbf{D}_{\mathfrak{E}}(u) = \iint_{\mathfrak{E}} \left[ \left(\frac{\partial u}{\partial x}\right)^2 + \left(\frac{\partial u}{\partial y}\right)^2 \right] dx\,dy \]
as the \textbf{Dirichlet integral}. \emph{It is invariant under conformal mapping.} Indeed: if the mapping (14) is conformal ($d \neq 0$ is then a consequence of the presupposed one-to-one character of the mapping), then in consequence of the Cauchy-Riemann differential equations
\[ \begin{gathered}
\left(\frac{\partial u}{\partial x}\right)^2 + \left(\frac{\partial u}{\partial y}\right)^2 = \left[ \left(\frac{\partial u}{\partial x'}\right)^2 + \left(\frac{\partial u}{\partial y'}\right)^2 \right] \left[ \left(\frac{\partial x'}{\partial x}\right)^2 + \left(\frac{\partial x'}{\partial y}\right)^2 \right], \\
d = \left(\frac{\partial x'}{\partial x}\right)^2 + \left(\frac{\partial x'}{\partial y}\right)^2,
\end{gathered} \]
thus
\[ \iint_{\mathfrak{E}} \left[ \left(\frac{\partial u}{\partial x}\right)^2 + \left(\frac{\partial u}{\partial y}\right)^2 \right] dx\,dy = \iint_{\mathfrak{E}'} \left[ \left(\frac{\partial u}{\partial x'}\right)^2 + \left(\frac{\partial u}{\partial y'}\right)^2 \right] dx'\,dy'. \]

\emph{This invariance law makes it possible to form the Dirichlet integral of a continuously differentiable function not only in the plane, but on an arbitrary Riemann surface.}

Let therefore $\mathfrak{F}$ be a closed Riemann surface, and $u$ a continuously differentiable function on it; we wish to explain what is to be understood by the Dirichlet integral $\mathbf{D}(u)$ of $u$ on $\mathfrak{F}$. If $\mathfrak{p}$ is a point on $\mathfrak{F}$ and $z$ a local uniformizing variable belonging to $\mathfrak{p}$ mapping a neighborhood of $\mathfrak{p}$ onto a domain of the $z$-plane containing the point $z = 0$, then a circle $|z| \leqq a$ belonging entirely to this domain is the image of a point set on $\mathfrak{F}$ which we call a $z$-\textbf{circle} around $\mathfrak{p}$. The points for which $|z| = a$ form the periphery of that circle. If to each point $\mathfrak{p}$ one arbitrarily assigns a local uniformizing variable $z$ and a $z$-circle $\mathrm{K}$, one can select among these $\mathrm{K}$ a finite number $\mathrm{K}_l$ ($l = 1, 2, \cdots, n$; belonging to the points $\mathfrak{p}_l$ with local uniformizing variables $z_l$) in such a way (Heine-Borel theorem) that every point of the surface is situated in the interior of one of the $n$ circles $\mathrm{K}_l$. (The Heine-Borel theorem for an arbitrary closed surface results if one forms on it a sequence of triangulations, each of which is a subdivision of the preceding one and which finally become arbitrarily fine.) The peripheries $\varkappa_l$ of the $\mathrm{K}_l$ are closed analytic lines and therefore intersect only in finitely many points. Consequently the closed set $\varkappa_1 + \varkappa_2 + \cdots + \varkappa_n$ determines on $\mathfrak{F}$ only finitely many domains, which, after each of them has been made into a closed set by addition

\origpage{81}
of its boundary, we denote by $\mathfrak{E}_1, \mathfrak{E}_2, \dots, \mathfrak{E}_r$. The boundary of each $\mathfrak{E}_h$ consists of finitely many analytic curve segments. Each $\mathfrak{E}_h$ lies entirely in one of the circles $\mathrm{K}_l$, and $z_l$ is then a \textbf{uniformizing variable belonging to $\mathfrak{E}_h$}, i.\,e.\ a function regular at all points of a certain domain containing $\mathfrak{E}_h$ which provides a one-to-one and domain-continuous conformal image of this domain in the $z_l$-plane. We imagine assigned in a definite manner to each $\mathfrak{E}_h$ a uniformizing variable $z_h = x_h + i y_h$ belonging to $\mathfrak{E}_h$ and then form
\[ \mathbf{D}(u) = \sum_{h=1}^r \iint_{\mathfrak{E}_h} \left[ \left(\frac{\partial u}{\partial x_h}\right)^2 + \left(\frac{\partial u}{\partial y_h}\right)^2 \right] dx_h\,dy_h. \]

Here $u$ in $\mathfrak{E}_h$ is to be conceived as expressed as a function of $x_h y_h$, and $\mathfrak{E}_h$ denotes simultaneously the planar image provided by $z_h$ of this piece of the Riemann surface. The number $\mathbf{D}(u)$ so ascertained is, according to the invariance law, independent of \emph{which uniformizing variable $z_h = x_h + i y_h$ is used in each piece $\mathfrak{E}_h$}. It is, however, also \emph{independent of the special manner of the dissection used into the uniformizable pieces $\mathfrak{E}_h$}. For let us choose in any other way points $\mathfrak{p}'_1, \dots, \mathfrak{p}'_{n'}$, associated uniformizing variables $z'$ and $z'_\nu$-circles $\mathrm{K}'_\nu$ with peripheries $\varkappa'$, in such a manner that the $\mathrm{K}'_\nu$ cover the entire surface! Let the $\varkappa'_1, \dots, \varkappa'_{n'}$ partition $\mathfrak{F}$ into the pieces $\mathfrak{E}'_1, \dots, \mathfrak{E}'_{r'}$, and for each $\mathfrak{E}'_{h'}$ let $z'_{h'} = x'_{h'} + i y'_{h'}$ be a uniformizing variable. We form anew
\[ \mathbf{D}'(u) = \sum_{h'=1}^{r'} \iint_{\mathfrak{E}'_{h'}} \left[ \left(\frac{\partial u}{\partial x'_{h'}}\right)^2 + \left(\frac{\partial u}{\partial y'_{h'}}\right)^2 \right] dx'_{h'}\,dy'_{h'}. \]

Then, as I assert,
\[ \mathbf{D}(u) = \mathbf{D}'(u). \]

Namely, apply the cuts $\varkappa_l$, $\varkappa'_\nu$ simultaneously. Since these lines intersect one another only in finitely many points, they partition the surface into finitely many pieces
\[ \mathfrak{E}_{hj} = \mathfrak{E}'_{h'j'} \left[ \begin{aligned} j &= 1, 2, \cdots, i_h; & h &= 1, 2, \cdots, r \\ j' &= 1, 2, \cdots, i'_{h'}; & h' &= 1, 2, \cdots, r' \end{aligned} \right]. \]
In the first designation each $\mathfrak{E}_{hj}$ is a part of $\mathfrak{E}_h$, in the second each $\mathfrak{E}'_{h'j'}$ a part of $\mathfrak{E}'_{h'}$. By the addition law,
\[ \mathbf{D}(u) = \sum_h \sum_j \iint_{\mathfrak{E}_{hj}} \left[ \left(\frac{\partial u}{\partial x_h}\right)^2 + \left(\frac{\partial u}{\partial y_h}\right)^2 \right] dx_h\,dy_h, \]

\origpage{82}
\[ \mathbf{D}'(u) = \sum_{h'} \sum_{j'} \iint_{\mathfrak{E}'_{h'j'}} \left[ \left(\frac{\partial u}{\partial x'_{h'}}\right)^2 + \left(\frac{\partial u}{\partial y'_{h'}}\right)^2 \right] dx'_{h'}\,dy'_{h'}, \]
by the invariance law,
\[ \iint_{\mathfrak{E}_{hj}} \left[ \left(\frac{\partial u}{\partial x_h}\right)^2 + \left(\frac{\partial u}{\partial y_h}\right)^2 \right] dx_h\,dy_h = \iint_{\mathfrak{E}'_{h'j'}} \left[ \left(\frac{\partial u}{\partial x'_{h'}}\right)^2 + \left(\frac{\partial u}{\partial y'_{h'}}\right)^2 \right] dx'_{h'}\,dy'_{h'}; \]
therefore
\[ \mathbf{D}(u) = \mathbf{D}'(u). \]

To every continuously differentiable function $u$ on $\mathfrak{F}$ there belongs thus a number $\mathbf{D}(u)$. It has the following properties:

1. It does not change if one replaces $u$ by $u + c$, where $c$ is a constant; it is never negative and $0$ only when $u$ is constant on the entire surface.

2. If $\mathfrak{p}$ is any point on the surface, $z = x + iy$ an associated local uniformizing variable, and $\mathrm{K}\colon |z| \leqq a$ a $z$-circle, then
\[ \mathbf{D}(u) = \iint_{\mathrm{K}} \left[ \left(\frac{\partial u}{\partial x}\right)^2 + \left(\frac{\partial u}{\partial y}\right)^2 \right] dx\,dy + R(u), \]
where $R(u) \geqq 0$, and does not change if one replaces $u$ by any continuously differentiable function on $\mathfrak{F}$ that agrees with $u$ for all points not lying in $\mathrm{K}$.

3. $\mathbf{D}(u)$ possesses quadratic character. This makes itself known in that: if $u$, $v$ are any two continuously differentiable functions on $\mathfrak{F}$ and $\lambda$, $\mu$ denote arbitrary constants, then
\[ \mathbf{D}(\lambda u + \mu v) = \lambda^2 \mathbf{D}(u) + 2\,\lambda\,\mu\,\mathbf{D}(uv) + \mu^2 \mathbf{D}(v) \]
is a quadratic form in $\lambda$, $\mu$.

On a non-closed Riemann surface as well, the Dirichlet integral of a continuously differentiable function $u$ can be formed, which here admittedly can also possess the value $\infty$. An unclosed surface will have to be dissected by peripheries $\varkappa_l$ of countably many ``circles'' $\mathrm{K}_l$ in such a way that every point belongs to only finitely many, but at least one of the circles $\mathrm{K}_l$.

\subsection*{§ 13. On the Poisson integral.}

As a tool for our investigations we require the solution of the so-called \emph{first boundary value problem of potential theory for the disc} with the aid of the \emph{Poisson integral}. If a continuous function is given on the unit circle of the complex $z$-plane, whose value at the position $\zeta = e^{i\varphi}$ may be denoted by $u(\varphi)$, then the equation

\origpage{83}
\[ 2\,\pi f(z) = \int_0^{2\pi} \frac{\zeta + z}{\zeta - z} u(\varphi)\,d\varphi \]
yields a regular-analytic function $f(z)$ for $|z| < 1$, whose real part assumes the values $u(\varphi)$ on the boundary; that is, the function $u$ which in the interior of the unit disc is set $=$ the real part of $f(z)$, but for values $\zeta = e^{i\varphi}$ on the unit circle\ednote{Misprint in the original: ``Einkeitskreise'' instead of ``Einheitskreise''.} is set $= u(\varphi)$, is a continuous function in the entire closed disc $|z| \leqq 1$.${}^{1)}$ If we use polar coordinates $z = r e^{i\varphi}$, we can write:
\[ 2\,\pi u = \int_0^{2\pi} \frac{1 - r^2}{1 + r^2 - 2r \cos(\varphi - \vartheta)} u(\vartheta)\,d\vartheta. \tag{15} \]

There is not, in addition to $u$, any other potential function $\bar{u}$ regular in the interior of the unit disc which in the sense explained assumes on the boundary of the disc the values $u(\varphi)$. For then $\bar{u} - u$ would be a function continuous in the closed disc, which as a regular potential function cannot assume a maximum or minimum anywhere in the interior, and which on the boundary is $= 0$. Accordingly both the maximum and the minimum of $\bar{u} - u$ (for $|z| \leqq 1$) would have to be $= 0$, that is, $\bar{u} \equiv u$.

The Poisson integral (15) provides an important convergence theorem: \emph{If $u_1$, $u_2$, $u_3$, \dots\ is a sequence of potential functions, all of which are regular for $|z| < 1$ and converge uniformly for $|z| \leqq q$ to a limit function $u$, where $q$ denotes any positive number $< 1$, then the limit function $u$ is likewise a potential function regular for $|z| < 1$.} For under the stated circumstances, according to the Poisson integral, we must have
\[ 2\,\pi u_n(r e^{i\varphi}) = \int_0^{2\pi} \frac{q^2 - r^2}{q^2 + r^2 - 2qr \cos(\varphi - \vartheta)} u_n(q e^{i\vartheta})\,d\vartheta \]
for $r < q$. Passing to the limit yields
\[ 2\,\pi u(r e^{i\varphi}) = \int_0^{2\pi} \frac{q^2 - r^2}{q^2 + r^2 - 2qr \cos(\varphi - \vartheta)} u(q e^{i\vartheta})\,d\vartheta, \]
and here on the right-hand side stands a potential function regular for $r < q$. If $q_1$ is a positive number $< 1$ and I choose $q > q_1$ and $< 1$, one further recognizes that in an arbitrary concentric disc whose radius is $< 1$ ($|z| \leqq q_1$), not only does $u_n$ converge uniformly to $u$, but

\textbf{1)} H. A. Schwarz, Gesammelte Abhandlungen, vol.~II, pp.~186--198. The proof is presented in almost all textbooks on the theory of functions.

\origpage{84}
also all derivatives of $u_n$ of arbitrarily high order converge uniformly to the corresponding derivatives of $u$.

\emph{The connection of potential theory with the Dirichlet integral emerges from the following theorem:}

\emph{If $v$ is a function continuous in the unit disc $\mathrm{K}\colon |z| \leqq 1$, continuously differentiable in its interior, with finite Dirichlet integral}
\[ \mathbf{D}(v) = \iint_{\mathrm{K}} \left[ \left(\frac{\partial v}{\partial x}\right)^2 + \left(\frac{\partial v}{\partial y}\right)^2 \right] dx\,dy \quad (z = x + iy) \]
\emph{and $u$ is that potential function regular for $|z| < 1$ whose boundary values on $|z| = 1$ coincide with those of $v$, then}
\[ \mathbf{D}(u) \leqq \mathbf{D}(v). \tag{16} \]

Here the Dirichlet integral is to be understood as an ``improper'' integral, namely as the limit of the integral
\[ \mathbf{D}_q(v) = \iint_{|z| \leqq q} \left[ \left(\frac{\partial v}{\partial x}\right)^2 + \left(\frac{\partial v}{\partial y}\right)^2 \right] dx\,dy \quad (q < 1) \]
as $\lim q = 1$. In general, when $0 \leqq q_1 < q_2 \leqq 1$, we use the notation
\[ \mathbf{D}_{q_1}^{q_2}(v) = \iint_{q_1 \leqq |z| \leqq q_2} \left[ \left(\frac{\partial v}{\partial x}\right)^2 + \left(\frac{\partial v}{\partial y}\right)^2 \right] dx\,dy. \]

The usual proof of (16) runs as follows: One sets $w = v - u$; integration by parts (Green's formula) yields:
\[ \mathbf{D}(wu) = \int_0^{2\pi} \left( w \frac{\partial u}{\partial r} \right)_{r=1} d\varphi - \iint_{\mathrm{K}} w\,\Delta u\,dx\,dy \quad \left( \Delta u = \frac{\partial^2 u}{\partial x^2} + \frac{\partial^2 u}{\partial y^2} \right), \]
and since in the interior of $\mathrm{K}\colon \Delta u = 0$, and on the boundary $w = 0$,
\[ \mathbf{D}(wu) = 0; \]
\[ \begin{aligned}
\mathbf{D}(v) = \mathbf{D}(u + w) &= \mathbf{D}(u) + 2\,\mathbf{D}(wu) + \mathbf{D}(w) \\
&= \mathbf{D}(u) + \mathbf{D}(w) \geqq \mathbf{D}(u).
\end{aligned} \]

If one wishes to make this inference rigorous, one must initially extend all Dirichlet integrals only over $|z| \leqq q\,(< 1)$ and then let $q$ converge to $1$. It then comes down to seeing that
\[ \mathbf{D}_q(wu) = q \int_0^{2\pi} \left( w \frac{\partial u}{\partial r} \right)_{r=q} d\varphi \]
tends to $0$ as $\lim q = 1$. Now one certainly knows that $w(q e^{i\varphi})$ converges uniformly in $\varphi$ to $0$ as $q$ approaches $1$, but by no means can it be said that

\origpage{85}
$\frac{\partial u}{\partial r}$ remains bounded in the entire unit disc; one can only assert that $(1 - r) \frac{\partial u}{\partial r}$ converges uniformly in $\varphi$ to $0$ as $\lim r = 1$, and that of course is not sufficient to conclude that $\lim_{q=1} \mathbf{D}_q(wu) = 0$. One will have to keep these difficulties in mind in order to properly appreciate the following proof due to Hadamard and Zaremba.${}^{1)}$

The boundary values $v(e^{i\varphi})$ may be called $v(\varphi)$ for short. If one expands the Poisson integral (15) in powers of $r$, one obtains
\[ \begin{aligned}
u &= \frac{a_0}{2} + a_1 r \cos\varphi + a_2 r^2 \cos 2\varphi + \cdots \\
&\qquad + b_1 r \sin\varphi + b_2 r^2 \sin 2\varphi + \cdots
\end{aligned} \]
\[ a_n = \frac{1}{\pi} \int_0^{2\pi} v(\varphi) \cos n\varphi\,d\varphi, \quad b_n = \frac{1}{\pi} \int_0^{2\pi} v(\varphi) \sin n\varphi\,d\varphi. \]

In the domain of potential functions, this expansion corresponds to the power series in the theory of analytic functions. In fact, $r^n \cos n\varphi$, $r^n \sin n\varphi$ are nothing other than the real and imaginary parts of $z^n$. The potential functions
\[ P_n = \frac{1}{\sqrt{\pi n}}\,r^n \cos n\varphi, \quad Q_n = \frac{1}{\sqrt{\pi n}}\,r^n \sin n\varphi \]
possess the following properties:
\[ \left\{
\begin{gathered}
\mathbf{D}_q(P_n, Q_m) = 0 \quad \text{without exception}, \\
\mathbf{D}_q(P_n, P_m) = 0, \quad \mathbf{D}_q(Q_n, Q_m) = 0 \quad \text{except for } n = m, \\
\mathbf{D}_q(P_n) = \mathbf{D}_q(Q_n) = q^{2n},
\end{gathered}
\right. \tag{I} \]
and if in general we set
\[ \iint_{|z| \leqq q} u v\,dx\,dy = \mathbf{J}_q(uv) \]
also these:
\[ \left\{
\begin{gathered}
\mathbf{J}_q(P_n, Q_m) = 0 \quad \text{without exception}, \\
\mathbf{J}_q(P_n, P_m) = \mathbf{J}_q(Q_n, Q_m) = 0 \quad \text{except for } n = m, \\
\mathbf{J}_q(P_n) = \mathbf{J}_q(Q_n) = \frac{q^{2n+2}}{2n(n+1)}.
\end{gathered}
\right. \tag{II} \]
We can write the expansion of $u$:

\textbf{1)} Hadamard, Sur le principe de Dirichlet, Bulletin de la Société mathématique de France, vol.~34 (1906), pp.~135--139. Zaremba, Sur le principe du minimum, Bulletin de l'Académie des sciences de Cracovie, July 1909, pp.~206\,ff.

\origpage{86}
\[ u - \frac{a_0}{2} = \sum_{n=1}^\infty [A_n P_n(xy) + B_n Q_n(xy)], \tag{17} \]
where $A_n = a_n \sqrt{\pi n}$, $B_n = b_n \sqrt{\pi n}$ are constants that can be represented as surface integrals in the form
\[ A_n = \mathbf{D}(v, P_n), \quad B_n = \mathbf{D}(v, Q_n) \tag{18} \]
Indeed, by Green's formula, for example,
\[ \mathbf{D}_q(v, P_n) = q \int_0^{2\pi} \left( v \frac{\partial P_n}{\partial r} \right)_{r=q} d\varphi = \sqrt{\frac{n}{\pi}}\,q^n \int_0^{2\pi} v(q e^{i\varphi}) \cos n\varphi\,d\varphi. \]
By passing to the limit $\lim q = 1$, it indeed follows that
\[ \mathbf{D}(v, P_n) = \sqrt{\frac{n}{\pi}} \int_0^{2\pi} v(\varphi) \cos n\varphi\,d\varphi = \sqrt{n\pi} \cdot a_n = A_n. \]

The series (17) converges uniformly and absolutely for $|z| \leqq q\,(< 1)$, together with all its derivatives. One can therefore calculate $\mathbf{D}_q(u)$ from it by termwise squaring followed by termwise integration. The formulas (I) then yield:
\[ \mathbf{D}_q(u) = \sum_{n=1}^\infty q^{2n}(A_n^2 + B_n^2). \tag{19} \]
On the other hand, on the basis of the same formulas and the expressions (18), it follows that
\[ \mathbf{D}\Big( v - \sum_{\nu=1}^n [A_\nu P_\nu + B_\nu Q_\nu] \Big) = \mathbf{D}(v) - \sum_{\nu=1}^n (A_\nu^2 + B_\nu^2). \]
The series
\[ \sum_{n=1}^\infty (A_n^2 + B_n^2) \]
must therefore, no matter how far one extends it, always remain $\leqq \mathbf{D}(v)$; this entails its convergence. If we combine this with equation (19), we have
\[ \mathbf{D}_q(u) \leqq \mathbf{D}(v), \]
and by passing to the limit $\lim q = 1$, the existence of $\mathbf{D}(u)$ follows, and that
\[ \mathbf{D}(u) \leqq \mathbf{D}(v) \tag{20} \]
holds.

\emph{Among all functions that coincide with $v$ in their boundary values, the potential function thus imparts to the Dirichlet integral its smallest value.} If we set $v - u = w$, then according to this theorem we must have not only for $\lambda = 1$, but for every real constant value of $\lambda$,

\origpage{87}
\[ \mathbf{D}(u) \leqq \mathbf{D}(u + \lambda w) \]
or written otherwise:
\[ 2\,\lambda\,\mathbf{D}(uw) + \lambda^2 \mathbf{D}(w) \geqq 0. \]
This entails the equation
\[ \mathbf{D}(uw) = 0 \]
as a consequence, which provides us in
\[ \mathbf{D}(v) = \mathbf{D}(u) + \mathbf{D}(w) = \mathbf{D}(u) + \mathbf{D}(v - u), \]
with a sharper statement than was contained in the inequality (20). In particular, it shows that the equality sign in (20) can hold only if $v - u = \mathrm{const.}$, or since this difference vanishes on the boundary, $v$ is identical with $u$.

From (19) it follows by passing to the limit $\lim q = 1$:
\[ \mathbf{D}(u) = \sum_{n=1}^\infty (A_n^2 + B_n^2). \tag{21} \]
For on the one hand,
\[ \mathbf{D}_q(u) \leqq \sum_{n=1}^\infty (A_n^2 + B_n^2), \quad \text{hence} \quad \mathbf{D}(u) \leqq \sum_{n=1}^\infty (A_n^2 + B_n^2), \]
and on the other hand,
\[ \mathbf{D}_q(u) \geqq \sum_{\nu=1}^n (A_\nu^2 + B_\nu^2)\,q^{2\nu} \]
yields the inequality valid for every $n$
\[ \mathbf{D}(u) \geqq \sum_{\nu=1}^n (A_\nu^2 + B_\nu^2). \]

$a_1 = \frac{A_1}{\sqrt{\pi}}$, $b_1 = \frac{B_1}{\sqrt{\pi}}$ are the values of $\frac{\partial u}{\partial x}$, $\frac{\partial u}{\partial y}$ at the origin. If we retain on the right-hand side of (21) only the first term, we accordingly find
\[ \left[ \left(\frac{\partial u}{\partial x}\right)^2 + \left(\frac{\partial u}{\partial y}\right)^2 \right]_{z=0} \leqq \frac{1}{\pi}\,\mathbf{D}(u). \tag{22} \]
If we apply this inequality, instead of to the unit disc, to the disc $\varkappa_z$ with center $z\;(\frac{1}{2} < |z| = r < 1)$ and radius $1 - r$, it follows that:
\[ \begin{aligned}
\left(\frac{\partial u}{\partial x}\right)^2 + \left(\frac{\partial u}{\partial y}\right)^2 &\leqq \frac{1}{\pi(1 - r)^2} \iint_{\varkappa_z} \left[ \left(\frac{\partial u}{\partial x}\right)^2 + \left(\frac{\partial u}{\partial y}\right)^2 \right] dx\,dy \\
&\leqq \frac{1}{\pi(1 - r)^2}\,\mathbf{D}_{2r-1}(u). \tag{23}
\end{aligned} \]
For the disc $\varkappa_z$ is entirely contained in the annulus $2r - 1 \leqq |z| \leqq 1$.

\origpage{88}
Since $\lim_{r=1} \mathbf{D}_{2r-1}(u) = 0$, this inequality implies a statement concerning the behavior of the derivatives $\frac{\partial u}{\partial x}$, $\frac{\partial u}{\partial y}$ at the boundary of the unit disc.

Just as (21) results from formulas (I) by way of (19), we obtain by using the system of formulas (II):
\[ \mathbf{J}\Big( u - \frac{a_0}{2} \Big) = \sum_{n=1}^\infty \frac{A_n^2 + B_n^2}{2\,n(n+1)}. \]
In conjunction with (21) this yields the inequality
\[ \mathbf{J}\Big( u - \frac{a_0}{2} \Big) \leqq \frac{1}{4}\,\mathbf{D}(u). \tag{24} \]

This inequality is closely connected with an interesting elementary problem of analysis. If one asks how well one can approximate a function $v$ of two arguments $xy$ in a region $\mathfrak{G}$ by a constant, the first mean value theorem of differential calculus\ednote{Misprint in the original: ``Differentialrechnng'' instead of ``Differentialrechnung''.} states that, if I take as approximation constant the value of $v$ at a fixed point $(x_0 y_0)$,
\[ | v(xy) - v(x_0 y_0) | \leqq \underset{\mathfrak{G}}{\mathrm{Max.}}\, \sqrt{\left(\frac{\partial v}{\partial x}\right)^2 + \left(\frac{\partial v}{\partial y}\right)^2} \times \text{distance } (xy;\,x_0 y_0) \]
results. Here the maximum of the absolute value of the difference $|v - v_0|$ is regarded as the error, and its magnitude proves to be essentially determined by the $\underset{\mathfrak{G}}{\mathrm{Max.}}\,\sqrt{\left(\frac{\partial v}{\partial x}\right)^2 + \left(\frac{\partial v}{\partial y}\right)^2}$. If, however, we judge the error, as corresponds to the method of least squares, by the square integral $\displaystyle\iint_{\mathfrak{G}} (v - v_0)^2 dx\,dy$, one may expect that the approximation achievable by a suitable choice of the constant $v_0$ is determined by the value of
\[ \iint_{\mathfrak{G}} \left[ \left(\frac{\partial v}{\partial x}\right)^2 + \left(\frac{\partial v}{\partial y}\right)^2 \right] dx\,dy. \]

This conjecture is confirmed by the above inequality for the case where the region $\mathfrak{G}$ is the unit disc and $v$ is a potential function.

We wish to free ourselves from the latter assumption by proving generally, using the notations introduced above:
\[ \mathbf{J}\Big( v - \frac{a_0}{2} \Big) \leqq \frac{1}{2}\,\mathbf{D}(v). \tag{25} \]
$\frac{a_0}{2}$ is in general not the value of $v$ at the center, but the mean value of $v$ on the boundary of the disc:

\origpage{89}
\[ \frac{a_0}{2} = \frac{1}{2\,\pi} \int_0^{2\pi} v(\varphi)\,d\varphi. \]
We make use of the Schwarz inequality
\[ \left( \int_0^1 fg\,dx \right)^2 \leqq \int_0^1 f^2 dx \cdot \int_0^1 g^2 dx, \]
for which one can also write
\[ \sqrt{\int_0^1 (f \pm g)^2 dx} \leqq \sqrt{\int_0^1 f^2 dx} + \sqrt{\int_0^1 g^2 dx} \]
[$f(x)$, $g(x)$ being any two continuous functions on the interval $0 \leqq x \leqq 1$]. It arises from the fact that the quadratic form in $\lambda\mu$:
\[ \int_0^1 (\lambda \cdot f(x) + \mu \cdot g(x))^2 dx = \mathrm{A}\,\lambda^2 + 2\,\mathrm{B}\lambda\,\mu + \mathrm{\Gamma}\,\mu^2 \]
cannot become negative, and therefore its discriminant
\[ \mathrm{A}\mathrm{\Gamma} - \mathrm{B}^2 \geqq 0 \]
must hold.

For a function continuous in the closed unit disc, vanishing on the boundary, and continuously differentiable in the interior, such as $w = v - u$, we obtain by means of the Schwarz inequality
\[ \begin{aligned}
[w(r_2 e^{i\varphi}) - w(r_1 e^{i\varphi})]^2 &= \left( \int_{r_1}^{r_2} \frac{\partial w(r e^{i\varphi})}{\partial r}\,dr \right)^2 \\
&\leqq \int_{r_1}^{r_2} \left(\frac{\partial w}{\partial r}\right)^2 r\,dr \cdot \int_{r_1}^{r_2} \frac{dr}{r} \qquad (0 < r_1 < r_2 < 1)
\end{aligned} \]
If one integrates this inequality with respect to $\varphi$ and then lets $r_2$ converge to $1$, since
\[ \mathbf{D}_{r_1}^{r_2}(w) = \int_{r_1}^{r_2} \int_0^{2\pi} \left[ \left(\frac{\partial w}{\partial r}\right)^2 r + \left(\frac{\partial w}{\partial \varphi}\right)^2 \frac{1}{r} \right] d\varphi\,dr \]
holds, we obtain the relation
\[ \int_0^{2\pi} [w(r e^{i\varphi})]^2 d\varphi \leqq \lg\frac{1}{r} \cdot \mathbf{D}_r^1(w). \tag{26} \]

We shall make another important application of this inequality shortly. For now, we replace $\mathbf{D}_r^1(w)$ on the right by the larger $\mathbf{D}(w)$, multiply by $r\,dr$, and integrate with respect to $r$ from $0$ to $1$. Then, since

\origpage{90}
\[ \int_0^1 r \lg\frac{1}{r}\,dr = \frac{1}{4}\colon \]
\[ \mathbf{J}(w) \leqq \frac{1}{4}\,\mathbf{D}(w). \tag{27} \]
We combine (27) with (24). We have
\[ \mathbf{J}\Big( v - \frac{a_0}{2} \Big) \leqq 2 \Big[ \mathbf{J}\Big( u - \frac{a_0}{2} \Big) + \mathbf{J}(w) \Big] \]
--- in view of the inequality $(\alpha + \beta)^2 \leqq 2\,[\alpha^2 + \beta^2]$ --- and hence by (24) and (27)
\[ \leqq \frac{1}{2}\,[\mathbf{D}(u) + \mathbf{D}(w)] = \frac{1}{2}\,\mathbf{D}(v). \]

With this, (25) is proven. We note once more that in (25) the factor $\frac{1}{2}$ can be replaced by the smaller $\frac{1}{4}$ if $v$ is a potential function or vanishes on the boundary.

If $v$ is continuously differentiable in a domain containing the unit disc in its interior, then replacing $v$ by $u$ in the unit disc reduces the value of the Dirichlet integral. However, since in the intended applications we can only use continuously differentiable functions, whereas the function $\bar{v}$ which coincides with $v$ outside the unit disc and with $u$ inside it does not in general possess this property across the boundary of the unit disc, we must ``smooth'' the course of the function $\bar{v}$ there in such a way that the value of the Dirichlet integral of $\bar{v}$ changes as little as one pleases. We wish therefore to show: there exists a continuously differentiable function $\tilde{v}$ which coincides with $v$ outside the unit disc, and whose Dirichlet integral $\mathbf{D}(\tilde{v})$ extended over the unit disc differs from $\mathbf{D}(\bar{v}) = \mathbf{D}(u)$ by as little as one pleases.

I choose $q > \frac{1}{2}$ and $< 1$ and form, with $\chi(r) = \left( \frac{r - q}{1 - q} \right)^2$, the function${}^{1)}$
\[ \tilde{v} = \begin{cases}
u & \text{for } |z| \leqq q \\
u + \chi(r)(v - u) & \text{for } q \leqq |z| = r \leqq 1.
\end{cases} \]
Since $\chi(r)$ vanishes to the 2nd order at $r = q$, this function $\tilde{v}$, together with its first derivatives, passes continuously across the circle $|z| = q$. Towards the boundary of the unit disc, $\tilde{v} - v$, together with both its first derivatives, converges uniformly to $0$. Thus, if one lets $\tilde{v}$ coincide with $v$ outside the unit disc, $\tilde{v}$ is continuous and continuously differentiable across the boundary of the unit disc as well. For example, we have
\[ \frac{\partial}{\partial x}(\tilde{v} - v) = [\chi(r) - 1]\,\frac{\partial(v - u)}{\partial x} + (v - u)\,\frac{d\chi}{dr}\,\frac{x}{r} \quad (\text{for } q < r < 1). \]

\textbf{1)} Following S. Zaremba, Krakauer Berichte, July 1909, pp.~246\,ff.

\origpage{91}
$\chi(r) - 1$ becomes zero of the 1st order at $r = 1$, and it therefore follows from inequality (23) that the first term on the right-hand side tends uniformly to $0$ as the periphery of the unit disc is approached from the interior; that the second term does likewise hardly needs mentioning.

We now estimate the Dirichlet integral to be extended over the unit disc
\[ \mathbf{D}(\tilde{v} - u) = \mathbf{D}_q^1 [\chi(r)\,(v - u)] \]
Since
\[ \frac{\partial}{\partial x}[\chi(r)\,(v - u)] = \chi(r)\,\frac{\partial(v - u)}{\partial x} + (v - u)\,\frac{2(r - q)}{(1 - q)^2}\,\frac{x}{r} \]
holds, one obtains for this integral a value $\leqq$ twice
\[ \iint_{q \leqq |z| \leqq 1} \chi^2(r) \left[ \left(\frac{\partial w}{\partial x}\right)^2 + \left(\frac{\partial w}{\partial y}\right)^2 \right] dx\,dy + \frac{4}{(1 - q)^4} \int_0^{2\pi} \int_q^1 w^2(r - q)^2 r\,dr\,d\varphi. \]
The first term here is, on account of $\chi^2 \leqq 1$:
\[ \leqq \mathbf{D}_q^1(w) \]
the second, since according to (26)
\[ \int_0^{2\pi} w^2\,d\varphi \leqq \left(\frac{1}{r} - 1\right) \mathbf{D}_q^1(w) \quad (q < r < 1)\text{ is:} \]
\[ \leqq \frac{4}{(1 - q)^4} \cdot \mathbf{D}_q^1(w) \cdot \int_q^1 (1 - r)(r - q)^2 dr = \frac{1}{3}\,\mathbf{D}_q^1(w). \]
In total we thus have
\[ \mathbf{D}(\tilde{v}) - \mathbf{D}(u) = \mathbf{D}(\tilde{v} - u) \leqq \frac{8}{3}\,\mathbf{D}_q^1(w), \]
and this can be made as small as one pleases by a suitable choice of $q$.

\subsection*{§ 14. Approach to the proof of the existence theorems. Setup of the elementary differentials.}

Instead of constructing an analytic function on the given closed Riemann surface $\mathfrak{F}$, we initially seek only the real part of such a function, that is, a real function harmonic on the surface. However, an everywhere regular potential function on the surface does not exist (apart from constants). We will therefore permit a singularity at one point for the harmonic function, namely the simplest that exists: it shall behave there like the real part of an analytic function possessing a pole of the 1st order. We thus set out to prove the following existence theorem:

\emph{If $O$ with the local uniformizing variable $z_0 = x_0 + i y_0$ is an arbitrary}

\origpage{92}
\emph{point on $\mathfrak{F}$, then there exists, at all points different from $O$, a regular potential function $U$ on $\mathfrak{F}$, which in the neighborhood of $O$ differs from $\frac{x_0}{x_0^2 + y_0^2}$ by a potential function regular at $O$.}

$U$ is obviously uniquely determined up to an additive constant by these properties. Even though $\mathfrak{F}$ may not be a real surface in space, but a Riemann surface in the abstract sense introduced in Chap.~I, it will nevertheless be permissible to regard $U$ as the potential of an incompressible, stationary, irrotational fluid flow on $\mathfrak{F}$, a flow which everywhere, apart from the single point $O$, is also source-free; but at the point $O$ there resides a ``double source'' of moment $2\pi$, whose direction coincides with the positive $x_0$-axis. It takes, to be sure, some boldness to infer the existence of $U$ from this hydrodynamic interpretation.${}^{1)}$

Could the Dirichlet integral of $U$ be formed, we would have to designate it as the energy of the fluid flow. Unfortunately, however, the integral becomes infinite due to the behavior of $U$ at the point $O$. We therefore proceed as follows: We draw about $O$ a $z_0$-disc $\mathrm{K}_0\colon |z_0| \leqq a_0$ and form in it the function
\[ \Phi = \frac{x_0}{x_0^2 + y_0^2} + \frac{x_0}{a_0^2}. \]
It has the property of possessing the normal derivative $0$ on the boundary of $\mathrm{K}_0$. Indeed, if we use polar coordinates $z_0 = r_0 e^{i\varphi_0}$, we have
\[ \Phi = \frac{\cos\varphi_0}{r_0} + \frac{r_0 \cos\varphi_0}{a_0^2}, \]
\[ \frac{\partial\Phi}{\partial r_0} = -\,\frac{\cos\varphi_0}{r_0^2} + \frac{\cos\varphi_0}{a_0^2}, \quad \text{for } r_0 = a_0\colon \frac{\partial\Phi}{\partial r_0} = 0. \]
We now consider
\[ u = \begin{cases}
U \text{ everywhere on } \mathfrak{F}\text{, except in } \mathrm{K}_0 \\
U - \Phi \text{ in } \mathrm{K}_0.
\end{cases} \]
This function has a jump along the circumference $\varkappa_0$ of $\mathrm{K}_0$, but is otherwise continuously differentiable. Moreover, the behavior of $u$ prevailing outside $\mathrm{K}_0$ can be continued inward a little way across $\varkappa_0$ as a continuously differentiable (even as a regular potential) function ($U$); likewise the $u$ prevailing inside $\mathrm{K}_0$ can be continued outward a little way across $\varkappa_0$ as a continuously differentiable (even as a regular potential) function ($U - \Phi$). Consequently, one can form the Dirichlet integral $\mathbf{D}(u)$ in the manner explained above, provided that among the circle

\textbf{1)} The presentation in Klein's aforementioned work of 1882 is based on these and similar physical intuitions; cf.\ the very instructive flow diagrams drawn there.

\origpage{93}
circumferences $\varkappa_l$ along which $\mathfrak{F}$ was dissected into the pieces $\mathfrak{E}_h$, one also always includes the circle $\varkappa_0$.

\emph{Among all functions $v$ that possess along $\varkappa_0$ the same jump as $u$, but are otherwise continuously differentiable, $u$ imparts to the Dirichlet integral $\mathbf{D}(v)$ its smallest value.} We wish to characterize the functions $v$ used here for comparison with $u$ somewhat more precisely. For this purpose I draw around $O$ a $z_0$-disc $\mathrm{K}_0^*\colon |z_0| \leqq b_0$, of somewhat larger radius $b_0$ than $a_0$. I call $\mathfrak{F} - \mathrm{K}_0$ the ``\emph{perforated}'' \emph{surface}, $\mathrm{K}_0$ the ``\emph{hole}'', $\mathrm{K}_0^*$ the ``\emph{cover}'', and the annulus $a_0 < |z_0| < b_0$, in which the cover overlaps the perforated surface, the ``\emph{sealing ring}''. A function $v$ is admitted to competition if it is continuously differentiable in the perforated surface and coincides in the hole with a function $v^*$ continuously differentiable in the entire cover, to which it is related in the sealing ring by $v = v^* + \Phi$. If $v$ is the symbol for any competing function, $v^*$ shall always denote the continuously differentiable function in the cover characterized here. The difference of two competing functions is continuously differentiable on the whole of $\mathfrak{F}$.

In order now to carry out the existence proof of $U$ or $u$, we pose the \emph{minimum problem} of finding, among all competing functions $v$, that function $u$ for which the Dirichlet integral $\mathbf{D}(v)$ attains its smallest value. The Thomson-Dirichlet principle asserts that such a minimizing function $u$ exists.${}^{1)}$ Riemann founded his existence proofs on this principle, though his formulation does not coincide entirely with ours.${}^{2)}$ Weierstrass's criticism showed that the self-evidence that seems to adhere to Dirichlet's principle is only illusory.${}^{3)}$ One therefore sought to secure Riemann's existence theorems by other methods, and this was first achieved in a brilliant manner by Schwarz and C.~Neumann, especially with the aid of the so-called alternating method.${}^{4)}$ Later, however, Dirichlet's principle was rigorously proved by Hilbert.${}^{5)}$ To his proof a larger series of papers by other authors is linked, by which the

\textbf{1)} Applied since 1847 by W.~Thomson (Lord Kelvin) in various papers. The name ``Dirichlet's principle'' originates from Riemann, who knew this mode of reasoning from Dirichlet's lectures.

\textbf{2)} The form of the formulation adopted here was presented by the author (aside from the lecture from which the present work arose) at the meeting of 16~Jan.~1912 of the Göttingen Mathematical Society.

\textbf{3)} Über das sog.\ Dirichletsche Prinzip, (1870), Weierstraß\ednote{Misprint in the original: ``Weierstraß'' without apostrophe.} Werke, vol.~2, pp.~49--54.

\textbf{4)} H.~A.~Schwarz, Berliner Berichte 1870 (Ges.\ Abhandlungen II, pp.~133--171). C.~Neumann, sächsische Berichte 1870; ibid., Vorlesungen über Riemanns Theorie der Abelschen Integrale, 2nd ed., Leipzig 1884, pp.~388--471.

\textbf{5)} Hilbert, ``Über das Dirichletsche Prinzip'', Festschrift zur Feier des 150jährigen Bestehens der K.~Gesellschaft der Wissenschaften zu Göttingen, 1901, reprinted in Math.\ Ann.\ vol.~59, pp.~161--186; ``Zur Theorie der konformen Abbildung'', Göttinger Nachrichten, 1909, pp.~314--323.

\origpage{94}
proof was simplified, not indeed in its basic idea, but in many details.${}^{1)}$ Dirichlet's principle has, since these investigations, once again belonged to the most powerful tools of analysis.

To justify the formulation of our problem as a minimum problem, it shall first be shown that the minimizing function $u$, if it exists (it is then itself an admitted competing function), possesses the desired properties, that is, $u$ is regular-harmonic at every point of the perforated surface, and the associated $u^*$ is regular-harmonic at every interior point of the cover.

1. If $\mathfrak{p}$ is a point of the perforated surface, $z$ an associated local uniformizing variable, $\mathrm{K}$ a $z$-disc about $\mathfrak{p}$ lying entirely within the perforated surface, I can replace $u$, if it is not harmonic in $\mathrm{K}$, by that potential function $\bar{u}$ in $\mathrm{K}$ which coincides with $u$ on the boundary of $\mathrm{K}$. For this function, the Dirichlet integral extended over $\mathrm{K}$ satisfies $\mathbf{D}_{\mathrm{K}}(\bar{u}) < \mathbf{D}_{\mathrm{K}}(u)$. By means of the smoothing procedure described at the end of §~13, I obtain a function $\tilde{u}$, continuously differentiable across the boundary of $\mathrm{K}$, which coincides with $u$ except in $\mathrm{K}$, but whose Dirichlet integral extended over $\mathrm{K}$ differs from $\mathbf{D}_{\mathrm{K}}(\bar{u})$ by as little as I please. In particular, I can accordingly ensure that $\mathbf{D}_{\mathrm{K}}(\tilde{u})$ is still $< \mathbf{D}_{\mathrm{K}}(u)$. Then $\tilde{u}$ is a competing function for which the integral $\mathbf{D}(\tilde{u})$ extended over all of $\mathfrak{F}$ possesses a smaller value than for $u$. This contradiction shows that $u$ must be a regular potential function in $\mathrm{K}$.

\rmfigure{figures/fig-23.png}{Fig.~23.}{Circle K about p with decomposition into K' and K'' by the hole boundary varkappa_0}

2. Let $\mathfrak{p}$ be an interior point of the cover and $\mathrm{K}$ in the $z_0$-plane a disc about $\mathfrak{p}$ belonging entirely to the cover. If $u^*$ is not harmonic at $\mathfrak{p}$, one can replace $u^*$ by a function $\widetilde{u^*}$ continuously differentiable in the entire cover, which coincides with $u^*$ outside $\mathrm{K}$ and for which $\mathbf{D}_{\mathrm{K}}(\widetilde{u^*}) < \mathbf{D}_{\mathrm{K}}(u^*)$. There is then a unique competing function $\tilde{u}$ which, except in $\mathrm{K}$, is $= u$ and whose associated $(\tilde{u})^*$ in the cover coincides with $\widetilde{u^*}$. For this function, as I now wish to show, the integral extended over the whole surface satisfies $\mathbf{D}(\tilde{u}) < \mathbf{D}(u)$, and the assumption that $u^*$ at $\mathfrak{p}$ was not a regular-harmonic function is thereby refuted. In the proof we assume for generality that $\mathrm{K}$ intersects the

\textbf{1)} I mention in particular the papers published in the Rendiconti del Circolo matematico di Palermo, vols.~22--24 (1906--07): B.~Levi, Sul principio di Dirichlet; Fubini, Il principio di minimo e i teoremi di esistenza per i problemi al contorno relativi alle equazione alle derivate parziali di ordini pari; Lebesgue, Sur le problème de Dirichlet; the paper by S.~Zaremba, Sur le principe du minimum, in the Krakauer Berichte, July 1909, and that by R.~Courant, Über die Methode des Dirichletschen Prinzips, in vol.~72 of Math.\ Ann.\ (1912).

\origpage{95}
boundary of the hole $\varkappa_0$. Then (Fig.~23) $\mathrm{K}$ is divided by $\varkappa_0$ into two parts $\mathrm{K}'$, $\mathrm{K}''$, of which $\mathrm{K}'$ belongs to $\mathrm{K}_0$. The equation to be proved,
\[ \mathbf{D}(u) - \mathbf{D}(\tilde{u}) = \mathbf{D}_{\mathrm{K}}(u^*) - \mathbf{D}_{\mathrm{K}}(\widetilde{u^*}) \quad (> 0) \]
is identical with the following:
\[ \mathbf{D}_{\mathrm{K}''}(u^* + \Phi) - \mathbf{D}_{\mathrm{K}''}(\widetilde{u^*} + \Phi) = \mathbf{D}_{\mathrm{K}''}(u^*) - \mathbf{D}_{\mathrm{K}''}(\widetilde{u^*}). \tag{28} \]
Since $\Phi$ is a regular potential function in $\mathrm{K}''$ and even beyond the boundary of $\mathrm{K}''$, we have for example
\[ \mathbf{D}_{\mathrm{K}''}(u^* + \Phi) = \mathbf{D}_{\mathrm{K}''}(u^*) + \mathbf{D}_{\mathrm{K}''}(\Phi) - 2 \int u^* \frac{\partial\Phi}{\partial n}\,ds, \]
where the last integral is to be extended over the boundary of $\mathrm{K}''$ ($ds$ denotes the arc element of this boundary, and the boundary normal $n$ points into the interior of $\mathrm{K}''$). But since $\frac{\partial\Phi}{\partial n}$ on $\varkappa_0$ is equal to $0$, that integral needs only to be extended over that part $\varkappa''$ of the boundary of $\mathrm{K}$ which runs outside $\mathrm{K}_0$. Our assertion then reduces to
\[ \int_{\varkappa''} (\widetilde{u^*} - u^*)\,\frac{\partial\Phi}{\partial n}\,ds = 0. \]
In this form it is self-evident, since along the boundary of $\mathrm{K}\colon \widetilde{u^*} = u^*$. One sees from this that for the correctness of our formulation it was decisive to add to the singular component $\frac{x_0}{x_0^2 + y_0^2}$ the term $\frac{x_0}{a_0^2}$, which caused the normal derivative of $\Phi$ on the boundary of the hole to become equal to $0$.

Since we shall prove the validity of Dirichlet's principle for open Riemann surfaces as well, we can likewise construct $u$, $u^*$, and $U$ on such a surface. The conditions of competition for the functions $v$ are then further restricted by the additional requirement that $\mathbf{D}(v)$ be finite. On an open surface the potential function $U$ can no longer be completely characterized by its singularity at the point $O$. It is, however, uniquely determined up to an additive constant if, in addition to specifying the characteristic singularity at $O$, we include in its description this fact: For every continuously differentiable function $w$ which is $= 0$ for all points in a neighborhood (no matter how small) of $O$ and which possesses a finite Dirichlet integral, we have
\[ \mathbf{D}(U, w) = 0. \]
For if we draw around $O$ a $z_0$-disc $\mathrm{K}_0^{**}\colon |z_0| \leqq c_0$ situated in $\mathrm{K}_0$,

\origpage{96}
so small that everywhere in $\mathrm{K}_0^{**}$ the function $w = 0$, then from the inequality valid for every constant $\lambda$
\[ \mathbf{D}(u) \leqq \mathbf{D}(u + \lambda w) \]
there first follows the relation $\mathbf{D}(u, w) = 0$. But the difference
\[ \mathbf{D}(U, w) - \mathbf{D}(u, w) \]
is
\[ \begin{aligned}
&\iint_{c_0 \leqq |z_0| \leqq a_0} \left( \frac{\partial\Phi}{\partial x_0}\,\frac{\partial w}{\partial x_0} + \frac{\partial\Phi}{\partial y_0}\,\frac{\partial w}{\partial y_0} \right) dx_0\,dy_0 \\
&\qquad = \int_0^{2\pi} \left[ r_0 w\,\frac{\partial\Phi}{\partial r_0} \right]_{r_0 = c_0}^{r_0 = a_0} d\varphi_0.
\end{aligned} \]
On $\mathrm{K}_0$ we have $\frac{\partial\Phi}{\partial r_0} = 0$, on $\mathrm{K}_0^{**}\colon w = 0$; consequently
\[ \mathbf{D}(U, w) = \mathbf{D}(u, w) = 0. \]

The further extensions of the method of Dirichlet's principle to be discussed in this section we undertake only for closed surfaces.

Instead of $\frac{\cos\varphi_0}{r_0}$, I can also prescribe as singularity at $O$
\[ \frac{\sin\varphi_0}{r_0} = \frac{y_0}{x_0^2 + y_0^2} \]
or more generally
\[ \frac{\cos n\varphi_0}{r_0^n}, \quad \frac{\sin n\varphi_0}{r_0^n} \]
(pole of the $n^{\text{th}}$ order; $n$ an integer and positive). For $\Phi$ I then have to make the ansatz
\[ \Phi = \frac{\substack{\cos\\\sin} n\varphi_0}{r_0^n} + r_0^n \cdot \frac{\substack{\cos\\\sin} n\varphi_0}{a_0^{2n}} \]
in all cases $\left(\frac{\partial\Phi}{\partial r_0}\right)_{r_0 = a_0} = 0$. Each of the potential functions thus obtained gives rise to a differential $d\tau$ (cf.\ p.~55). This differential is everywhere regular except at $O$, and becomes infinite at $O$ like
\[ -\,\frac{n\,dz_0}{z_0^{n+1}}, \quad \text{resp.} -\; i\,\frac{n\,dz_0}{z_0^{n+1}}, \]
while the real part${}^{1)}$ of the integral $\displaystyle\int d\tau$ extended over any closed curve is always $0$. The differentials found are of the \textbf{second kind}, that is, they nowhere possess a residue different from $0$. Those possessing residues different from $0$ are designated as \textbf{differentials of the third kind}; but a differential is said to be of the \textbf{first kind} if it nowhere possesses a pole, let alone a residue different from $0$.

\textbf{1)} For the real and imaginary parts of a complex quantity $c = a + ib$ we shall use the notation $a = \Re c$, $b = \Im c$.

\origpage{97}
Let $1$ and $2$ be two points inside $\mathrm{K}_0$, and $1', 2'$ their mirror images${}^{1)}$ in the $z_0$-plane with respect to $\varkappa_0$ constructed by the method of reciprocal radii. Let $r_1, r_2; r'_1, r'_2$ be the lengths of the radius vectors from the points $1, 2; 1', 2'$ to a variable point of the $z_0$-plane, and $\varphi_1, \varphi_2, \varphi'_1, \varphi'_2$ the associated azimuths (angles with the direction of the positive $x_0$-axis, determined only up to integer multiples of $2\pi$). If we wish to find a potential function which is regular except at $1$ and $2$, but becomes infinite at $1$ like $\lg r_1$, and at $2$ like $-\lg r_2$, we take
\[ \Phi = (\lg r_1 - \lg r_2) + (\lg r'_1 - \lg r'_2). \]
The normal derivative of $\Phi$ on $\varkappa_0$ is $0$. The sealing ring is to be taken so narrow that the points $1', 2'$ do not come to lie in the cover. The resulting potential $U$ gives rise to a differential $d\omega_{12}$ which possesses at each of $1$ and $2$ a pole of the 1st order with residue $+ 1$ and $- 1$, respectively. If $\mathfrak{p}, \mathfrak{q}$ are any two points on $\mathfrak{F}$, one can interpolate along a curve connecting $\mathfrak{p}$ with $\mathfrak{q}$ such points $\mathfrak{p} = 1, 2, 3, \dots, \mathfrak{q} = n$ that
\[ d\omega_{12}, \, d\omega_{23}, \, \dots, \, d\omega_{n-1, n} \]
can be formed by our method.
\[ d\omega_{\mathfrak{p}\,\mathfrak{q}} = d\omega_{12} + d\omega_{23} + \dots + d\omega_{n-1, n} \]
is then a differential which possesses only at $\mathfrak{p}, \mathfrak{q}$ a pole of the 1st order with residues $+ 1$ and $- 1$, respectively. If $\mathfrak{p}_1, \mathfrak{p}_2, \dots, \mathfrak{p}_n$; $\mathfrak{p}_0$ are any $n + 1$ distinct points on the surface and $A_1, A_2, \dots, A_n$ arbitrary complex numbers with sum $0$, then
\[ \sum_{l=1}^n A_l\,d\omega_{\mathfrak{p}_l\,\mathfrak{p}_0} \]
is a differential that possesses at each of the points $\mathfrak{p}_1, \mathfrak{p}_2, \dots, \mathfrak{p}_n$ a pole of the 1st order with residues $A_1, A_2, \dots, A_n$, while it behaves regularly at all other points (including $\mathfrak{p}_0$). One sees that the residues of a differential regular except for poles can be prescribed arbitrarily, subject to the condition that their sum must be equal to $0$.

Let us return to the disc $\mathrm{K}_0$ with the two eccentrically situated points $1, 2$ within it.
\[ \Phi = (\varphi_1 - \varphi_2) - (\varphi'_1 - \varphi'_2) \]

\textbf{1)} We can assume $\mathrm{K}_0$ such that neither $1$ nor $2$ lies at the center of this disc.

\origpage{98}
is a single-valued regular potential function in the sealing ring (and even somewhat beyond its boundaries). According to Dirichlet's principle there therefore exists a regular potential function $u$ in the perforated surface and a regular potential function $u^*$ in the cover, which in the sealing ring stand in the relation $u = u^* + \Phi$ to one another. The associated $U$ is not single-valued, but acquires branch points of infinitely high order at $1$ and $2$. If, however, we cut open the surface along any continuous curve $\sigma_0$ leading within $\mathrm{K}_0$ from $1$ to $2$, $U$ is single-valued in the dissected surface. This follows immediately from the fact that with respect to a closed curve $\gamma$ lying in the cover which does not intersect $\sigma_0$, the points $1$ and $2$ in the $z_0$-plane have the same ``order'' (cf.\ p.~56). In spite of its multi-valuedness, $U$ gives rise to a differential $d\omega'_{12}$, single-valued on the undissected surface, which possesses poles of the 1st order only at $1$ and $2$ with residues $\mp i$. If the line $\sigma_0$ is an elementary segment on the triangulated surface $\mathfrak{F}$, then $\Re \displaystyle\int_\gamma d\omega'_{12} = 0$ will hold for every closed curve $\gamma$ which does not meet this segment; but for a polygon $\beta$ that crosses $\sigma_0$ at a single point from right to left, one will have
\[ \Re \int_\beta d\omega'_{12} = 2\pi. \]

If $\alpha$ is a closed curve on $\mathfrak{F}$, we can distribute points $1, 2, 3, \dots, n$ on $\alpha$ so densely that each individual subcurve $12, 23, \dots, n\,1$ into which $\alpha$ is thereby divided comes to lie entirely within a disc such as $\mathrm{K}_0$. By then forming
\[ d\omega_\alpha = \frac{1}{2\pi}\,(d\omega'_{12} + d\omega'_{23} + \dots + d\omega'_{n1}) \]
we obtain a differential $d\omega_\alpha$ regular \emph{everywhere without exception}; it has the property that for every closed curve $\gamma$ that does not meet $\alpha$, the real part of $\displaystyle\int_\gamma d\omega_\alpha = 0$. If, however, $\alpha$ was a polygon and $\beta$ is a polygon that crosses $\alpha$ at a single point from right to left, then
\[ \Re \int_\beta d\omega_\alpha = 1. \]

The canonical dissection serves to demonstrate that every arbitrary differential of the first kind can be generated by linear combination from the differentials $d\omega_\alpha$. Namely, if
\[ \alpha_1 + \alpha_2, \, \alpha_3 + \alpha_4, \, \dots, \, \alpha_{2p-1} + \alpha_{2p} \]
are the pairs of retrosections of a canonical dissection, and if for abbreviation we write $d\omega_{\alpha_h} = d\omega_h$, then

\origpage{99}
\[ \Re \int_{\alpha_{2l}} d\omega_{2l-1} = 1, \quad \Re \int_{\alpha_{2l-1}} d\omega_{2l} = -1 \quad [l = 1, \dots, p], \]
and all other expressions of the form $\Re \displaystyle\int_{\alpha_h} d\omega_i$ are $= 0$. If now the real part of $\displaystyle\int_{\alpha_h} d\omega$ is equal to $c_h$, then
\[ \Re \int \bigl[ d\omega - (c_2\,d\omega_1 - c_1\,d\omega_2) - (c_4\,d\omega_3 - c_3\,d\omega_4) - \dots \bigr] \tag{29} \]
vanishes ($= 0$) upon integration over any closed curve, since this ``integral function'' vanishes for the $2p$ curves $\alpha_h$ that form a basis of the closed paths. Consequently, (29), integrated from a fixed initial point to a variable endpoint, is a single-valued and indeed regular-harmonic function of this endpoint; but since no such function exists apart from constants, there results the identity:
\[ dw = (c_2\,d\omega_1 - c_1\,d\omega_2) + (c_4\,d\omega_3 - c_3\,d\omega_4) + \dots, \]
that is, \emph{from the $2p$ differentials $d\omega_h$, every differential of the first kind can be composed in one and only one way with the help of constant real factors: the $d\omega_h$ form a ``real'' basis of the differentials of the first kind.} --- The transition from \emph{one} real basis of the differentials of the first kind to another is mediated by linear transformation with arbitrary real coefficients and a determinant different from $0$. From the $2p$ differentials $d\omega_h$ of a real basis one can select such differentials (say $d\omega_1, \dots, d\omega_q$) that between them there exists no linear relation with arbitrary constant \emph{complex} coefficients, but indeed between any $q+1$ differentials of the basis. Since then
\[ d\omega_1, \dots, d\omega_q; \; i\,d\omega_1, \dots, i\,d\omega_q \tag{30} \]
are linearly independent in the real sense, we must have $2q \leqq 2p$. But since every differential of the basis, and thus every differential of the first kind in general, can be linearly composed with complex coefficients from $d\omega_1, \dots, d\omega_q$, hence linearly with real coefficients from the differentials (30), we must have $2q \geqq 2p$, and consequently $q=p$. From the differentials of a real basis one can accordingly extract a ``complex'' basis $d\omega_1, \dots, d\omega_p$ consisting of $p$ differentials. Every differential of the first kind $dw$ admits one and only one representation of the form
\[ dw = C_1\,d\omega_1 + \dots + C_p\,d\omega_p \]
with complex constants $C_1, \dots, C_p$.

If the power expansion of an arbitrary differential $dv$ at a point $\mathfrak{p}_1$, to which the local uniformizing variable $z_1$ belongs, is given by
\[ dv = \left( \frac{A_{-m}}{z_1^m} + \dots + \frac{A_{-1}}{z_1} + A_0 + A_1 z_1 + \dots \right) dz_1 \]

\origpage{100}
then we call
\[ \left( \frac{A_{-m}}{z_1^m} + \dots + \frac{A_{-1}}{z_1} \right) dz_1 \]
the \textbf{principal part} of $dv$ at the point $\mathfrak{p}_1$. If, on a closed Riemann surface, one prescribes arbitrarily the poles and the principal parts at these poles of a differential $dv$, but in such a way that the sum of the residues is $= 0$, and in addition the real part of $\displaystyle\int_\gamma dv$ for $2p$ linearly independent closed paths $\gamma$, then there exists one and only one differential $dv$ satisfying the stated requirements; it can be obtained by additive combination from the \textbf{elementary differentials} constructed above by us with the aid of Dirichlet's principle.

\subsection*{§ 15. Proof of Dirichlet's principle.${}^{1)}$}

Let the conditions of competition for the function $v$ be established as in the preceding section, from which we also adopt all other notations. It is first to be noted that competing functions exist at all (also in the case of non-closed surfaces), for example, those that are $= 0$ everywhere except at the points of $\mathrm{K}_0$. Let $d$ be the greatest lower bound of the Dirichlet integral $\mathbf{D}(v)$ for all competing functions; then we will always have $\mathbf{D}(v) \geqq d$, but, if $\varepsilon$ is an arbitrarily small positive quantity, there will still always exist competing functions $v$ for which $\mathbf{D}(v) < d + \varepsilon$. If $v$ is a competing function, then $v + a$, where $a$ is an arbitrary real constant, is also always such a function, which moreover imparts to the Dirichlet integral the same value as $v$. We can consequently further restrict the conditions of competition by the requirement that the boundary integral of $v^*$ around the circle $\mathrm{K}_0$
\[ \int_0^{2\pi} v^*_{r_0 = a_0} \cdot d\varphi_0 = 0 \]
shall hold. If to each competing function $v$ there is somehow assigned a number $\vartheta(v)$, a limit equation such as
\[ \lim_v \vartheta(v) = \vartheta_0 \]

\textbf{1)} In carrying out the proof, the paper by Zaremba cited on p.~94 was mainly taken as a basis. That paper, however, refers only to planar regions (namely to the solution of the first boundary value problem of potential theory for such regions). How Dirichlet's principle is to be grounded for arbitrary Riemann surfaces is shown by \emph{R.~König}, Konforme Abbildung der Oberfläche einer räumlichen Ecke, Leipzig Habilitationsschrift 1911, reprinted in Math.\ Ann.\ vol.~71, pp.~184---205. \emph{Our} formulation, differing from the Riemann-Hilbert one, brings extensive simplifications.

\origpage{101}
will mean that for every positive number $\varepsilon$ a positive number $\delta$ can be specified such that for all competing functions $v$ for which $\mathbf{D}(v) < d + \delta$, the difference $|\vartheta(v) - \vartheta_0| < \varepsilon$. The necessary and sufficient condition for a $\lim_v \vartheta(v)$ to exist in this sense is that for every positive $\varepsilon$ a positive $\delta$ can be found such that for any two competing functions $v_1, v_2$ that impart to the Dirichlet integral values $< d + \delta$,
\[ |\vartheta(v_1) - \vartheta(v_2)| < \varepsilon \]
holds.

The theorem we have to prove states: \emph{There exists a competing function $u$ with the property $\mathbf{D}(u) = d$.}

We begin this proof with the \emph{B.~Levi inequality}${}^{1)}$, which states that two competing functions $v_1, v_2$ whose Dirichlet integral comes close to the greatest lower bound $d$ possess a difference $v_1 - v_2$ whose Dirichlet integral is very small:
\[ \sqrt{\mathbf{D}(v_1 - v_2)} \leqq \sqrt{\mathbf{D}(v_1) - d} + \sqrt{\mathbf{D}(v_2) - d}. \]
Indeed, if $\lambda_1, \lambda_2$ are any two constants, $\lambda_1 + \lambda_2 \neq 0$, then with $v_1, v_2$ also
\[ \frac{\lambda_1 v_1 + \lambda_2 v_2}{\lambda_1 + \lambda_2} \]
is a competing function, and therefore
\[ \mathbf{D}(\lambda_1 v_1 + \lambda_2 v_2) \geqq d \cdot (\lambda_1 + \lambda_2)^2. \]
This inequality does not lose its validity for $\lambda_1 + \lambda_2 = 0$. The quadratic form in $\lambda_1, \lambda_2$:
\[ \lambda_1^2 [\mathbf{D}(v_1) - d] + 2\lambda_1 \lambda_2 [\mathbf{D}(v_1 v_2) - d] + \lambda_2^2 [\mathbf{D}(v_2) - d] \]
is therefore always $\geqq 0$, and for that reason
\[ [\mathbf{D}(v_1) - d][\mathbf{D}(v_2) - d] \geqq [\mathbf{D}(v_1 v_2) - d]^2 \]
must hold. But we have
\[ \begin{aligned}
0 \leqq \mathbf{D}(v_1 - v_2) &= \mathbf{D}(v_1) - 2\,\mathbf{D}(v_1 v_2) + \mathbf{D}(v_2) \\
&= [\mathbf{D}(v_1) - d] + [\mathbf{D}(v_2) - d] - 2\,[\mathbf{D}(v_1 v_2) - d] \\
&\leqq [\mathbf{D}(v_1) - d] + [\mathbf{D}(v_2) - d] + 2\sqrt{[\mathbf{D}(v_1) - d][\mathbf{D}(v_2) - d]} \\
&= \bigl\{ \sqrt{\mathbf{D}(v_1) - d} + \sqrt{\mathbf{D}(v_2) - d} \bigr\}^2.
\end{aligned} \]
From Levi's inequality one can already infer something similar to the fact that $v$ itself, if one varies it in such a way that $\mathbf{D}(v)$ is driven towards its greatest lower bound $d$, converges to a limit function which will then be the sought-after minimizing function. For if

\textbf{1)} B.~Levi, Sul principio di Dirichlet, Rendiconti del Circolo Matematico di Palermo, vol.~22 (1906), pp.~293---360, §~7.

\origpage{102}
$\mathfrak{p}$ is any point on the surface, $z = x + iy$ a local uniformizing variable for $\mathfrak{p}$, $\mathrm{K} : |z| \leqq a$ a $z$-disc, and $\mathfrak{E}$ an arbitrary closed set in $\mathrm{K}$ possessing a definite area in the $z$-plane, then the \emph{Schwarz} inequality applied to $\frac{\partial(v_1 - v_2)}{\partial x}$ and $1$ yields:
\[ \left| \iint_{\mathfrak{E}} \left\{ \frac{\partial v_1}{\partial x} - \frac{\partial v_2}{\partial x} \right\} dx\,dy \right| \leqq a\sqrt{\pi} \cdot \sqrt{\mathbf{D}(v_1 - v_2)}. \]
Thus there exist, uniformly for all regions $\mathfrak{E}$ situated in $\mathrm{K}$, the limit values
\[ \lim_v \iint_{\mathfrak{E}} \frac{\partial v}{\partial x}\,dx\,dy, \qquad \lim_v \iint_{\mathfrak{E}} \frac{\partial v}{\partial y}\,dx\,dy. \]
That is: although I cannot infer from the derivatives $\frac{\partial v}{\partial x}, \frac{\partial v}{\partial y}$ themselves that they converge, this does indeed hold for the \emph{areally integrated} derivatives. It is more convenient, instead of the derivatives, to return to the function $v$ itself; we will then see that in the same sense
\[ \lim_v \iint_{\mathfrak{E}} v\,dx\,dy \]
exists. If the \textbf{minimizing function} $u$ exists, this limit must moreover be $=\iint_{\mathfrak{E}} u\,dx\,dy$. If we take for $\mathfrak{E}$ the disc $\mathrm{K}$ itself and observe that $u$ will be a potential function, so that (at any rate, when $\mathrm{K}$ lies entirely in the perforated surface)
\[ \begin{aligned}
\frac{\iint\limits_{\mathrm{K}} u\,dx\,dy}{\pi a^2} &= \text{mean value of } u \text{ in } \mathrm{K} \\
&= u(\mathfrak{p})
\end{aligned} \]
holds, one is led to provide the desired existence proof by assigning the value
\[ \lim_v \frac{1}{\pi a^2} \iint_{\mathrm{K}} v\,dx\,dy \]
to the center of the disc as function value $u$, and then showing that the function thus obtained is indeed the minimizing function: Because the limit function $u$ will be a potential function, the convergence-generating integration over discs performed has no influence in the limit.

Proceeding to the execution of this line of thought, we must first replace the Dirichlet integral of $v_1 - v_2$ by the square integral of $v_1 - v_2$. For that purpose we show: If $\mathfrak{p}$ is a point, $z$ an associated

\origpage{103}
local uniformizing variable, $\mathrm{K} : |z| \leqq a$ a $z$-disc, then there exists a number $C$ such that for every continuously differentiable function $w$ on the surface whose boundary integral $\displaystyle\int_0^{2\pi} w\,d\varphi_0$ around $\mathrm{K}_0$ has the value $0$,
\[ \iint_{\mathrm{K}} w^2 dx\,dy \leqq C \cdot \mathbf{D}(w) \]
holds.

We form a chain of points $\mathfrak{p}_0 = O$, $\mathfrak{p}_1, \mathfrak{p}_2, \dots, \mathfrak{p}_n = \mathfrak{p}$ with associated local uniformizing variables $z_0, z_1, z_2, \dots, z_n = z$ and draw about each point $\mathfrak{p}_i$ a $z_i$-disc $\mathrm{K}_i : |z_i| \leqq a_i$ (about $\mathfrak{p}_0$ the disc $\mathrm{K}_0$, about $\mathfrak{p} = \mathfrak{p}_n$ the disc $\mathrm{K} = \mathrm{K}_n$); care is taken that $\mathrm{K}_{i+1}$ always overlaps $\mathrm{K}_i$ (that is, $\mathrm{K}_i$ and $\mathrm{K}_{i+1}$ have interior points in common). Then first of all, according to §~13,
\[ \iint_{\mathrm{K}_0} w^2 dx_0\,dy_0 \leqq \frac{a_0^2}{2} \iint_{\mathrm{K}_0} \left[ \left(\frac{\partial w}{\partial x_0}\right)^2 + \left(\frac{\partial w}{\partial y_0}\right)^2 \right] dx_0\,dy_0 \leqq \frac{a_0^2}{2}\,\mathbf{D}(w). \]
For $\mathrm{K}_0$ our assertion is thus true; we will conclude from this that it is likewise valid for $\mathrm{K}_1$. By the theorem just used, there is at any rate a constant $c$ such that
\[ \iint_{\mathrm{K}_1} (w - c)^2 dx_1\,dy_1 \leqq \frac{a_1^2}{2}\,\mathbf{D}(w) \]
holds. If $\mathfrak{E}_{01}$ denotes the point set formed by the points lying both in $\mathrm{K}_0$ and $\mathrm{K}_1$, then in $\mathfrak{E}_{01}$ the variables $z_1 = x_1 + i y_1$, $z_0 = x_0 + i y_0$ are connected by a transformation which is still regular analytic even beyond the boundary of $\mathfrak{E}_{01}$. The Jacobian
\[ \frac{\partial(x_1, y_1)}{\partial(x_0, y_0)} = \left| \frac{dz_1}{dz_0} \right|^2 \]
therefore has an upper bound $M_{01}$ in $\mathfrak{E}_{01}$, and we have:
\[ \begin{aligned}
\iint_{\mathfrak{E}_{01}} w^2 dx_1\,dy_1 &= \iint_{\mathfrak{E}_{01}} w^2\,\frac{\partial(x_1, y_1)}{\partial(x_0, y_0)}\,dx_0\,dy_0 \\
&\leqq M_{01} \iint_{\mathfrak{E}_{01}} w^2 dx_0\,dy_0 \leqq M_{01} \iint_{\mathrm{K}_0} w^2 dx_0\,dy_0 \leqq M_{01} \frac{a_0^2}{2}\,\mathbf{D}(w),
\end{aligned} \tag{31} \]
and on the other hand
\[ \iint_{\mathfrak{E}_{01}} (w - c)^2 dx_1\,dy_1 \leqq \frac{a_1^2}{2}\,\mathbf{D}(w). \tag{32} \]
If $J_{01}$ denotes the area of the image set of $\mathfrak{E}_{01}$ in the $z_1$-plane:

\origpage{104}
\[ J_{01} = \iint_{\mathfrak{E}_{01}} dx_1\,dy_1 > 0, \]
then addition of (31), (32) yields, on account of
\[ c^2 \leqq 2\,[w^2 + (c - w)^2]\colon \]
\[ c^2 \cdot J_{01} \leqq (a_0^2\,M_{01} + a_1^2)\,\mathbf{D}(w) \]
and then
\[ \begin{aligned}
\iint_{\mathrm{K}_1} w^2 dx_1\,dy_1 &\leqq 2 \left[ \iint_{\mathrm{K}_1} (w - c)^2 dx_1\,dy_1 + c^2 a_1^2\pi \right] \\
&\leqq a_1^2 \left( 1 + \frac{2\,\pi(a_1^2 + a_0^2\,M_{01})}{J_{01}} \right) \cdot \mathbf{D}(w).
\end{aligned} \]
With this the proof for $\mathrm{K}_1$ is accomplished. From $\mathrm{K}_1$ we now infer to $\mathrm{K}_2$ etc., and finally arrive at $\mathrm{K}_n = \mathrm{K}$. Applying the result to the difference of two competing functions $v_1, v_2$, we have
\[ \iint_{\mathrm{K}} (v_1 - v_2)^2 dx\,dy \leqq C \bigl\{ \sqrt{\mathbf{D}(v_1) - d} + \sqrt{\mathbf{D}(v_2) - d} \bigr\}^2. \tag{33} \]
The mean value
\[ \frac{1}{\pi a^2} \iint_{\mathrm{K}} f dx\,dy \]
of a function continuous in $\mathrm{K}$ I denote in general by $\underset{\mathrm{K},\,z}{\mathbf{M}} f$. I consider separately the points on the perforated surface and the points in the cover. If $\mathfrak{p}$ lies on the perforated surface, let $\mathrm{K}$ be taken such that $\mathrm{K}$ also lies entirely in the perforated surface. Then according to the last inequality there exists
\[ \lim_v \underset{\mathrm{K},\,z}{\mathbf{M}} v = u. \]
If $\mathfrak{p}$ is a point in the cover, $z_0(\mathfrak{p}) = c_0$, let $z_0 - c_0$ be adopted as local uniformizing variable $z$, and let the $z$-disc $\mathrm{K}$ lie entirely in the cover. There exists
\[ \lim_v \underset{\mathrm{K},\,z}{\mathbf{M}} v^* = u^*. \]
\emph{We assert: $u, u^*$ depend only on $\mathfrak{p}$, but not on the choice of the local uniformizing variable $z$ nor on the $z$-disc $\mathrm{K}$.}

I carry out the proof for points $\mathfrak{p}$ in the perforated surface. Let $z' = x' + iy'$ thus be another local uniformizing variable for $\mathfrak{p}$, $\mathrm{K}' : |z'| \leqq a'$ a $z'$-disc. I first assume that $\mathrm{K}'$ is situated entirely in $\mathrm{K}$. Let $\bar{v}$ be that potential function in $\mathrm{K}$ which coincides with $v$ on the boundary. We have $\mathbf{D}_{\mathrm{K}}(\bar{v}) \leqq \mathbf{D}_{\mathrm{K}}(v)$, and there exists a competing function $\tilde{v}$ which coincides with $v$ except in $\mathrm{K}$ and for which

\origpage{105}
\[ \mathbf{D}_{\mathrm{K}}(\tilde{v} - \bar{v}) = \mathbf{D}_{\mathrm{K}}(\tilde{v}) - \mathbf{D}_{\mathrm{K}}(\bar{v}) \]
is as small as one pleases. We have
\[ \mathbf{D}_{\mathrm{K}}(\tilde{v} - v) = \mathbf{D}(\tilde{v} - v) \leqq \bigl( \sqrt{\mathbf{D}(v) - d} + \sqrt{\mathbf{D}(\tilde{v}) - d} \bigr)^2; \]
consequently
\[ \mathbf{D}_{\mathrm{K}}(\bar{v} - v) = \mathbf{D}_{\mathrm{K}}(v) - \mathbf{D}_{\mathrm{K}}(\bar{v}) \leqq 4\,[\mathbf{D}(v) - d] \]
must hold. Since $\bar{v} - v$ vanishes on the boundary of the disc $\mathrm{K}$, it follows that
\[ \iint_{\mathrm{K}} (\bar{v} - v)^2 dx\,dy \leqq a^2 [\mathbf{D}(v) - d] \]
and by means of the \emph{Schwarz} inequality
\[ \left| \underset{\mathrm{K},\,z}{\mathbf{M}} v - \underset{\mathrm{K},\,z}{\mathbf{M}} \bar{v} \right| \leqq \frac{1}{\sqrt{\pi}} \sqrt{\mathbf{D}(v) - d}. \]
Since $\bar{v}$ is a potential function, $\underset{\mathrm{K},\,z}{\mathbf{M}} \bar{v} = \bar{v}(\mathfrak{p})$, and so we have:
\[ \lim_v \underset{\mathrm{K},\,z}{\mathbf{M}} v = \lim_v \bar{v}(\mathfrak{p}), \qquad \lim_v \bar{v}(\mathfrak{p}) = u. \]
In $\mathrm{K}'$, $z$ and $z'$ are related by a transformation analytic even beyond the boundaries of $\mathrm{K}'$; if $M$ is an upper bound for $\left|\frac{dz'}{dz}\right|^2$ in $\mathrm{K}'$, then
\[ \iint_{\mathrm{K}'} (\bar{v} - v)^2 dx'\,dy' \leqq a^2 M [\mathbf{D}(v) - d]\colon \]
\[ \lim_v \underset{\mathrm{K}',\,z'}{\mathbf{M}} \bar{v} = \lim_v \underset{\mathrm{K}',\,z'}{\mathbf{M}} v. \]
But since $\bar{v}$ is also a potential function in the variables $x', y'$, we must have $\underset{\mathrm{K}',\,z'}{\mathbf{M}} \bar{v} = \bar{v}(\mathfrak{p})$; therefore
\[ \lim_v \underset{\mathrm{K}',\,z'}{\mathbf{M}} v = \lim_v \bar{v}(\mathfrak{p}) = \lim_v \underset{\mathrm{K},\,z}{\mathbf{M}} v. \]

If $\mathrm{K}'$ does not lie in $\mathrm{K}$, we use as an intermediate link a $z'$-disc $\mathrm{K}'_*$ lying both in $\mathrm{K}'$ and in $\mathrm{K}$. Our reasoning then yields:
\[ \text{1)}\ \lim_v \underset{\mathrm{K},\,z}{\mathbf{M}} v = \lim_v \underset{\mathrm{K}'_*,\,z'}{\mathbf{M}} v; \qquad \text{2)}\ \lim_v \underset{\mathrm{K}',\,z'}{\mathbf{M}} v = \lim_v \underset{\mathrm{K}'_*,\,z'}{\mathbf{M}} v, \]
from which also\ednote{Misprint in the original: K, z and K', z' appear in print under lim instead of under M; the variable symbol v under lim is missing.}
\[ \lim_{\mathrm{K},\,z} \mathbf{M} v = \lim_{\mathrm{K}',\,z'} \mathbf{M} v \]
results in this case. If one wishes to carry out the analogous inferences for the points of the cover, one must in the case where $\mathrm{K}$ reaches across the boundary of the hole take into account the relation resulting from the discussion on p.~95,

\origpage{106}
\[ \mathbf{D}_{\mathrm{K}}(v) - \mathbf{D}_{\mathrm{K}}(\tilde{v}) = \mathbf{D}_{\mathrm{K}}(v^*) - \mathbf{D}_{\mathrm{K}}(\tilde{v}^*) \]
by which the equation $\left(\frac{\partial\Phi}{\partial r_0}\right)_{r_0 = a_0} = 0$ is utilized.

We assign the number $u$, resp.\ $u^*$, as function value to the point $\mathfrak{p}$; we then have a function $u(\mathfrak{p})$ defined in the perforated surface and a function $u^*(\mathfrak{p})$ in the cover. To investigate the function $u(\mathfrak{p})$ in the perforated surface, I consider all points $\mathfrak{q}$ that lie in the $z$-disc $\mathrm{K}$ about $\mathfrak{p}$, and carry out the following considerations on the image of $\mathrm{K}$ in the $z$-plane (in which $\mathfrak{p}$ is the origin). Around $\mathfrak{q}$ I draw (in the $z$-plane) the disc $\mathrm{K}_\mathfrak{q}$ tangent to $\mathrm{K}$ from the inside; its radius is $q = a - |z(\mathfrak{q})|$. If $v_1$ is a second competing function alongside $v$ and $\bar{v}$ has the previous meaning, formula (33), in which we replace $v_2$ by $\bar{v}$, yields
\[ \iint_{\mathrm{K}_\mathfrak{q}} (\bar{v} - v_1)^2 dx\,dy \leqq C \bigl( \sqrt{\mathbf{D}(v) - d} + \sqrt{\mathbf{D}(v_1) - d} \bigr)^2, \]
hence
\[ \left| \bar{v}(\mathfrak{q}) - \underset{\mathrm{K}_\mathfrak{q}}{\mathbf{M}} v_1 \right| \leqq \frac{\sqrt{C}}{q\sqrt{\pi}} \bigl( \sqrt{\mathbf{D}(v) - d} + \sqrt{\mathbf{D}(v_1) - d} \bigr). \]
In this I carry out with $v_1$ the passage to the limit $\lim\mathbf{D}(v_1) = d$:
\[ |u(\mathfrak{q}) - \bar{v}(\mathfrak{q})| \leqq \frac{\sqrt{C}}{q\sqrt{\pi}} \sqrt{\mathbf{D}(v) - d}. \]
Thus $\bar{v}(\mathfrak{q})$ converges \emph{uniformly to $u(\mathfrak{q})$ in any concentric disc lying entirely within $\mathrm{K}$}. Accordingly, $u$ is a potential function at $\mathfrak{p}$ (p.~83), and also the first derivatives\ednote{Misprint in the original: ``Abteilungen'' instead of ``Ableitungen''.} of $\bar{v}$ converge in the same sense uniformly to the respective derivatives of $u$.

Likewise $u^*$ is a regular potential function in the cover. If $\mathfrak{p}$ denotes a point inside the sealing ring and we draw in the $z_0$-plane around $\mathfrak{p}$ a disc $\mathrm{K}$ lying \emph{entirely in the sealing ring}, we have
\[ \begin{gathered}
u(\mathfrak{p}) = \lim_v \underset{\mathrm{K}}{\mathbf{M}} v; \\
u^*(\mathfrak{p}) = \lim_v \underset{\mathrm{K}}{\mathbf{M}} v^* = \lim_v \underset{\mathrm{K}}{\mathbf{M}} v - \underset{\mathrm{K}}{\mathbf{M}} \Phi.
\end{gathered} \]
Since $\Phi$ is a regular potential function in the sealing ring, the equation $\underset{\mathrm{K}}{\mathbf{M}} \Phi = \Phi(\mathfrak{p})$ holds; we thus conclude
\[ u(\mathfrak{p}) = u^*(\mathfrak{p}) + \Phi(\mathfrak{p}). \]
\emph{We have thus succeeded in constructing a regular-harmonic function $u$ in the perforated surface and a regular-harmonic function $u^*$ in the cover which in the sealing ring stand in the relation $u = u^* + \Phi$ to one another.} It remains for us to demonstrate that the function

\origpage{107}
$u$ which is $= u(\mathfrak{p})$ in the perforated surface and $= u^*(\mathfrak{p})$ in the hole satisfies the minimum equation $\mathbf{D}(u) = d$ --- although for applications to the theory of \emph{closed} Riemann surfaces the properties of $u$, $u^*$ established so far would obviously suffice. For the uniformization theory, however, with which the final sections of this work will deal, the minimal property is essential.

The dissection of the surface by the circumferences $\varkappa_l$ of discs $\mathrm{K}_l$ necessary for the calculation of the Dirichlet integral can be arranged in such a way that each disc $\mathrm{K}_l$ either lies entirely in the perforated surface or belongs entirely to the cover (and then appears as a disc in the $z_0$-plane). Let $\mathfrak{E}_h$ be one of the pieces into which the surface is thereby dissected and which belongs to a disc $\mathrm{K}_l$ of the first kind. We use the uniformizing variable $z_l = z = x + iy$, with respect to which $\mathrm{K}_l$ is a disc, for the uniformization of $\mathfrak{E}_h$. Let $\mathrm{K}'_l$ be a $z_l$-disc somewhat larger than $\mathrm{K}_l$ but still lying entirely in the perforated surface, and let $\bar{v}$, for each competing function $v$, be that potential function in $\mathrm{K}'_l$ which coincides with $v$ on the boundary of $\mathrm{K}'_l$. Then $\frac{\partial\bar{v}}{\partial x}$, $\frac{\partial\bar{v}}{\partial y}$ converge uniformly in $\mathrm{K}_l$ to $\frac{\partial u}{\partial x}$, $\frac{\partial u}{\partial y}$; hence
\[ \lim_v \mathbf{D}_{\mathrm{K}_l}(\bar{v} - u) = 0. \]

On the other hand,
\[ \mathbf{D}_{\mathrm{K}_l}(v - \bar{v}) \leqq \mathbf{D}_{\mathrm{K}'_l}(v - \bar{v}) \leqq 4[\mathbf{D}(v) - d], \]
hence
\[ \begin{gathered}
\lim_v \mathbf{D}_{\mathrm{K}_l}(v - u) = 0, \\
\lim_v \mathbf{D}_{\mathfrak{E}_h}(v - u) = 0, \\
\mathbf{D}_{\mathfrak{E}_h}(u) = \lim_v \mathbf{D}_{\mathfrak{E}_h}(v).
\end{gathered} \]

The same is proved for a piece $\mathfrak{E}_h$ belonging to a disc of the second kind [initially in the form
\[ \lim_v \mathbf{D}_{\mathfrak{E}_h}(v^* - u^*) = 0]. \]

If one forms $\displaystyle\sum\mathbf{D}_{\mathfrak{E}_h}(u)$ over any finite number of pieces $\mathfrak{E}_h$, this sum is
\[ = \lim_v \sum\mathbf{D}_{\mathfrak{E}_h}(v) \leqq \lim_v \mathbf{D}(v) = d. \]
Therefore $\mathbf{D}(u)$ is finite and $\leqq d$. Since it cannot be $< d$, we must have $\mathbf{D}(u) = d$, q.~e.~d.

\origpage{108}

\subsection*{§ 16. Relations between the differentials on a closed Riemann surface.}

Instead of setting up the elementary differentials of the 2nd and 3rd kind independently of one another with the help of Dirichlet's principle as in § 14, one can also derive all differentials regular except for poles that belong to a given closed Riemann surface $\mathfrak{F}$ from the simplest of those denoted by $d\tau$ in § 14,${}^{1)}$ and in this way one gains at the same time a very complete overview of the relations existing between the individual differentials.

Let $q$ be a point on $\mathfrak{F}$, and $\zeta = \xi + i\eta$ a local uniformizing variable belonging to $q$. The potential function which becomes infinite only at $q$ like
\[ \Re \frac{1}{\zeta} = \frac{\xi}{\xi^2 + \eta^2}, \]
we denote --- as the real part of the Abelian integral of the 2nd kind $\tau_q$ which is single-valued only on the covering surface $\tilde{\mathfrak{F}}$ and possesses a pole of the 1st order at $q$ --- by $\Re\tau_q$. The meaning of this function symbol accordingly depends not only on the point $q$, but also on the choice of the local uniformizing variable $\zeta$ belonging to $q$; where this is to be emphasized, we write $q\,\zeta$ instead of the index $q$ alone. If we use $i\zeta$ in place of $\zeta$ as local uniformizing variable, we obtain that potential function
\[ \Re\tau'_q = \Re\tau'_{q\,\zeta} = \Re\tau_{q,\,i\zeta}, \]
which becomes infinite only at the point $q$ and indeed like $\frac{-\eta}{\xi^2 + \eta^2}$.

To investigate the dependence of these functions on the point of infinity $q$, we proceed as follows: We delimit around $q$ a $\zeta$-disc $\mathrm{K}\colon |\zeta| \leqq \alpha$, denote by $q^*$ a point variable in $\mathrm{K}$, and understand by
\[ \Re\tau_{q^*}, \quad \Re\tau'_{q^*} \]
those potential functions becoming infinite at $q^*$ that correspond to the local uniformizing variable $\zeta - \zeta(q^*)$ at $q^*$. We shall presently show that these functions --- with an appropriate normalization of the additive constants still at our disposal for each position of $q^*$ --- are continuous with respect to $q^*$; for points $\mathfrak{p}$ lying in the neighborhood $\mathrm{K}$ of $q$, we will have
\[ \Re\tau_{q^*}(\mathfrak{p}) = \Re \frac{1}{\zeta(\mathfrak{p}) - \zeta(q^*)} + R_{q^*}(\mathfrak{p}), \quad \Re\tau'_{q^*}(\mathfrak{p}) = \Re \frac{-i}{\zeta(\mathfrak{p}) - \zeta(q^*)} + R'_{q^*}(\mathfrak{p}) \tag{34} \]
where $R$, $R'$, as long as $q^*$ and $\mathfrak{p}$ lie in $\mathrm{K}$, are continuous functions of $q^*$ and entirely regular potential functions in $\mathfrak{p}$. Let $q_1$, $q_2$ be two points within $\mathrm{K}$ connected by a rectifiable path of integration $\gamma$

\textbf{1)} Cf.\ Klein-Fricke, \emph{Theorie der elliptischen Modulfunktionen}, vol.~I, pp.~518\,ff.

\origpage{109}
within $\mathrm{K}$. By \emph{longitudinal} juxtaposition of double sources along $\gamma$, that is, by juxtaposition of double sources of constant moment whose direction coincides with the direction of $\gamma$, we obtain in
\[ \int\limits_{q_1}^{q_2} \{ \Re\tau_{q^*}(\mathfrak{p})\,d\xi_{q^*} - \Re\tau'_{q^*}(\mathfrak{p})\,d\eta_{q^*} \} = \Re\omega_{q_1 q_2}(\mathfrak{p}) \quad [\xi_{q^*} + i\eta_{q^*} = \zeta(q^*)] \tag{35} \]
a potential function which possesses at each of $q_1$, $q_2$ a sink and a source, respectively, that is, a logarithmic singularity with residues $+1$ and $-1$, respectively [formation of a magnet from elementary magnets!], and thus the Abelian integral of the 3rd kind $\omega_{q_1 q_2}$. For, for points $\mathfrak{p}$ not lying in $\mathrm{K}$, the left-hand side of (35) is obviously a regular potential function; but if $\mathfrak{p}$ lies in $\mathrm{K}$, it is by (34), apart from an additively entering regular potential function, $= \lg r_1 - \lg r_2$, where $r_1, r_2$ denote the rectilinear distances from $q_1, q_2$ to $\mathfrak{p}$ in the $\zeta$-plane. --- By \emph{transverse} juxtaposition of the double sources we obtain in
\[ \int\limits_{q_1}^{q_2} \{ \Re\tau_{q^*}(\mathfrak{p})\,d\eta_{q^*} + \Re\tau'_{q^*}(\mathfrak{p})\,d\xi_{q^*} \} = \Re\omega'_{q_1 q_2}(\mathfrak{p}) \tag{35'} \]
a potential function that has at $q_1, q_2$ two vortex points with opposite sense of rotation [equivalence of a magnetic double layer with an electric current!], which namely for points $\mathfrak{p}$ within $\mathrm{K}$, apart from an additive regular potential function, is equal to the angle that the straight connecting lines $q_1\mathfrak{p}$, $q_2\mathfrak{p}$ form with each other in the $\zeta$-plane. The function $\Re\omega'_{q_1 q_2}$ is single-valued only on the surface $\mathfrak{F}$ cut open along the integration path $\gamma$.

For the complete justification of this derivation, the continuity of $\Re\tau_{q^*}$, $\Re\tau'_{q^*}$ with respect to the ``parameter'' $q^*$, which is essential for the possibility of the integrations performed in (35), (35'), remains to be proved. For this we must return to the construction of these potential functions. Understanding by $\bar{\zeta}(q^*)$ the reflection of $\zeta(q^*)$ in the $\zeta$-plane with respect to the circle $\mathrm{K}$, we form
\[ \Phi_{q^*}(\zeta) = \Re \left[ \frac{1}{\zeta - \zeta(q^*)} - \left( \frac{\bar{\zeta}(q^*)}{\alpha} \right)^2 \left( \frac{1}{\zeta - \bar{\zeta}(q^*)} + \frac{1}{\bar{\zeta}(q^*)} \right) \right] \]
and set
\[ u_{q^*} = \begin{cases}
\Re\tau_{q^*} \text{ outside } \mathrm{K} \\
\Re\tau_{q^*} - \Phi_{q^*} \text{ inside } \mathrm{K}.
\end{cases} \]
Since $\Phi_{q^*}$ has normal derivative $0$ on the boundary of $\mathrm{K}$, $u_{q^*} - u_q$ imparts to the Dirichlet integral $\mathbf{D}(v)$ extended over the whole of $\mathfrak{F}$ its smallest value among all functions $v$ that have on the boundary of $\mathrm{K}$ the jump $\Phi_{q^*} - \Phi_q$, but are otherwise continuous. As a comparison function can serve

\origpage{110}
in particular that function $v$ which is equal to $0$ outside $\mathrm{K}$, but inside $\mathrm{K}$ is defined as the regular potential function coinciding in its boundary values with $\Phi_q - \Phi_{q^*}$ :
\[ -v = \frac{2}{\alpha^2} \, \Re \left[ \frac{\bar{\zeta}^2(q^*)}{\bar{\zeta}(q^*) - \zeta} - \zeta - \bar{\zeta}(q^*) \right]\colon \]
\[ \mathbf{D}(u_{q^*} - u_q) \leqq \mathbf{D}(v) = \frac{4\pi}{\alpha^2} \frac{q^2(2 - q^2)}{(1 - q^2)^2} \left( q = \frac{|\zeta(q^*)|}{\alpha} \right). \tag{36} \]

The Dirichlet integral of $u_{q^*} - u_q$ thus converges to $0$ as $q^*$ approaches $q$; from this results the asserted continuity. If $\mathfrak{p} \neq q$, we take $\mathrm{K}$ so small that $\mathfrak{p}$ comes to lie outside $\mathrm{K}$, and choose a local uniformizing variable $z$ for $\mathfrak{p}$ and a disc $K\colon |z| \leqq a$ about $\mathfrak{p}$ that has no point in common with $\mathrm{K}$. If $\mathfrak{p}^*$ is a variable point in $K$, $z(\mathfrak{p}^*) = x + iy$, $\frac{1}{a}\,|z(\mathfrak{p}^*)| = r$, then from (36) and (22) it follows:
\[ \left( \frac{\partial u_{q^*}}{\partial x} - \frac{\partial u_q}{\partial x} \right)^2 + \left( \frac{\partial u_{q^*}}{\partial y} - \frac{\partial u_q}{\partial y} \right)^2 < 8 \left[ \frac{q}{a\alpha(1 - q^2)(1 - r)} \right]^2, \]
that is,
\[ \left| \frac{d\tau_{q^*}}{dz} - \frac{d\tau_q}{dz} \right| < 2\sqrt{2} \, \frac{q}{a\alpha(1 - q^2)(1 - r)}, \]
hence
\[ \lim_{q^* = q} \frac{d\tau_{q^*}}{dz} = \frac{d\tau_q}{dz} \tag{37} \]
uniformly with respect to $\mathfrak{p}^*$ if $\mathfrak{p}^*$ is restricted to a concentric disc somewhat smaller than $K$ ($r \leqq r_0 < 1$). --- If $\mathfrak{p}$ lies in $\mathrm{K}$, then from (36) it follows:
\[ \left( \frac{\partial u_{q^*}(\zeta)}{d\xi} - \frac{\partial u_q(\zeta)}{\partial\xi} \right)^2 + \left( \frac{\partial u_{q^*}(\zeta)}{\partial\eta} - \frac{\partial u_q(\zeta)}{\partial\eta} \right)^2 < 8 \left[ \frac{q}{\alpha^2(1 - q^2)(1 - r)} \right]^2, \quad \left( r = \frac{|\zeta|}{\alpha} \right), \]
consequently\ednote{Misprint in the original: $d\xi$ instead of $\partial\xi$ in the denominator of the first term of the above formula.} uniformly for $|\zeta| \leqq r_0\alpha$ ($r_0$ a fixed number $< 1$):
\[ \lim_{q^* = q} \frac{\partial u_{q^*}(\zeta)}{\partial\xi}, \, \frac{\partial u_{q^*}(\zeta)}{\partial\eta} = \frac{\partial u_q(\zeta)}{\partial\xi}, \, \frac{\partial u_q(\zeta)}{\partial\eta}. \]

If we carry out the normalization of the additive constants by the equation $R_{q^*}(q^*) = 0$, there results from this by integration:

1) $\lim\limits_{q^* = q} R_{q^*}(\mathfrak{p}) = R_q(\mathfrak{p})$ uniformly for all $\mathfrak{p}$ in a certain neighborhood of $q$;

2) $\lim\limits_{q^* = q} \Re\tau_{q^*}(\mathfrak{p}) = \Re\tau_q(\mathfrak{p})$ uniformly for all $\mathfrak{p}$ on $\mathfrak{F}$, if an arbitrarily small neighborhood of the point $q$ is excluded.

With this the asserted continuity properties are established for the (arbitrary) point $q^* = q$.

We combine (35), (35') into the single equation

\origpage{111}
\[ \int\limits_{q_1}^{q_2} \{ \Re\tau_{q^*}(\mathfrak{p}) + i\,\Re\tau'_{q^*}(\mathfrak{p}) \} \, d\zeta(q^*) = \Re\omega_{q_1 q_2}(\mathfrak{p}) + i\,\Re\omega'_{q_1 q_2}(\mathfrak{p}). \tag{38} \]

Let $\mathfrak{p}$ lie outside $\mathrm{K}$, and we use the symbols $K$, $\mathfrak{p}^*$, and $z = z(\mathfrak{p}^*) = x + iy$ as above. If $f(z)$ is an analytic function of the complex variable $z$, then
\[ \frac{\partial}{\partial x}\,\Re f(z) = \Re \frac{df}{dz}, \quad \frac{\partial}{\partial y}\,\Re f(z) = -\,\Im \frac{df}{dz}. \]

If we thus differentiate equation (38), in which $\mathfrak{p}$ is first replaced by the variable point $\mathfrak{p}^*$, with respect to $x$ and $y$, we obtain:
\[ \int\limits_{q_1}^{q_2} \left\{ \Re \frac{d\tau_{q^*}(z)}{dz} + i\,\Re \frac{d\tau'_{q^*}(z)}{dz} \right\} d\zeta(q^*) = \Re \frac{d\omega_{q_1 q_2}(z)}{dz} + i\,\Re \frac{d\omega'_{q_1 q_2}(z)}{dz}, \tag{39} \]
\[ \int\limits_{q_1}^{q_2} \left\{ \Im \frac{d\tau_{q^*}(z)}{dz} + i\,\Im \frac{d\tau'_{q^*}(z)}{dz} \right\} d\zeta(q^*) = \Im \frac{d\omega_{q_1 q_2}(z)}{dz} + i\,\Im \frac{d\omega'_{q_1 q_2}(z)}{dz}. \tag{39\,i} \]

That the differentiation under the integral sign is permissible follows from the fact that $\frac{d\tau_{q^*}(z)}{dz}$, $\frac{d\tau'_{q^*}(z)}{dz}$ depend continuously on $q^*$ according to (37).

Now we carry out the essential inference of this investigation: from (39) it follows: the integrand
\[ \Re \frac{d\tau_{q^*}(z)}{dz} + i\,\Re \frac{d\tau'_{q^*}(z)}{dz} \]
is in $\mathrm{K}$ such a function of $\zeta = \zeta(q^*)$ that its integral with respect to $\zeta$ is independent of the path, that is, \emph{is a regular analytic function of $\zeta$}. In place of $\frac{d}{dz}$ I also write $\frac{d}{d\mathfrak{p}^*}$. On the local uniformizing variable $\zeta$ belonging to $q$,
\[ \Re \frac{d\tau_{q\,\zeta}(z)}{dz} + i\,\Re \frac{d\tau'_{q\,\zeta}(z)}{dz} \]
depends, as one sees directly or concludes from (39), in the manner in which the value of a differential at a point depends on the local uniformizing variable used. There thus exists a differential $dv_{\mathfrak{p}^*}(q)$, regular-analytic for all $q$ on $\mathfrak{F}$ except for $q = \mathfrak{p}^*$, such that
\[ \Re \frac{d\tau_{q^*}(z)}{dz} + i\,\Re \frac{d\tau'_{q^*}(z)}{dz} = \frac{dv_{\mathfrak{p}^*}(\zeta)}{d\zeta} \left( = \frac{dv_{\mathfrak{p}^*}(q^*)}{dq^*} \right) \]
holds. (39) now writes itself as:

\origpage{112}
\[ \int\limits_{q^* = q_1}^{q_2} dv_{\mathfrak{p}^*}(q^*) = \Re \frac{d\omega_{q_1 q_2}(\mathfrak{p}^*)}{d\mathfrak{p}^*} + i\,\Re \frac{d\omega'_{q_1 q_2}(\mathfrak{p}^*)}{d\mathfrak{p}^*}. \]

This formula holds for an arbitrarily extended integration path $\gamma = (q_1\,q_2)$, which now no longer needs to be restricted to the neighborhood $\mathrm{K}$ of $q$; only $d\omega'_{q_1 q_2}$ must be understood as that differential of the 3rd kind which arises in the manner described in § 14 by interpolation of points and summation along $\gamma$. On the other hand, $d\omega_{q_1 q_2}$ is independent of the manner of interpolation, and therefore $\Re \displaystyle\int\limits_{q_1}^{q_2} dv_{\mathfrak{p}^*}$ is independent of the integration path, or $\Re v_{\mathfrak{p}^*}$ is a single-valued potential function on $\mathfrak{F}$ which becomes singular only at the point $\mathfrak{p}^*$. The behavior at the point $\mathfrak{p}^*$ is surveyed if, for points $q^*$ lying in the disc $K$ about $\mathfrak{p}$, one uses $z - z(q^*)$ as local uniformizing variable:
\[ \Re\tau_{q^*}(z) = \Re \frac{1}{z - z(q^*)} + R_{q^*}(z), \]
where $R_{q^*}(z)$, $\frac{\partial R_{q^*}(z)}{\partial x}$, $\frac{\partial R_{q^*}(z)}{\partial y}$ are continuous functions of $q^*$, even at the point $z$; therefore, apart from an additive term that remains between fixed bounds as $q^*$ approaches $\mathfrak{p}^*$ and is consequently a regular function of $q^*$ even for $q^* = \mathfrak{p}^*$,
\[ \Re \frac{d\tau_{q^*}(\mathfrak{p}^*)}{d\mathfrak{p}^*} + i\,\Re \frac{d\tau'_{q^*}(\mathfrak{p}^*)}{d\mathfrak{p}^*} \sim \frac{-1}{[z(\mathfrak{p}^*) - z(q^*)]^2} \]
that is, $dv_{\mathfrak{p}^*}(q^*)$ is a differential of the 2nd kind whose principal part at the point $\mathfrak{p}^*$ is given by
\[ \frac{-\,dz(q^*)}{[z(q^*) - z(\mathfrak{p}^*)]^2} \]
and whose integral possesses a single-valued real part on $\mathfrak{F}$:
\[ dv_{\mathfrak{p}^*} = d\tau_{\mathfrak{p}^*} = d\tau_{\mathfrak{p}^* z}. \]

We summarize the results obtained. \emph{Let $\mathfrak{p}$, $q$ be two distinct points on $\mathfrak{F}$, $z$ a local uniformizing variable for $\mathfrak{p}$, $\zeta$ one for $q$; these local uniformizing variables serve to uniquely fix the function symbols $\tau_\mathfrak{p}$, $\tau_q$ and the differentiation symbols $\frac{d}{d\mathfrak{p}}$, $\frac{d}{dq}$. Then the fundamental reciprocity laws hold:}
\[ \left\{ \begin{aligned}
\Re \frac{d\tau_q(\mathfrak{p})}{d\mathfrak{p}} &= \Re \frac{d\tau_\mathfrak{p}(q)}{dq}, & \Re \frac{d\tau'_q(\mathfrak{p})}{d\mathfrak{p}} &= \Im \frac{d\tau_\mathfrak{p}(q)}{dq}, \\
\Im \frac{d\tau_q(\mathfrak{p})}{d\mathfrak{p}} &= \Re \frac{d\tau'_\mathfrak{p}(q)}{dq}, & \Im \frac{d\tau'_q(\mathfrak{p})}{d\mathfrak{p}} &= \Im \frac{d\tau'_\mathfrak{p}(q)}{dq}.
\end{aligned} \right. \tag{I} \]

\origpage{113}
\emph{The differentials of the 3rd kind arise from those of the 2nd kind according to the formulas:}
\[ \left\{ \begin{aligned}
\Re \int\limits_{q_1}^{q_2} d\tau_\mathfrak{p} &= \Re \frac{d\omega_{q_1 q_2}(\mathfrak{p})}{d\mathfrak{p}}, & \Im \int\limits_{q_1}^{q_2} d\tau_\mathfrak{p} &= \Re \frac{d\omega'_{q_1 q_2}(\mathfrak{p})}{d\mathfrak{p}}, \\
\Re \int\limits_{q_1}^{q_2} d\tau'_\mathfrak{p} &= \Im \frac{d\omega_{q_1 q_2}(\mathfrak{p})}{d\mathfrak{p}}, & \Im \int\limits_{q_1}^{q_2} d\tau'_\mathfrak{p} &= \Im \frac{d\omega'_{q_1 q_2}(\mathfrak{p})}{d\mathfrak{p}}.
\end{aligned} \right. \tag{II} \]

We obtain a special case of the equations in the second column when we let $(q_1\,q_2)$ expand into a closed path $\alpha$:
\[ \int\limits_\alpha d\tau_\mathfrak{p} = 2\,\pi i\,\Re \frac{dw_\alpha(\mathfrak{p})}{d\mathfrak{p}}, \quad \int\limits_\alpha d\tau'_\mathfrak{p} = 2\,\pi i\,\Im \frac{dw_\alpha(\mathfrak{p})}{d\mathfrak{p}}. \tag{$\mathrm{II}_0$} \]

\emph{They show that the differentials of the 1st kind are the periods of the elementary integrals of the 2nd kind.}

We further read off from them the following three facts ($\alpha$, $\beta$, $\gamma$ denote closed curves on $\mathfrak{F}$):
\[ \begin{aligned}
&1. & &\text{if } \alpha \sim 0, & &\text{then } dw_\alpha \text{ is identically } = 0; \\
&2. & &\text{if } \alpha \sim \beta, & &\text{then } dw_\alpha = dw_\beta; \\
&3. & &\text{if } \alpha + \beta \sim \gamma, & &\text{then } dw_\alpha + dw_\beta = dw_\gamma.
\end{aligned} \]

To each deck transformation $S$ of the covering surface $\hat{\mathfrak{F}}$ of the integral functions into itself corresponds a class of mutually homologous closed paths $\alpha$ on $\mathfrak{F}$: they are the traces of all curves on $\hat{\mathfrak{F}}$ that lead from any point $\hat{\mathfrak{p}}$ to the point $\hat{\mathfrak{p}}S$ situated above it. To all these $\alpha$ there belongs, by the 2nd of the facts just noted, the same differential $dw_\alpha$, which, since it is uniquely determined by $S$, we may fittingly denote also by $dw_S$. If we extend the integrations to be performed from an initial point $\mathfrak{p}_0$ chosen once and for all, then
\[ -\,\Re \int\limits_{\mathfrak{p}_0}^\mathfrak{p} dw_S = u_S(\mathfrak{p}) \]
depends on the path of integration used only in so far as a change of this path (with fixed endpoint $\mathfrak{p}$) can increase or decrease $u_S$ by an integer. (And by this property the integrals $w_S$ are also distinguished within the totality of all integrals of the 1st kind.) The quantity
\[ \chi_S(\mathfrak{p}) = e^{2\pi i u_S(\mathfrak{p})} \]

\origpage{114}
of absolute value 1 is thus a single-valued function of $\mathfrak{p}$. If $S$, $T$ are any two deck transformations of $\hat{\mathfrak{F}}$ into itself, then
\[ \chi_S(\mathfrak{p}) \cdot \chi_T(\mathfrak{p}) = \chi_{S\,T}(\mathfrak{p}). \]

In general: if in any way to each deck transformation $S$ of $\hat{\mathfrak{F}}$ there is assigned a number $\chi_S$ of absolute value 1 such that the multiplication rule
\[ \chi_S \cdot \chi_T = \chi_{S\,T} \]
holds, one says that a \textbf{system of characters} (of the commutative group of all deck transformations of $\hat{\mathfrak{F}}$ into itself) is defined. Accordingly we will call the quantities $\chi_S(\mathfrak{p})$ corresponding to the various transformations $S$ the system of \textbf{integral characters} of $\mathfrak{p}$. \emph{To give all the integral characters of $\mathfrak{p}$ means as much as: simultaneously to give the values of all integrals of the 1st kind at the point $\mathfrak{p}$ --- or more precisely: to give what is invariant in this simultaneous system of values under deformation of the integration path.} --- For a later application we further note that, if one considers a system of characters as an individual, a ``point'', the totality of all possible systems of characters forms a closed $2p$-dimensional\ednote{Misprint in the original: ``dimensonale'' instead of ``dimensionale''.} manifold. For if
\[ S_1, \, S_2, \, \dots, \, S_{2p} \]
is a basis for the group of deck transformations of $\hat{\mathfrak{F}}$ (cf.\ p.~74), then one obtains all systems of characters once and only once if one lets
\[ \chi_{S_1}, \, \chi_{S_2}, \, \dots, \, \chi_{S_{2p}} \]
traverse independently of one another in a $\chi$-plane the circumference of the unit circle $|\chi| = 1$.

\emph{With the help of formulas (II), the reciprocity law (I) can be transferred from the differentials of the 2nd kind to the integrals of the 3rd kind:}
\[ \left\{ \begin{aligned}
\Re [\omega_{q_1 q_2}]_{\mathfrak{p}_1}^{\mathfrak{p}_2} &= \Re [\omega_{\mathfrak{p}_1 \mathfrak{p}_2}]_{q_1}^{q_2}, & \Im [\omega_{q_1 q_2}]_{\mathfrak{p}_1}^{\mathfrak{p}_2} &= \Re [\omega'_{\mathfrak{p}_1 \mathfrak{p}_2}]_{q_1}^{q_2}, \\
\Re [\omega'_{q_1 q_2}]_{\mathfrak{p}_1}^{\mathfrak{p}_2} &= \Im [\omega_{\mathfrak{p}_1 \mathfrak{p}_2}]_{q_1}^{q_2}, & \Im [\omega'_{q_1 q_2}]_{\mathfrak{p}_1}^{\mathfrak{p}_2} &= \Im [\omega'_{\mathfrak{p}_1 \mathfrak{p}_2}]_{q_1}^{q_2}.
\end{aligned} \right. \tag{III} \]

These equations, which are usually referred to as the \emph{theorem on the interchange of argument and parameter}, are valid when the paths used $\mathfrak{p}_1\,\mathfrak{p}_2$, $q_1\,q_2$ do not cross each other. For the proof we assume that\ednote{Misprint in the original: ``das'' instead of ``daß''.} the path $\mathfrak{p}_1\,\mathfrak{p}_2$ lies entirely in the $z$-disc $K$ about $\mathfrak{p}$, and $q_1\,q_2$ entirely in the $\zeta$-disc $\mathrm{K}$ about $q$. We can assume (which may perhaps only be achieved by rotating the coordinate axes, which however has no influence on the meaning of the differentials $d\omega_{\mathfrak{p}_1 \mathfrak{p}_2}$, $d\omega_{q_1 q_2}$) that the straight segment $s = (\mathfrak{p}_1\,\mathfrak{p}_2)$ in the $z$-plane is parallel to the real ($x$-) axis, and the segment $\sigma = (q_1\,q_2)$ in the $\zeta$-plane is parallel to the real ($\xi$-) axis. If we now replace $\mathfrak{p}$, $q$ in the 1st equation (I) by the variable points $\mathfrak{p}^*$, $q^*$ within

\origpage{115}
$K$ and $\mathrm{K}$, respectively (where to uniquely fix the notation the uniformizing variables $z$ and $\zeta$ are to be used as above), and then integrate with respect to $z = z(\mathfrak{p}^*)$ along $s$, and with respect to $\zeta = \zeta(q^*)$ along $\sigma$, we obtain
\[ \int\limits_{q_1}^{q_2} \int\limits_{\mathfrak{p}_1}^{\mathfrak{p}_2} \Re \frac{d\tau_{q^*}(z)}{dz} \cdot dz\,d\zeta = \int\limits_{\mathfrak{p}_1}^{\mathfrak{p}_2} \int\limits_{q_1}^{q_2} \Re \frac{d\tau_{\mathfrak{p}^*}(\zeta)}{d\zeta} \cdot d\zeta\,dz. \]
Since the inner integration on the left-hand side runs parallel to the real axis, its execution yields
\[ \int\limits_{\mathfrak{p}_1}^{\mathfrak{p}_2} \Re \frac{d\tau_{q^*}(z)}{dz} \cdot dz = \Re \int\limits_{\mathfrak{p}_1}^{\mathfrak{p}_2} d\tau_{q^*} = \Re \frac{d\omega_{\mathfrak{p}_1 \mathfrak{p}_2}(\zeta)}{d\zeta} \]
and for the same reason
\[ \int\limits_{q_1}^{q_2} \Re \frac{d\omega_{\mathfrak{p}_1 \mathfrak{p}_2}(\zeta)}{d\zeta} \cdot d\zeta = \Re \int\limits_{q_1}^{q_2} d\omega_{\mathfrak{p}_1 \mathfrak{p}_2} = \Re [\omega_{\mathfrak{p}_1 \mathfrak{p}_2}]_{q_1}^{q_2}. \]
Treating the right-hand side in the same way, the first equation (III) results. One thus obtains equations (III) from (I) by integrating the latter both with respect to the argument ($\mathfrak{p}$) and with respect to the parameter ($q$), interchanging the order of integrations on one side, and then carrying out the inner integrations with the help of equations (II).

In order to survey the modifications that arise if the paths $\mathfrak{p}_1\,\mathfrak{p}_2$, $q_1\,q_2$ intersect, we have to consider the case where both paths run in the $z$-disc $K$. We then deform the path $q_1\,q_2$ in such a way that it no longer meets $\mathfrak{p}_1\,\mathfrak{p}_2$, but still runs within $K$. Thereby the value of the quantities appearing in the 1st and 4th equations of (III) does not change; these two formulas are therefore valid without any restriction. On the other hand, the value of the right-hand side of the 2nd and 3rd equations may by this deformation be increased or decreased by an integer multiple of $2\pi$, while the value of the left-hand sides is also preserved here. Without any restriction these two formulas are therefore valid only in the modified form
\[ \begin{gathered}
\Im [\omega_{q_1 q_2}]_{\mathfrak{p}_1}^{\mathfrak{p}_2} = \Re [\omega'_{\mathfrak{p}_1 \mathfrak{p}_2}]_{q_1}^{q_2} - 2n\pi, \quad \Re [\omega'_{q_1 q_2}]_{\mathfrak{p}_1}^{\mathfrak{p}_2} = \Im [\omega_{\mathfrak{p}_1 \mathfrak{p}_2}]_{q_1}^{q_2} - 2n\pi \\
\{\, n \text{ an integer} \,\}
\end{gathered} \]

If $q_1\,q_2$ is in particular a closed path $\alpha$, one has
\[ \int\limits_\alpha d\omega_{\mathfrak{p}_1 \mathfrak{p}_2} = 2\pi i\,\Re \int\limits_{\mathfrak{p}_1}^{\mathfrak{p}_2} dw_\alpha, \quad \Im \int\limits_\alpha d\omega'_{\mathfrak{p}_1 \mathfrak{p}_2} = 2\pi\,\Im \int\limits_{\mathfrak{p}_1}^{\mathfrak{p}_2} dw_\alpha. \tag{$\mathrm{III}_0$} \]

\origpage{116}
In the 1st equation it is assumed that $\mathfrak{p}_1\,\mathfrak{p}_2$ does not meet the curve $\alpha$. The latter yields, if we take also for $\mathfrak{p}_1\,\mathfrak{p}_2$ a closed path $\beta$, the symmetry law
\[ \Im \int\limits_\alpha dw_\beta = \Im \int\limits_\beta dw_\alpha. \tag{$\mathrm{III}_{00}$} \]

It finally remains to derive from the differentials $d\tau_q$, $d\tau'_q$ those integrals of the 2nd kind that possess poles of higher than 1st order. From (I) there results by differentiation with respect to $\xi = \xi(q^*)$, when $q$ is first replaced by $q^*$:
\[ \Re \frac{d^2\tau_p(q)}{dq^2} + i\,\Re \frac{d^2\tau'_p(q)}{dq^2} = \left[ \frac{\partial}{\partial\xi} \left( \frac{d\tau_{q^*}(p)}{dp} \right) \right]_{\xi(q^*) = 0}. \]
It follows from this that the left-hand side can be set
\[ = \frac{d\tau_q^2(p)}{dp} \]
where $d\tau_q^2$ denotes a differential everywhere regular on $\mathfrak{F}$ that becomes infinite only at the point $q$, and indeed, as the expression on the left-hand side shows, with the principal part $-\,\frac{2\,d\zeta}{\zeta^3} = d\left(\frac{1}{\zeta^2}\right)$. The periods of this integral $\tau_q^2$ of the 2nd kind are purely imaginary. We thus have the relations:
\[ \left\{ \begin{aligned}
\Re \frac{d\tau_q^2(p)}{dp} &= \Re \frac{d^2\tau_p(q)}{dq^2}, & \Im \frac{d\tau_q^2(p)}{dp} &= \Re \frac{d^2\tau'_p(q)}{dq^2}, \\
\Re \frac{d{\tau_q^2}'(p)}{dp} &= \Im \frac{d^2\tau_p(q)}{dq^2}, & \Im \frac{d{\tau_q^2}'(p)}{dp} &= \Im \frac{d^2\tau'_p(q)}{dq^2}.
\end{aligned} \right. \tag{$\mathrm{I}^2$} \]
If in (II), ($\mathrm{II}_0$) one replaces $p$ by $p^*$ and differentiates with respect to $x = \Re z(p^*)$, one gets
\[ \left\{ \begin{aligned}
\Re \int\limits_{q_1}^{q_2} d\tau_p^2 &= \Re \frac{d^2\omega_{q_1 q_2}(p)}{dp^2}, & \Im \int\limits_{q_1}^{q_2} d\tau_p^2 &= \Re \frac{d^2\omega'_{q_1 q_2}(p)}{dp^2}, \\
\Re \int\limits_{q_1}^{q_2} d{\tau_p^2}' &= \Im \frac{d^2\omega_{q_1 q_2}(p)}{dp^2}, & \Im \int\limits_{q_1}^{q_2} d{\tau_p^2}' &= \Im \frac{d^2\omega'_{q_1 q_2}(p)}{dp^2}.
\end{aligned} \right. \tag{$\mathrm{II}^2$} \]
\[ \left\{ \int\limits_\alpha d\tau_p^2 = 2\pi i\,\Re \frac{d^2 w_\alpha(p)}{dp^2}, \quad \int\limits_\alpha d{\tau_p^2}' = 2\pi i\,\Im \frac{d^2 w_\alpha(p)}{dp^2}. \right. \tag{$\mathrm{II}_0^2$} \]
Corresponding equations ($\mathrm{I}^n$), ($\mathrm{II}^n$), ($\mathrm{II}_0^n$) hold for integrals with a pole of the $n^{\text{th}}$ order.

One sees how all these relations (which otherwise, following Riemann and C.~Neumann, are usually derived by the method of ``integrating around'' the boundary of the dissected surface) are an immediate consequence of the intuitive and well-known fact that by

\origpage{117}
longitudinal juxtaposition of double sources one obtains two simple sources with opposite signs, but by transverse juxtaposition a pair of oppositely rotating vortices.

\subsection*{§~17. Single-valued functions on $\mathfrak{F}$ as a subclass of additive and multiplicative functions on $\mathfrak{F}$. The Riemann-Roch theorem and Abel's theorem.}

By a transcendentally normalized Abelian integral of the 2nd kind or an \textbf{``additive function''} we understand a function $v(\hat{\mathfrak{p}})$, single-valued on the covering surface $\hat{\mathfrak{F}}$ of the integral functions and regular except for poles, which behaves additively under every deck transformation $S$ of $\hat{\mathfrak{F}}$:
\[ v(\hat{\mathfrak{p}}S) = v(\hat{\mathfrak{p}}) + c\,i, \]
where $c\,i$ denotes a purely imaginary constant belonging to the transformation $S$ and independent of $\hat{\mathfrak{p}}$. There are only finitely many points $\mathfrak{p}$ on $\mathfrak{F}$ over which poles of such an additive function $v$ lie, and if $z$ denotes a local uniformizing variable belonging to $\mathfrak{p}$ and therefore also to all points of $\hat{\mathfrak{F}}$ lying above it, then $v$ has at all points lying above $\mathfrak{p}$ the same ``principal part'':
\[ \frac{A_{-1}}{z} + \frac{A_{-2}}{z^2} + \dots + \frac{A_{-r}}{z^r}. \]
The poles and associated principal parts can be prescribed arbitrarily and determine the function $v$ uniquely up to an additive constant:
\[ \begin{aligned}
v &= \sum\limits_\mathfrak{p} \left[ (a\tau_\mathfrak{p} + a'\tau'_\mathfrak{p}) + (a^2\tau_\mathfrak{p}^2 + a^{2\prime}{\tau_\mathfrak{p}^2}') + \dots + (a^r\tau_\mathfrak{p}^r + a^{r\prime}{\tau_\mathfrak{p}^r}') \right] \\
&\qquad + \text{arbitrary const.} \quad (a^h = \Re A_{-h}, \ a^{h\prime} = -\,\Im A_{-h})
\end{aligned} \tag{40} \]
(the summation extends over the distinct poles $\mathfrak{p}$). This representation is the analogue of the \emph{partial fraction decomposition} of rational functions. To obtain this analogue, we must accordingly extend the class of single-valued functions on the surface $\mathfrak{F}$ having at most poles to the more comprehensive totality of ``additive functions''.

Alongside the \emph{partial fraction decomposition}, the representation of rational functions as a \emph{product of linear factors} is important. In order to transfer it from the sphere to arbitrary closed Riemann surfaces, we have to speak of \textbf{``multiplicative functions''}, that is, of those single-valued functions $\Theta$ on the covering surface $\hat{\mathfrak{F}}$, regular except for poles, which behave multiplicatively under every deck transformation $S$ of $\hat{\mathfrak{F}}$:

\origpage{118}
\[ \Theta(\hat{\mathfrak{p}}S) = \mu_S \cdot \Theta(\hat{\mathfrak{p}}), \]
where the multiplier $\mu_S$ is a constant belonging to $S$ of absolute value 1.${}^{1)}$ The $\mu_S$ form a system of characters. $\Theta$ has the same order at all points lying above a position $\mathfrak{p}$ of $\mathfrak{F}$; this order can be positive (then $\Theta$ has a zero over $\mathfrak{p}$), negative ($\Theta$ has a pole over $\mathfrak{p}$), or zero. \emph{The points on $\mathfrak{F}$ over which the zeros and poles of a multiplicative function are to lie can be prescribed arbitrarily in finite number according to their position and multiplicity --- under the single restriction that the number of zeros must equal the number of poles, provided every zero and every pole is counted with the correct multiplicity. By prescribing the zeros and poles, however, $\Theta$ is uniquely determined up to an arbitrary constant factor $(\neq 0)$.} For there can be no multiplicative function $\Theta$ apart from constants that possesses neither zeros nor poles; from such a function $\lg |\Theta|$ would yield a potential function single-valued on $\mathfrak{F}$ and everywhere regular, which cannot exist. That the number of zeros must equal that of the poles results from the fact that $\frac{d\Theta}{\Theta}$ is a differential of the 3rd kind on $\mathfrak{F}$ whose sum of residues is
\[ = \textit{number of zeros minus number of poles of } \Theta \]
If $1$, $2$ are any two points on $\mathfrak{F}$, then
\[ e^{\omega_{12}} = \Theta_{12} \quad (e = \text{base of natural log.}) \]
is a multiplicative function which possesses over $1$ a zero and over $2$ a pole of the first order (and by this property of having only one zero and one point of infinity corresponds to the linear factor $z - a$ or $\frac{z-a}{z-b}$ on the $z$-sphere). If
\[ \begin{gathered}
\mathfrak{p}_1 \, \mathfrak{p}_2 \, \dots \, \mathfrak{p}_r \text{ are the prescribed zeros,} \\
q_1 \, q_2 \, \dots \, q_r \text{ the prescribed poles}
\end{gathered} \]
of a function $\Theta$ (where every zero and every pole is listed as many times as its multiplicity requires), then one obtains the sought-after $\Theta$ from the equation
\[ \Theta = \text{Const.} \cdot \Theta_{\mathfrak{p}_1 q_1} \cdot \Theta_{\mathfrak{p}_2 q_2} \cdots \Theta_{\mathfrak{p}_r q_r}, \]
\emph{in which we have before us the analogue of the product representation of rational functions.} We write this equation also in the symbolic form

\textbf{1)} Cf.\ Riemann, Theorie der Abelschen Funktionen, Werke (2nd ed.), p.~140; but especially Appell, Journal de mathématiques, 3rd ser., vol.~9 (1883), and Acta Mathematica vol.~13 (1890).

\origpage{119}
\[ \Theta = \frac{\mathfrak{p}_1 \, \mathfrak{p}_2 \, \dots \, \mathfrak{p}_r}{q_1 \, q_2 \, \dots \, q_r}. \tag{41} \]
Here it was taken for granted that all $\mathfrak{p}_h$ are distinct from all $q_h$. If that is not the case, the symbol on the right-hand side is to be given the meaning that results if, as with ordinary fractions, one ``cancels'' the points occurring in both numerator and denominator.

The theorem that a multiplicative function becomes $0$ as often as $\infty$ holds in particular also for every single-valued function $f$ on $\mathfrak{F}$ free from essential singularities. If we apply it, instead of to $f$, to $f - a$, where $a$ is an arbitrary constant, we see: \emph{A single-valued function on $\mathfrak{F}$ regular except for poles assumes every value, including $\infty$, equally often.}

In the case $p = 0$, the integral of the 2nd kind $\tau_q$ is single-valued on $\mathfrak{F}$ itself; consequently, since it possesses only a single pole of the 1st order, it assumes every value once and only once and maps $\mathfrak{F}$ conformally onto the sphere: \emph{The sphere is the only Riemann surface of genus $0$.}

\emph{We now have to investigate in general how the single-valued functions on $\mathfrak{F}$ are classified as special cases under the additive and multiplicative functions}, that is, we want to establish the conditions under which the representations (40) and (41) yield single-valued functions on $\mathfrak{F}$. Thereby we will be led to the \emph{Riemann-Roch theorem} and to \emph{Abel's theorem}, respectively. We begin with (40) and initially assume only simple poles for ease of presentation. Let $\mathfrak{p}_1$, $\mathfrak{p}_2$, \dots, $\mathfrak{p}_m$ thus be any $m$ mutually distinct points on $\mathfrak{F}$. Single-valued functions on $\mathfrak{F}$ which have poles of the 1st order at most at $\mathfrak{p}_1$, $\mathfrak{p}_2$, \dots, $\mathfrak{p}_m$, but are otherwise regular, must be of the form:
\[ f = (a_1\tau_{\mathfrak{p}_1} + a'_1\tau'_{\mathfrak{p}_1}) + (a_2\tau_{\mathfrak{p}_2} + a'_2\tau'_{\mathfrak{p}_2}) + \dots + (a_m\tau_{\mathfrak{p}_m} + a'_m\tau'_{\mathfrak{p}_m}) + (b + i b'), \tag{42} \]
where the $a_h$, $a'_h$; $b$, $b'$ denote real constants. \emph{The condition that the $2p$ periods of $f$ corresponding to the $2p$ retrosections $\alpha_\iota$ of a canonical dissection must vanish yields $2p$ homogeneous linear equations for $a_h$, $a'_h$ with real coefficients}, equations that can certainly be satisfied by constants $a_h$, $a'_h$ not all vanishing if $m > p$.

\emph{If, therefore, more than $p$ points on $\mathfrak{F}$ are prescribed arbitrarily, there always exist non-constant single-valued functions on $\mathfrak{F}$ having poles of the 1st order at most at these points, but otherwise regular. Thereby in particular the existence of non-constant single-valued functions regular except for poles is secured on every closed Riemann surface.}

Let $\mathfrak{p}_1$, $\mathfrak{p}_2$, \dots, $\mathfrak{p}_r$ be any $r$ mutually distinct points on

\origpage{120}
$\mathfrak{F}$ and $m_1$, $m_2$, \dots, $m_r$ associated integers ($> 0$, $= 0$, or $< 0$), then we designate a function or a differential on $\mathfrak{F}$ that has at the point $\mathfrak{p}_h$ at least the order $m_h$ ($h = 1, \, 2, \, \dots, \, r$), but is otherwise regular, as a \textbf{multiple of the divisor}${}^{1)}$
\[ \mathfrak{d} = \mathfrak{p}_1^{m_1} \mathfrak{p}_2^{m_2} \dots \mathfrak{p}_r^{m_r}. \]
It is convenient to agree to calculate with such divisors --- symbols consisting of a collection of finitely many points with associated order numbers (exponents) --- with respect to multiplication and division as with ordinary fractions. The sum of exponents $\displaystyle\sum_{h=1}^r m_h = m$ we call the \textbf{total order} of the divisor. The theorem just stated can then be generalized as follows:

\emph{If the total order $m$ of a divisor $\mathfrak{d}$ is not smaller than $p$, then among the single-valued functions on $\mathfrak{F}$ there exist multiples of $\frac{1}{\mathfrak{d}}$, and indeed at least $m + 1 - p$ linearly independent ones (in the complex sense).}

For if among the exponents $m_h$ the first $s$ are positive and the others negative, then those additive functions having over $\mathfrak{p}_h$ a pole of at most the $m_h^{\text{th}}$ order ($h = 1, \, 2, \, \dots, \, s$) and otherwise regular depend linearly and homogeneously on $2\left(1 + \displaystyle\sum_{h=1}^s m_h\right)$ real constants. The requirement that all periods of the additive function shall vanish entails $2p$ linear homogeneous real equations for these constants; the requirement that the function shall vanish at the points $\mathfrak{p}_{s+1}$, \dots, $\mathfrak{p}_r$ to the orders $-m_{s+1}$, \dots, $-m_r$, respectively, entails a further $2\displaystyle\sum_{h=s+1}^r (-m_h)$ such equations. The excess of the number of unknowns over the number of equations,
\[ 2\left(1 + \sum_{h=1}^s m_h\right) - 2p - 2\sum_{h=s+1}^r (-m_h) = 2(m + 1 - p) \]
gives a lower bound for the number of linearly independent solution systems (in the real sense).

The most complete answer to these questions, however, is provided only by the \emph{Riemann-Roch theorem}${}^{2)}$. As the most transparent example we take again

\textbf{1)} This terminology originates from the arithmetic theory of algebraic functions of Hensel and Landsberg; cf.\ Hensel and Landsberg, Theorie der algebraischen Funktionen einer Variablen, Leipzig (Teubner) 1902.

\textbf{2)} Riemann considers only the so-called ``general'' case ($R = p$ in the notation of the text); the result valid without exception for every integral divisor $\mathfrak{d}$ is due to G.~Roch, Crelles Journal, vol.~64 (1865), pp.~372---376; fractional divisors were considered by Klein (Riemannsche Flächen I, autographed lecture, Göttingen 1892, pp.~110---111), E.~Ritter (Math.\ Ann.\ vol.~44, 1894, p.~314), Hensel and Landsberg (op.\ cit., cf.\ especially pp.~362---364).

\origpage{121}
an ``integral'' divisor composed entirely of simple points
\[ \mathfrak{d} = \mathfrak{p}_1 \mathfrak{p}_2 \dots \mathfrak{p}_m. \]
The $2p$ linear homogeneous equations demanding the vanishing of the periods of (42) for the retrosections $\alpha_h$ read according to ($\mathrm{II}_0$) §~16:
\[ \begin{gathered}
\sum_{\iota=1}^m \left( a_\iota \Re \frac{dw_h(\mathfrak{p}_\iota)}{d\mathfrak{p}_\iota} + a'_\iota \Im \frac{dw_h(\mathfrak{p}_\iota)}{d\mathfrak{p}_\iota} \right) = 0. \\
(h = 1, \, 2, \, \dots, \, 2p).
\end{gathered} \tag{43} \]
If the ``rank'' of this system of equations is $= 2R$ (that is, if $2R$ is the number of linearly independent equations among these $2p$ equations for the $2m$ unknowns $a_\iota$, $a'_\iota$), then there are exactly
\[ A = m + 1 - R \]
linearly independent functions (in the complex sense) on $\mathfrak{F}$ that are multiples of $\frac{1}{\mathfrak{d}}$. The ``transposed'' system of equations belonging to (43) (in which rows have become columns and columns have become rows), whose real unknowns we denote by $b_h$ ($h = 1, \, 2, \, \dots, \, 2p$), reads:
\[ \left\{ \begin{aligned}
\Re \left( \sum_{h=1}^{2p} b_h \frac{dw_h(\mathfrak{p}_l)}{d\mathfrak{p}_l} \right) &= 0, \\
\Im \left( \sum_{h=1}^{2p} b_h \frac{dw_h(\mathfrak{p}_l)}{d\mathfrak{p}_l} \right) &= 0.
\end{aligned} \right. \qquad [l = 1, \, 2, \, \dots, \, m] \]
Its rank is the same: $2R$, and the number of its linearly independent solution systems $\{b_h\}$ is accordingly $2p - 2R$. But that means: the number $B$ of (in the complex sense) linearly independent differentials of the 1st kind which vanish at the points $\mathfrak{p}_l$ ($l = 1, \, 2, \, \dots, \, m$) is
\[ B = p - R. \]
\emph{The relation thus established,}
\[ A = B + (m + 1 - p) \]
\emph{forms the content of the Riemann-Roch theorem.} By it the \emph{problem} of determining the number $A$ is, to be sure, not actually solved, but nevertheless reduced to a substantially simpler one. For the linear family of differentials of the 1st kind to which the number $B$ refers is obviously by no means as complicated a system as that of the

\origpage{122}
functions on the surface. Moreover, we shall see presently that whenever $m > 2p - 2$, one must set $B = 0$, since the total order of the zeros of a differential of the 1st kind cannot exceed $2p - 2$. And finally, the reciprocity that exists according to the Riemann-Roch theorem between functions and differentials is in itself a remarkable and interesting fact. In any case, this theorem is one of the most important starting points for a deeper insight into the nature of functions on a closed Riemann surface.

If the divisor $\mathfrak{d}$ contains a denominator (negative exponents) and multiple points occur in it (exponents $> 1$ or $< -1$), then, while the proof procedure otherwise remains the same, one must also draw upon equations (II) and possibly ($\mathrm{II}_0^n$), ($\mathrm{II}^n$) of the preceding section.${}^{1)}$ The result is analogous. We can state it as follows:

(\emph{Riemann-Roch theorem.}) \emph{Between the number $B$ of (in the complex sense) linearly independent differentials which are multiples of the arbitrary divisor $\mathfrak{d}$ of total order $m$, and the number $A$ of linearly independent functions single-valued on $\mathfrak{F}$ which are multiples of the reciprocal divisor $\frac{1}{\mathfrak{d}}$, there holds the relation}
\[ A = B + (m + 1 - p). \]

We have seen: a multiple of $\frac{1}{\mathfrak{p}_1 \, \mathfrak{p}_2 \dots \mathfrak{p}_m}$ exists among the single-valued, non-constant functions on $\mathfrak{F}$ whenever $m > p$. We add: If $m \leqq p$, such a multiple is \emph{in general}, namely except for special positions of the $m$ points $\mathfrak{p}_\iota$, \emph{not} present. Thereby the function-theoretic significance of the genus $p$ is once again brought into a bright light. A non-constant multiple of $\frac{1}{\mathfrak{p}_1 \, \mathfrak{p}_2 \dots \mathfrak{p}_p}$ exists among the functions on $\mathfrak{F}$ according to the Riemann-Roch theorem if and only if there exists a differential of the 1st kind $dw$ which vanishes (but not identically) at the points $\mathfrak{p}_1$, $\mathfrak{p}_2$, \dots, $\mathfrak{p}_p$. If
\[ dw_1^*, \, dw_2^*, \, \dots, \, dw_p^* \]
denotes a complex basis for the differentials of the 1st kind, then it must be possible to find

\textbf{1)} Cf.\ the passages cited in Klein and Hensel-Landsberg. E.~Ritter concludes loc.\ cit.\ as follows: If $f_0 = \frac{\mathfrak{d}'}{\mathfrak{d}}$ is a function that is a multiple of $\frac{1}{\mathfrak{d}}$ ($\mathfrak{d}'$ integral), I obtain all such functions $f$ from: $f = f_0 \cdot f'$, $f'$ a multiple of $\frac{1}{\mathfrak{d}'}$, and thereby return to the case of an integral divisor $\mathfrak{d}'$. But here the proof is missing that [when $B + (m + 1 - p) > 0$] an $f_0$ exists at all.

\origpage{123}
complex numbers $c_1^*, \, c_2^*, \, \dots, \, c_p^*$, not all vanishing, such that
\[ c_1^* dw_1^* + c_2^* dw_2^* + \dots + c_p^* dw_p^* \]
possesses zeros at those $p$ points. We choose any $p$ mutually distinct points $\mathfrak{p}_1^0$, $\mathfrak{p}_2^0$, \dots, $\mathfrak{p}_p^0$, associated local uniformizing variables $z_h$, and for each $\mathfrak{p}_h^0$ a $z_h$-disc $\mathrm{K}_h$: $|z_h| \leqq a_h$ ($h = 1, \, 2, \, \dots, \, p$) such that these discs do not mutually encroach upon one another. If for every position of the points $\mathfrak{p}_h$ within $\mathrm{K}_h$ there existed a differential of the 1st kind\ednote{Misprint in the original: ``eine'' instead of ``ein''.} vanishing simultaneously at the points $\mathfrak{p}_h$, then identically in $z_1$, $z_2$, \dots, $z_p$ ($|z_h| \leqq a_h$) we would have
\[ \Delta \equiv \begin{vmatrix}
\dfrac{dw_1^*(z_1)}{dz_1}, & \dots, & \dfrac{dw_p^*(z_1)}{dz_1} \\
\cdot & \cdot & \cdot & \cdot & \cdot & \cdot & \cdot & \cdot & \cdot & \cdot & \cdot & \cdot \\
\cdot & \cdot & \cdot & \cdot & \cdot & \cdot & \cdot & \cdot & \cdot & \cdot & \cdot & \cdot \\
\dfrac{dw_1^*(z_p)}{dz_p}, & \dots, & \dfrac{dw_p^*(z_p)}{dz_p}
\end{vmatrix} = 0 \]
This, however, contradicts the linear independence of the differentials $dw_h^*$. For if we imagine $z_1$ in $\Delta = 0$ as variable, but $z_2, \, \dots, \, z_p$ held fixed for a moment, we have before us a homogeneous linear relation between the differentials $dw_h^*(z_1)$ which is satisfied identically in $z_1$; hence the coefficients of this relation, that is, the subdeterminants formed from the last $p - 1$ rows of $\Delta$, must all vanish individually, and indeed identically in $z_2, \, \dots, \, z_p$. If we repeat this inference by applying it now, instead of to $\Delta$, to each of the aforementioned $(p - 1)$-rowed subdeterminants, and continue thus step by step, we finally ascend to the absurd assertion that the elements of the last row, that is, the differentials $dw_h^*$ themselves, are identically $= 0$.

For every point $\mathfrak{p}$ we obtain by the formula
\[ d\tau_{\mathfrak{p}} - i d\tau'_{\mathfrak{p}} = dw_{\mathfrak{p}} \]
a differential of the 1st kind.${}^{1)}$ This is a very natural method for generating such differentials from those of the 2nd kind; the only question is whether it is comprehensive enough to provide a full basis for the differentials of the 1st kind. We can now answer this question in the affirmative. Namely, if I choose $p$ points $\mathfrak{p}_1, \, \mathfrak{p}_2, \, \dots, \, \mathfrak{p}_p$ such that no single-valued function $\neq \mathrm{const.}$ exists on $\mathfrak{F}$ that is a multiple of $\frac{1}{\mathfrak{p}_1 \, \mathfrak{p}_2 \dots \mathfrak{p}_p}$, then the associated $dw_{\mathfrak{p}_1}, \, dw_{\mathfrak{p}_2}, \, \dots, \, dw_{\mathfrak{p}_p}$ form a com-

\textbf{1)} A changed choice of the local uniformizing variable for $\mathfrak{p}$ results only in its multiplication by a constant factor, so that we can regard it as essentially determined by $\mathfrak{p}$ alone.

\origpage{124}
plex basis of the differentials of the 1st kind. For if there were a relation
\[ C_1 dw_{\mathfrak{p}_1} + C_2 dw_{\mathfrak{p}_2} + \dots + C_p dw_{\mathfrak{p}_p} = 0, \]
then for every closed path $\alpha$ on $\mathfrak{F}$ one would have:
\[ \sum_{h=1}^p C_h \left( \int_\alpha d\tau_{\mathfrak{p}_h} - i \int_\alpha d\tau'_{\mathfrak{p}_h} \right) = 0. \]
If we retain from this equation, say, only the imaginary part ($c_h = \Re C_h$, $c'_h = \Im C_h$):
\[ \sum_{h=1}^p \left( c_h \int_\alpha d\tau_{\mathfrak{p}_h} + c'_h \int_\alpha d\tau'_{\mathfrak{p}_h} \right) = 0, \]
then it turns out that
\[ \sum_{h=1}^p (c_h \tau_{\mathfrak{p}_h} + c'_h \tau'_{\mathfrak{p}_h}) \]
would have to be a single-valued function on the surface.

Two paths are open to us for obtaining functions on the surface: either to combine additive functions linearly with such constants that the periods are made to vanish, --- or to divide two differentials by one another.${}^{1)}$ By combining both possibilities, one can arrive at important theorems. First we note: since the quotient of two differentials is a function on $\mathfrak{F}$,
1) every differential is uniquely determined up to a constant factor by its representing divisor, in whose numerator the zeros of the differential are recorded in their multiplicity and whose denominator contains the poles --- and
2) the total order $d$ of this representing divisor is the same for all differentials; $d$ is the number of zeros of every differential of the 1st kind and thus for $p > 0$ certainly $\geqq 0$. We will prove that this important invariant of every Riemann surface depends solely on the genus $p$, and indeed according to the simple equation:
\[ d = 2p - 2. \]

Let $\mathfrak{d}$, $\mathfrak{e}$ be two divisors of total order $m$ and $n$, respectively, whose product $\mathfrak{d}\mathfrak{e}$ represents a differential $dv$ on $\mathfrak{F}$: $m + n = d$. Let there be exactly $M$ linearly independent differentials that are multiples of

\textbf{1)} This latter method is also applicable to open Riemann surfaces and yields here immediately the existence proof for functions on the surface that do not reduce to mere constants. Cf.\ also Koebe, Crelles Journal vol.~139, 1911, pp.~291\,f.

\origpage{125}
$\mathfrak{d}$, and $N$ such differentials that are multiples of $\mathfrak{e}$. Among the single-valued functions on $\mathfrak{F}$ we consider all multiples $f$ of $\frac{1}{\mathfrak{d}}$. If we generate them by combination of additive functions, the Riemann-Roch theorem yields the result that the number of linearly independent $f$
\[ = m + 1 - p + M \]
For every $f$, the differential $f\,dv$ is a multiple of $\mathfrak{e}$; that is, in a second way we obtain all functions $f$ by dividing all differentials that are multiples of $\mathfrak{e}$ by $dv$. Paying attention only to the numbers, this yields the relation
\[ N = m + 1 - p + M. \]
Analogously,
\[ M = n + 1 - p + N. \]
Addition of these two equations yields
\[ d = 2p - 2, \]
subtraction the \emph{Brill-Noether reciprocity theorem}:${}^{1)}$
\[ M - N = n - m. \]\ednote{Misprint in the original: From the two equations it follows that $2(M - N) = n - m$; the factor 2 is missing in the original.}

Yet enough of these considerations, which all link themselves as if naturally to the ``partial fraction decomposition'' (40)! We turn to the ``product representation'' (41). To calculate the multiplier $\mu_S$ belonging to $S$ of
\[ \Theta = \frac{\mathfrak{p}_1 \, \mathfrak{p}_2 \dots \mathfrak{p}_r}{\mathfrak{q}_1 \, \mathfrak{q}_2 \dots \mathfrak{q}_r} \]
we lay a closed curve $\alpha$ on $\mathfrak{F}$ that does not pass through the $\mathfrak{p}_h$, $\mathfrak{q}_h$ and leads on $\hat{\mathfrak{F}}$ from a point $\hat{\mathfrak{p}}$ to the point $\hat{\mathfrak{p}}S$, and draw integration paths $\mathfrak{p}_h\mathfrak{q}_h$ that do not meet $\alpha$. Then according to ($\mathrm{III}_0$), §~16:
\[ \mu_S = \prod_{h=1}^r e^{\int\limits_\alpha d\omega_{\mathfrak{p}_h \mathfrak{q}_h}} = \prod_{h=1}^r e^{-2\pi i \cdot \Re \int\limits_{\mathfrak{q}_h}^{\mathfrak{p}_h} dw_\alpha} = \prod_{h=1}^r \frac{\chi_S(\mathfrak{p}_h)}{\chi_S(\mathfrak{q}_h)}. \tag{44} \]
It is natural to introduce as the \textbf{integral characters} of an arbitrary divisor
\[ \mathfrak{d} = \mathfrak{p}_1^{m_1} \, \mathfrak{p}_2^{m_2} \dots \mathfrak{p}_r^{m_r} \]
the quantities
\[ \chi_S(\mathfrak{d}) = \chi_S(\mathfrak{p}_1)^{m_1} \cdot \chi_S(\mathfrak{p}_2)^{m_2} \dots \chi_S(\mathfrak{p}_r)^{m_r} \]
Then equation (44) reads in words${}^{2)}$:

\emph{The multipliers of a function $\Theta$ coincide with the integral characters of its representing divisor.}

\textbf{1)} A.~Brill and M.~Noether, Mathematische Annalen, vol.~7 (1874), p.~283.

\textbf{2)} Cf.\ Appell, Acta Mathematica vol.~13 (1890), pp.~13--14 of the prize essay.

\origpage{126}
In particular ($\mu_S = 1$):

\emph{Prescribed points $\mathfrak{p}_h$, $\mathfrak{q}_h$ ($h = 1, \, 2, \, \dots, \, r$) are the totality of zeros and poles, respectively, of a single-valued function on $\mathfrak{F}$ if and only if the product of the integral characters of the zeros is equal to the product of the integral characters of the poles:}
\[ \prod_{h=1}^r \chi_S(\mathfrak{p}_h) = \prod_{h=1}^r \chi_S(\mathfrak{q}_h). \]

If we return to the meaning of the integral characters and use a particular canonical dissection $\alpha_h$ together with an associated real basis $dw_{\alpha_h} = dw_h$ of the integrals of the 1st kind, we can replace this requirement by the following $2p$ equations:
\[ \begin{gathered}
\sum_{l=1}^r \Re \int\limits_{\mathfrak{p}_0}^{\mathfrak{p}_l} dw_h - \sum_{l=1}^r \Re \int\limits_{\mathfrak{q}_0}^{\mathfrak{q}_l} dw_h = n_h \qquad (n_h = \text{integer}) \\
[h = 1, \, 2, \, \dots, \, 2p].
\end{gathered} \]
The paths of integration $\mathfrak{p}_0 \mathfrak{p}_l$, $\mathfrak{q}_0 \mathfrak{q}_l$ are initially arbitrary. If, however, we subsequently deform one of these integration paths, say $\mathfrak{p}_0 \mathfrak{p}_1$, in such a way that its old and its new position together form a closed line which is homologous to
\[ (n_1 \alpha_2 - n_2 \alpha_1) + \dots + (n_{2p-1} \alpha_{2p} - n_{2p} \alpha_{2p-1}) \]
the integers $n_h$ on the right-hand side of the last equation are thereby made into $0$. There then results:

(\emph{Abel's theorem.})${}^{1)}$ \emph{The points $\mathfrak{p}_h$, $\mathfrak{q}_h$ [$h = 1, \, 2, \, \dots, \, r$] are the totality of zeros and poles, respectively, of a function on $\mathfrak{F}$ if and only if --- with a suitable choice, independent of $w$, of the integration paths leading to}

\textbf{1)} The name, as I use it here, is not entirely accurate. Due to Abel is only the theorem that the stated condition is \emph{necessary}; on the other hand, the theorem stated by Abel [developed by him in magnificent simplicity in the short note ``Démonstration d'une propriété générale d'une certaine classe de fonctions transcendantes'', Crelles Journal vol.~4 (1829), pp.~200---201 = Œuvres complètes, Nouvelle édition (1881), vol.~I, pp.~515---517] is more general, in so far as it does not concern integrals of the 1st kind alone. Cf.\ p.~136 of this work. A memoir submitted by Abel in 1826 to the Paris Academy, ``Mémoire sur une propriété générale d'une classe très-étendue de fonctions transcendantes'', on this subject was long lost through Cauchy's fault and only published after Abel's death: Mémoires présentés par divers savants, vol.~VII (1841) = Abel, Œuvres complètes, Nouvelle édition (1881), vol.~I, pp.~145---241. Furthermore the theorem is contained in Abel's posthumous manuscript ``Sur la comparaison des fonctions transcendantes'', Œuvres complètes, vol.~II, pp.~55---66. The converse of Abel's theorem for integrals of the 1st kind stands between the lines in Riemann, but was explicitly stated first by Clebsch (without a completely adequate proof) (Crelles Journal vol.~63, 1864, p.~198), and utilized by him most extensively in the theory of algebraic curves.

\origpage{127}
$\mathfrak{p}_h$, $\mathfrak{q}_h$ --- for every integral of the 1st kind $w$ the relation
\[ \sum_{h=1}^r w(\mathfrak{p}_h) = \sum_{h=1}^r w(\mathfrak{q}_h) \]
holds.

If $dv_1$, $dv_2$ are any two differentials on $\mathfrak{F}$ regular except for poles, and $\mathfrak{d}_1$, $\mathfrak{d}_2$ their representing divisors, then, since $\frac{dv_1}{dv_2} = \frac{\mathfrak{d}_1}{\mathfrak{d}_2}$ is a function on $\mathfrak{F}$, according to Abel's theorem
\[ \chi_S\left( \frac{\mathfrak{d}_1}{\mathfrak{d}_2} \right) = \frac{\chi_S(\mathfrak{d}_1)}{\chi_S(\mathfrak{d}_2)} = 1 \quad\text{or}\quad \chi_S(\mathfrak{d}_1) = \chi_S(\mathfrak{d}_2). \]
We thus have this corollary of Abel's theorem:

\emph{For all differentials on $\mathfrak{F}$ and their representing divisors, respectively, the integral characters have the same values $\chi_S^0$. A divisor $\mathfrak{d}$ of total order $2p - 2$ is a representative of a differential if and only if the system of its integral characters coincides with the character system $\chi_S^0$.}

We make an immediate application of Abel's theorem to the case $p = 1$. Here there exists, up to a constant factor, only a single differential of the 1st kind $dw$. The values corresponding to the different integration paths that $w(\mathfrak{p})$ can assume at a point $\mathfrak{p}$ form a parallelogrammatic point lattice in the complex $w$-plane. All lattices that thus appear assigned to the individual points $\mathfrak{p}$ of $\mathfrak{F}$ arise from one another by translation, or have, as we shall say, the same ``position and shape'' $\Lambda$. By conceiving each lattice of this position and shape as a ``point'', the $w$-plane is transformed into a closed Riemann surface $\mathfrak{F}_\Lambda$ (cf.\ p.~28). If the same lattice ever corresponded to two distinct points $\mathfrak{p}$, $\mathfrak{q}$ of $\mathfrak{F}$, there would exist according to Abel's theorem a function on $\mathfrak{F}$ possessing a pole of the 1st order only at $\mathfrak{p}$ (and vanishing at $\mathfrak{q}$); this function would have to map the surface bijectively and conformally onto the sphere, which is incompatible with $p = 1$. To distinct points $\mathfrak{p}$ there therefore always correspond distinct lattices $w = w(\mathfrak{p})$. For this reason $dw$ can nowhere be $0$, as also follows from our general formula $d = 2p - 2$ for $p = 1$. Thus $\mathfrak{F}$ is mapped bijectively and bicontinuously (moreover conformally) onto a domain $\mathfrak{G}$ of the manifold $\mathfrak{F}_\Lambda$. Since $\mathfrak{F}$ is closed, $\mathfrak{G}$ must also be closed and consequently identical with the whole of $\mathfrak{F}_\Lambda$: $\mathfrak{F}_\Lambda$ \emph{is conformally equivalent to the given Riemann surface and can be addressed as the normal form of Riemann surfaces of genus $1$.} For $\mathfrak{F}_\Lambda$, the $w$-plane is the ``covering

\origpage{128}
surface of the integral functions'' (for $w$ is itself an integral function), but at the same time, since the plane is simply connected, the ``universal'' covering surface. For the functions on the base surface $\mathfrak{F}$, $w$ is a \emph{uniformizing variable}, that is, all these functions represent themselves as single-valued functions of the complex variable $w$ varying in the schlicht plane, and indeed this representation takes place here --- in accordance with the fact that to all points $w$ of a $\Lambda$-lattice by virtue of $w = w(\mathfrak{p})$ belongs the same point $\mathfrak{p}$ --- by \emph{doubly periodic}, so-called \emph{elliptic functions}. Two Riemann surfaces of genus 1 in their normal form $\mathfrak{F}_{\Lambda'}$, $\mathfrak{F}_{\Lambda''}$ are conformally equivalent if and only if the lattices of shape $\Lambda'$ and $\Lambda''$ are Euclidean-similar.

The \emph{``inversion problem''} solved in the case $p = 1$ by elliptic functions consists for a surface $\mathfrak{F}$ of arbitrary genus $p$ in the following: If $w_h^* [h = 1, \, 2, \, \dots, \, p]$ is a complex basis of the differentials of the 1st kind on $\mathfrak{F}$, then for arbitrarily prescribed numbers $\mathfrak{f}_1, \, \dots, \, \mathfrak{f}_p$, points $\mathfrak{p}_1, \, \dots, \, \mathfrak{p}_p$ on $\mathfrak{F}$ are to be found such that (with appropriate choice of the integration paths)
\[ \sum_{l=1}^p w_h^*(\mathfrak{p}_l) = \mathfrak{f}_h \quad [h = 1, \, 2, \, \dots, \, p] \tag{45} \]
holds. Following Abel's theorem, this inversion problem was posed by Jacobi and, after important preparatory work by Göpel and Rosenhain, was solved in general by Riemann and Weierstrass with the aid of the $\vartheta$-functions, certain entire transcendental functions of the $p$ arguments $\mathfrak{f}_h$.${}^{1)}$ The problem can obviously also be formulated as follows: For a given system of characters $\chi_S$, points $\mathfrak{p}_1, \, \mathfrak{p}_2, \, \dots, \, \mathfrak{p}_p$ are to be determined for which
\[ \chi_S(\mathfrak{p}_1) \cdot \chi_S(\mathfrak{p}_2) \dots \chi_S(\mathfrak{p}_p) = \chi_S \tag{46} \]
holds; and it is then, by Abel's theorem, equivalent to the problem of finding a function
\[ \Theta = \frac{\mathfrak{p}_1 \, \mathfrak{p}_2 \dots \mathfrak{p}_p}{\mathfrak{p}_0^p} \]

\textbf{1)} Riemann, Theorie der Abelschen Funktionen, Crelles Journal, vol.~54 (1857) = Werke, 2nd ed., pp.~88---142; Über das Verschwinden der Theta-Funktionen, Crelles Journal, vol.~65 (1865) = Werke, pp.~212---224. Weierstrass, Vorlesungen über die Theorie der Abelschen Transzendenten, Werke, vol.~4. Stahl, Theorie der Abelschen Funktionen, Leipzig 1896. Prym und Rost, Theorie der Prymschen Funktionen erster Ordnung, Leipzig 1911, Part 2, Section 7. Krazer, Lehrbuch der Thetafunktionen, Leipzig 1903. --- The great significance of the inversion problem lies for us today not only (and probably not even predominantly) in its value in itself, but in the magnificent series of ideas to whose creation Riemann and Weierstrass were impelled by their efforts to solve it.

\origpage{129}
with prescribed multipliers $\chi_S$ which possesses a pole of at most the $p^{\text{th}}$ order only at the point $\mathfrak{p}_0$. It thereby subordinates itself naturally under the following general question:

\emph{The functions $\Theta$ with prescribed multiplier system $\chi_S$ form a linear family $G_\chi$, that is, one does not leave the domain of these functions if one combines arbitrary ones among them linearly with constant coefficients. If $\mathfrak{d}$ is any divisor of total order $m$, the question is asked for all functions from $G_\chi$ that are multiples of $\frac{1}{\mathfrak{d}}$, and in particular for the number $A$ of linearly independent ones among them.}

For the general character system $\chi_S$ this is the same question that we treated on pp.~120---122 for the \textbf{``principal character''} $\chi_S = 1$. The results obtained there can be completely transferred; first of all the inequality
\[ A \geqq m + 1 - p, \]
\emph{in which in particular ($\mathfrak{d} = \mathfrak{p}_0^p$, $m = p$) the solvability of the inversion problem is explicitly contained.} For if $\Theta_0 = \mathfrak{d}_0$ is any function from $G_\chi$ --- that such functions exist is a lemma whose proof we shall supply presently ---, then the multiples of $\frac{1}{\mathfrak{d}}$ lying in $G_\chi$ arise by multiplying all multiples of $\frac{1}{\mathfrak{d}\mathfrak{d}_0}$ existing among the single-valued functions on $\mathfrak{F}$ by $\Theta_0$. But the number of linearly independent ones among the latter is, since $\mathfrak{d}\mathfrak{d}_0$ like $\mathfrak{d}$ has total order $m$ (see p.~120), $\geqq m + 1 - p$.

Moreover, the formula of the full Riemann-Roch theorem can also be transferred immediately. For this purpose we must draw, alongside multiplicative functions, upon the multiplicative differentials introduced into the theory by Prym.${}^{1)}$ Such a \textbf{multiplicative} or \textbf{Prym differential} $dT$ is single-valued on $\hat{\mathfrak{F}}$, analytic, free from essential singularities, and behaves under the deck transformations $S$ of $\hat{\mathfrak{F}}$ according to the equation
\[ dT(\hat{\mathfrak{p}}S) = \chi_S \cdot dT(\hat{\mathfrak{p}}), \]
where the $\chi_S$ are constant multipliers of absolute value 1. Every Prym differential can be represented by a divisor $\mathfrak{t}$, in whose numerator are recorded with the correct multipli-

\textbf{1)} Crelles Journal vol.~70 (1869), pp.~354---262\ednote{Misprint in the original: 354---262 instead of 354---362.}; vol.~71 (1870), pp.~223---236 and pp.~305---315. Cf.\ the papers by Appell cited on p.~118, further: Prym and Rost, Theorie der Prymschen Funktionen erster Ordnung, Leipzig 1911; R.~König, Berichte der kgl.\ Sächsischen Gesellschaft der Wissenschaften (mathematisch-physikalische Klasse), vol.~63 (1911), pp.~348---368.

\origpage{130}
city those points over which the zeros of $dT$ lie, and whose denominator contains the poles of $dT$. From the fact that the quotient of a Prym differential and an arbitrary Abelian (that is, single-valued on $\mathfrak{F}$) differential is a function $\Theta$ with the same multipliers as the Prym differential, it follows that by the representing divisor $\mathfrak{t}$ --- which can be prescribed arbitrarily provided only its total order is $2p - 2$ --- a Prym differential $dT$ is uniquely determined up to a constant factor; the multipliers of $dT$ are $= \chi_S(\mathfrak{t}) : \chi_S^0$. If we now understand by $B$ the number of linearly independent Prym differentials with multipliers $\chi_S^{-1}$ that are multiples of $\mathfrak{d}$, the \emph{generalized Riemann-Roch theorem} holds${}^{1)}$:
\[ A = B + (m + 1 - p). \]
For $B$ is identical with the number of linearly independent multiples of $\mathfrak{d}\mathfrak{d}_0$ found among the Abelian differentials (from which the Prym differentials in question arise by division by $\Theta_0$). If we write $\chi_S^{-1}$ instead of $\chi_S$, this relation thus holds between

the number $B$ of those linearly independent \emph{differentials} which are multiples of $\mathfrak{d}$ and possess the multiplier system $\chi_S$, on the one hand,

and the number $A$ of those linearly independent \emph{functions} which are multiples of the reciprocal divisor $\frac{1}{\mathfrak{d}}$ and possess the reciprocal multiplier system $\frac{1}{\chi_S}$, on the other hand.

To state at least one concrete special case of this theorem, let us take for $\mathfrak{d}$ the divisor ``1'' containing not a single point with an exponent different from 0; then, since apart from constants there is no function $\Theta$ without poles, we must set $A = 0$, and $A = 1$ only in the case of the principal character. Thus: \emph{the Prym differentials of the 1st kind with prescribed multipliers form a linear family of degree $p - 1$ (of degree $p$ only in the case of the principal character).}

In these considerations we had to rely on the

\emph{Lemma: If $\chi_S$ is a prescribed system of characters, there always exist functions $\Theta$ whose multipliers are $= \chi_S$.}

\textbf{1)} First formulated by E.~Ritter, Math.\ Ann.\ vol.~44, p.~314. Cf., however, the footnote on p.~122. Ritter uses, instead of divisor symbolism, the theory of forms on Riemann surfaces founded by Klein, which enables a real\ednote{Misprint in the original: ``rcale'' instead of ``reale''.} representation of divisors with the aid of multiplicative\ednote{Misprint in the original: ``mnltiplikativer'' instead of ``multiplikativer''.} forms, and in Math.\ Ann.\ vol.~47 (1896), pp.~157---221, guided by the principles introduced by Riemann into the theory of linear differential equations, generalizes his investigations to systems of $n$ forms that transform homogeneously and linearly under the execution of a deck transformation $S$.

\origpage{131}
According to Abel's theorem, finding such a function
\[ \Theta_0 = \frac{q_1 q_2 \dots q_n}{q_1^0 q_2^0 \dots q_n^0} \]
amounts to finding points $q_l, q_l^0 [l = 1, 2, \dots, n]$ such that
\[ \sum_{l=1}^n \{ w_h^*(q_l) - w_h^*(q_l^0) \} = \text{prescribed values } \mathfrak{F}_h \quad [h = 1, 2, \dots, p] \tag{47} \]
holds. Let $q_1^0, \dots q_p^0$ be $p$ mutually distinct points situated such that no differential of the 1st kind vanishes simultaneously at these $p$ points (p.~123), $z_l$ associated local uniformizing variables, and $\mathrm{K}_l$ such $z_l$-discs about $q_l^0$ that nowhere mutually overlap. The equations
\[ \sum_{l=1}^p \int\limits_{q_l^0}^{q_l} dw_h^* = \mathfrak{F}_h \quad [h = 1, 2, \dots, p] \tag{48} \]
for the unknowns $q_l$ have, when the prescribed values $\mathfrak{F}_h$ are sufficiently small, one and only one solution for which $q_l$ (together with the integration path $q_l^0 q_l$) comes to lie within $\mathrm{K}_l$, since the Jacobian
\[ \left| \frac{dw_h^*(z_l)}{dz_l} \right|_{z_1 = 0, \dots, z_p = 0} \neq 0. \]
If, however, $\mathfrak{F}_h$ are arbitrary values, one first chooses a positive integer $N$ so large that such a solution $q_l [l = 1, 2, \dots, p]$ of equations (48) is possible when on the right $\mathfrak{F}_h$ is replaced by $\frac{\mathfrak{F}_h}{N}$. Then there is also a solution of the (unaltered) equations (47) in which $n = Np$ and in the sequence $q_1 q_2 \dots q_n$ each of the points $q_1 \dots q_p$, in the sequence $q_1^0 q_2^0 \dots q_n^0$ each of the points $q_1^0 \dots q_p^0$, occurs $N$ times in all.

The \emph{inversion problem} (46) \emph{is always solvable}, as we saw above; \emph{it also possesses in general only one solution}. For, according to the generalized Riemann-Roch theorem, the solution is multi-valued (infinitely many-valued) only if there exists a Prym differential
\[ dT = \mathfrak{p}_0^p \cdot (q_1 q_2 \dots q_{p-2}) = \mathfrak{t} \]
with multipliers $\chi_S^{-1}$ which vanishes at the point $\mathfrak{p}_0$ to at least the $p^{\text{th}}$ order; but then
\[ \chi_S = \frac{\chi_S^0}{\chi_S(\mathfrak{t})} = \chi_S^0 : \prod_{i=1}^{p-2} \chi_S(q_i). \]
Within the $2p$-dimensional manifold of all possible character systems, those obtained from this formula when $q_1, \dots, q_{p-2}$ independently traverse the surface $\mathfrak{F}$ form

\origpage{132}
a manifold $\mathrm{X}$ of only $(2p - 4)$ dimensions, which takes on the role of a \emph{singular locus} for the inversion problem.${}^{1)}$

If $\chi_S$ is any character system distinct from the principal character, one can choose the points $q_1, q_2, \dots, q_{p-2}$ such that (up to a constant factor) there exists only a single Prym differential of the 1st kind
\[ dT = (q_1 \dots q_{p-2}) \cdot (p_1 p_2 \dots p_p) \]
with multipliers $\chi_S$ that vanishes at those points, --- in accordance with the theorem that the Prym differentials of the 1st kind form a $(p - 1)$-dimensional family. This choice has the consequence that the equation
\[ \prod_{l=1}^p \chi_S(p_l) = \chi_S \cdot \left( \chi_S^0 : \prod_{i=1}^{p-2} \chi_S(q_i) \right) \]
possesses only a single solution $p_1 p_2 \dots p_p$, or: that the character system standing on the left-hand side of this equation is \emph{non-singular}. On the other hand, the fraction placed in parentheses on the right belongs to $\mathrm{X}$. That there are only $p - 1$ Prym differentials of the 1st kind is accordingly, if we make use of the language of the inversion problem, equivalent to the theorem:

\emph{The singular locus $\mathrm{X}$ never transforms into itself when all character systems belonging to it are subjected to multiplication by a fixed character system $\chi_S$} (the trivial case $\chi_S = 1$ of course excepted).

This may suffice to make clear the connection of the inversion problem with the theory of multiplicative functions and differentials. \emph{Analytically}, one will attack the solution of Jacobi's problem by calling to aid an arbitrary function $f(\mathfrak{p})$ single-valued on $\mathfrak{F}$ and regular except for poles, and then seeking to determine, instead of the points $\mathfrak{p}_1, \mathfrak{p}_2, \dots, \mathfrak{p}_p$ themselves, the values of the function $f$ at the points $\mathfrak{p}_1, \mathfrak{p}_2, \dots, \mathfrak{p}_p$ from (45) in their dependence on $\mathcal{F}_1, \mathcal{F}_2, \dots, \mathcal{F}_p$. But since these quantities $f(\mathfrak{p}_1), f(\mathfrak{p}_2), \dots, f(\mathfrak{p}_p)$ are determined only up to their order, one more properly replaces them by their \emph{elementary symmetric} functions, that is, by the coefficients of that equation of the $p^{\text{th}}$ degree
\[ \lambda^p + A_1 \lambda^{p-1} + \dots + A_p = 0, \]
whose roots are the numbers
\[ \lambda = f(\mathfrak{p}_1), f(\mathfrak{p}_2), \dots, f(\mathfrak{p}_p) \]
These coefficients $A_h$, expressed in $\mathcal{F}_1, \mathcal{F}_2, \dots, \mathcal{F}_p$, are what, following Jacobi's proposal, are called \textbf{Abelian functions}. Except

\textbf{1)} In the case $p = 1$, $\mathrm{X}$ is not present at all (p.~127); in the case $p = 2$ it consists of the single character system $\chi_S^0$.

\origpage{133}
for singular systems of values $(\mathcal{F}) = (\mathcal{F}_1, \mathcal{F}_2, \dots, \mathcal{F}_p)$, which in the entire $2p$-dimensional $\mathcal{F}$-space constitute a locus of only $(2p - 4)$ dimensions, the Abelian functions are \emph{single-valued and regular-analytic}. They are furthermore \emph{$(2p)$-fold periodic}. For if $\alpha_l [l = 1, 2, \dots, 2p]$ is a basis of the closed paths on $\mathfrak{F}$ and
\[ a_h^l = \int\limits_{\alpha_l} dw_h^*, \]
then the numbers $a_1^l, a_2^l, \dots, a_p^l$ form for each value of the index $l$ a period system of the Abelian functions $A((\mathcal{F}))$:
\[ A((\mathcal{F} + a^l)) = A((\mathcal{F})) \text{ identically in } \mathcal{F}_1, \mathcal{F}_2, \dots, \mathcal{F}_p, \]
and these $2p$ period systems are linearly independent in the sense that the determinant whose $l^{\text{th}}$ row is $=$
\[ \Re a_1^l, \, \dots, \, \Re a_p^l, \, \Im a_1^l, \, \dots, \, \Im a_p^l \]
turns out to be different from 0. If one carries out explicitly the proof given here for the solvability of the inversion problem, one arrives at the result that an Abelian function $A((\mathcal{F}))$ can be represented \emph{in every finite piece of the $\mathcal{F}$-space}
\[ |\mathcal{F}_1| < M, \quad |\mathcal{F}_2| < M, \, \dots, \quad |\mathcal{F}_p| < M \]
as the quotient of two functions which in this finite piece are entirely regular-analytic.${}^{1)}$ The indeterminacy for the singular systems of values $(\mathcal{F})$ comes about because for these the numerator and denominator of the aforementioned representation both vanish. The point set formed by the points of indeterminacy permits no translations into itself other than those which the periodicity of the Abelian functions brings with it as a matter of course. --- Beyond this, it was shown by Riemann and Weierstrass through a much more penetrating analysis that the Abelian functions, even without restriction to a finite region, can be represented as quotients of entire transcendental functions, namely the $\vartheta$-functions, and they gave an explicit analytic expression for the $\vartheta$-function in the form of a rapidly convergent infinite series. This $\vartheta$-function also suffices to ground the general theory of $2p$-fold periodic functions of $p$ independent complex arguments, of which the theory of Abelian functions is only a special case.${}^{2)}$ ---

If we look back upon the insights gained in this long section, we cannot above all close our eyes to the \emph{conviction}

\textbf{1)} Cf.\ Weierstrass, Werke, vol.~4, pp.~451---456.

\textbf{2)} Cf.\ Chapter 4 in Krazer's Lehrbuch der Thetafunktionen (Leipzig 1903), where further literature is also indicated.

\origpage{134}
of the \emph{high function-theoretic significance of the number $p$}; yet in its root and its essence this number is and remains an \emph{analysis-situs quantity}. The simple basic outlines shimmering through everywhere, by which the legal order of the function-theoretic commonwealth --- intricate in its details and not quite easy to survey --- is governed, all point back to analysis situs as to an invisible ``divine lawgiver'' hovering above function-theoretic reality.

As for what belongs properly to function theory, however, we were concerned with \emph{two circles of ideas}, which may perhaps be contrasted with one another once more through the catchwords:

additive function, partial fraction decomposition, Riemann-Roch theorem ---

multiplicative function, product representation, Abel's theorem.

These two circles of ideas interpenetrate in the \emph{theory of single-valued functions on $\mathfrak{F}$}, whose construction from here proceeds effortlessly by means of the reciprocity laws obtained in §~16. The \emph{significance} of the edifice thus erected, however, is set in its proper light only when we come to know the system of singularity-free single-valued functions on a closed Riemann surface, as is to happen in the next section, from a third point of view as an \emph{algebraic function field}; for only from this standpoint do those functions appear most closely bound up with our other interests, the \emph{algebraic} and the \emph{geometric} --- in so far as these relate to the theory of algebraic curves in the plane and in higher-dimensional spaces.

\subsection*{?~18. The algebraic function field.}

Let $z$ be a function on the given closed Riemann surface $\mathfrak{F}$ which does not reduce to a constant, but assumes every value $n$-times. One can make $\mathfrak{F}$ into an $n$-sheeted covering surface spreading over the $z$-sphere by saying of a point $\mathfrak{p}$ on $\mathfrak{F}$ at which $z$ possesses the value $a$ that it lies over the point $z = a$ of the $z$-sphere. Over those finitely many values $a$ for which the $n$ points $\mathfrak{p}$ at which $z$ assumes the value $a$ are not all distinct from one another, there lie, to be sure, fewer than $n$ points of the covering surface. If at a point $\mathfrak{p}_0$ the function $z$ assumes the finite value $a$ $r$-times, then $\sqrt[r]{z - a}$ is a local uniformizing variable for $\mathfrak{p}_0$, the point $\mathfrak{p}_0$ lying over $a$ is a branch point of order $r - 1$. If $\mathfrak{U}$ is any neighborhood of $\mathfrak{p}_0$, there is a disc $|z - a| \leqq \varepsilon$ on the $z$-sphere such that over every point of this disc (except over the center) there lie exactly $r$ points of $\mathfrak{F}$ belonging to the neighborhood $\mathfrak{U}$. In $\mathfrak{p}_0$, $dz$ has a zero of the $(r - 1)$th or-

\origpage{135}der. If at the point $\mathfrak{p}_0$, $z$ assumes the value $\infty$ $s$-times, then $\sqrt[s]{\frac{1}{z}}$ is a local uniformizing variable for $\mathfrak{p}_0$, and we have a branch point of the $(s - 1)$th order; in $\mathfrak{p}_0$, $dz$ possesses a pole of order $(s + 1)$. We denote the sum of the orders of all branch points on the surface, its ``\textbf{branching number}'', by $V$. It is
\[ \text{\emph{number of zeros}} - \text{\emph{number of poles of }} dz = \Sigma(r - 1) - \Sigma(s + 1), \]
where the first sum on the right refers to all points of the surface except those lying over $z = \infty$, while the second sum concerns precisely these points. The right-hand side of the equation is
\[ = \Sigma(r - 1) + \sum_\infty (s - 1) - 2 \cdot \sum_\infty s = V - 2n, \]
the left-hand side $= 2p - 2$, we thus find:
\[ V = 2(p + n - 1). \]

This formula, which shows that $V$ is always even and allows the genus of the surface to be calculated from the number of sheets and the branching number, can also be proved purely topologically as follows. One triangulates the $z$-sphere in such a way that $z = \infty$ and those points over which branch points lie all appear as vertices, and also that over at most one of the three vertices of any triangle there lies a branch point. By considerations of a similar kind to those in ?~6, one sees that this triangulation transfers to the surface $\mathfrak{F}$, regarded as a covering surface of the $z$-sphere. Over each elementary triangle of the triangulated sphere there lie $n$ triangles of $\mathfrak{F}$, over each edge of the sphere $n$ edges of $\mathfrak{F}$. For the total number $D$ and $K$ of triangles and edges of $\mathfrak{F}$ we accordingly have
\[ D = n D_0, \quad K = n K_0, \]
if $D_0, K_0, E_0$ denote the numbers of triangles, edges, and vertices of the triangulated sphere. If over a vertex $z = a$ of the sphere there lie $h$ points with branching orders $r_1 - 1, \dots, r_h - 1$, this gives
\[ h = n - [(r_1 - 1) + \dots + (r_h - 1)] \]
vertices on $\mathfrak{F}$. The total number of vertices on $\mathfrak{F}$ is therefore $E = n E_0 - V$. Since the sphere is simply connected, Euler's polyhedron formula gives:
\[ E_0 + D_0 - K_0 = 2. \]
Therefore
\[ K - E - D = n(K_0 - E_0 - D_0) + V = V - 2n, \]
and since
\[ K - E - D + 2 = 2p \]
(?~12), as above:

\origpage{136}\[ V = 2(p + n - 1). \]
From this it can now be concluded (conversely to what was done in the preceding paragraph) that the total order of $dz$ (and thus of every differential on the surface) is $= V - 2n = 2p - 2$.

The representation of a given Riemann surface $\mathfrak{F}$ as a covering surface over the $z$-sphere can also serve to prove very simply that part of \emph{Abel's theorem} which states that the condition
\[ \sum_{k=1}^n w(\mathfrak{p}_k^0) = \sum_{k=1}^n w(\mathfrak{q}_k^0) \]
is necessary in order that
\[ \frac{\mathfrak{p}_1^0 \mathfrak{p}_2^0 \dots \mathfrak{p}_n^0}{\mathfrak{q}_1^0 \mathfrak{q}_2^0 \dots \mathfrak{q}_n^0} \]
be a function $z$ on the surface. For if in general $\mathfrak{p}_1, \mathfrak{p}_2, \dots, \mathfrak{p}_n$ are the points of $\mathfrak{F}$ lying according to this view over the point $z$ of the $z$-sphere, then, if $dw$ is an arbitrary differential of the 1st kind on $\mathfrak{F}$,
\[ \frac{dW}{dz} = \frac{dw(\mathfrak{p}_1)}{dz} + \frac{dw(\mathfrak{p}_2)}{dz} + \dots + \frac{dw(\mathfrak{p}_n)}{dz} \]
defines an everywhere regular differential $dW$ on the $z$-sphere, and such a differential does not exist except: $dW \text{ identically} = 0$. If we draw any curve on the $z$-sphere, e.g.\ a meridian leading from the south pole $z = 0$ to the north pole $z = \infty$, the points of $\mathfrak{F}$ lying above it form curves $\gamma_1, \dots, \gamma_n$ (which may meet at branch points), and by integrating $dW = 0$ it follows that:
\[ \int\limits_{\gamma_1} dw + \dots + \int\limits_{\gamma_n} dw = 0. \tag{49} \]
This equation is equivalent to the necessary condition stated in Abel's theorem. However, one does not arrive at the insight that this condition is sufficient by this route; but nothing prevents us, following Abel, from replacing the differential of the 1st kind $dw$ by any differential on the surface provided with poles but otherwise regular, and deriving for such a differential a relation similar to (49). We shall, however, not pursue this further. ---

\emph{Let $z$ again be a function on the surface $\mathfrak{F}$ which assumes every value $n$\,times. If $f$ is then any function on the surface, $f$ satisfies identically an equation of the $n^{\text{th}}$ degree}
\[ f^n + r_1(z) f^{n-1} + \dots + r_{n-1}(z) f + r_n(z) = 0, \tag{50} \]
\emph{in which the $r_i(z)$ are rational functions of $z$.} This means: If $\mathfrak{p}_0$ is any point on $\mathfrak{F}$, $t$ an associated local uniformizing variable, and $z = z(t)$, $f = f(t)$ are the expansions of $z$ and $f$ in integral powers in the neighborhood of $\mathfrak{p}_0$, then the left-hand side of (50) becomes iden-

\origpage{137}tically 0 when one substitutes the power series $z(t)$ for $z$ and the power series $f(t)$ for $f$. To prove this equation (50), we first exclude on the $z$-sphere the point $z = \infty$ and those points over which branch points lie or poles of the function $f$. For every other $z$-value we form the number
\[ r_1(z) = f(\mathfrak{p}_1) + f(\mathfrak{p}_2) + \dots + f(\mathfrak{p}_n), \]
where the sum on the right extends over the $n$ points lying over $z$. It is independent of the order of their numbering. Since we can arrange this numbering so that in the neighborhood of a non-excluded point
\[ z_0 \colon \mathfrak{p}_1^0, \, \mathfrak{p}_2^0, \, \dots, \, \mathfrak{p}_n^0 \]
$f(\mathfrak{p}_1), f(\mathfrak{p}_2), \dots, f(\mathfrak{p}_n)$ become power series in $z - z_0$, $r_1(z)$ is a regular-analytic function at all non-excluded points. At the excluded points this function can have no essential singularities, hence only poles; therefore it is a rational function of $z$. In the same way one recognizes that the remaining elementary symmetric functions of $f(\mathfrak{p}_1), \dots, f(\mathfrak{p}_n)$ are expressed rationally in terms of $z$.

If over $z_0$ there lie the $n$ distinct points $\mathfrak{p}_1^0, \dots, \mathfrak{p}_n^0$, we can choose the function $f$ so that it possesses distinct values at these $n$ places. If $d\tau_k (k = 1, \dots, n)$ is a differential which possesses a pole only at the point $\mathfrak{p}_k^0$ and has there the principal part $-\frac{dz}{(z - z_0)^2}$, one can set, for example,
\[ f = (z - z_0)^2 \left( C_1 \frac{d\tau_1}{dz} + \dots + C_n \frac{d\tau_n}{dz} \right) \]
by choosing for $C_1, \dots, C_n$ any $n$ distinct constants. The equation of the $n$\,th degree for $f$ is then \textbf{irreducible}; that is, its left-hand side, the polynomial of the $n$\,th degree in the variable $u$
\[ F_z(u) \equiv u^n + r_1(z) u^{n-1} + \dots + r_n(z) \]
cannot be split into two polynomials $F_z^{(1)}(u) \cdot F_z^{(2)}(u)$ whose coefficients are likewise rational functions of $z$. For suppose this were possible; $f$ can be expanded in the neighborhood of $\mathfrak{p}_1^0$ into a power series in terms of the local uniformizing variable $z - z_0$ belonging to $\mathfrak{p}_1^0$. Suppose this power series, substituted for $u$, satisfies the equation $F_z^{(1)}(u) = 0$. If we connect $\mathfrak{p}_1^0$ with any of the points $\mathfrak{p}_k^0$ by a curve on $\mathfrak{F}$ whose trace on the $z$-sphere passes through none of the excluded points, the equation $F_z^{(1)}(u) = 0$ must persist for all function elements $(z, f)$ along this curve. Consequently $f(\mathfrak{p}_k^0)$ is also a root of the equation $F_{z_0}^{(1)}(u) = 0$; the latter thus has $n$ distinct roots and must therefore be of degree $n$

\origpage{138}; $F_z^{(2)}(u)$ can therefore only be of degree 0 in $u$; and a genuine factorization of the hypothesized kind does not occur.

To each point $\mathfrak{p}$ of $\mathfrak{F}$ belongs a function element $(z, f)$ satisfying the equation $F_z(u) = 0$. For two distinct points these two function elements are always distinct, and the elements belonging to all points $\mathfrak{p}$ exhaust the totality of those satisfying that equation. In other words: \emph{The totality of those function elements which satisfy the irreducible algebraic equation $F_z(u) = 0$ constitutes a single analytic configuration in the Weierstrassian sense. This analytic configuration, considered as a Riemann surface, is conformally equivalent to the given Riemann surface; the given surface is the Riemann surface belonging to the algebraic configuration defined by the equation $F_z(u) = 0$.}${}^{1)}$

Every function $f^*$ on the surface regular except for poles can be expressed rationally in terms of $f$ with coefficients that are rational functions of $z$:
\[ f^* = R_0(z) + R_1(z) f + \dots + R_{n-1}(z) f^{n-1}, \]
and every function formed rationally from $z$ and $f$ is regular except for poles on the surface. If one regards $z$ as the independent variable and $f$ as the algebraic function of $z$ defined by the equation $F_z(u) = 0$, then those functions which can be expressed rationally in terms of $f$ with coefficients rational in $z$ form an \textbf{algebraic function field}. Our assertion can therefore be stated as follows: \emph{If all functions regular except for poles on a closed Riemann surface are regarded as functions of an arbitrary one among them (which need only not be a mere constant), they form an algebraic function field.}

For the proof we set (\emph{Lagrange's interpolation formula})
\[ \frac{G_z(u)}{F_z(u)} = \frac{f^*(\mathfrak{p}_1)}{u - f(\mathfrak{p}_1)} + \frac{f^*(\mathfrak{p}_2)}{u - f(\mathfrak{p}_2)} + \dots + \frac{f^*(\mathfrak{p}_n)}{u - f(\mathfrak{p}_n)}. \]
$\mathfrak{p}_1, \mathfrak{p}_2, \dots, \mathfrak{p}_n$ again denote the $n$ points of $\mathfrak{F}$ lying over the point $z$ of the $z$-sphere. By the same reasoning as above, one finds that the coefficients of the polynomial $(n - 1)^{\text{th}}$ degree $G_z(u)$ are rational functions of $z$. For all values $z$ for which $f(\mathfrak{p}_1), \dots, f(\mathfrak{p}_n)$ are distinct from one another, it follows that
\[ f^*(\mathfrak{p}_1) = \left[ \frac{G_z(u)}{F_z'(u)} \right]_{u = f(\mathfrak{p}_1)} \quad \left[ F_z'(u) = \frac{dF_z(u)}{du} \right]. \]
One can (say by applying the ``Euclidean division algorithm'';

\origpage{139}cf.\ e.g.\ Weber, Algebra, 2nd ed., Braunschweig 1899, vol.~1, p.~41) determine two polynomials $H_z(u), L_z(u)$ with coefficients rational in $z$ such that
\[ H_z(u) F_z'(u) + L_z(u) F_z(u) = 1 \]
holds. The equation
\[ f^*(\mathfrak{p}) = [G_z(u) H_z(u)]_{u = f(\mathfrak{p})} \]
is then an identity on the surface.

If instead of $z$ one takes any other function on the surface, $\bar{z}$, as the independent variable, one can furthermore determine in infinitely many ways a function $\bar{f}$ on the surface of such a kind that all functions can be expressed rationally in terms of $\bar{z}$ and $\bar{f}$. Between $\bar{z}$ and $\bar{f}$ there subsists an irreducible algebraic equation $\bar{F}_{\bar{z}}(\bar{f}) = 0$. $\bar{z}$ and $\bar{f}$ are rational functions of the variables $z, f$ connected by the equation $F_z(f) = 0$; and conversely $z$ and $f$ are rational functions of the variables $\bar{z}, \bar{f}$ connected by the equation $\bar{F}_{\bar{z}}(\bar{f}) = 0$. By the ``\textbf{birational transformation}'' $(z, f) \longleftrightarrow (\bar{z}, \bar{f})$ the equations
\[ F_z(u) = 0 \quad \text{und} \quad \bar{F}_{\bar{z}}(\bar{u}) = 0 \]
pass into one another. The degree of this equation is of course by no means an invariant under birational transformation, but the genus $p$ certainly is.

If there exists on $\mathfrak{F}$ a function $z$ that assumes every value only once, the associated algebraic field is the field of rational functions of $z$ (and $p = 0$). If on $\mathfrak{F}$ there exists no function assuming every value only once, but there is one, $z$, that assumes every value exactly twice, we can assume the algebraic equation determining the associated field (which must be quadratic) to be of the form:
\[ u^2 = (z - e_1)(z - e_2) \dots (z - e_l), \]
where the $e_i$ are all distinct from one another. These $l$ points and, if $l$ is odd, also the point $\infty$ are branch points of the 1st order, and the genus $p$ is accordingly $= \frac{l}{2} - 1$ if $l$ is even, $= \frac{l - 1}{2}$ if $l$ is odd. We see: $l = 1$ or $= 2$ leads back again to the rational field $p = 0$; $l = 3$ or $= 4$ entails $p = 1$, which is the elliptic case; if $l > 4$, we obtain the so-called \textbf{hyperelliptic} function fields. --- In an arbitrary algebraic function field of genus $p = 1$ there always exists a function with two prescribed poles, hence a function assuming every value only twice; in every function field of genus 2 we obtain a function of the same kind by dividing two linearly inde-

\origpage{140}pendent Abelian differentials of the 1st kind by one another. Only from $p = 3$ onwards is the hyperelliptic case no longer the general one.${}^{1)}$

If one does not, as was done here, take the Riemann surface as one's point of departure, but rather a specific algebraic equation $F_z(u) = 0$ (as is done in the Weierstrassian and other theories), one may regard it as a natural demand to construct in a purely algebraic way all those functions and differentials whose existence was inferred in the preceding sections with the aid of Dirichlet's principle, as rational expressions in $z, f$ [the differentials in the form $R(z, f)dz$]. But as soon as what is given is not an algebraic equation but the Riemann surface, the function-theoretic route discussed here following Riemann must on the contrary appear as the natural one. However important those algebraic construction principles may be, especially with regard to special applications of the theory, --- one may nevertheless designate Riemann's standpoint as the higher one, since from it a more comprehensive and deeper insight into the peculiar laws governing this realm of mathematical knowledge becomes possible than can be gained by any other route. Even if one follows Weierstrass's conceptual formations, one will, as I mentioned earlier, not be able to circumvent the conception of the analytic configuration as a two-dimensional manifold without doing violence to the subject, and it is then only a short step to place the topological properties belonging to this manifold --- whose incisive function-theoretic significance has meanwhile emerged sufficiently --- ahead of all others as the most primitive. Beyond this, it is characteristic of Riemann's manner of treat-

\textbf{1)} The generalization of this theorem to arbitrary (non-closed) surfaces was proved by P.~Koebe, Comptes Rendus, 1~July 1909. Cf.\ E.~Freundlich, Funktionen mit vorgeschriebenem unendlichbl?ttrigen Existenzbereich, G?ttingen dissertation 1910.

\textbf{1)} The outline of the theory of algebraic functions that we could give here is only incomplete. More detailed accounts will be found by the reader, in addition to the works already cited of Riemann, Weierstrass, Klein, C.~Neumann, Stahl, Hensel-Landsberg, in the following presentations: Clebsch and Gordan, Theorie der Abelschen Funktionen, Leipzig 1866 (curve-theoretic). Brill and Noether, ?ber die algebraischen Funktionen und ihre Anwendung in der Geometrie, Math.\ Ann.\ vol.~7 (1874), pp.~269---310 (curve-theoretic); Die Entwicklung der Theorie der algebraischen Funktionen, Bericht der Deutschen Mathematiker-Vereinigung vol.~3, Berlin 1894. Dedekind and Weber, Theorie der algebraischen Funktionen einer Ver?nderlichen, Crelles Journal vol.~92 (1882), pp.~181---290 (arithmetic); also presented in Weber, Algebra, vol.~III, 2nd ed., Braunschweig 1908, pp.~623\,ff. Klein-Fricke, Theorie der elliptischen Modulfunktionen (1890---92), vol.~I, Section~III, Chaps.~1, 2, and vol.~II, Section~VI, Chap.~1. Klein, Riemannsche Fl?chen I, II, autographed lectures, G?ttingen 1892/93. Appell et Goursat, Th?orie des fonctions alg?briques, Paris 1895. Baker, Abel's theorem and the allied theory incl.\ the theory of the Thetafunctions, Cambridge 1897. Fields, Theory of the algebraic functions of a complex variable, Berlin 1906. Stahl, Abri? einer Theorie der algebraischen Funktionen einer Ver?nderlichen in neuer Fassung (posthumous work, edited by L?ffler and Noether), Leipzig 1911.

\origpage{141}ment that in it everywhere not the analytic configuration, but the Riemann surface is viewed as the given entity, and the construction of an associated analytic configuration forms precisely a principal component of the problems to be solved. In Riemann's own presentation, to be sure, this standpoint does not yet stand out with that complete clarity with which we can now extract it on the basis of the works of Prym, Dedekind${}^{1)}$, C.~Neumann, and especially of Klein.

Every Riemann surface of genus $p$ can, as we saw, be represented as a many-sheeted covering surface over the sphere (with finitely many branch points, but without boundaries). This ``normal form'' can, however, even if one normalizes the number of sheets $n$ by the condition $n = p + 1$ (which can always be achieved), still be produced in the most varied ways. A very much higher fundamental significance belongs to the essentially uniquely determined normal form of Riemann surfaces of arbitrary genus which is provided by uniformization theory (theory of automorphic functions).

\subsection*{?~19. Uniformization.}

In the theory of uniformization, the Weierstrassian and Riemannian circles of ideas fuse into a complete unity. While in Weierstrass the analytic configuration $(z, u)$ is described \emph{at each individual point} by a \emph{particular} representation with the aid of a parameter $t$ (the ``local uniformizing variable''): $z = z(t)$, $u = u(t)$, Riemann, to be sure, obtains a \emph{unified} representation $z = z(\mathfrak{p})$, $u = u(\mathfrak{p})$ of the whole configuration, but is thereby forced to conceive the parameter $\mathfrak{p}$ as a point on a Riemann surface (not as a complex variable in the ordinary sense); in uniformization theory, the concern is to produce for an analytic configuration a unified representation $z = z(t)$, $u = u(t)$ with the aid of a parameter $t$, the \textbf{uniformizing variable}, varying in a domain of the schlicht complex plane. As the true founders of the \emph{theory of automorphic functions}, to which this problem leads, F.~Klein and H.~Poincar?${}^{2)}$ are to be named, whose general conceptions and results in the literature

\textbf{1)} Prym represents this standpoint in his papers from 1869 onwards; of Dedekind, the paper on the elliptic modular function from the year 1877 cited on p.~35 comes into consideration here.

\textbf{2)} By Poincar? see, besides numerous Comptes-Rendus notes from the years 1881/82, especially the memoirs in the Acta Mathematica, vols.~1, 3, 4, 5 (1882/84); by Klein the papers in the Math.\ Ann., vols.~19, 20, 21 (1882/83), further the comprehensive presentation completed recently: Fricke and Klein, Vorlesungen ?ber die Theorie der automorphen Funktionen, Leipzig 1897---1912.

\origpage{142}appear prepared by important, albeit more special investigations, especially by Riemann, Schwarz, Fuchs, Dedekind, Klein, and Schottky. The proof of the possibility of uniformization on the basis of the idea of the covering surface was provided completely only quite recently (1907) by P.~Koebe and H.~Poincar?.${}^{1)}$ It is above all thanks to Klein, Poincar?, and Koebe that today the theory of uniformization, which may claim a central position within complex function theory, stands before us as a mathematical edifice of exceptional harmony and grandeur.${}^{2)}$ --- The fundamental idea of the proof given in the following, of deducing the existence of the uniformizing variable from Dirichlet's principle, is due to Hilbert.${}^{3)}$

The uniformizing variable $t$ which we seek shall be so constituted that it is suitable as a local uniformizing variable at every place of the given surface $\mathfrak{F}$. It will therefore have to be a single-valued, regular-analytic function, apart from poles of the 1st order, on the universal covering surface $\widetilde{\mathfrak{F}}$. If we seek that $t$ to which belongs the strongest uniformizing power, we will seek to determine $t$ in such a way that it never assumes the same value at two distinct points of the surface $\widetilde{\mathfrak{F}}$, hence maps $\widetilde{\mathfrak{F}}$ bijectively and conformally onto a domain of the $t$-sphere. Then not only will the functions on the base surface $\mathfrak{F}$ be representable as single-valued functions of $t$, but the much more comprehensive totality of those functions (in general infinitely many-valued on $\mathfrak{F}$) which arise from a function element on $\mathfrak{F}$ that can be continued without branching and unrestrictedly along all paths in $\mathfrak{F}$. And since $\widetilde{\mathfrak{F}}$ (in contrast to $\mathfrak{F}$) is simply connected, the possibility of such a mapping does not conflict with the topological properties of $\widetilde{\mathfrak{F}}$. The assumption adopted in the preceding sections, that $\mathfrak{F}$ is closed, we can now gladly drop, since it brings with it no simplification whatsoever for uniformization theory.

\textbf{1)} Poincar?, Acta Mathematica vol.~31 (1908), pp.~1---63; Koebe, Nachrichten der K.~Ges.\ d.~Wissensch.\ zu Göttingen 1907, pp.~191---210 and pp.~633---669.

\textbf{2)} Cf.\ the compilation of the recent literature in Koebe, ?ber die Uniformisierung der algebraischen Kurven, I [Math.\ Ann.\ vol.~67, 1909, pp.~146 to 149], II [Math.\ Ann.\ vol.~69, 1910, pp.~2---3] and III [Math.\ Ann.\ vol.~72, 1912, pp.~438---439]; further Koebe, ?ber die Uniformisierung beliebiger analytischer Kurven I [Crelles Journal vol.~138, 1910, p.~195] and II [Crelles Journal vol.~139, 1911, pp.~251\,ff.]. For a general orientation regarding the results and problems of this part of function theory, the report on the Karlsruhe proceedings (1911) in the Jahresbericht der Deutschen Mathematiker-Vereinigung vol.~21, 1912, pp.~153---166 serves admirably.

\textbf{3)} Zur Theorie der konformen Abbildung, G?ttinger Nachrichten, 1909, pp.~314 to 323.

\origpage{143}We obtain the sought-after uniformizing variable simply by applying Dirichlet's principle not to $\mathfrak{F}$, but to the covering surface $\widetilde{\mathfrak{F}}$. We choose on $\widetilde{\mathfrak{F}}$ a point $O$ with the local uniformizing variable $z_0$ and construct with the aid of Dirichlet's principle that potential function $U$, regular on the whole of $\widetilde{\mathfrak{F}}$ apart from the point $O$, which behaves at $O$ like $\Re \frac{1}{z_0}$ and possesses the property that

1.\ the Dirichlet integral of $U$ extended over the entire surface $\widetilde{\mathfrak{F}}$ with the exclusion of an arbitrarily small $z_0$-disc about $O$ is finite,

2.\ for every continuously differentiable function $w$ on $\widetilde{\mathfrak{F}}$ with finite Dirichlet integral which vanishes in the neighborhood of $O$, the variation
\[ \mathbf{D}(U, w) = 0 \]
becomes.

$U$ gives rise to a differential $d\tau$ on $\widetilde{\mathfrak{F}}$, and this must, since $\widetilde{\mathfrak{F}}$ is simply connected, be the differential of a specific function
\[ \tau = U + iV \]
whose real part coincides with $U$ and which is regular-analytic everywhere, apart from the point $O$, but possesses at $O$ a pole of the 1st order. $\tau$ is then a uniformizing variable such as we seek. The proof of this fact succeeds in a very elegant manner with the aid of the following deduction due to Mr.~Koebe.${}^{1)}$

We show first:

\emph{If $V_0$ is any real constant, those points on $\widetilde{\mathfrak{F}}$ at which $V > V_0$ form a single domain, as do those at which $V < V_0$.}

For $O$, $\frac{1}{\tau}$ is a local uniformizing variable, and let $\mathrm{K}_0 \colon \left|\frac{1}{\tau}\right| \leqq a_0$ be a $\frac{1}{\tau}$-disc about $O$. If $\mathfrak{E}(V_0)$ is that closed set on $\widetilde{\mathfrak{F}}$ consisting of all points at which $V = V_0$, then certainly only two of the domains determined by $\mathfrak{E}(V_0)$ have points lying in $\mathrm{K}_0$. If our assertion were false, there would among the domains determined by $\mathfrak{E}(V_0)$ be one, let it be called $\mathfrak{G}$, that does not penetrate into the neighborhood $\mathrm{K}_0$ of $O$. Let $\varphi(u), \psi(u)$ now be any two continuous and continuously differentiable functions defined for all real $u$-values. We form the following function $w$ on the surface $\widetilde{\mathfrak{F}}$:

\textbf{1)} Über die Hilbertsche Uniformisierungsmethode, Nachrichten der K.~Ges.\ d.~Wissensch.\ zu Göttingen 1910, pp.~61---65.

\origpage{144}\[ w = \left\{ \begin{aligned} &\varphi(U)\psi(V) &&\text{for all points within } \mathfrak{G}, \\ &0 &&\text{for all points not belonging to } \mathfrak{G} \text{.} \end{aligned} \right. \]
It will be everywhere continuously differentiable if
\[ \psi(V_0) = 0, \quad \psi'(V_0) = 0 \quad \left[ \psi' = \frac{d\psi}{du} \right] \]
holds. In the neighborhood $\mathrm{K}_0$ of $O$, $w$ vanishes identically. If $\mathfrak{p}$ is any point in $\mathfrak{G}$ and $z = x + iy$ a local uniformizing variable for $\mathfrak{p}$, then
\[ \begin{aligned} \frac{\partial w}{\partial x} &= \varphi'(U)\psi(V)\,\frac{\partial U}{\partial x} + \varphi(U)\psi'(V)\,\frac{\partial V}{\partial x} \\ &= \varphi'(U)\psi(V)\,\frac{\partial U}{\partial x} - \varphi(U)\psi'(V)\,\frac{\partial U}{\partial y}, \\ \frac{\partial w}{\partial y} &= \varphi'(U)\psi(V)\,\frac{\partial U}{\partial y} + \varphi(U)\psi'(V)\,\frac{\partial U}{\partial x}. \end{aligned} \]
If $\varphi, \psi, \varphi', \psi'$ are bounded functions, the Dirichlet integral of $w$ extended over the whole of $\widetilde{\mathfrak{F}}$ will thus be finite. Under the stated assumptions $\mathbf{D}(U, w) = 0$ would therefore have to hold. Now:
\[ \frac{\partial w}{\partial x} \frac{\partial U}{\partial x} + \frac{\partial w}{\partial y} \frac{\partial U}{\partial y} = \varphi'(U)\psi(V)\left[ \left(\frac{\partial U}{\partial x}\right)^2 + \left(\frac{\partial U}{\partial y}\right)^2 \right]. \]
If we therefore take care that $\varphi', \psi$ are positive for all values of their argument ($\psi$ except for $u = V_0$), we arrive at a contradiction.${}^{1)}$

From the fact thus proved and the circumstance that $\widetilde{\mathfrak{F}}$ is simply connected, one can draw the following conclusions regarding the behavior of $\tau$:

1.\ \emph{$d\tau$ has nowhere a zero.} For if at a point $\mathfrak{p}_0$ of $\widetilde{\mathfrak{F}}$, where $\tau$ possesses the value $\tau_0$, $d\tau = 0$ were to hold, then not $\tau - \tau_0$, but $\sqrt[r]{\tau - \tau_0}$ ($r$ integer and $\geqq 2$) would be a local uniformizing variable for $\mathfrak{p}_0$; let us assume that $r = 2$ (for higher $r$ the proof proceeds analogously). I set $\tau - \tau_0 = \sigma^2$ and draw in the complex $\sigma$-plane a circle $\mathrm{K}$ with center $\sigma = 0$, so small that it appears as the conformal image produced by the function $\sigma$ of a certain neighborhood of the point $\mathfrak{p}_0$ on $\widetilde{\mathfrak{F}}$. I take four points $\mathfrak{p}_1, \mathfrak{p}_2, \mathfrak{q}_1, \mathfrak{q}_2$ lying in $\mathrm{K}$, as indicated by Figure~24, which are connected criss-cross by two rectilinear segments $\alpha, \beta$ intersecting at the origin. In $\mathfrak{p}_1$ and

\textbf{1)} An ansatz satisfying all requirements is obtained, for example, if with the aid of the functions
\[ \alpha(u) = \operatorname{arctg} u \quad \left( -\frac{\pi}{2} < \alpha(u) < \frac{\pi}{2} \right), \quad \beta(u) = \frac{u^2}{1 + u^2}; \]
\[ \varphi(u) = \alpha(u), \quad \psi(u) = \beta(u - V_0) \]
we form.

\origpage{145}$\mathfrak{p}_2$, $V > V_0$, in $\mathfrak{q}_1, \mathfrak{q}_2$ on the contrary $V < V_0$. This figure I imagine transferred to the surface $\widetilde{\mathfrak{F}}$, which let be triangulated in a definite manner${}^{1)}$; then I can connect $\mathfrak{p}_1\,\mathfrak{p}_2$ by a polygonal path $\alpha'$ in all points of which $V > V_0$, likewise $\mathfrak{q}_1\,\mathfrak{q}_2$ by a polygonal path $\beta'$ in all points of which $V < V_0$.${}^{2)}$ The two closed curves $\alpha + \alpha'$, $\beta + \beta'$ intersect only at the point $\mathfrak{p}_0$, and from this it follows, exactly as on p.~46, that $\beta + \beta'$ does not divide the surface $\widetilde{\mathfrak{F}}$. If I cross a segment belonging to $\beta'$ by an elementary segment $\alpha^*$, I can connect its endpoints with each other by a polygonal path that nowhere meets $\beta + \beta'$. Together with $\alpha^*$ this forms a polygon that leaves $\widetilde{\mathfrak{F}}$ undivided. Such a polygon, however, cannot exist because $\widetilde{\mathfrak{F}}$ is simply connected.

\rmfigure{figures/fig-24.png}{Fig.~24.}{Circle K with polygonal paths alpha, beta and alpha', beta'}

2. By saying of a point of the surface $\widetilde{\mathfrak{F}}$ at which $\tau$ possesses the value $\tau_0$ that it lies over the point $\tau_0$ of the $\tau$-sphere, $\widetilde{\mathfrak{F}}$ becomes a covering surface $\widetilde{\mathfrak{F}}_\tau$ over the $\tau$-sphere or the $\tau$-plane. This is, by what was just proved under 1., unramified. Over $\tau = \infty$ lies the single point $O$. Starting from $\tau = \infty$, I trace the line $V = V_0$ in the $\tau$-plane in the direction from smaller to larger $U$-values, and on $\widetilde{\mathfrak{F}}_\tau$ a continuously moving point $\tilde{\mathfrak{p}}$ starting from $O$, which always remains over the point describing this line in the $\tau$-plane. If before returning to $\infty$ I encounter no boundary, $\tilde{\mathfrak{p}}$, since over $\infty$ lies only the one point $O$, describes a closed curve $\tilde{\gamma}$ simply covering the line $V = V_0$ of the $\tau$-plane. If, however, I encounter a boundary, I obtain a line $\tilde{\gamma}_1$ on $\widetilde{\mathfrak{F}}_\tau$ simply covering a certain segment $U < U_1$ of the straight line $V = V_0$. Then I trace the straight line $V = V_0$ proceeding further from larger to smaller values of $U$ and obtain on $\widetilde{\mathfrak{F}}_\tau$ a line $\tilde{\gamma}_2$ simply covering a segment $U > U_2$ of this straight line. Together $\tilde{\gamma}_1 + \tilde{\gamma}_2$ form an open curve without end $\tilde{\gamma}$ passing through $O$ on $\widetilde{\mathfrak{F}}_\tau$.${}^{3)}$ In each of the two cases the line $\tilde{\gamma}$ divides the surface $\widetilde{\mathfrak{F}}_\tau$, since the latter is simply con-

\textbf{1)} On $\widetilde{\mathfrak{F}}$ I may then no longer designate $\alpha$, $\beta$ as ``segments''.

\textbf{2)} Naturally I can assume that $\beta'$ has no other points in common with $\beta$ other than the two endpoints; cf.\ p.~46.

\textbf{3)} A \textbf{curve without end} has a representation
\[ \tilde{\mathfrak{p}} = \tilde{\mathfrak{p}}(\lambda) \quad [0 < \lambda < 1], \]
and the points $\tilde{\mathfrak{p}}$ belonging to it form, even though no curve point corresponds to the parameter values $\lambda = 0$ and $\lambda = 1$, a \emph{closed} set.

\origpage{146}nected, into two domains $\mathfrak{G}'$, $\mathfrak{G}''$. If there were further points on $\widetilde{\mathfrak{F}}$ at which $V = V_0$ other than on $\tilde{\gamma}$, say in $\mathfrak{G}'$, there would in $\mathfrak{G}'$ be both points at which $V < V_0$ and points at which $V > V_0$. The point set $\mathfrak{E}(V_0)$ would accordingly have to determine at least three domains. With the line $\tilde{\gamma}$, therefore, the points at which $V = V_0$ are exhausted. As a covering surface of the $\tau$-sphere $\widetilde{\mathfrak{F}}$ is thus everywhere at most two-sheeted. \emph{A value $U_0 + iV_0$ is on $\widetilde{\mathfrak{F}}$ certainly assumed once and only once if $V = V_0$ on $\widetilde{\mathfrak{F}}$ is a closed line.}

We wish to show more precisely that $\tilde{\gamma}$ must always divide the surface $\widetilde{\mathfrak{F}}$. For if that is not the case, we construct a two-sheeted unramified unbounded covering surface over $\widetilde{\mathfrak{F}}$, which we obtain by slitting $\widetilde{\mathfrak{F}}$ along $\tilde{\gamma}$, producing two copies of the surface $\widetilde{\mathfrak{F}}$ thus cut open, and joining their cut edges criss-cross to one another. Expressed more abstractly this means (cf.\ the analogous construction on pp.~31\,f.): To each point $\tilde{\mathfrak{p}}$ of $\widetilde{\mathfrak{F}}$ we assign two ``points lying above it'' $\tilde{\mathfrak{p}}^1, \tilde{\mathfrak{p}}^2$. If $\tilde{\mathfrak{p}}_0$ is a point not lying on $\tilde{\gamma}$, $\tau_0 = \tau(\tilde{\mathfrak{p}}_0)$, $\mathrm{K}$ an arbitrary $(\tau - \tau_0)$-disc containing no point of $\tilde{\gamma}$, then those points $\tilde{\mathfrak{p}}^1$ (with the upper index $1$) which lie over the points $\tilde{\mathfrak{p}}$ situated in the interior of $\mathrm{K}$ form a ``neighborhood'' of $\tilde{\mathfrak{p}}_0^1$, and those points $\tilde{\mathfrak{p}}^2$ which lie over the same points $\tilde{\mathfrak{p}}$ form a ``neighborhood'' of $\tilde{\mathfrak{p}}_0^2$. If on the other hand $\tilde{\mathfrak{p}}_0$ lies on $\tilde{\gamma}$, let $\mathrm{K}$ denote an arbitrary $(\tau - \tau_0)$-disc. Those points $\tilde{\mathfrak{p}}^1$ whose trace points $\tilde{\mathfrak{p}}$ lie within $\mathrm{K}$ and satisfy the condition $V \geqq V_0$ shall, together with all points $\tilde{\mathfrak{p}}^2$ lying over the interior points $\tilde{\mathfrak{p}}$ of $\mathrm{K}$ satisfying the condition $V < V_0$, form a ``neighborhood'' of $\tilde{\mathfrak{p}}_0^1$, and analogously the neighborhood of $\tilde{\mathfrak{p}}_0^2$ shall be declared. Since in the latter case all points situated in $\mathrm{K}$ at which $V = V_0$ certainly belong to $\tilde{\gamma}$, this definition of the concept of neighborhood is in harmony with all requirements to be made of such a definition. If $\tilde{\gamma}$ does not divide, it is clear that the manifold just declared also satisfies the condition that any two of its points can be connected by a continuous curve. The existence of such a covering surface, however, contradicts the fact that $\widetilde{\mathfrak{F}}$ is simply connected.

\rmfigure{figures/fig-25.png}{Fig.~25.}{Rectangular strip in the tau plane with U = U_0, V = V_0' and V = V_0''}

The final step of the proof consists in demonstrating \emph{the theorem that there can be at most a single real number $V_0$ for which the associated line $V = V_0$ on $\widetilde{\mathfrak{F}}$ is open.} For if there were two such lines $\tilde{\gamma}'$, $\tilde{\gamma}''$: $V = V_0'$, or $V = V_0''$, one would call to aid a closed line $U = U_0$ running entirely in $\mathrm{K}_0$ (such a large positive value

\origpage{147}being taken for the constant $U_0$). Let $U > U_2'$ be the one segment $\tilde{\gamma}_2'$ of the curve $\tilde{\gamma}'$; $U > U_2''$ the segment $\tilde{\gamma}_2''$ of $\tilde{\gamma}''$. Over the heavily drawn line of the $\tau$-plane in Figure~25 lies a curve without end simply covering it on $\widetilde{\mathfrak{F}}_\tau$. On account of the simple connectivity, this divides $\widetilde{\mathfrak{F}}$ into two domains, and let $\mathfrak{G}$ be that one of the two domains which does not contain the point $O$. We again set
\[ w = \begin{cases} \varphi(U)\psi(V) & \text{innerhalb } \mathfrak{G}, \\ 0 & \text{au?erhalb } \mathfrak{G}. \end{cases} \]
In order for this function to be continuously differentiable on $\widetilde{\mathfrak{F}}$, we must have
\[ \varphi(U_0) = \varphi'(U_0) = 0, \quad \begin{gathered} \psi(V_0') = \psi'(V_0') = 0, \\ \psi(V_0'') = \psi'(V_0'') = 0 \end{gathered} \]
hold. Let $\varphi$, $\varphi'$, $\psi$, $\psi'$ be bounded and $\varphi'$ except for $u = U_0$, $\psi$ except for $u = V_0'$ and $u = V_0''$ positive.${}^{1)}$ Then a contradiction with the equation $\mathbf{D}(U, w) = 0$ comes about. Therewith is proved:

\emph{$\tau$ maps the surface $\widetilde{\mathfrak{F}}$ bijectively and conformally either onto the full sphere (1st Case)}

or

\emph{onto the sphere with the exception of one point $\tau_0$ (2nd Case)}

or

\emph{onto the sphere with the exception of a slit $V = V_0$, $U_1 \leqq U \leqq U_2$ (3rd Case).}

We replace $\tau$ by a somewhat different uniformizing variable $t$. In the first case, to be sure, we shall have $\tau = t$. In the second we ensure that the one point of the sphere which is omitted is the point at infinity; one therefore sets $t \doteq \frac{1}{\tau - \tau_0}$; $t$ maps the covering surface onto the entire plane (without the point at infinity). In the third case we can first achieve by an entire linear transformation that the slit is given by
\[ V = 0, \quad -1 \leqq U \leqq +1 \]
with the aid of the formula
\[ \tau = \frac{1}{2} \left( t + \frac{1}{t} \right) \]
the slit $\tau$-sphere is then mapped conformally onto the\ednote{Misprint in the original: ``des'' instead of ``das''.} interior of the unit disc $|t| < 1$ of the $t$-plane.

The described construction of the uniformizing variable $t$ takes place in two clearly separated steps. First the surface $\widetilde{\mathfrak{F}}$ solves the pro-

\textbf{1)} All requirements are satisfied by [cf.\ footnote to p.~144]:
\[ \varphi = \alpha(u - U_0) \cdot \beta(u - U_0), \quad \psi = \beta(u - V_0') \cdot \beta(u - V_0''). \]

\origpage{148}blem in so far as it belongs to analysis situs, and then the function-theoretic theorem \emph{that every simply connected surface can be conformally mapped onto a domain of the sphere} (when one applies this theorem to $\widetilde{\mathfrak{F}}$) yields the uniformizing variable. By a slight modification of the chain of thought one can show that every \emph{schlichtartig} surface can also be mapped conformally onto a domain of the sphere${}^{1)}$; for non-schlichtartig surfaces, however, this is ruled out on topological grounds: not even a bijective and domain-continuous mapping onto a domain of the sphere can then take place. Every uniformizing variable belonging to $\mathfrak{F}$ (unramified relative to $\mathfrak{F}$) will map a certain unramified unbounded schlichtartig covering surface $\bar{\mathfrak{F}}$ over $\mathfrak{F}$ conformally onto a planar domain. Two uniformizing variables which map the same covering surface $\bar{\mathfrak{F}}$ onto a sub-domain of the plane are counted in the same \textbf{class} $\{\bar{\mathfrak{F}}\}$. The setting up of all uniformizing variables then requires the solution of two problems:
1.\ of the topological problem: to find for a given surface all unramified unbounded schlichtartig covering surfaces $\bar{\mathfrak{F}}$;

2.\ of the problem of conformal mapping: to map a schlichtartig surface $\bar{\mathfrak{F}}$ in every possible way conformally onto a planar domain.

That the latter is always possible in at least one (and thus also in infinitely many) ways, hence that every class $\{\bar{\mathfrak{F}}\}$ of uniformizing variables possible according to the topological condition of schlichtartigkeit contains uniformizing variables really present function-theoretically, is the content of the \emph{general Koebe uniformization principle}.${}^{2)}$ This, to be sure, carries even further in so far as, besides those unramified relative to $\mathfrak{F}$, it also encompasses the whole multitude of ramified uniformizing variables. Beyond question, however, among all possible uniformizing variables the \emph{greatest} fundamental significance belongs to the one constructed above, designated by $t$.

\subsection*{?~20. Riemann surfaces and non-Euclidean motion groups. Fundamental domains. Poincar? $\Theta$-series.}

To what extent is the uniformizing variable $t$ determined by its properties mentioned at the conclusion of the preceding section, that is, in how manifold a way can one map the surface $\widetilde{\mathfrak{F}}$ conformally onto the

\textbf{1)} Koebe, ?ber die Hilbertsche Uniformisierungsmethode, G?ttinger Nachrichten 1910, pp.~67---74.

\textbf{2)} See especially P.~Koebe, ?ber die Uniformisierung beliebiger analytischer Kurven, Erster Teil: Das allgemeine Uniformisierungsprinzip, Crelles Journal vol.~138 (1910), pp.~192---253.

\origpage{149}sphere, the plane, or the interior of the unit disc? This question obviously amounts to: in how manifold a way can one map the sphere, plane, or disc interior conformally onto one of these three regions itself? First, it is clear: the sphere, already by virtue of its closedness, can be mapped neither onto the plane nor onto the disc interior. But a conformal mapping of the plane onto the disc interior is also impossible; a function $t^*(t)$ which transformed the $t$-plane conformally into the interior of the unit disc of the $t^*$-plane would namely be an entire transcendental function whose absolute value remained below the bound $1$ for all values of the argument, and such a function does not exist by Liouville's theorem (apart from a constant, which is completely meaningless here). Furthermore the following simple theorems now hold:

(Case~1.) \emph{All conformal mappings of the sphere (represented in the ordinary way by a complex variable)\ednote{Misprint in the original: ``Kuge'' instead of ``Kugel''.} onto itself are given by linear transformations.}

(Case~2.) \emph{The complex plane can be mapped conformally onto itself only by an entire linear transformation.}

(Case~3.) \emph{The interior of the unit disc can likewise be mapped conformally into itself only by linear transformations.}

Case~1. We only have to show: If the $t$-sphere is mapped conformally onto the $t^*$-sphere, $t^* = t^*(t)$, such that $t = 0$ passes into $t^* = 0$ and $t = \infty$ into $t^* = \infty$, then $t^* = ct$ must hold ($c$ constant). Indeed $\frac{1}{t^*}$ is then regular at the north pole of the $t$-sphere ($t = \infty$) and has a zero there, so $\frac{t}{t^*}$ is also still regular there. $t^*$ becomes $0$ only for $t = 0$, and indeed of the 1st order, since the relation $t \to t^*$ is bijective; accordingly $\frac{t}{t^*}$ is regular on the whole $t$-sphere, and thus a constant.

Case~2. Again we can assume that $t = 0$ is mapped into $t^* = 0$. One then has first to prove that
\[ \lim \frac{1}{t^*} = 0 \quad \text{for} \quad \lim \frac{1}{t} = 0. \tag{51} \]

\rmfigure{figures/fig-26.png}{Fig.~26.}{Circles in the $t$- and $t^*$-plane}

To the circle $|t^*| = R^*$ corresponds in the $t$-plane a closed curve $\mathfrak{C}$; let $R_0$ be the greatest distance of a point on $\mathfrak{C}$ from the origin, $R$ an arbitrary number $> R_0$. In the disc $|t| \leqq R$, $|t^*|$ assumes its maximum, which must be $> R^*$, on the boundary, say for $t = t_0$. Since the image curve $\mathfrak{K}^*$ of the circle $|t| = R$ in the $t^*$-plane cannot meet the circle $|t^*| = R^*$

\origpage{150}--- for the circle $|t| = R$ and the curve $\mathfrak{C}$, of which those two curves are the images, do not meet in the $t$-plane --- and since the point $t^*(t_0)$ on $\mathfrak{K}^*$ possesses an absolute value $> R^*$, we must have $|t^*| > R^*$ for all points on $\mathfrak{K}^*$; that is, as soon as $|t| > R_0$, we have $|t^*| > R^*$; that is assertion (51). $\frac{1}{t^*(t)}$ is accordingly regular at the north pole of the $t$-sphere and has a zero there, and one can continue the deduction exactly as in Case~1.

Case~3 we settle with the aid of the so-called Schwarz lemma.${}^{1)}$ Again we can assume that $t = 0$ passes into $t^* = 0$. If I consider the regular function $\frac{t^*}{t}$ in the disc $|t| \leqq q\,(< 1)$, its absolute value must reach its maximum on the boundary and is therefore $< \frac{1}{q}\,(|t^*| < 1, |t| = q)$. Since I can choose $q$ as close to $1$ as I please, $\left| \frac{t^*}{t} \right| \leqq 1$ must hold for all $t$ in the interior of the unit disc. Likewise there results $\left| \frac{t}{t^*} \right| \leqq 1$, consequently $\left| \frac{t^*}{t} \right| = 1$. This relation is possible only by $\frac{t^*}{t}$ being a constant of absolute value $1$.

Therewith the assertions made are demonstrated: \emph{the local uniformizing variable $t$ is in any case always uniquely determined up to a linear transformation.} The linear transformations which map the interior of the unit disc into itself have the form:
\[ t^* = \frac{(a + ib)t + (c - id)}{(c + id)t + (a - ib)} \tag{52} \]
\[ \{ a, b, c, d \text{ real; determinant } (a^2 + b^2) - (c^2 + d^2) = 1 \}. \]

In particular it results, since the deck transformations of $\widetilde{\mathfrak{F}}$ are bijective conformal mappings of $\widetilde{\mathfrak{F}}$ onto itself, that these deck transformations must be mapped in the $t$-plane as linear transformations. To the group of deck transformations there thus corresponds a certain group $\varGamma$ of linear transformations isomorphic to it. No transformation belonging to the group $\varGamma$ (except the identity) may possess a \textbf{fixed point} in the image domain (sphere, plane, or disc interior), that is, a point which passes into itself under the transformation. Since on the sphere every linear transformation has a fixed point, there can in Case~1 be no other deck transformation than the identity; $\widetilde{\mathfrak{F}}$ is then identical with $\mathfrak{F}$ itself and $\mathfrak{F}$ as a Riemann surface with the

\textbf{1)} Schwarz, Gesammelte Abhandlungen vol.~II, pp.~109---111. The application of the lemma to the present problem following Poincar?, Acta Mathematica vol.~4, pp.~231---232.

\origpage{151}sphere. \emph{Case~1 can occur only when the given Riemann surface $\mathfrak{F}$ is equivalent to the sphere.}

The points on the surface $\widetilde{\mathfrak{F}}$ lying over one and the same point of $\mathfrak{F}$ appear in the mapping as a system of points $t$ equivalent with respect to $\varGamma$. Such a system $\Sigma$ has the property that any two of its points pass into one another by a transformation belonging to $\varGamma$, but also that every point into which a point of $\Sigma$ is transformed by a transformation belonging to $\varGamma$ belongs for its part to $\Sigma$ (cf.\ p.~27). The single-valued functions on $\mathfrak{F}$ regular except for poles appear, when expressed in terms of $t$, as \textbf{automorphic functions} of that variable belonging to $\varGamma$, that is, as functions $z(t)$ behaving invariantly under the transformations of the group $\varGamma$:
\[ z \left( \frac{\alpha t + \beta}{\gamma t + \delta} \right) = z(t), \]
if $t^* = \frac{\alpha t + \beta}{\gamma t + \delta}$ is a transformation of the group $\varGamma$. The group $\varGamma$ must be \textbf{discontinuous}; that shall mean: a system of points equivalent with respect to $\varGamma$ may never exhibit a condensation point within the image domain. If we define a Riemann surface $\mathfrak{F}_\varGamma$ by regarding each system of points belonging to the image domain equivalent with respect to $\varGamma$ as a ?point? of $\mathfrak{F}_\varGamma$ and transferring the angle measurement from the $t$-plane directly to $\mathfrak{F}_\varGamma$, this $\mathfrak{F}_\varGamma$ as a Riemann surface is equivalent to the given $\mathfrak{F}$ and \emph{represents the most suitable normal form to which every Riemann surface can be brought.}

In Case~2, $\varGamma$ must consist of entire linear transformations possessing no fixed point in the finite plane. Only translations $t' = t + \alpha$ satisfy this requirement. A discontinuous group of translations can consist${}^{1)}$

1.\ only of the identity; then $\widetilde{\mathfrak{F}}$ is identical with $\mathfrak{F}$, and $\mathfrak{F}$ is equivalent to the plane [the ?\emph{punctured}? sphere, i.e.\ the sphere without the north pole];

or

2.\ of the repetitions of a single translation
\[ \varGamma\colon \quad t^* = t + n\alpha \quad [n = 0, \pm 1, \pm 2, \dots]; \]
then $\mathfrak{F}$ is obviously equivalent to the infinitely long straight \emph{circular cylinder}, and there-

\textbf{1)} That only the three enumerated cases are possible follows from the procedure of adapting a (two-dimensional) number lattice with respect to an included lattice; see pp.~73\,f.\ of this work or G.~Kowalewski, Die komplexen Ver?nderlichen und ihre Funktionen, Leipzig 1911, pp.~57\,f.

\origpage{152}with by the Mercator projection to the \emph{doubly punctured sphere} (sphere without north and south poles) as a Riemann surface;

or can

3.\ be of the form:
\[ t^* = t + m\alpha + n\beta \quad \begin{pmatrix} m = 0, \pm 1, \pm 2, \dots \\ n = 0, \pm 1, \pm 2, \dots \end{pmatrix}, \]
where $\alpha$, $\beta$ are two translations in different directions. Then $\mathfrak{F}$ is closed, $dt$ a single-valued everywhere regular differential on $\mathfrak{F}$ possessing nowhere a zero, and $\mathfrak{F}$ thus a \emph{closed Riemann surface of genus 1.} Every such surface --- whose type (but not whose most general form) is the torus --- in fact also permits a uniformization by the integral of the 1st kind, in which the universal covering surface is mapped onto the plane.

Apart from the few exceptional cases enumerated above ($\mathfrak{F} = {}$full sphere, singly or doubly punctured sphere, closed surface of genus $1$), Case~3 always obtains, in which \emph{the interior of the unit disc} appears as the image domain. Since the periphery of the unit disc will in general be a natural boundary for the automorphic functions of $t$ representing the functions on the base surface $\mathfrak{F}$, one designates $t$ as the \textbf{limit-circle uniformizing variable}.

The group $\varGamma$ is not completely uniquely determined by the given surface. For $t$ can be replaced by any variable $t'$ arising from $t$ by means of a linear transformation $T_0$ mapping the interior of the unit disc into itself; thereby $\varGamma$ is transformed into the group
\[ \varGamma' = T_0^{-1} \cdot \varGamma \cdot T_0. \]
We introduce the following terminology:

As a ?point of $\mathfrak{E}$? we designate every point $t$ situated in the interior of the unit disc, as a ?straight line on $\mathfrak{E}$? that part situated in the interior of the unit disc of every circle which stands perpendicular to the periphery of the unit disc (the diameters of the unit disc also count among these ?straight lines?). As a ?motion of $\mathfrak{E}$? counts every linear transformation mapping the interior of the unit disc into itself, and as ?congruent? such point sets on $\mathfrak{E}$ which can be transformed into one another by a ?motion?. The ordinary angle measure is retained. Then for these ?points? and ?straight lines? the \emph{complete Bolyai-Lobachevskian geometry} holds, hence that geometry whose axioms differ from those of Euclidean geometry only by the omission of the parallel axiom${}^{1)}$, so that we can designate $\mathfrak{E}$ as the \emph{non-Euclidean plane}.

\textbf{1)} The model of plane non-Euclidean geometry resulting here is connected most closely with that discovered by Klein in the year 1871 based on Cayley's

\origpage{153}projective metrics (?ber die sogenannte Nicht-Euklidische Geometrie, Math.\ Ann.\ vol.~4, pp.~573---625); cf.\ concerning this Fricke and Klein, Vorlesungen ?ber die Theorie der automorphen Funktionen, vol.~I, Leipzig 1897, pp.~3---59. A good orientation regarding non-Euclidean geometry is given by the book of R.~Bonola, which appeared in German translated by Liebmann as vol.~4 of the collection ``Wissenschaft und Hypothese'' (Die Nicht-Euklidische Geometrie; Leipzig 1908).

\textbf{1)} Cf.\ the concluding remarks of Koebe's first communication on the uniformization of arbitrary analytic curves, G?ttinger Nachrichten 1907, pp.~209---210.

\textbf{2)} N.-E. $=$ non-Euclidean.

The complex variable $t$ (bound to the inequality $|t| < 1$) is then to be conceived as a coordinate suitable for representing the points of the Lobachevskian plane in the same way as, say, rectangular Cartesian coordinates $x$, $y$ or their complex combination $z = x + iy$ represent the points of the Euclidean plane. Linear transformation of $t$ mapping the interior of the unit disc into itself yields another coordinate system of the Lobachevskian plane on an equal footing with the original one. In the present interpretation $\varGamma$ appears as a group of motions of the non-Euclidean plane, or, more precisely, as a \emph{representation} of such a group with the aid of a specific coordinate $t$. That group of linear substitutions $\varGamma'$ which we obtain from $\varGamma$ when we replace $t$ by $t'$ by means of a transformation mapping the interior of the unit disc into itself is in this view to be regarded merely as another representation of the same motion group of the non-Euclidean plane (namely with the aid of another coordinate $t'$). $\varGamma$ contains no rotations, that is, no motion of $\mathfrak{E}$ which leaves a point of $\mathfrak{E}$ fixed.

\emph{To every Riemann surface there accordingly corresponds (apart from the four exceptional cases enumerated previously) a single, completely determined, rotation-free discontinuous motion group $\varGamma$ of the Lobachevskian plane. Two Riemann surfaces are conformally equivalent (belong, as Riemann expresses it, to the same class or are realizations of one and the same ideal Riemann surface) if and only if the associated non-Euclidean motion groups are congruent in the sense of Lobachevskian geometry. Conversely, to every discontinuous non-Euclidean motion group without rotations there also belongs a specific class of Riemann surfaces (as whose representative the surface $\mathfrak{F}_\varGamma$ constructed on p.~151 can serve).}${}^{1)}$

To represent a discontinuous motion group $\varGamma$ of the N.-E.${}^{2)}$ plane intuitively, one uses, following Klein and Poincar?, a simple and gapless covering of the plane by N.-E.\ congruent re-

\origpage{154}gions (\textbf{fundamental domains}) which arise from one another by the motions contained in $\varGamma$. A partition of this kind as simple as possible is suggested by the following considerations.

\rmfigure{figures/fig-27.png}{Fig.~27.}{Non-Euclidean segment}

(\emph{Definition of N.-E.\ distance.}) If $t_1$, $t_2$ are any two distinct points in the interior of the unit disc, we connect them by a N.-E.\ straight line, that is, in the Gaussian $t$-plane by a circle which intersects the unit circle perpendicularly at the places $t_{\infty_1}$ and $t_{\infty_2}$. This circle is uniquely determined by $t_1$, $t_2$, and the points shall lie on it one after another in the order $t_{\infty_1}, t_1, t_2, t_{\infty_2}$ (Fig.~27). The cross-ratio
\[ d(t_1\,t_2) = \frac{t_1 - t_{\infty_2}}{t_1 - t_{\infty_1}} : \frac{t_2 - t_{\infty_2}}{t_2 - t_{\infty_1}} \]
has a real positive value $> 1$, and its real logarithm
\[ r(t_1\,t_2) = \lg d(t_1\,t_2) \]
is therefore positive. $r(t_1\,t_2)$ is invariant under all linear transformations mapping the interior of the unit disc into itself, and the N.-E.\ segment $t_1\,t_2$ is N.-E.\ congruent to another such segment $t_1^*\,t_2^*$ if and only if $r(t_1\,t_2) = r(t_1^*\,t_2^*)$. Furthermore, if three points $t_1$, $t_2$, $t_3$ lie on a N.-E.\ line in the order of their numbering, then
\[ r(t_1\,t_2) + r(t_2\,t_3) = r(t_1\,t_3). \]
By virtue of these two circumstances we have to designate the number $r(t_1\,t_2)$ as the \textbf{non-Euclidean distance} of the two points $t_1$, $t_2$, and are thereby enabled to measure in the N.-E.\ plane not only angles, but also segments.

If $\Sigma_0$ is a system of points equivalent with respect to $\varGamma$, $t_0$ an individual point of $\Sigma_0$, and $t$ an arbitrary point of the N.-E.\ plane, then because of the discontinuity of $\varGamma$ there are only finitely many points in $\Sigma_0$ whose distance${}^{1)}$ from $t$ does not exceed an arbitrarily prescribed magnitude $R$. There are thus among the points of $\Sigma_0$ one or more (but certainly only finitely many) $t_h$ such that the distance $r(t\,t_h)$ is smallest among all distances from $t$ to points of the system $\Sigma_0$; $t$ then lies, as I shall briefly express it, nearest to $t_h$. If the distance $r(t\,t_h)$ is strictly smaller than the distance to all points of the system $\Sigma_0$ distinct from $t_h$, I can delimit an entire neighborhood about $t$ such that all points of this neighborhood also still lie nearest to $t_h$. If $t_h$ is any of the points of $\Sigma_0$, I u-

\textbf{1)} The word ``distance'' and all other geometric expressions are until further notice to be understood in the N.-E.\ sense.

\origpage{155}nite all those $t$ which lie nearest to $t_h$ into a point set $\mathfrak{P}_h$, as whose ?center? the point $t_h$ counts. The $\mathfrak{P}_h$ thus coming into being about all the centers belonging to $\Sigma_0$ are throughout distinct in their interior points and together cover the entire N.-E.\ plane. $\mathfrak{P}_h$ arises from $\mathfrak{P}_0$ by that motion contained in the group $\varGamma$ which transforms $t_0$ into $t_h$, and is thus N.-E.\ congruent to $\mathfrak{P}_0$. To each point $t$ there is found in $\mathfrak{P}_0$ one equivalent with respect to $\varGamma$, and two \emph{interior} points of $\mathfrak{P}_0$ are never equivalent to each other: $\mathfrak{P}_0$ has the properties of a fundamental domain. Moreover $\mathfrak{P}_0$ is convex, that is, if $t'$, $t''$ are any two points situated in $\mathfrak{P}_0$, then all points of the N.-E.\ straight connecting segment $t'\,t''$ also belong to $\mathfrak{P}_0$. The perpendicular bisector $\mathfrak{g}_h$ of the segment $t_0\,t_h$ is the locus of all points having equal distance from $t_0$ and $t_h$; for all centers $t_h$ distinct from $t_0$ we obtain such a line $\mathfrak{g}_h$. The boundary of the closed set $\mathfrak{P}_0$ is composed of segments of these lines, and the interior of $\mathfrak{P}_0$ consists of all those points lying with respect to all the lines $\mathfrak{g}_h$ on the same side as the point $t_0$. Since all points of the line $\mathfrak{g}_h$ have a distance $\geqq \frac{1}{2} r(t_0\,t_h)$ from $t_0$, and among the centers $t_h$ only finitely many are found whose distance from $t_0$ does not exceed an arbitrarily prescribed bound $2R$, it turns out that only finitely many of the straight segments belonging to the boundary of $\mathfrak{P}_0$ have a distance $\leqq R$ from the main center $t_0$. $\mathfrak{P}_0$ is accordingly a convex polygon with finitely or infinitely many sides, whose vertices (that is, those points having equal distance from three or more points of the system $\Sigma_0$) nowhere accumulate in the finite region. Following Fricke, $\mathfrak{P}_0$ is called a \textbf{normal polygon} belonging to the group $\varGamma$. The sides of the normal polygon are related to one another in pairs by the fact that any two sides of a pair can be transformed into one another by a motion belonging to $\varGamma$. It never happens here that two sides with a common vertex are related to one another; for this vertex would have to be a fixed point of the motion transforming one side into the other. By repetition and composition one can obtain from the motions transforming the paired sides of $\mathfrak{P}_0$ into one another [these are those motions which transform $\mathfrak{P}_0$ into equivalent fundamental domains $\mathfrak{P}_h$ adjacent to $\mathfrak{P}_0$ along a side] all the motions of the group $\varGamma$. The proof for this is provided by connecting an arbitrary point $t_h$ belonging to $\Sigma_0$ with $t_0$ by a polygonal path which passes through no vertices of the polygonal partition $\{\mathfrak{P}_h\}$, and then paying attention to \emph{which} polygons of the partition this path successively passes through. If at the vertex $1$ of $\mathfrak{P}_0$ in all $e$ polygons $\mathfrak{P}_h$ equivalent to $\mathfrak{P}_0$ meet ($\mathfrak{P}_0$ itself is to be included), there are exactly $e$ centers to which $1$ lies near-

\origpage{156}est, and among the points of $\mathfrak{P}_0$ there are $e$ which are equivalent to $1$ with respect to $\varGamma$, or, as one says following Poincar?, form a \textbf{cycle} of vertices of the polygon $\mathfrak{P}_0$. The sum of the angles of $\mathfrak{P}_0$ situated at the vertices of a single such cycle is $= 2\pi$.

The Riemann surface $\mathfrak{F}_\varGamma$ belonging to the group $\varGamma$ (p.~151) is obviously closed if and only if a positive number $R$ exists such that for every point of the N.-E.\ plane at least one equivalent with respect to $\varGamma$ is present which has a N.-E.\ distance $\leqq R$ from the center $t_0$. Then every point has from that nearest to it in the system $\Sigma_0$ a distance $\leqq R$; in particular $\mathfrak{P}_0$ lies entirely in the N.-E.\ circle of radius $R$ about $t_0$. Since the vertices of $\mathfrak{P}_0$ nowhere accumulate in the finite region, under these circumstances only finitely many such vertices can exist at all: $\mathfrak{P}_0$ \emph{has finitely many sides.} Let the number of its sides, which must be even because of the pairwise assignment, be set $= 2s$, the number of distinct vertex cycles $= c$, so that $2\pi c$ denotes the angle sum of $\mathfrak{P}_0$. One can decompose $\mathfrak{P}_0$ into finitely many N.-E.\ rectilinear triangles in such a way that a \emph{triangulation} of $\mathfrak{F}_\varGamma$ thereby comes about. If $D_0$ is the number of triangles, $K_0$ the number of edges appearing within $\mathfrak{P}_0$ in this decomposition, $E_0$ the number of vertices appearing within $\mathfrak{P}_0$, $2e_0$ the number of vertices newly added on the sides of $\mathfrak{P}_0$ to the polygon vertices (which will be pairwise equivalent and divide the periphery of $\mathfrak{P}_0$ into $2(s + e_0)$ individual segments), then, since a convex polygon is simply connected,
\[ D_0 + E_0 - K_0 = 1. \]
If, however, we regard the decomposition of $\mathfrak{P}_0$ as a triangulation of $\mathfrak{F}_\varGamma$, let the numbers of triangles, edges, and vertices be respectively $= D, K, E$. We have
\[ D = D_0, \quad K = K_0 + (s + e_0), \quad E = E_0 + e_0 + c. \]
The genus $p$ of $\mathfrak{F}_\varGamma$ is determined from the equation
\[ 2p - 2 = K - E - D = s - c - 1. \]

\emph{It shows how one can find the genus of a Riemann surface when this is given in its normal form (that is, when the associated non-Euclidean motion group $\varGamma$ is given), namely by constructing a normal polygon belonging to $\varGamma$.} One can give the equation a still more elegant form if one considers that the angle sum of a triangle in non-Euclidean geometry is smaller than $\pi$, and indeed by as much as the area of the triangle amounts to. The N.-E.\ area $J$ of $\mathfrak{P}_0$, if one decomposes $\mathfrak{P}_0$ from $t_0$ into $2s$ triangles, results accordingly as $2\pi(s - c - 1)$:
\[ J = 4\pi\,(p - 1). \]

\origpage{157}Since each cycle consists of at least\ednote{Misprint in the original: ``wenigsten'' instead of ``wenigstens''.} three vertices,
\[ 3c \leqq 2s, \quad c \leqq \frac{2s}{3}, \quad s - c \geqq \frac{s}{3}, \]
hence
\[ \begin{gathered} \frac{s}{3} \leqq 2p - 1; \\ 2s \leqq 12p - 6. \end{gathered} \]
\emph{This upper bound depending only on the genus $p$ holds for the number of sides and vertices of a normal polygon.}

The division of the N.-E.\ plane into normal polygons serves not only to visualize N.-E.\ motion groups, but also provides the means to construct such groups directly.${}^{1)}$

In the motion groups $\varGamma$ of the non-Euclidean plane (or in the associated manifolds $\mathfrak{F}_\varGamma$), the purest embodiment of the idea of the Riemann surface, freed from all accidental features, confronts us. As the crowning achievement of the whole structure of this part of Riemann's function theory, the solution of the following problem would be required: \emph{When a Riemann surface is given in its normal form (that is, the associated motion group of the non-Euclidean plane), to express all single- or multiple-valued unramified functions on the surface by means of closed analytic formulas in the complex coordinate $t$ representing the points of the Lobachevskian plane.} This problem has not been solved completely hitherto; an important start towards it are the $\Theta$-series introduced by Poincar?.${}^{2)}$ If the prescribed group $\varGamma$ mapping the interior of the unit disc $|t| < 1$ into itself consists of the substitutions:
\[ \begin{gathered} t_i^* = tS_i = \frac{\alpha_i t + \beta_i}{\gamma_i t + \delta_i} \quad (\alpha_i\delta_i - \beta_i\gamma_i = 1) \\ [i = 0, 1, 2, 3, \dots; S_0 \text{ the identity}], \end{gathered} \]
then, for example,

\textbf{1)} Besides the normal polygons, the ``canonical polygons'' are of great importance, whose theory was developed by Fricke (see Fricke-Klein, Vorlesungen ?ber die Theorie der automorphen Funktionen, Leipzig 1897 to 1912) and with whose aid this investigator succeeded in formulating and proving rigorously the assertion already stated by Riemann that the Riemann surfaces of genus $p\,(> 1)$ form a $(6p - 6)$-dimensional manifold. These are things connected most closely with the ``continuity method'' recently coming again to the foreground of consideration, by which Klein and Poincar? tried to prove the uniformizability of algebraic configurations right at the beginning of the eighties. Cf.\ the report on the Karlsruhe proceedings cited on p.~142, Deutsche Mathematiker-Vereinigung 1912.

\textbf{2)} Poincar?, M?moire sur les fonctions fuchsiennes, Acta Mathematica, vol.~1 (1882), p.~207.

\origpage{158}\[ \Theta(t) = \sum_{i=0}^\infty \frac{1}{(\alpha_i t + \beta_i)(\gamma_i t + \delta_i)^3}, \quad \Theta'(t) = \sum_{i=0}^\infty \frac{1}{(\alpha_i t + \beta_i)^3(\gamma_i t + \delta_i)} \]
are Poincar? $\Theta$-series. To prove the absolute and uniform convergence of these series, one describes about an arbitrary point $t_0\,(|t_0| < 1)$ an ordinary circle $\varkappa$ of radius $a\colon |t - t_0| \leqq a$, so small that the equivalent circles $\varkappa S_i$ do not mutually meet. The surface integral $J_i$ of
\[ \left| \frac{dt_i^*}{dt} \right|^2 = \frac{1}{|\gamma_i t + \delta_i|^4} \]
extended over $\varkappa$ is the (Euclidean) area of $\varkappa S_i$; consequently $\displaystyle\sum_{i=0}^\infty J_i$ is convergent. If $t$ has from the point $t_0$ a Euclidean distance $|t - t_0| \leqq \frac{a}{2}$, then [formula (22), (23) on p.~87]
\[ \left( \frac{a}{2} \right)^2 \pi \left| \frac{dt_i^*}{dt} \right|^2 \leqq J_i. \]
Consequently
\[ \sum \frac{1}{|\gamma_i t + \delta_i|^4} \]
is uniformly convergent in the neighborhood of $t = t_0$. Since furthermore uniformly for $|t| \leqq q\,(< 1)$
\[ \lim_{i = \infty} \left| \frac{\alpha_i t + \beta_i}{\gamma_i t + \delta_i} \right| = 1 \]
holds (for the equivalent points $t_i^*$ accumulate only at the boundary of the unit disc), there results from this the absolute and uniform convergence of the two series $\Theta$, $\Theta'$. $\Theta$ has at $t = 0$ and all places equivalent to it with respect to the group $\varGamma$ a pole of the 1st order and is otherwise regular, $\Theta'$ has there a pole of the 3rd order and is otherwise regular. From the equation
\[ \Theta(t)\,(dt)^2 = \sum_{i=0}^\infty \frac{(dt_i^*)^2}{t_i^*} \]
it emerges that $\Theta(t)\,(dt)^2$ passes into itself under an arbitrary transformation $S_i$ of the group:
\[ \Theta(tS_i) = (\gamma_i t + \delta_i)^4 \cdot \Theta(t). \]
Likewise follows
\[ \Theta'(tS_i) = (\gamma_i t + \delta_i)^4 \cdot \Theta'(t). \]
$\frac{\Theta(t)}{\Theta'(t)}$ is accordingly a function automorphic with respect to the group $\varGamma$; it does not vanish identically, but possesses at the point $t = 0$ (and all those equivalent to it) a zero of the 2nd order; in short: is a

\origpage{159}single-valued, non-constant function on the base surface $\mathfrak{F}_\varGamma$ free from essential singularities.

One sees from this example how it must be possible, once the Riemann surface is given in its normal form, to derive all theorems concerning the existence of functions and integrals on it with the aid of analytic formulas in a manner similar to what has long been done for $p = 1$ by the theory of elliptic functions. Only when this goal has been attained in general will a completely unified construction of Riemann's function theory be possible. One will then by transcendental means (with the aid of the method of the alternating procedure or of Dirichlet's principle) no longer prove the existence of functions or integrals, but solely the existence of the limit-circle uniformizing variable, that is, first produce by transcendental means the normal form of the Riemann surface, in case this is not yet given in its normal form; but then descend with the aid of analytic formulas to all other functions unramified relative to the given surface. Here there still lie rewarding problems for the future.

\subsection*{?~21. Conformal mapping of a Riemann surface onto itself.}

Of a motion group of the N.-E.\ plane we say that it contains \textbf{infinitesimal transformations} if there are motions in the group that differ arbitrarily little from the identity. A discontinuous motion group, e.g.\ $\varGamma$, certainly contains no infinitesimal transformations. Of this theorem, however, the converse also holds:

\emph{A N.-E.\ motion group is discontinuous if and only if it contains no infinitesimal motions.}

In the proof we make use of the complex coordinate $t$ as hitherto. A linear transformation $T\colon t \to t^*$ has on the $t$-sphere in general two distinct fixed points $\tau'$, $\tau''$ and can then be written as follows:
\[ \frac{t^* - \tau''}{t^* - \tau'} = \mu \cdot \frac{t - \tau''}{t - \tau'}{}^{1)}; \tag{53} \]
the constant $\mu$ is called the \textbf{multiplier} of $T$. If $T$ maps the interior of the unit disc into itself, two possibilities are present:

\emph{either} $\tau'$, $\tau''$ both lie on the boundary of the unit disc, and $\mu$ is then positive (\textbf{hyperbolic transformation}),

\emph{or} $\tau'$, $\tau''$ are mirror images of each other with respect to the unit circle, and $\tau'$ lies in the interior of the same; then $|\mu| = 1$ (\textbf{elliptic}

\textbf{1)} If $\tau'' = \infty$, this means: $\displaystyle\frac{1}{t^* - \tau'} = \mu \cdot \frac{1}{t - \tau'}$.

\origpage{160}\textbf{transformation}). In the N.-E.\ plane $T$ is a rotation about $\tau'$, and the real number $\varphi$ determined by the equation $\mu = e^{i\varphi}$ up to integral multiples of $2\pi$ is the rotation angle of $T$.

As a transitional case between the hyperbolic and elliptic transformation there interpolates the \textbf{parabolic} transformation, which possesses only \emph{one} fixed point $\tau'$. If it maps the interior of the unit disc into itself, it has the form
\[ \frac{1}{t^* - \tau'} = \frac{1}{t - \tau'} + \frac{ib}{\tau'} \quad (|\tau'| = 1, b \text{ real}). \tag{54} \]

We base our proof on the following \emph{lemma}: If
\[ t_n, t_n^* \quad (n = 1, 2, 3, \dots) \]
are two sequences of points in the N.-E.\ plane which converge to the same point $t_0\,(|t_0| < 1)$; if furthermore $T_n$ is a N.-E.\ motion transforming $t_n$ into $t_n^*$, and there exists a neighborhood of $t_0$ containing no fixed point of any of the motions $T_n\,(n = 1, 2, 3, \dots)$, then $T_n$ converges to the identity.

If $\tau_n'\,(|\tau_n'| \leqq 1)$, $\tau_n''$ are the two fixed points of $T_n$ (which may also coincide), there shall thus exist a positive number $l$ such that for all $n$
\[ |t_0 - \tau_n'|, \quad |t_0 - \tau_n''| > \frac{3l}{2} \]
holds. If I set
\[ |t_n - t_n^*| = \varepsilon_n, \]
I can assume that for all $n$
\[ |t_n - t_0|, \quad |t_n^* - t_0| < \frac{l}{2}, \quad \text{hence} \quad \varepsilon_n < l \]
holds. The differences
\[ |t_n - \tau_n'|, \quad |t_n - \tau_n''|, \quad |t_n^* - \tau_n'|, \quad |t_n^* - \tau_n''| \]
are then all four $> l$. If $T_n$ is not parabolic, it has the form
\[ \frac{t^* - \tau_n''}{t^* - \tau_n'} = \mu_n \cdot \frac{t - \tau_n''}{t - \tau_n'}. \tag{55} \]
If I set herein $t = t_n$, hence $t^* = t_n^*$, there comes
\[ \mu_n = \frac{t_n^* - \tau_n''}{t_n - \tau_n''} : \frac{t_n^* - \tau_n'}{t_n - \tau_n'} = \left( 1 + \frac{t_n^* - t_n}{t_n - \tau_n''} \right) : \left( 1 + \frac{t_n^* - t_n}{t_n - \tau_n'} \right). \]
This yields the estimate
\[ |\mu_n - 1| \leqq \frac{2\varepsilon_n}{l - \varepsilon_n}. \tag{56} \]
In place of (55) I can write
\[ \frac{1}{t^* - \tau_n'} = \lambda_n + \mu_n \frac{1}{t - \tau_n'} \quad (\lambda_n \text{ constant}); \tag{57} \]
for this form must be had by every linear substitution with the fixed point

\origpage{161}$\tau_n'$. Even when the transformation $T_n$ is parabolic, I can put it into the form (57); specifically we will then have $\mu_n = 1$, so (56) likewise applies. We gain the limit equation
\[ \lim_{n = \infty} \mu_n = 1. \]
Setting in (57) again $t = t_n$, $t^* = t_n^*$, we find
\[ \lambda_n = \frac{t_n - t_n^*}{(t_n - \tau_n')(t_n^* - \tau_n')} - \frac{\mu_n - 1}{t_n - \tau_n'}, \]
\[ |\lambda_n| \leqq \frac{\varepsilon_n}{l^2} + \frac{|\mu_n - 1|}{l}, \quad \lim_{n = \infty} \lambda_n = 0. \]
Bringing (57) finally into its natural form
\[ t^* = \frac{\alpha_n t + \beta_n}{\gamma_n t + \delta_n}, \]
there result the values
\[ \begin{aligned} \alpha_n &= 1 + \lambda_n \tau_n', &\quad \beta_n &= \tau_n'(\mu_n - 1 - \lambda_n \tau_n'), \\ \gamma_n &= \lambda_n, &\quad \delta_n &= \mu_n - \lambda_n \tau_n', \end{aligned} \]
which differ from the corresponding coefficients of the identical substitution $\begin{pmatrix} 1 & 0 \\ 0 & 1 \end{pmatrix}$ by less than
\[ |\mu_n - 1| + |\lambda_n|. \]
The lemma is therewith proved.

Let now $\varGamma^*$ be a N.-E.\ motion group \emph{without infinitesimal transformations} and $\tau$ a point of the N.-E.\ plane which is a fixed point of a rotation belonging to $\varGamma^*$. By $S_\tau$ I denote all rotations contained in $\varGamma^*$ with fixed point $\tau$, which form a group in themselves. If I normalize the rotation angles $\varphi$ of the $S_\tau$ by the condition $0 \leqq \varphi < 2\pi$, then among the $S_\tau$, on account of the absence of infinitesimal transformations, there is one, $S_\tau^0$, with \emph{smallest} rotation angle $\varphi_0 > 0$. The rotation angle of every $S_\tau$ is an integral multiple of $\varphi_0$, since otherwise by suitable choice of an integer $k$ one could generate a transformation $S_\tau \cdot (S_\tau^0)^{-k}$ with smaller rotation angle than $S_\tau^0$; $S_\tau^0$ is thus in the group of the $S_\tau$ a ``primitive'' transformation from which all others are obtained by exponentiation. If one determines an integer $h$ by the condition
\[ h\varphi_0 \leqq 2\pi < (h + 1)\varphi_0, \]
then $(S_\tau^0)^{h+1}$ certainly acquires a smaller rotation angle $(h + 1)\varphi_0 - 2\pi$ than $S_\tau^0$, unless $h\varphi_0 = 2\pi$. We must therefore have $\varphi_0 = \frac{2\pi}{h}$, and $h$ is the \emph{order} of the finite cyclic group $(S_\tau)$ consisting of the rotations
\[ (S_\tau^0)^1, (S_\tau^0)^2, (S_\tau^0)^3, \dots, (S_\tau^0)^h = 1 \]

\origpage{162}The fixed points of the rotations occurring in $\varGamma^*$ can nowhere accumulate in the finite region. For if $\tau\,(|\tau| < 1)$ were such an accumulation point and
\[ S_n \left( \text{angle of rotation} = \frac{2\pi}{h_n} \right) \qquad [n = 1, 2, 3, \dots] \]
a sequence of primitive rotations from $\varGamma^*$ whose fixed points $\tau_n$ are all distinct and converge to $\tau$, then a priori two possibilities exist:

1. The orders $h_n$ remain below a fixed bound for all $n$. Then among the $S_n$ infinitely many are found to which belongs the same order $h$; let these be the transformations $S_1', S_2', S_3', \dots$. $S_n' S_{n+1}'^{-1}$ is infinitesimal for infinitely large $n$. This case is therefore excluded.

2. The orders $h_n$ do not remain below a fixed bound. Then from the $S_n$ one can select a sequence $S_n'$ such that the associated order numbers converge to $\infty$ and the rotation angles consequently to $0$; this likewise contradicts the absence of infinitesimal operations in $\varGamma^*$.

Having established this, we can now also show that a system of points equivalent with respect to $\varGamma^*$ can possess no accumulation point in the finite region. For if $t_0\,(|t_0| < 1)$ were the limit of a sequence of equivalent points $t_n\,(n = 1, 2, 3, \dots)\colon$
\[ \lim_{n = \infty} t_n = t_0, \]
then let $S_n$ denote that motion belonging to $\varGamma^*$ which transforms $t_n$ into $t_{n+1}$. About $t_0$ I delimit a neighborhood $\mathfrak{U}_0$ containing, except possibly at $t_0$, no fixed point of a rotation belonging to $\varGamma^*$. By the lemma, only for finitely many indices $n$ can the fixed point $\tau_n'$ of $S_n\,(|\tau_n'| \leqq 1)$ lie outside $\mathfrak{U}_0$; from a certain $n$ onwards we must consequently have $\tau_n' = t_0$, and it involves no restriction to assume that this already holds from $n = 1$ onwards. All $t_n$ then arise from $t_1$ by such rotations about $t_0$ as are contained in the group $\varGamma^*$; by such rotations, however, I can bring $t_1$ into only finitely many distinct positions. The possibility of a condensation point is thus refuted, and the proof of the theorem stated at the beginning (p.~159) is accomplished.

Having premised this, we come now to the proper subject of this concluding section of our presentation. The concern is the question of those bijective conformal mappings which a Riemann surface $\mathfrak{F}$ can undergo onto itself. These mappings form a group which obviously possesses for the theory of the Riemann surface $\mathfrak{F}$ an importance analogous to that of, say, the group of motions for metric geometry, and we

\origpage{163}therefore pronounce a fact of fundamental importance when we formulate the following theorem, due essentially to Klein:

\emph{The group of conformal mappings of a Riemann surface onto itself is always discontinuous, apart from the following seven exceptional cases: $\mathfrak{F} = {}$full sphere, singly or doubly punctured sphere, spherical cap, punctured spherical cap (i.e.\ cap without its center), spherical zone (between two circles of latitude)${}^{1)}$, closed Riemann surface of genus 1.}

\emph{Proof}${}^{2)}$: Let $\widetilde{\mathfrak{F}}$ denote the universal covering surface over $\mathfrak{F}$, $t$ the limit-circle uniformizing variable mapping $\widetilde{\mathfrak{F}}$ conformally onto the interior of the unit disc, $\varGamma$ the N.-E.\ motion group belonging to the Riemann surface $\mathfrak{F}$; the cases in which instead of the unit disc the $t$-sphere or the infinite $t$-plane appears we can leave aside here, since in them our theorem is invalid in any case (exceptional cases 1, 2, 3, 7). Let $C$ be a conformal mapping of $\mathfrak{F}$ onto itself transforming the point $\mathfrak{p}_0$ into $\mathfrak{p}_0'$. Let $\tilde{\mathfrak{p}}_0$ be a point on $\widetilde{\mathfrak{F}}$ over $\mathfrak{p}_0$, $\tilde{\mathfrak{p}}_0'$ a point on $\widetilde{\mathfrak{F}}$ over $\mathfrak{p}_0'$. Because of the simple connectivity of $\widetilde{\mathfrak{F}}$ there is a unique bijective and domain-continuous mapping $\widetilde{C}$ of $\widetilde{\mathfrak{F}}$ onto itself in which
1.~$\tilde{\mathfrak{p}}_0$ passes into $\tilde{\mathfrak{p}}_0'$,

2.~each point $\tilde{\mathfrak{p}}$ on $\widetilde{\mathfrak{F}}$ passes into a point $\widetilde{\mathfrak{p}C}$ whose trace point $\mathfrak{p}C$ arose by the mapping $C$ from the trace point $\mathfrak{p}$ of $\tilde{\mathfrak{p}}$.
$\widetilde{C}$ is likewise conformal and thus appears in the $t$-plane as a linear transformation $T$ mapping the interior of the unit disc into itself, and indeed as a transformation which maps each system of points equivalent with respect to $\varGamma$ into a like point system. This latter property is expressed in the equation
\[ T \cdot \varGamma \cdot T^{-1} = \varGamma, \]
which states that $T$ is \emph{permutable} with $\varGamma$, or: that for every substitution $S$ belonging to $\varGamma$ the ``transform'' $TST^{-1}$ likewise belongs to $\varGamma$. Conversely it is also clear that every linear substitution $T$ permutable with $\varGamma$ mapping $|t| < 1$ into itself yields a conformal mapping of $\mathfrak{F}_\varGamma$ onto itself. All motions of the N.-E.\ plane permutable with $\varGamma$ evidently form a group $\varGamma_c$ containing $\varGamma$ as a part. The discontinuity of this group $\varGamma_c$ is to be proved; it suffices for this, by the

\textbf{1)} These are, according to the width of the zone, infinitely many essentially distinct Riemann surfaces.

\textbf{2)} This proof is due to Poincar?: Acta Mathematica vol.~7, 1885, pp.~16---19. Poincar? there, to be sure, considers only closed Riemann surfaces, but his proof remains valid verbatim also for open surfaces.

\origpage{164}theorem at the beginning of this section, to recognize the absence of infinitesimal operations in $\varGamma_c$.

If $S, T$ are any two linear transformations mapping the interior of the unit disc into itself, those two points into which $T$ casts the fixed points of $S$ are the fixed points of the transformation
\[ TST^{-1} = S'. \]
In order for $S'$ to have the \emph{same} fixed points as $S$, $T$ must therefore either map each of the two fixed points of $S$ into itself [that is, if $S$ is not parabolic, possess the same fixed points as $S$, and if $S$ is parabolic, have at least one fixed point coinciding with the fixed point of $S$] or $T$ must interchange the two fixed points $\sigma', \sigma''$ of the (non-parabolic) transformation $S$ with each other, that is, map $\sigma'$ into $\sigma''$ and $\sigma''$ into $\sigma'$. To effect the latter, as follows from (54), a parabolic $T$ is unable; a non-parabolic $T$, (53), can do so only if $\mu^2 = 1$, hence $\mu = -1$, that is, if $T$ is an elliptic transformation of rotation angle $\pi$. If in particular $S$ is to be permutable with $T$:
\[ TS = ST, \quad TST^{-1} = S, \]
this is consequently possible only if

I. both transformations $S$ and $T$ are non-parabolic and possess the same fixed points, or

II. both transformations $S$ and $T$ are parabolic and possess the same fixed point, or

III. $S$ and $T$ are elliptic transformations\ednote{Misprint in the original: ``Transformaticnen'' instead of ``Transformationen''.} of rotation angle $\pi$, or

IV. one of the two transformations $S, T$ is the identity.

If $\varGamma$ consists solely of the identity, $\mathfrak{F}$ is conformally equivalent to the interior of the unit disc or, what is the same thing, to a spherical cap [exceptional case 4].

If $S$ is any substitution from $\varGamma$ distinct from the identity and we assume, in contradiction to our assertion, that a sequence of operations $T_n \neq 1$ from $\varGamma_c$ converges to the identity $1$, then along with $T_n S T_n^{-1}$ also $T_n S T_n^{-1} S^{-1}$ would be contained in $\varGamma$. This operation, however, just like $T_n$, converges to the identity with indefinitely growing $n$, and since $\varGamma$ contains no infinitesimal transformations, we must have from a finite $n$ onwards
\[ T_n S T_n^{-1} S^{-1} = 1, \]
that is, $S$ must be permutable with $T_n$. Of the previously enumerated four cases in which this is possible, only cases I. and II. come into consideration here, since $S$ is not elliptic. Applying the result resulting therefrom to all possible operations $S \neq 1$ of the group $\varGamma$, we recognize that only two possibilities remain open to us:

\origpage{165}I. All operations of $\varGamma$ are hyperbolic and have the same two fixed points, which one can transfer to the places $+i$ and $-i$;

II. All operations of $\varGamma$ are parabolic and have the same fixed point, which may lie at the place $i$.

Case~I.: By the mapping
\[ z = \lg \frac{i - t}{i + t} \]
the interior of the unit disc of the $t$-plane passes into the parallel strip $-\frac{\pi}{2} < y < +\frac{\pi}{2}$ of the $z = (x + iy)$-plane, and $\varGamma$ into a group of translations parallel to the $x$-axis, that is, $\varGamma$ acquires the form
\[ z^* = z + n \cdot \frac{2\pi}{a} \qquad \begin{bmatrix} a \text{ a positive constant;} \\ n = 0, \pm 1, \pm 2, \dots \end{bmatrix}. \tag{58} \]
$\mathfrak{F}_\varGamma$ is then by means of
\[ w = e^{aiz} \]
mapped bijectively onto the circular annulus
\[ e^{-\frac{1}{2}a\pi} < |w| < e^{\frac{1}{2}a\pi}, \]
which one can finally transform by stereographic projection into a spherical zone whose midline is the equator: exceptional case 6.

Case~II.: By the mapping
\[ z = \frac{2}{t - i} - i \qquad (z = x + iy) \]
the interior of the unit disc of the $t$-plane passes into the upper half-plane $y > 0$ and $\varGamma$ into a group of translations parallel to the $x$-axis; $\varGamma$ thus again acquires the form (58), and $w = e^{aiz}$ transforms $\mathfrak{F}_\varGamma$ into the punctured unit disc of the $w$-plane:
\[ 0 < |w| < 1\colon \qquad \text{exceptional case 5.} \]

The main theorem is therewith completely proved, and we have at the same time found occasion to become acquainted with the applicability of the normal form $\mathfrak{F}_\varGamma$ of Riemann surfaces in an important example. That in the seven excluded cases a continuous group of conformal mappings of the surface onto itself really exists is trivial. It is also easy to specify these groups completely.

\emph{A closed Riemann surface of genus $p > 1$ admits only finitely many conformal mappings onto itself.}${}^{1)}$ For a system of points equivalent with respect to this group may exhibit no accumulation point on the closed surface, and can therefore consist of only finitely many points.

\textbf{1)} For this more special theorem, first stated by H.~A.~Schwarz, other more algebraic proofs exist, one by Weierstrass (1875), published only in 1895 (Werke vol.~II, pp.~235---244), one by Noether (Math.\ Ann.\ vol.~20, pp.~59---62, and vol.~21, pp.~138---140; 1882), and one by Hurwitz (Math.\ Ann.\ vol.~41, pp.~403---411; 1893).

\origpage{166}

\section*{Index of Technical Terms.}

[The numbers indicate the page on which the respective technical term is introduced.]

\subsection*{A.}

Mapping 19.

---, conformal 36.

Abelian function 132.

Abel's theorem 126, 136.

Closed point set 18.

Adaptation of a lattice to an included lattice 73.

Additive function 117.

Algebraic configuration 138.

--- function field 138.

Analysis situs 20.

Analytic differential 55.

Analytic continuation, immediate 2.

--- ---, mediate 2.

--- --- along a curve 2.

--- --- on a Riemann surface 43.

Analytic function in the Weierstrassian sense 4.

--- --- on a Riemann surface 36.

Analytic configuration 11.

Analytic curve 39.

Analytic neighborhood 9.

Analytically connected series of function elements 11.

Equivalence in the sense of analysis situs 20.

--- in the sense of conformal mapping 36.

Equivalent pairs of power series 6.

--- points 27, 151.

Automorphic function 151.

\subsection*{B.}

Basis of closed paths 74.

--- of a linear family 68.

Boundary of a closed set 78.

Arbitrarily fine partition 31.

Domain determined (by a closed point set) 58.

Motion of the non-Euclidean plane 152.

Image sense of rotation 60.

Image curve 19.

Image set 19.

Birational transformation 139.

Bolyai-Lobachevskian geometry 152.

\subsection*{C.}

Character system (of an Abelian group) 114.

\subsection*{D.}

Representation of a function element 6.

Cap 93.

Deck transformation 50.

(a point) covers (another) 50.

Differential 55.

---, multiplicative 129.

---, Prym 129.

---, regular-analytic 55.

--- of the 1st, 2nd, 3rd kind 96, 97.

Dirichlet integral 86.

Discontinuous group 151.

Divisor 120.

Rotation in the non-Euclidean plane 153.

Sense of rotation, uniform 61.

---, continuous 61.

Triangle 21.

Star of triangles 22.

\subsection*{E.}

Plane, non-Euclidean 152.

---, projective 25.

Vertex 22.

One-sheeted covering surface 47.

Simple chain of triangles 23\ednote{Misprint in the original: after 23 the period is missing.}

Simple polygonal path 44.

\origpage{167}
Simply connected 47.

Uniform sense of rotation 61.

One-sided surface 61.

Elementary differentials of the 1st, 2nd, 3rd kind 100.

Elementary triangle 23.

Elementary segment 31, 43.

End (curve without end) 145.

Distance, non-Euclidean 154.

\subsection*{F.}

Fixed point 150.

Surface 23.

---, simply connected 47.

---, one-sided 61.

---, perforated 93.

---, closed 24.

---, open 24.

---, Riemann 36.

---, schlichtartig 45.

---, two-sided 61.

Continuation, analytic 2, 43.

Fundamental domain 154.

Function, Abelian 132.

---, additive 117.

---, analytic 4, 36.

--- on a two-dimensional manifold 18.

--- on a Riemann surface (in the narrower sense) 43.

---, automorphic 151.

Function field, algebraic 138.

---, hyperelliptic 139.

Function, harmonic 38.

---, multiplicative 117.

---, potential 38.

---, regular-analytic, on a Riemann surface 36.

Function element (in the narrower sense) 1.

--- (in the wider sense) 6.

---, regular 8.

---, branched and unbranched 9.

Function, continuously differentiable 39.

\subsection*{G.}

Domain 18.

Domain-continuous 19.

Configuration, algebraic 138.

---, analytic 11.

Perforated surface 93.

Geometry, Bolyai-Lobachevskian 152.

Directed segment 44.

Total order of a divisor 120.

Genus 76.

Closed surface 24.

--- polygonal path 44.

--- polyhedron 52.

Separated banks 65.

Lattice 28, 72, 127.

Equation, irreducible 137.

Degree of a linear family 69.

Boundary of a domain 58.

Limit-circle uniformizing variable 152.

Base surface 47.

Group, discontinuous 151.

Validity of a representation 9.

\subsection*{H.}

Harmonic function 38.

Accumulation point 18.

Principal character 129.

Principal part of an additive function 117.

--- of a differential 100.

Homology of closed paths 71.

--- of integral functions 68.

Hyperbolic transformation 154.

Hyperelliptic 139.

\subsection*{I.}

Indicatrix 62.

---, coherent 62.

Induced 16, 62.

Infinitesimal transformation 159.

Area 78.

Interior point 18.

Interior edge 52.

Integral characters of a divisor 125.

--- of a point 114.

Integral, Dirichlet 80.

---, Poisson 83.

Integral function 68.

---s, linearly dependent 68.

Irreducible equation 137.

Isolated point 18.

\subsection*{K.}

Canonical dissection 76.

Edge 22.

---, interior, and boundary edge 52.

Edge path 44.

Chain of triangles 23.

Class of Riemann surfaces 153.

--- of uniformizing variables 148.

Coherent indicatrices 62.

Conformally equivalent 36.

Conformal mapping 36.

\origpage{168}
Circle of convergence 1.

Coordinate ratio 21.

Circle 80.

Critical point 4.

Curve 18.

---, analytic 39.

--- without end 145.

---, continuously differentiable 39.

Curve function 68.

---, linear 68.

\subsection*{L.}

(a point) lies over (another) 47.

Limit 100.

Linearly dependent integral functions 68.

Linear curve function 68.

Left 76.

Lobachevskian geometry 152.

Hole 93.

\subsection*{M.}

Manifold, two-dimensional 17.

Center of a function element 1.

M?bius strip 26.

Modulus of Riemann surfaces 41.

Multiplicative differential 129.

Multiplicative function 117.

Multiplier of a transformation 159.

Multiple of a divisor 120.

\subsection*{N.}

Non-Euclidean motion 152.

--- plane 152.

--- distance 154.

Normal representation of a function element 8.

Normal form of a Riemann surface 151.

Normal polygon 155.

Zero of a differential 55.

--- of a function 38.

\subsection*{O.}

Open surface 24.

--- polyhedron 52.

Order of the zero of a differential 55.

--- --- --- of a function 38.

--- of the pole of a differential 55.

--- --- --- of a function 38.

--- of a differential at a point 55.

--- of a function at a point 38.

--- of a group 161.

Order of a point with respect to a curve 56.

--- --- (total order) of a divisor 120.

---, branching 9.

Local uniformizing variable 36.

\subsection*{P.}

Parabolic transformation 160.

Periods of an integral function 77.

Poincar? $\Theta$-series 158.

Pole of a differential 55.

--- of a function 38.

Polyhedron 52.

---, closed and open 52.

Polygon 44.

(Normal) polygon 155.

Potential function 38.

Projective plane 25.

Prym differential 129.

Point of a two-dimensional manifold 17.

Point lattice 28, 72, 127.

Point, interior 18.

---, isolated 18.

---, critical 4.

\subsection*{Q.}

Crosscut 52.

\subsection*{R.}

Boundary edge of an open polyhedron 52.

Right 76.

Regular-analytic differential 55.

--- --- function on a Riemann surface 36.

Regular function element 8.

--- covering surface 50.

Residue 56.

Riemann surface 36.

Riemann-Roch theorem 122.

--- ---, generalized 130.

Pair of loop cuts 75.

\subsection*{S.}

Scheme of a surface 29.

Schlichtartig surface 45.

Singular point of a function 38.

Jump of an integral function 77.

Trace point 47.

Continuous mapping 19.

--- sense of rotation 61.

--- function 18.

Continuously differentiable function 39.

--- --- curve 39.

\origpage{169}
Star 22.

Segment, directed 44.

Polygonal path 44.

---, simple 44.

---, closed 44.

Symbol of a differential 120, 129.

--- of a function 119.

System of equivalent points 27, 151.

\subsection*{T.}

Sub-triangle 31.

$\Theta$-series 158.

Topology 20.

Torus 27.

Transformation, birational 139.

---, elliptic 160.

---, hyperbolic 159.

---, infinitesimal 159.

---, parabolic 160.

Triangulation 22.

$t$-neighborhood 10.

\subsection*{U.}

Covering surface 47.

--- of the integral functions 74.

---, one-sheeted 47.

---, regular 50.

---, unbounded 47.

---, universal 50.

---, unramified 47.

(a polygonal path) does not self-intersect 44.

Bank 45.

---s, separated 65.

Neighborhood, analytic 9.

--- on a surface 17.

Inversion problem 128.

Unbounded covering surface 47.

Uniformizing variable 141.

(Limit-circle) uniformizing variable 152.

(Local) uniformizing variable 36.

(belonging to $\mathfrak{E}$) uniformizing variable 81.

Uniformization principle 148.

Universal covering surface 50.

Immediate analytic continuation 2.

Subdivision 31.

Unbranched function element 9.

--- covering surface 47.

\subsection*{V.}

Variable, uniformizing 141.

Condensation point 18.

Closing ring 93.

Branched function element 9.

Branching order 9.

Branch point 32.

Branching number 135.

\subsection*{W.}

Value of an analytic function 5.

--- of a function on a surface 18.

--- of a curve function 68.

Essential singularity 38.

Angle 39.

\subsection*{X. Y.}

(vacant.)

\subsection*{Z.}

Decomposition into elementary triangles 23.

Decomposition of a domain by a set 44.

Dissection, canonical 76.

(Analytically) connected series of function elements 11.

(Simply) connected surface 47.

Degree of connectivity 69.

Two-dimensional manifold 17.

Two-sided surface 61.

Cycle of vertices 156.
\end{document}
